\documentclass[11pt]{amsart}
\usepackage{lmodern}
\usepackage[expansion=false]{microtype}
\usepackage[a4paper,textwidth=152mm,textheight=242mm,centering]{geometry}
\usepackage{amsmath,amssymb,mathtools,booktabs,enumitem,needspace}
\usepackage[colorlinks,linkcolor=blue!50!black,citecolor=blue!50!black,urlcolor=blue!50!black,filecolor=blue!50!black]{hyperref}
\usepackage{xcolor}
\setlength{\emergencystretch}{1.4em}
\numberwithin{equation}{section}
\newtheorem{theorem}{Theorem}[section]
\newtheorem{lemma}[theorem]{Lemma}
\newtheorem{proposition}[theorem]{Proposition}
\newtheorem{corollary}[theorem]{Corollary}
\theoremstyle{definition}
\newtheorem{definition}[theorem]{Definition}
\newtheorem{example}[theorem]{Example}
\newtheorem{remark}[theorem]{Remark}
\theoremstyle{plain}
\newtheorem{mainthm}{Theorem}
\newtheorem{maincor}[mainthm]{Corollary}
\renewcommand{\themainthm}{\Alph{mainthm}}
\DeclareMathOperator{\lcm}{lcm}
\DeclareMathOperator{\Tr}{Tr}
\DeclareMathOperator{\FPdim}{FPdim}
\DeclareMathOperator{\Irr}{Irr}
\DeclareMathOperator{\Gal}{Gal}
\newcommand{\QQ}{\mathbb Q}
\newcommand{\ZZ}{\mathbb Z}
\newcommand{\C}{\mathcal C}
\newcommand{\supfile}[2]{\href{run:rank-bound-supplement/#1}{#2}}
\allowdisplaybreaks[2]
\title[The rank of simple integral modular data]{The rank of simple integral modular data}
\author{Sebastian Burciu}
\address[S. Burciu]{Simion Stoilow Institute of Mathematics of the Romanian Academy, Bucharest, Romania}
\email{sebastian.burciu@imar.ro}
\author{Jingcheng Dong}
\address[J. Dong]{College of Mathematics and Statistics, Nanjing University of Information Science and Technology, Nanjing, China}
\email{jcdong@nuist.edu.cn}
\author{S\'ebastien Palcoux}
\address[S. Palcoux]{Beijing Institute of Mathematical Sciences and Applications, Beijing, China}
\email{sebastien.palcoux@gmail.com}
\date{}
\hypersetup{pdftitle={The rank of simple integral modular data},pdfauthor={Sebastian Burciu, Jingcheng Dong and Sebastien Palcoux}}
\begin{document}
\lefthyphenmin=2 \righthyphenmin=3\relax

\begin{abstract}
A non-pointed based-simple integral fusion ring with symmetric Drinfeld
duality, cyclotomic modular data and compatible rotation traces has
rank at least $45$, at least $217$ in odd dimension, and at least $219$
in odd nonsquare dimension. These bounds
apply to non-pointed simple integral modular fusion categories.
An actual or effective row of odd prime-square codegree has degree
prime to that prime. In square dimension a prime denominator $p$ occurs
with multiplicity $pj$ for a proper divisor $j$ of $p-1$, and its $j$
neutral labels form one Galois orbit. A projection onto a phase fiber
with a single eligible neutral duality orbit gives related restrictions
at two distinct primes. In odd dimension one diagonal character value
determines the twist, and local phase moments bound both ordinary mass
and Gauss sums. We also classify the nontrivial perfect weakly integral modular categories
through rank $44$. At rank $44$ they are the twisted doubles of $A_6$
and the $A_5$-gaugings of the pointed modular category on the deleted
binary permutation module.
The latter give the least rank at which a perfect weakly integral modular
category can fail to be group-theoretical.
\end{abstract}

\maketitle

\section*{Introduction}
Let $R$ be a commutative fusion ring with distinguished basis
$\{x_i\}_{i\in I}$, unit $x_0$ and Frobenius--Perron degrees
$d_i=\FPdim(x_i)$. Its rank is $|I|$. The ring is \emph{integral} if
all $d_i$ are integers, \emph{pointed} if all $d_i=1$, and
\emph{based-simple} if each nonunit basis element generates $R$ under
products, duality and passage to constituents. For example, the
representation ring of $A_5$ is non-pointed, based-simple and integral
of rank five \cite[\S5]{LPR}. Its degrees are $1,3,3,4,5$ and
$9\nmid60$, so its probability basis cannot admit symmetric Drinfeld
duality (Proposition~\ref{alg:duality}). We ask how small the rank of a non-pointed based-simple ring can be when its character
table and twists satisfy the arithmetic of modular and rotation data.

\begin{mainthm}\label{alg:rank}
Let $R$ be a non-pointed based-simple commutative integral fusion ring.
Suppose its probability basis $\{x_i/d_i\}$ admits a symmetric
algebraic Drinfeld map, and $R$ carries the cyclotomic modular data
and rotation certificates of Section~\ref{alg:modular}. Then
$\operatorname{rk}(R)\ge45$. If its global dimension is odd, then
$\operatorname{rk}(R)\ge217$, and if it is odd and nonsquare, then
$\operatorname{rk}(R)\ge219$.
\end{mainthm}

The probability basis identifies the fusion basis with its characters. The extra hypotheses concern the congruence representation of $\mathrm{SL}_2(\ZZ)$ and the traces of rotations on tensor powers. Proposition~\ref{formal:criterion} gives finite sufficient tests for the rank theorem. The positive spherical structure and generalized Frobenius--Schur rotations supply these data in a modular fusion category (Proposition~\ref{alg:realization}).
\begin{maincor}\label{main:rank}
Every non-pointed simple integral modular fusion category over
$\mathbb C$ has rank at least $45$, at least $217$ in odd dimension,
and at least $219$ in odd nonsquare dimension.
\end{maincor}

Whether such a category exists is open. Liu, Palcoux and Ren showed
that existence is equivalent to the existence of a simple integral
fusion category that is not weakly group-theoretical
\cite[Question~1.3 and Theorem~5.8]{LPR}. This connects the rank
problem to the question of Etingof, Nikshych and Ostrik on fusion
categories of integer dimension \cite[Question~2]{ENOWGT}.

A non-pointed based-simple ring is \emph{perfect}, meaning that the
unit is the only basis element of degree one. In a fusion category
this means that the unit is the only invertible simple object.
Alekseyev, Bruns, Palcoux and Petrov excluded nontrivial perfect
integral modular categories through rank $13$, and through rank $24$ in odd dimension \cite[Theorems~1.5 and~12.2]{ABPP}.
Dong and Palcoux obtained lower bounds of $15$ in general and $27$
in odd dimension \cite[Corollaries~1.2 and~1.3]{DP}, completing in
odd dimension the rank-$25$ analysis of Czenky, Gvozdjak and Plavnik
\cite[Theorem~1.1]{CGP}. Berean-Dutcher, C\'ordova, Parra-Martinez
and Tuddenham searched numerically for modular data through rank $22$
\cite[Introduction and \S3]{Atlas}.

Burciu, Dong and Palcoux established categorical lower bounds of $32$ in general and $57$ in odd dimension \cite[Theorem~22.2]{BDPcore}.

For a nontrivial perfect weakly integral modular category of rank at
most $44$, Corollary~\ref{main:rank} forces a minimal nontrivial
fusion subcategory to be Tannakian. Counting simples in the resulting
equivariantization yields the classification in Corollary~\ref{ext:perfect}.
The least such rank is $22$, answering the rank part of
\cite[Question~11.8]{BPBurnside}; the examples there are twisted doubles
of $A_5$. The rank-$37$ family occurs in the icosahedral orbifold
constructions of Kac--Todorov and Mason--Ng
\cite[\S6B]{KacTodorov}, \cite[\S4, Example~3]{MasonNgOrbifolds}.
At rank $44$ there are $24$ twisted doubles of $A_6$ and $18$
permutation gaugings of dimension $57\,600$. In particular, $44$ is
the least rank of a perfect weakly integral modular category that is
not group-theoretical.

\paragraph{From characters to rank.}
Write $D=\sum_i d_i^2=qM^2$ with $q$ square-free. The formal codegrees
$f_i=D/d_i^2$ are integers. Thus $b_i=M/d_i$ are integers and
$\sum_i b_i^{-2}=q$. This is the square-denominator form of the rank-finiteness argument
\cite[Proposition~8.38]{ENOannals}, \cite[Lemma~2.3]{HR}, \cite[\S4]{ABPP}.
The least denominator measures how much of this sum one label can supply.

The primes impose restrictions in two different ways. Reduction of a character at a prime defines its block. The primary factors of the modular representation determine which twists can occur in that block. A square-free codegree is impossible in a nontrivial perfect ring (Theorem~\ref{thm:squarefree-codegree}). Across each partition of the dimension primes, at least three distinct nonunit denominators meet both sides and their least common multiple is $M$ (Theorem~\ref{new:prime-cut}). In square dimension the complete class at a prime denominator $p$ has size $pj$, where $j$ is a proper divisor of $p-1$. Its $j$ members of trivial $p$-phase form a single Galois orbit (Theorem~\ref{new:one-packet}). We use these facts to bound the finite denominator search.

Balancing and character orthogonality give further restrictions on a
degree multiset. They bound the multiplicities carried by one twist and
the sums of principal-block entries over each degree class. Common row
support also constrains primary phases.

Odd dimension affords a different approach. Duality pairs every nonunit label. Pairing terms in the modular relation at two recovers the odd-order twist from a single diagonal character value. Pairing terms in a partial Weyl operator detects the components of the row field at every activity level. Section~\ref{sec:odd-new} combines these facts with the prime support and phase-mass bounds. In particular a phase fiber containing only one eligible neutral dual
pair gives the trace constraints of a prime-square row (Proposition~\ref{new:one-neutral-fiber}).
On three-prime support this excludes a primitive fifth-phase fiber with a
single neutral pair of denominator fifteen.

\paragraph{The finite input.}
The covering denominator census through rank $44$ contains the square-
and nonsquare-dimensional multisets in the first two columns of
Table~\ref{tab:counts}. Its purpose is to contain every possible degree
multiset in the stated scope; a list in this census need not be
realizable. The next five columns count those remaining
after the primary mass bound, the prime-square gcd condition, defect-one
block tiling, principal-block support, and the ternary trace congruence,
in that order. Section~\ref{sec:finish} proves the five tests and
specifies the separate checks of search coverage and terminal exclusion.

\begin{table}[htbp]
\caption{Degree multisets in the covering census and survivors after each successive test.}
\label{tab:counts}
\centering\footnotesize\setlength{\tabcolsep}{5.3pt}
\begin{tabular}{@{}lrrrrrrr@{}}\toprule
 & \multicolumn{2}{c}{Covering census} & \multicolumn{5}{c}{Successive survivors} \\
\cmidrule(lr){2-3}\cmidrule(lr){4-8}
Rank & Square & Nonsquare & Mass & Gcd & Tiling & Block & Trace \\ \midrule
$1$--$38$ (each) & $0$ & $0$ & $0$ & $0$ & $0$ & $0$ & $0$ \\
39 & $1$ & $0$ & $1$ & $1$ & $1$ & $0$ & $0$ \\
40 & $2$ & $0$ & $2$ & $2$ & $2$ & $0$ & $0$ \\
41 & $17$ & $4$ & $17$ & $15$ & $15$ & $0$ & $0$ \\
42 & $111$ & $65$ & $119$ & $98$ & $98$ & $0$ & $0$ \\
43 & $914$ & $723$ & $955$ & $801$ & $801$ & $30$ & $0$ \\
44 & $6\,269$ & $10\,093$ & $7\,113$ & $5\,277$ & $5\,249$ & $214$ & $0$ \\
\midrule Total & $7\,314$ & $10\,885$ & $8\,207$ & $6\,194$ & $6\,166$ & $244$ & $0$ \\
\bottomrule
\end{tabular}
\end{table}

\section{Symmetric integral character data}\label{alg:framework}

\begin{definition}\label{alg:based}\label{alg:probability}\label{alg:drinfeld}
Let $R=\bigoplus_{i\in I}\ZZ x_i$ be a commutative fusion ring, with
unit $x_0$, involution $x_i^*=x_{i^*}$, fusion coefficients $N_{ij}^k$,
degrees $d_i=\FPdim x_i$, rank $r=|I|$ and $D=\sum_i d_i^2$.
A based subring is closed under duality and constituents. The ring is
based-simple if only $\ZZ 1$ and $R$ occur, and perfect if $d_i=1$ only
at the unit. Write $c_i=x_i/d_i$. A symmetric algebraic Drinfeld map
identifies the $c_i$ with the characters $\phi_i$ so that
$\phi_i(c_j)=\phi_j(c_i)$. Write $\chi_i(j)=\phi_i(x_j)$,
$S_{ij}=d_i\chi_i(j)$ and $s=S/\sqrt D$.
\end{definition}

\begin{proposition}[Normalization and duality]\label{alg:normalization}\label{alg:duality}
In the probability basis $c_i c_j=\sum_k p_{ij}^k c_k$,
$d_i=(p_{ii^*}^0)^{-1/2}$ and $N_{ij}^k=d_i d_j p_{ij}^k/d_k$.
A probability algebra with this normalization gives a fusion ring
when all $N_{ij}^k$ are integral. The symmetric map identifies its probability
basis with the dual basis under pointwise multiplication.
Then $S_{i0}=d_i$, $s$ is unitary, and
$f_i=\sum_j|\chi_i(j)|^2=D/d_i^2$. If all degrees are integral, $f_i$ is an integer and
$S_{ij}/\lcm(d_i,d_j)$ is an algebraic integer. If, in addition,
the character values are cyclotomic, then
$\sigma(S_{ij})=S_{\sigma i,j}=S_{i,\sigma j}$, where
$i\mapsto\sigma i$ is a degree-preserving permutation.
\end{proposition}
\begin{proof}
For rescaling and orthogonality see \cite[\S\S2.1--2.2]{BPBurnside}.
Evaluation at $c_l$ and symmetry identify the structure constants on the two bases.
The diagonal $SS^*=S^*S=\operatorname{diag}(d_i^2 f_i)$ commutes with $S$,
so its nonzero vacuum column makes it scalar. Its vacuum entry is $D$.
Under the integrality hypothesis the formal codegrees
$f_i=D/d_i^2$ are rational algebraic integers, hence integers.
B\'ezout applied to $S_{ij}/d_i$ and $S_{ij}/d_j$ gives the lcm
assertion. Galois permutes the character rows. Their formal
codegrees satisfy $f_{\sigma i}=\sigma(f_i)=f_i$, so
$d_{\sigma i}=d_i$. Symmetry now gives the displayed covariance.
\end{proof}

In a non-pointed based-simple ring, weak integrality already implies integrality by the dimension grading \cite[Theorem~3.10]{GelakiNikshych}: the integral component contains a nonunit constituent of $x_ix_i^*$ for $d_i>1$.

\section{Modular and rotation data}\label{alg:modular}

\begin{definition}[Cyclotomic modular data]\label{alg:mod-data}
We use the usual $S,T$ formulation of modular data
\cite[Definition~1]{GannonData} with the following arithmetic hypotheses.
On an integral fusion ring with symmetric Drinfeld duality, choose
roots of unity $\theta_i$ with $\theta_0=1$ and a root $\gamma$ of
two-power order. Put $T=\operatorname{diag}(\theta_i)$,
$N=\operatorname{ord}(T)$ and $L=\operatorname{ord}(T/\gamma)$.
Require that $N,D$ have the same prime divisors and that
\[
 W=\begin{pmatrix}0&-1\\1&0\end{pmatrix}\longmapsto s,\qquad
 U=\begin{pmatrix}1&1\\0&1\end{pmatrix}\longmapsto T/\gamma
\]
give a representation $\rho$ of $\mathrm{SL}_2(\ZZ/L)$ satisfying
\begin{equation}\label{alg:covariance}
 \sigma_a\rho(g)=\rho(D_a gD_a^{-1}),\qquad D_a=\operatorname{diag}(a,1),
 \quad \sigma_a(\zeta_L)=\zeta_L^a\quad(a\in(\ZZ/L)^\times).
\end{equation}
It suffices to check covariance on $W$, as it holds on $U$.
\end{definition}

Here are the identities implicit in this definition
\cite[\S1.3]{DLNCongruence}. The unit entry gives $N\mid L$.
Symmetric duality gives $s^2=C=\rho(-I)$, the charge-conjugation
permutation, so $\theta_{i^*}=\theta_i$. From
$WUW=U^{-1}WU^{-1}$ and its vacuum entry,
\[
 sTs=\xi T^{-1}sT^{-1},\qquad
 \xi=\gamma^3=D^{-1/2}\sum_i d_i^2\theta_i .
\]
For $v_i=d_i\theta_i$ the vacuum column gives
$(s^{-1}v)_\ell=\sqrt D\,\xi s_{0\ell}/\theta_\ell$.
Diagonalizing fusion and inverting the modular relation gives balancing:
\[
 \sum_kN_{i^*j}^k v_k
 =\sum_\ell\frac{\overline{s_{i\ell}}s_{j\ell}}{s_{0\ell}}(s^{-1}v)_\ell
 =\sqrt D\,\xi(sT^{-1}s^{-1})_{ji}=\theta_i\theta_j S_{ij}.
\]
In particular $S,\chi_i$ have entries in $\QQ(\zeta_N)$ and
$\sqrt D\in\QQ(\zeta_L)$, so covariance takes place over $\QQ(\zeta_L)$.
A local factor at a prime dividing $L$ but not $D$ acts by scalars:
its upper unipotent is scalar, as is the conjugate lower unipotent.

\begin{definition}[Rotation certificates]\label{alg:rot-data}
For $m\ge1$ define $L_m^i$ linearly by
\begin{equation}\label{alg:rotation}
 L_m^i(x_c)=\sum_{a,b}s_{ia}s_{ib}(\theta_a/\theta_b)^m N_{ab}^c .
\end{equation}
A certificate for $i,m,c$ is a multiset of $m$th roots $\lambda_j$
with $\sum_j\lambda_j^t=L_{m/t}^i(x_c^t)$ for every positive $t\mid m$.
A ring satisfying the preceding modular hypotheses and admitting all
these certificates is \emph{rotation-admissible}.
\end{definition}

\begin{lemma}\label{alg:rotation-arithmetic}
Every nonnegative integral combination $V$ of basis elements admits
the same certificates. All power traces belong to $\QQ(\zeta_N)$;
in particular $L_m^i(x_c)$ is integral over $\ZZ$ and belongs to
$\QQ(\zeta_m)\cap\QQ(\zeta_N)$. For $p\mid m$ and $y$ with
algebraic-integer coefficients,
$L_m^i(y)^p\equiv L_{m/p}^i(y^p)\pmod{\mathfrak p}$ at every place
above $p$. Also $L_{N_{p'}}^0(x_c)\in\ZZ$, where
$N_{p'}=N/p^{v_p(N)}$.
\end{lemma}
\begin{proof}
Regard copies of summands of $V$ as distinct letters. For each cyclic
orbit of length-$m$ words of least period $b$, let $Y$ be the product
of one period. Induce its certificate from $C_{m/b}$ to $C_m$; it
exists for $b=1$ by hypothesis and for $b>1$ by induction on $m$.
At the $t$th power ($t\mid m$) the induced trace vanishes unless
$b\mid t$, when it equals $bL_{m/t}^i(Y^{t/b})$. Summing counts each
word of length $t$ in $V^t$ once.

Since $s_{ia}s_{ib}=S_{ia}S_{ib}/D$, \eqref{alg:rotation} puts
divisor-power traces in $\QQ(\zeta_N)$; any other positive power
trace is Galois conjugate to that at $\gcd(t,m)$. The congruence on
basis elements follows by reducing the sum of roots above $p$, and
extends by the binomial theorem and commutativity. At $i=0,m=N_{p'}$,
\eqref{alg:rotation} involves only $p$-power roots with rational
weights. Its value also lies in $\QQ(\zeta_{N_{p'}})$, hence in $\ZZ$.
\end{proof}

\subsection{Number-theoretic conventions}\label{alg:arithmetic}
We normalize $v_p(p)=1$ and take congruences in residue fields.
The conductor of an abelian field $K$ is the least $n$ such that
$K\subseteq\QQ(\zeta_n)$. We use $\mathcal O_{\QQ(\zeta_n)}=\ZZ[\zeta_n]$,
trivial intersection for coprime conductors
\cite[Ch.~2, Theorem~2.6]{Washington}, and total ramification
(resp.\ unramifiedness) over $\QQ_p$ of $p$-power (resp.\ prime-to-$p$)
cyclotomic extensions \cite[Ch.~IV, \S4]{SerreLocal}. Over a complete
discrete valuation ring the local factors of a finite algebra lift
from its residue algebra \cite[Tag 04GH]{Stacks}.

For a cyclotomic integer $\alpha$, write $\langle\alpha\rangle$ for
its Galois mean and $C(\alpha)=\langle\alpha\overline\alpha\rangle$.
The norm and arithmetic--geometric mean inequality give $C(\alpha)\ge1$
when $\alpha\ne0$. For nonzero $\alpha$, equality or the bound $|\sigma\alpha|\le1$ at
every embedding forces a root of unity by Kronecker's theorem.
A root of order $k$ has mean $\mu(k)/\varphi(k)$
\cite[Theorem~8.6]{Apostol}. The integral power basis of a dyadic
cyclotomic field is orthonormal for $(v,w)\mapsto\langle v\overline w\rangle$:
the quotient of distinct basis powers has order at least four.
The mean of an element of this field is its constant coefficient.

Write $D=qM^2$ with $q$ square-free. Then $d_i\mid M$. If $D$ is
\emph{square-full} (every prime exponent is at least two), its prime
divisors are precisely those of $M$.

\begin{lemma}[Cyclotomic norms with prescribed reduction]\label{nt:residue-norm}\label{r6:dyadic-mean}
Suppose the cyclotomic integer $\alpha$ satisfies $\alpha\equiv s$ at every place above $p$, $0\le s<p$. Then $(p-1)C(\alpha)\ge s(p-s)$. Let $e$ be the size of the orbit of $\alpha$ under the subgroup of order $p-1$ of cyclotomic inertia (the trivial subgroup if $p=2$), and put $h=(p-1)/e$. In the fixed field of the stabilizer the sharp bound is
\[
 (p-1)C(\alpha)\ \ge p\min_{k\in\mathbb Z}\big((s-hk)^2+h|k|\big)-s^2 .
\]
At $e=p-1$ this recovers the first bound, the coset minimum of the scaled dual root lattice $A_{p-1}^*$ in the prime cyclotomic field \cite[p.~86]{ConwaySloaneLattice}.
If $\alpha\ne0$ lies in a $p$-power cyclotomic field and $s=0$,
then $C(\alpha)\ge2$. For a sum of $n$ roots of $p$-power order,
the residue is $s\equiv n\pmod p$.
\end{lemma}
\begin{proof}
At level $p$ work over $K=\QQ(\zeta_n)$, $p\nmid n$, with $\alpha\in K(\zeta_p)$. Write $\alpha=\sum_{i=0}^{p-1}c_i\zeta_p^i$, $c_i\in\mathcal O_K$. The sum of the coefficients equals $s$ modulo $p\mathcal O_K$ by unramifiedness, so subtracting a common coefficient gives a unique expansion with $\sum c_i=s$. The subgroup of order $h$ fixing $\alpha$ permutes its coefficients. Let $b_1,\dots,b_e$ be the common values on its nonzero exponent orbits, and put $z=\sum b_j$. Trace over the $p$-part gives
\[
 (p-1)C(\alpha)=p\big(C(s-hz)+h\sum_{j=1}^e C(b_j)\big)-s^2 .
\]
For an algebraic integer $v$ put $t=\Re\langle v\rangle$, with brackets denoting the Galois mean. Then $C(v)\ge|t|$, by Cauchy--Schwarz and the norm bound $C(v)\ge1$ for $v\ne0$. Applying this to $z-k$ shows that $C(z)$ dominates the piecewise linear interpolation of the integer squares at $\Re\langle z\rangle$. Consequently $C(s-hz)+h\sum C(b_j)$ dominates the piecewise linear interpolation of $(s-hk)^2+h|k|$ there. A minimizing $k$ lies in $[0,e]$. Giving $k$ of the $b_j$ the value one and the rest zero attains the bound in the field fixed by this subgroup.

For arbitrary ramification work in $K(\zeta_{p^a})$ and expand over $K(\zeta_p)$ in the first $p^{a-1}$ powers of $\zeta_{p^a}$, an integral basis \cite[Theorem~2.6]{Washington}. Their trace orthogonality is the cyclotomic mean-square decomposition of Cassels \cite{CasselsRobinson}; see \cite[Lemma~5.2.1]{CMSCyclotomic}. The basis exponents can be any representatives modulo $p^{a-1}$. Choose them equivariantly for the stabilizer $H$ of order $h$: it acts freely on every nonzero orbit modulo $p^{a-1}$, so lift one representative and take its $H$-translates. Replace the corresponding powers on each nonzero orbit by the $H$-conjugates of an independent uniform $p$th root, retaining the power one at exponent zero. The resulting cyclotomic integer $\beta$ is fixed by $H$ and still reduces to $s$ at every place above $p$. The replacements have zero mixed moments, also within a single orbit since its multipliers are distinct modulo $p$. Thus $\mathbb E C(\beta)=C(\alpha)$, and the level-$p$ estimate applies.
Finally suppose $\alpha\ne0$ lies in a $p$-power cyclotomic field,
$s=0$ and $C(\alpha)<2$. Its orthogonal integral expansion over
$\QQ(\zeta_p)$ has exactly one nonzero coefficient $c$, since each
costs at least one. This coefficient reduces to zero above $p$.
Its normalized expansion $c=\sum_{i=0}^{p-1}a_i\zeta_p^i$ has
$a_i\in\ZZ$ and $\sum a_i=0$, so
$C(\alpha)=C(c)=p\sum a_i^2/(p-1)\ge2p/(p-1)>2$, a contradiction.
\end{proof}

\begin{lemma}[Dyadic cancellation]\label{r7:coupled-gap}
Let $\alpha$ be a sum of an even number of roots of two-power order,
let $\beta$ be a sum of at most $k$ arbitrary roots of unity, and
let $u\in\QQ$ satisfy $u>k/2$. If $(\alpha,\beta)\ne(0,0)$, then
\[
 C(u\alpha+\beta)\ \ge\
 \min\{1,\,2(u-k/2)^2\}.
\]
Consequently, for $k=3$, a nonzero $9\alpha+4\beta$ has mean
squared modulus at least sixteen, and its vanishing forces
$\alpha=\beta=0$.
\end{lemma}
\begin{proof}
Project by relative trace onto a common dyadic cyclotomic field $K$
containing $\alpha$ and all dyadic parts of the roots in $\beta$.
The normalized trace of a root is its dyadic part times the
Ramanujan mean of its odd part, of absolute value at most one.
In the orthonormal power basis of $K$, the projection of $\beta$
therefore has coefficient $\ell^1$-norm at most $k$. The coefficient
vector of $\alpha$ is integral and has even sum, since every dyadic
root is a signed basis vector.

Suppose $\alpha\ne0$ and its coefficient support has size $m$.
For $m\ge2$, Cauchy--Schwarz on this support bounds the squared
norm of the projected sum below by
$(um-k)^2/m\ge2(u-k/2)^2$; the last expression is minimized at
$m=2$, since $u>k/2$. For $m=1$ the sole coefficient has magnitude
at least two, giving the stronger bound $(2u-k)^2$.
If $\alpha=0$ and $\beta\ne0$, the norm inequality gives
$C(\beta)\ge1$. Projection cannot increase squared norm.
\end{proof}


\begin{lemma}[Dyadic projection of a short Gauss relation]\label{r7:gauss-projection}
Let $A$ be a sum of an even number of dyadic roots and let $B$ be a sum of roots whose orders involve only two and three. Suppose a cyclotomic integer $Z$ satisfies
\[
 Z=\epsilon+\sum_j c_j\zeta_j,\qquad
 \epsilon,c_j\in\mathbb Q,\qquad L=|\epsilon|+\sum_j|c_j|,
\]
where the $\zeta_j$ are arbitrary roots of unity.
\begin{enumerate}[label=\textup{(\roman*)}]
\item If $B$ has at most eight terms, $Z=9A+4B$ and $L<2$, then $A=B=0$.
\item If $B$ has at most five terms, $E$ has zero or two dyadic terms, $Z=36A+16B+9E$ and $L<6$, then $(A,B,E)=(0,0,0)$ or $(0,-\eta,2\eta)$ for a dyadic root $\eta$.
\end{enumerate}
\end{lemma}
\begin{proof}
Choose $K$ containing all dyadic parts in the relation. The relative-trace argument of Lemma~\ref{r7:coupled-gap} bounds the coefficient $\ell^1$-norm of the projection of $Z$ by $L$.

If $B\notin K$, its squared distance from $K$ is at least $3/4$: the cyclotomic orthogonal decomposition \cite[Lemma~5.2.1]{CMSCyclotomic} either has a nonzero component above $K(\zeta_3)$, of mean square at least one, or its constant component is $P+Q\zeta_3$ with $Q\ne0$, of squared distance $3C(Q)/4\ge3/4$. In (i) this gives $C(Z)\ge12>L^2$, and in (ii) $C(Z)\ge192>L^2$, both impossible. Hence $B,Z\in K$.

Write their integer coefficient vectors in its orthonormal power basis as $b,z$, and that of $A$ as $a$. The sum of the entries of $a$ is even. In (i), $z=9a+4b$ has even coefficient sum, so $\|z\|_1<2$ forces $z=0$. Coprimality gives $A=4\delta$, $B=-9\delta$ for an algebraic integer $\delta$. Every conjugate of $\delta$ has modulus at most $8/9$, so $\delta=0$.

For (ii), $\|b\|_1\le5$ by relative trace, and the vector $e$ of $E$ is zero, one entry $\pm2$, or two entries $\pm1$. Whenever $e_j=z_j=0$, $9a_j+4b_j=0$ forces $a_j=b_j=0$ under this budget. If $e=0$, $\|z\|_1<6$ allows at most one nonzero entry $\pm4$, which would make the sole nonzero entry of $a$ odd. If $e$ has one entry, change its sign to two. The bound allows only $z_j=\pm2$ there and zero elsewhere. The even entry $a_j$ then solves $9a_j+4b_j=-4$ or $-5$; the second is odd, and the first, with $|b_j|\le5$, gives $a_j=0$. If $e$ has two entries, their absolute $z$-values are $(1,1)$ or $(1,3)$ in either order. A value one requires $|b_j|\ge4$, and a value three requires $|b_j|\ge3$, by reducing $36a_j+16b_j+9e_j=z_j$ modulo sixteen. This exceeds the budget five. The preceding cases give exactly the two asserted triples.
\end{proof}


\section{Dimensions and character fields}\label{sec:dimensions}
Throughout this section $R$ is integral with symmetric Drinfeld
duality and cyclotomic characters. Orthogonality and symmetry give
\begin{equation}\label{eq:foundorth}
 \sum_j\chi_i(j)\overline{\chi_k(j)}=\delta_{ik}f_i,\qquad
 \sum_{j\ne0}d_j\chi_i(j)=-1\quad(i\ne0).
\end{equation}

\begin{lemma}\label{lem:basic-table}
In a nontrivial perfect $R$, every nonunit row has a zero and $D$ is
square-full. If $x_i$ generates $R$, then $S_{ij}=0$ for nonunit $j$
with $\gcd(d_i,d_j)=1$, and at least two primes divide $d_i$.
A non-pointed based-simple $R$ is perfect and at least three primes
divide $D$.
\end{lemma}
\begin{proof}
For Burnside vanishing and square-fullness use \cite[Theorem~1.2]{BPBurnside} and the algebraic argument of \cite[Lemma~2.5]{BDPcore}. The same lemma gives strictness $|S_{ij}|<d_i d_j$ for nonunit $j$ when $x_i$ generates. Indeed the equality labels form a based subring \cite[Lemma~2.35]{BPBurnside}. For coprime degrees $S_{ij}/(d_i d_j)$ is integral, so the strict bound on all its conjugates forces zero. Orthogonality in \eqref{eq:foundorth} then rules out a prime-power generator. In a based-simple non-pointed ring the invertibles span a proper based subring, and all nonunit degrees have gcd one by the same orthogonality. This forces three dimension primes. For the categorical version see \cite[Lemma~7.1 and Corollary~7.2]{ENOWGT}.
\end{proof}

In the non-pointed based-simple case, Blichfeldt's divisibility
argument also applies (compare Cameron and Kiyota
\cite{CameronKiyota} for groups). For a nonunit degree $x$ of
multiplicity $m$ put $A=\sum_{d_i=x}x_i$,
$n=mx$ and $L=\{\chi_j(A):j\ne0\}$. Galois invariance and the
strict bound above give $L\subset\ZZ\cap(-n,n)$. For
$P(T)=\prod_{t\in L}(T-t)$, evaluation on the characters gives
$P(A)=D^{-1}P(n)\sum_i d_ix_i$. Its integral unit coefficient yields
$D\mid P(n)\mid(2n-1)!$.

\begin{proposition}[Fusion multiplicities from the twist of a constituent]\label{new:fusion-phase-cap}
Assume that $R$ is non-pointed and based-simple and carries the
cyclotomic modular data of Definition~\ref{alg:mod-data}.
Let $m$ be the least nonunit degree and $\ell$ the least prime
dividing $D$. For nonunit labels $i,j$ of degrees $a,b$, put
$w_\eta=\sum_{\theta_k=\eta}N_{ij}^k d_k$.
At least two $w_\eta$ are positive. If $\gcd(a,b)=1$, then
\[
 w_\eta\le ab/\ell\qquad\text{for every }\eta.
\]
Equality at one phase forces exactly $\ell$ equally weighted phases, forming a coset of the $\ell$th roots. In particular the product of two coprime degrees uses at least $\ell$ distinct twists.

For nonunit $i,j,k$ of degrees $a,b,c$, set $B(a,b)=\lfloor ab/\ell\rfloor$ for coprime $a,b$, $B(a,b)=ab-m$ for noncoprime $a\ne b$, and $B(a,a)=a^2-1$. Then
\[
 N_{ij}^{k^*}\le\min\{\lfloor B(a,b)/c\rfloor,\lfloor B(b,c)/a\rfloor,
 \lfloor B(c,a)/b\rfloor\}.
\]
For a specified nondual input pair one may also replace $B(a,b)$ by $\min\{B(a,b),ab-m\}$.
\end{proposition}
\begin{proof}
Balancing expresses $\sum_\eta w_\eta\eta$ as a unit scalar times $S_{i^*j}$. The strict character bound in the proof of Lemma~\ref{lem:basic-table} gives $|S_{i^*j}|<ab$, while $\sum w_\eta=ab$. This proves the first assertion. At coprime degrees $S_{i^*j}=0$ by that lemma. Divide the relation by a phase $\eta$ present and take its normalized cyclotomic trace. For a root of exact order $h>1$ the trace is $\mu(h)/\varphi(h)$ by the Ramanujan sum formula. Each prime of $h$ divides $D$, so this is at least $-1/(\ell-1)$, with equality only when $h=\ell$. Thus $0\ge w_\eta-(ab-w_\eta)/(\ell-1)$. Equality confines the relation to an $\ell$-gon; the minimal polynomial $1+t+\cdots+t^{\ell-1}$ forces equality of all its coefficients. Finally a second constituent contributes at least $m$ for a nondual pair and at least one otherwise. The uncapped three-way bound is \cite[Lemma~21.7]{BDPcore}. Apply the phase bounds to each input pair using the symmetry of $N_{ij}^{k^*}=[1](x_ix_jx_k)$.
\end{proof}


\begin{lemma}[Deleting a Galois orbit]\label{lem:deletion}
In a nontrivial perfect ring, the nonunit degrees remaining after
deletion of any one Galois orbit of nonunit labels have gcd one.
In particular at least two such orbits have degrees prime to each
prime $p$.
\end{lemma}
\begin{proof}
Choose a column in the orbit. By Lemma~\ref{lem:basic-table} some
nonunit row vanishes there. Its entries on the entire orbit are
Galois conjugate and hence vanish. Vacuum orthogonality
\eqref{eq:foundorth} expresses $-1$ as a sum over the remaining
nonunit degrees; compare \cite[proof of Theorem~1.1]{DP}.
\end{proof}

\paragraph{Enumerating denominators.}
Order $b_i=M/d_i$ for $i\ne0$ nondecreasingly. Perfection gives
\begin{equation}\label{eq:recip}
 q=M^{-2}+\sum_{i=1}^{r-1}b_i^{-2},\qquad qb_1^2<r .
\end{equation}
We use the ordered recurrence and terminal divisor tests of
\cite[\S18.2, Propositions~18.2--18.3]{BDPcore}. After a nonempty prefix of length $j$ let $L=\lcm(b_1,\ldots,b_j)$,
$h=r-1-j$ and $E=L^2(q-\sum_{i\le j}b_i^{-2})$. In the simple
case the next $b$ satisfies
\begin{equation}\label{eq:interval}
 b\ge b_j,\qquad L^2<Eb^2\le(h+1/36)L^2.
\end{equation}
Indeed the residual sum exceeds $1/b^2$ and is at most $h/b^2+1/M^2$,
with $M\ge6b$ by Lemma~\ref{lem:basic-table} ($M\ge2b$ and $1/4$
suffice under perfection). For $g=\gcd(L,b)$, $u=b/g$, the next
pair is $(uL,u^2E-(L/g)^2)$.

With two nonunit degrees $d\le e$ left, write $M=kL$. The prefix
degrees are multiples of $k$, so Lemma~\ref{lem:deletion} forces
$\gcd(k,de)=1$. Thus $d,e\mid L$, $Ek^2=1+d^2+e^2$ and $kL/e\ge b_j$.
Test divisor pairs, or enumerate $b=kL/e$ by \eqref{eq:interval} and
recover $e=L/\gcd(L,b)$, then $k,d$. If $2\mid k$ the right side
is $3\bmod4$; if $5\mid k$ its two squares cannot sum to $-1\bmod5$.
Hence $\gcd(k,10)=1$, $k^2\equiv1\bmod8$, and
$E\not\equiv0,4,7\bmod8$. With one degree left the deletion lemma
gives $k=1$ and $E=1+d^2$.

\begin{lemma}[Combining denominator reservations]\label{new:reservation-calculus}
Let $b_1\le\cdots\le b_n$ be denominators, all at least $B>0$.
Suppose each of several families reserves distinct positions with
specified denominator floors at least $B$. Pad each family's floors to length $n$
with $B$ and sort them increasingly, obtaining $a^{(\lambda)}$.
Then
\[
 b_i\ge a_i:=\max_\lambda a_i^{(\lambda)},\qquad
 \sum_i b_i^{-2}\le\sum_i a_i^{-2}.
\]
Families proved disjoint may first be merged. For undecided alternative
cases, take the coordinatewise minimum of their resulting bounds.
If floors $c_1\ge\cdots\ge c_v$ are charged to labels, each label
paying at most $h\in\mathbb Z_{>0}$ charges, they reserve distinct positions with floors
$c_1,c_{h+1},c_{2h+1},\ldots$.

In a closed degree class, suppose the labels eligible for a quota
form a Galois-stable union of size at most $m\in\mathbb Z_{\ge0}$,
and every orbit in it has size divisible by one of the positive
integers $d_1,\ldots,d_k$. The class can supply at most
\[
 Q(m;d_1,\ldots,d_k)
 =\max\bigl(\{0,\ldots,m\}\cap
               (d_1\mathbb Z_{\ge0}+\cdots+d_k\mathbb Z_{\ge0})\bigr)
\]
labels to that quota, with $Q(m;\varnothing)=0$.
Credits for separate quotas are upper bounds;
their residual demands may be added only when disjointness is proved.
\end{lemma}
\begin{proof}
A reservation matches its floors to distinct denominators no smaller
than those floors; sorting preserves the coordinatewise inequalities.
All simultaneous families hold, whereas at least one alternative case
holds. Among the first $jh+1$ charges at least $j+1$ labels are needed,
giving the charged floors. Finally, an eligible Galois-stable union is
a disjoint union of whole orbits, so its cardinality belongs to the
displayed additive semigroup.
\end{proof}

\begin{lemma}[Row conductor]\label{conduct}
The conductor $n$ of $K_i=\QQ(\chi_i(j):j)$ divides $f_i$.
\end{lemma}
This is \cite[Theorem~2.9]{BDPcore}. Proposition~\ref{alg:duality} gives precisely the symmetric integral character data of Definition~2.2 there.

\begin{lemma}[Odd dimension and duality]\label{alg:odd-parity}
If $D=qM^2$ is odd, every nonunit label has a distinct dual of the
same degree, every nonunit character field is nonreal, and
$D\equiv r\pmod{16}$ and $r\equiv q\pmod8$. Dual labels also have the same twist
when modular data are present.
\end{lemma}
\begin{proof}
The coefficient trace and the regular trace are congruent modulo two in odd dimension. Applying this to $x_i^2$ and $x_i$ gives $\delta_{i,i^*}=\delta_{i,0}$, as in \cite[Lemma~2.5]{BDPcore}; see also \cite[Corollary~8.2]{NSfs} in categories. Conjugation acts by duality, so every nonunit row field is nonreal. Since each odd square is one modulo eight and each nonunit occurs in a pair, $D-r=2\sum_{\{i,i^*\}}(d_i^2-1)$ is divisible by sixteen. Finally $M^2\equiv1\pmod8$ and $s^2=C$ commutes with $T$.
\end{proof}

\begin{lemma}[Quadratic incidence and its weighted norm]
\label{odd:quadratic-incidence}\label{new:quadratic-weighted-norm}
Let $K=\QQ(\sqrt\delta)$ be a quadratic field, with $\delta$
square-free, and let $h_1,\ldots,h_n$ be the lengths of the Galois label orbits whose
character fields contain $K$. Put $u=D/|\delta|$.
If $u$ is not a rational square, then $n=2k$ and
\[
 (-1)^k\prod_{i=1}^{2k}h_i
       \in N_{\QQ(\sqrt u)/\QQ}\bigl(\QQ(\sqrt u)^\times\bigr).
\]
In particular, $n$ is always even in odd square dimension.
If $n\equiv2\pmod4$ and $\prod h_i$ is a rational square,
then $u$ is a sum of two rational squares.
\end{lemma}
\begin{proof}
The assertion is immediate for $n=0$. Otherwise take a common
cyclotomic field of the character values and let $\eta$ be the
quadratic Galois character of $K$. Its isotypic space in the
rational permutation space on the labels has dimension $n$:
each transitive orbit contributes one exactly when its stabilizer
lies in $\ker\eta$. The orbit sign vectors have Gram matrix
$H=\operatorname{diag}(h_i)$. Symmetric Galois covariance gives
$S=\sqrt\delta B$ on this space, with $B$ rational and
$H$-self-adjoint. Charge conjugation acts by
$\operatorname{sgn}(\delta)$, so $S^2=DC$ gives $B^2=uI$.

For nonsquare $u$ the space is $k$-dimensional over
$E=\QQ(\sqrt u)$, hence $n=2k$. For an $E$-basis $(e_i)$,
the rational basis $(e_i,Be_i)$ has Gram matrix
$\left(\begin{smallmatrix}F&G\\G&uF\end{smallmatrix}\right)$,
where $F,G$ are symmetric. Its determinant is
$(-1)^kN_{E/\QQ}\det(G+\sqrt uF)$. A rational change of basis
multiplies the determinant by a rational square, itself a norm.
This proves the formula. In odd square dimension the only possible
square-$u$ exception is $K=\QQ(i)$, excluded by the odd conductor.
The last conclusion is immediate if $u$ is a square; otherwise
the hypotheses make $-1$ a norm:
$x^2-uy^2=-1$ with $x,y\in\QQ$ and $y\ne0$, so
$u=(x/y)^2+(1/y)^2$. Orbits are counted separately even when
their character fields coincide.
\end{proof}

\paragraph{Projection estimates.}\label{lin:projection}
We use Bessel's inequality and positive-semidefinite compression in
one normalization. If $v_a$ are orthogonal rows of squared norms
$f_a$, all orthogonal to the degree vector $d$, then
\begin{equation}\label{eq:projection-master}
 P=I-\sum_a f_a^{-1}v_av_a^*\succeq D^{-1}dd^*.
\end{equation}
Every coordinate compression of $P$ is positive semidefinite; a
restriction of the normalized row matrix is a contraction, and its
squared Hilbert--Schmidt norm is at most its rank. We use the
standard Schur-complement criterion for its small principal blocks
\cite[\S7.7]{HornJohnson}. The arithmetic restrictions on entries
and the equality cases are established at each application below.

\begin{proposition}[A rank inequality for character fields]\label{r5:field-rank}
Let $O$ be a column orbit under Galois conjugation. Take mutually
orthogonal, nonzero rows with algebraic integer entries, equivariant
under conjugation of the column, and partition them into families
$A_1,\ldots,A_k$. Assume that every selected row is nonzero at a
representative $Y$ of $O$. Let $F_i$ be the field generated by the
entries at $Y$ of the rows in $A_i$, and $h_i=[F_i:\QQ]$. Then
\[
 [F_1\cdots F_k:\QQ]\le |O|,\qquad
 |O|\sum_{v\in\cup A_i}\|v\|^{-2}\le 1+\sum_{i=1}^k(h_i-1).
\]
If the selected rows are nonunit character rows, or their effective
rows from Proposition~\ref{new:effective}, then the second bound
strengthens to
\[
 |O|\left(\frac1{f_Y}+\sum_{v\in\cup A_i}\|v\|^{-2}\right)
 \le 1+\sum_{i=1}^k(h_i-1).
\]
\end{proposition}
\begin{proof}
The stabilizer of $Y$ fixes all the displayed entries, giving the
first inequality. On $O$, every row in $A_i$ is the pullback of a
function on the embeddings of $F_i$. These $k$ function spaces have
a common constant line, giving the stated upper bound for the rank of
the restricted row matrix. Apply the contraction estimate in
\eqref{eq:projection-master} to the normalized rows. Each entry at $Y$ occurs on $O$ with all its conjugates
equally often; a nonzero algebraic integer has mean squared modulus
at least one. Thus each row contributes at least $|O|/\|v\|^2$.
For the stronger bound, append the normalized degree row. It is
orthogonal to the selected rows and constant on $O$, so adds no
dimension to the function space, but contributes $|O|d_Y^2/D$.
\end{proof}

\subsection{Principal blocks}\label{blk:principal}
Write $[i]=x_i$ and let $e_i$ be the idempotent of $\chi_i$ over a
splitting field. The unit coefficient $\tau(a)=\sum_i\chi_i(a)/f_i$
satisfies $\tau([i][j^*])=\delta_{ij}$, so $R$ is its own trace dual.
For $p\mid D$ the \emph{$p$-blocks} are the fibers of character
reduction above $p$. The partition is independent of the place:
$\sigma\chi_i(j)=\chi_i(\sigma j)$ permutes the arguments. The block
$P$ of the vacuum is Galois stable; its idempotent $e_P=\sum_{j\in P}e_j$
is the unit of the lifted dimension factor of $R\otimes\ZZ_p$
\cite[Tag 04GH]{Stacks}. More generally, over the completed valuation
ring $\mathcal O$ of a splitting field above $p$, every full block $B$
has a lifted idempotent $e_B\in R\otimes\mathcal O$. Hence
\begin{equation}\label{eq:block-trace}
 \sum_{j\in B}d_j^2\chi_j(y)=D\tau(e_By),\qquad
 \tau(e_By)\in\mathcal O\quad(y\in R\otimes\mathcal O).
\end{equation}

\begin{lemma}\label{blk:support}
The principal-block norm $A_i=\sum_{j\in P}|\chi_i(j)|^2$ is an
integer divisible by $p^{v_p(f_i)}$. Moreover
\begin{equation}\label{eq:principal-trace}
 \sum_{j\in P}d_j\chi_i(j)=\frac D{d_i}\tau(e_Px_i)
 \quad\text{is an integer divisible by }p^{v_p(D/d_i)}.
\end{equation}
A row of $p$-power codegree is supported in $P$, with support equal
to $P$ when $p\nmid d_i$.
\end{lemma}
\begin{proof}
Symmetry and \eqref{eq:block-trace} give
$A_i=f_i\tau(e_Px_ix_i^*)$ and \eqref{eq:principal-trace}.
Their left sides are Galois-invariant algebraic integers, and both
traces are $p$-integral. If $f_i$ is a $p$-power, then
$0<A_i\le f_i$ forces $A_i=f_i$, proving support in $P$.
For $j\in P$ and $p\nmid d_i$, principality gives
$\chi_j(i)\equiv d_i\ne0$, proving equality of supports.
\end{proof}

\begin{lemma}\label{blk:height}
If $p^2\mid D$, then $p^2\mid f_i$ for every $i\in P$.
\end{lemma}
\begin{proof}
This is the principal-block height bound of
\cite[Corollary~4.1]{BDPcore}. Its hypotheses are the integral
principal idempotent and the symmetric order at $p$, as above.
\end{proof}

\begin{proposition}[Degree-class moments and simultaneous block traces]\label{new:block-census}\label{new:cross-block-traces}
Let $R$ be non-pointed and based-simple, and let $Y$ have nonunit
degree $y$. On a nonunit degree class $x$ of multiplicity $m_x$, put
$g_x=\gcd(x,y)$, $u_x=x/g_x$ and
\[
 A_x=\sum_{d_i=x}\frac{\chi_Y(i)}{u_x},\qquad
 Q_x=\sum_{d_i=x}\left|\frac{\chi_Y(i)}{u_x}\right|^2.
\]
These are rational integers, with $|A_x|<m_xg_x$, and
\[
 \sum_x xu_x A_x=-1,\qquad \sum_xu_x^2Q_x=f_Y-1.
\]
If $A_x=m_xt+r$, $0\le r<m_x$, then
$Q_x\ge(m_x-r)t^2+r(t+1)^2\ge A_x^2/m_x$.
If $Y\in P_p$ and $p\nmid u_x$, then
\[
 A_x\equiv m_xg_x,\quad Q_x\equiv m_xg_x^2\pmod p,
 \qquad
 Q_x-2jA_x+m_xj^2\ge m_xC_p([g_x-j]_p)\quad(j\in\mathbb Z),
\]
where $C_p$ is any applicable bound in Lemma~\ref{nt:residue-norm}.
When $p\nmid x$, one also has $g_x>1$ and
$Q_x\ge m_xs_x(p-s_x)/(p-1)$, with $s_x=[g_x]_p$.
If $\min(m_x,m_y)<p$, the inertia orbit in the sharper bound has
size $e\mid p-1$, $e\le\min(m_x,m_y)$.

The classwise bounds hold on each Galois-stable part of a degree
class, using its cardinality in place of $m_x$. The two moment
identities hold after summing over a partition of all nonunit
labels. For a Galois-stable partition into such parts $A$, write
$x_A,u_A,a_A$ for their
degree, quotient and trace. If the partition refines a principal
$\ell$-block, then also
\begin{equation}\label{eq:simultaneous-block-trace}
 1+\sum_{A\subseteq P_\ell}x_Au_Aa_A
       \equiv0\pmod{\ell^{v_\ell(D/d_Y)}}.
\end{equation}
These congruences hold simultaneously for any collection of primes.
\end{proposition}
\begin{proof}
The quotients defining $A_x,Q_x$ are algebraic integers by
Proposition~\ref{alg:duality}; Galois permutes the degree class.
This also holds on any Galois-stable part. Orthogonality gives
the two sums, and Lemma~\ref{lem:basic-table} gives the strict bound.
The inequality $C(z)\ge|\Re\langle z\rangle|$ used in
Lemma~\ref{nt:residue-norm}, applied after translation by $t$, gives
the interpolation of consecutive integer squares.

Principality gives residue $g_x$ whenever $p\nmid u_x$, at every
place above $p$ by Galois stability. This proves the congruences and
the translated norm bounds. For $p\nmid x$ the entries are nonzero,
so coprime vanishing gives $g_x>1$. The inertia orbit of an entry
has at most $m_x$ members by column conjugation and at most $m_y$
by row conjugation. If their minimum is less than $p$, wild inertia
fixes it and its orbit size divides $p-1$.
For a block-refining partition, \eqref{eq:principal-trace} identifies
the left side of \eqref{eq:simultaneous-block-trace} with
$(D/d_Y)\tau(e_{P_\ell}x_Y)$. Its rational integrality and the
$\ell$-adic integrality of the trace prove the congruence.
Intersecting degree classes with the principal blocks and their
complements gives one partition for all the primes.
\end{proof}

Coprime vanishing and principal support may force $A_x=Q_x=0$.
The \supfile{census/engine_src/degree_rows.py}{class-sum test}
enumerates the remaining integer $A_x$ using the displayed strict
bound and $u_x^2A_x^2\le m_x(f_Y-1)$, and tests their simultaneous
trace equation and norm budget. The
\supfile{census/arithmetic/block_trace.py}{simultaneous-block test} also
retains all possible intersection cardinalities and minimizes the
norm cost at each attainable residue in
\eqref{eq:simultaneous-block-trace}. A resource limit retains the
degree multiset; it never certifies an exclusion.

\begin{proposition}[Integral spheres for principal degree slices]\label{new:degree-slices}
Assume that $R$ is non-pointed and based-simple. Let $P$ be the
principal block at $p$ and let $I$ be a Galois-stable set of $a>0$
nonunit labels of a common degree $y$ and codegree $f$.
Set $u_x=x/\gcd(x,y)$ and $g_x=\gcd(x,y)$.
There are ordinary integers $z_j$, one at each nonunit label $j$,
satisfying, with $x=d_j$ in each summand,
\[
 \sum_{j\ne0}u_x^2 z_j^2=a(f-a),\qquad
 \sum_{j\ne0}x u_x z_j=-a,\qquad |z_j|<a g_x.
\]
They vanish wherever every chosen row vanishes.
If $I\subset P$, then $z_j\equiv a g_x\pmod p$ when $p\nmid u_x$.
This applies, in particular, with $a$ the entire cardinality of the
intersection of $P$ with the degree-$y$ class.
\end{proposition}
\begin{proof}
Put $v_j=\sum_{i\in I}\chi_i(j)$ and $z_j=v_j/u_x$.
Lcm integrality and Galois stability make $z_j$ an ordinary integer.
As in Proposition~\ref{new:block-census}, strictness and the residue
follow pointwise from based simplicity and principality.
Orthogonality gives $\|v\|^2=af$, $\sum_jd_jv_j=0$ and $v_0=a$;
removing the vacuum gives both identities. Vanishing is immediate.
\end{proof}

For each allowed coordinate let $s$ be the spacing of its congruence
class ($p$, or~$1$ in the unrestricted case) and choose a nearest
representative $z_j^0$ to zero. Discard a proposed $a$ if any
$|z_j^0|\ge a g_x$, or if
$L=\sum u_x^2(z_j^0)^2>a(f-a)$.
Set $B=a(f-a)-L$. A coordinate can change only if an adjacent
representative $z_j^0\pm s$ satisfies the strict bound and costs at
most $B$ more in squared norm. All changes of $\sum xu_xz_j$ are
therefore divisible by the gcd of $xu_xs$ over those coordinates.
A zero gcd means that the sum is fixed. Finally the admissible
cardinalities must permit principal-block mass
$1+\sum_y a_y y^2\equiv0\pmod{p^{v_p(D)}}$ and gcd one among the
degrees used. These are the integer tests in
\supfile{census/degree_slices.py}{the degree-slice checker}.

At a fixed denominator prefix, if the degree-$y$ class of denominator
$b=M/y$ is already closed, the coefficient $u_x$ for an observed
denominator $c=M/x$ equals $b/\gcd(b,c)$. Thus its contribution
to $L$ can be charged before knowing the final degrees whenever the
relevant residue of $g_x=M/\lcm(b,c)$ is fixed. If that residue is
undetermined, taking the least cost over its possibilities preserves
the necessary inequality. The nonnegative contributions of the suffix
can be omitted.


\section{Character blocks and primary twists}\label{blk:section}
Assume rotation-admissibility. At $p\mid D$ let $w,u,h_a$ be the
images of $W,U,\operatorname{diag}(a,a^{-1})$ in the $p$-factor
of $\rho$ under Chinese remaindering. Write $s=w'w$ for the
complementary and local Weyl factors. On the coordinate vector $v_i$,
$u$ has eigenvalue $\theta_{i,p}/\gamma_p$, with subscripts denoting
primary components. A label is \emph{active} at $p$ if $\theta_{i,p}\ne1$.
A \emph{phase fiber} fixes any primary components; its mass is the
sum of the squared degrees of its labels. Complementary phase
projectors commute with the $p$-factor. If $\sigma_a$ has multiplier
$a$ at $p$ and one elsewhere, \eqref{alg:covariance} gives
$h_as=\sigma_a(s)$ and $h_av_{\sigma_a i}=\varepsilon(a)v_i$, where
$\varepsilon(a)=\sqrt D/\sigma_a(\sqrt D)$. Since $h_auh_a^{-1}=u^{a^2}$,
Galois raises the normalized local twists to the square-unit powers.

\subsection{Local invariance}
\begin{lemma}\label{blk:local}
Every full $p$-block $B$ spans an invariant subspace for the $p$-factor.
Equivalently, $\sum_{z\in B}s_{xz}\overline{s_{yz}}=0$ if
$\theta_{x,p}\ne\theta_{y,p}$.
\end{lemma}
\begin{proof}
For $A_j=sT^js^{-1}$, substitution of
$e_z=\sum_c\overline{\chi_z(c)}[c]/f_z$ in \eqref{alg:rotation} and
multiplicativity give the first identity below (replace $b$ by $b^*$,
using $\theta_{b^*}=\theta_b$); Lemma~\ref{alg:rotation-arithmetic}
gives the second for integral coefficients:
\begin{equation}\label{eq:rotation-entry}
 L_j^i(e_z)=|(A_j)_{iz}|^2,\qquad
 L_j^i(y)^p\equiv L_{j/p}^i(y^p)\pmod{\mathfrak p}\quad(p\mid j).
\end{equation}
Put $m=N_{p'}$. The reduced character of $B$ defines a residue field
$k=\mathbb F_{p^h}$ of $R/pR$; its embeddings give conjugate blocks
$B=B_0,\ldots,B_{h-1}$. Choose $b\ge v_p(N)$ divisible by $h$
with $p^b\equiv1\bmod m$. For $V$ a nonnegative integral combination
and $n\ge2b$, the certificate at $mp^n$ gives
\begin{equation}\label{eq:block-limit}
 L_m^i(V^{p^n})-L_{mp^b}^i(V^{p^{n-b}})
       \in p^{n-2b+1}\overline{\ZZ}_p .
\end{equation}
Indeed these are the traces of the $p^n$th and $p^{n-b}$th powers.
The powers agree on roots of order dividing $mp^{n-b}$. At exact
$p$-part $p^c>p^{n-b}$, the multipliers $1\bmod m$, $1\bmod p^b$
fix $\QQ(\zeta_N)$ and give orbits of size $p^{c-b}$, on which both
powers are constant. Multiplicities are constant too, by Fourier
inversion from the power traces in $\QQ(\zeta_N)$.

The Teichm\"uller limit and separation of the residue-field embeddings in \cite[Theorem~3.2, proof of equation~(3.5)]{BDPcore} now apply to \eqref{eq:block-limit}. They give $\sum_{z\in B}|(A_m)_{iz}|^2=\mathbf1_{i\in B}$. Hence the block projector commutes with $A_m$ and with $u$. The two conjugate generating unipotents preserve the block, as in \cite[Lemma~3.3]{BDPcore}.
\end{proof}

\begin{lemma}\label{blk:witness}
The nonunit degrees in the principal $p$-block $P$ are nonempty
and have gcd one. In the non-pointed based-simple case they are
pairwise noncoprime. Consequently the nonunit denominators divisible
by $p$ include three distinct values and have lcm $M$.
In this case the set $A_p$ of nonunit degrees $d_j$ satisfying
$p^2\mid f_j$ and $\gcd(d_j,d_i)>1$ for every $p$-free nonunit
degree $d_i$ also has gcd one. If exactly three primes divide $D$,
then $p^2\mid M$ for every $p\mid D$.
\end{lemma}
\begin{proof}
Choose a $p$-active $X$, since $p\mid N$. Block orthogonality gives
$\sum_{0\ne j\in P}d_j\chi_X(j)=-1$. In the simple case,
principality makes the entry between a row in $P$ and a $p$-free
degree nonzero; coprime vanishing then applies. Two $p$-divisible
degrees already share $p$. Thus the nonunit degrees in $P$ are
pairwise noncoprime, and their gcd one forces three distinct values.
Lemmas~\ref{lem:basic-table} and~\ref{blk:height} put these degrees
in $A_p$ and their denominators among the multiples of $p$.
Their lcm is $M$ because their degrees have gcd one.
With three dimension primes, some nonunit degree in $P$ is divisible
by $p$: otherwise all contain both other primes, contradicting gcd
one. Its denominator also contains $p$ by Lemma~\ref{blk:height},
so $p^2\mid M$.
\end{proof}

\begin{theorem}[Three denominators across every prime cut]
\label{new:prime-cut}
Assume the manuscript's non-pointed, based-simple,
rotation-admissible hypotheses and write $D=qM^2$, with $q$
square-free. Partition the prime divisors of $M$ into at least
two nonempty parts. Let $H$ consist of the nonunit labels whose
denominators meet at least two parts.
The degrees occurring on $H$ have gcd one and include at least
three distinct values. Equivalently, the distinct denominators
on $H$ number at least three and have least common multiple $M$.
\end{theorem}
\begin{proof}
For $p\mid M$, Lemma~\ref{blk:height} gives $p^2\mid qb_Y^2$,
hence $p\mid b_Y$, on $P_p$. For every $t\mid M$, choose $p$ in
another part. Lemma~\ref{blk:witness} supplies a nonunit $Y\in P_p$
with $t\nmid d_Y$, so $p,t\mid b_Y$ and $Y\in H$.
Thus the degrees on $H$ have gcd one.

If they had at most two distinct values, they would be two coprime
integers $d,e>1$. Fix a part $J$ and $p\in J$. Every nonunit
degree in $P_p$ outside $H$ is divisible by $M_{J'}>1$.
The gcd-one and pairwise-noncoprime assertions of
Lemma~\ref{blk:witness} therefore put exactly one of $d,e$, say $a$,
in $P_p$, with $(a,M_{J'})=1$; otherwise a common prime would
divide every nonunit degree in that block. Hence $a$ is supported
on $J$. Applying the same argument in another part produces the
other degree $a'$, supported there. Principality makes the entry
between the row of degree $a$ in $P_p$ and a label of degree $a'$
nonzero, since $p\nmid a'$. Coprime vanishing makes it zero,
a contradiction. Finally
$\gcd_{Y\in H}d_Y=M/\lcm_{Y\in H}b_Y$.
\end{proof}

\begin{theorem}[Principal blocks and products of character rows]
\label{rev:principal-product}
Let $p\mid D$ and let $P$ be the principal $p$-block in a nontrivial perfect integral
ring with symmetric Drinfeld duality and cyclotomic characters.
Put $e=v_p(D)$. For any list $X_1,\ldots,X_m$ (repetition is
allowed), write $t_i=v_p(d_{X_i})$. If $\sum_i t_i<e$, there is a
nonunit $Y\in P$ with $S_{X_iY}\ne0$ for every $i$ and
\[
  2v_p(d_Y)\le\sum_i\min\{t_i,v_p(d_Y)\}.
\]
In particular, every row hits a nonunit of $p$-free degree in $P$
which also hits every row of $p$-free degree. The $p$-free nonunits
of $P$ occupy at least two Galois orbits. If $D$ is odd and $3\mid D$,
the principal three-block contains at least eight nonunit labels
of degree prime to three.
\end{theorem}
\begin{proof}
Let $\tau$ extract the unit coefficient and $e_P$ denote the block
idempotent. By \eqref{eq:block-trace} and symmetry,
\[
 \frac{D}{\prod_i d_{X_i}}\tau\left(e_P\prod_i x_{X_i}\right)
 =1+\sum_{0\ne Y\in P}
       d_Y^{2-m}\prod_i\chi_{X_i}(Y).
\]
The trace on the left is $p$-integral. Its prefactor has positive
valuation. Hence a nonzero summand on the right beyond the unit has
valuation at most zero. Put $s=v_p(d_Y)$. Lcm integrality gives
$v_p(\chi_{X_i}(Y))\ge s-\min(s,t_i)$, so that summand has valuation
at least $2s-\sum_i\min(s,t_i)$. This proves the inequality.

Append every row of $p$-free degree to the one specified row $X$.
The budget holds since $d_X^2\mid D$, and $2s\le\min(s,v_p(d_X))$
forces $s=0$. To delete any one Galois orbit of $p$-free principal
nonunits, take for $X$ a row vanishing on that orbit by Burnside
vanishing and Galois covariance. Thus at least two such orbits occur.
In odd dimension each contains dual pairs. In $P_3$ the count of
three-free nonunits is $2\bmod3$, by \eqref{eq:block-trace} at $1$.
It is even and at least four, hence at least eight.
\end{proof}

\begin{theorem}[Detecting equal primary twists from degrees]
\label{rev:phase-overlap}
Assume the non-pointed, based-simple, rotation-admissible hypotheses.
For $p\mid D$ let $C_p$ contain the degrees of all nonunits in $P_p$.
For two distinct labels of degrees $a,b$ let $H$ consist of the
$c\in C_p$ at which both rows are allowed to be nonzero. Write
$u_a(c)=c/\gcd(a,c)$. Then
\[
  \gcd_{c\in H}\bigl(u_a(c)u_b(c)\bigr)\ne1
  \quad\Longrightarrow\quad \theta_{a,p}=\theta_{b,p},
\]
where the twists refer to the chosen labels and $\gcd\varnothing=0$.
An allowed support may include extra columns. The vacuum row allows
every column. For $a,c>1$ one may require $\gcd(a,c)>1$ and
$\lcm(a,c)^2<D$.
\end{theorem}
\begin{proof}
Different primary twists give orthogonality on $P_p$ by
Lemma~\ref{blk:local}. Removing the vacuum expresses $-1$ as a sum
of terms $\chi_a(Y)\overline{\chi_b(Y)}$ divisible as algebraic
integers by $u_a(d_Y)u_b(d_Y)$. Their common rational divisor
divides one. Coprime vanishing gives the first support condition.
For the second, if $S_{aY}\ne0$, Galois averaging the row norm
gives $D\ge a^2+\lcm(a,d_Y)^2$, using the nonzero algebraic integer
$S_{aY}/\lcm(a,d_Y)$ and reserving the vacuum entry. Compare
\cite[Theorem~1.1]{DP}.
\end{proof}

A nonunit $c$ can be placed in $C_p$ only if $p^2\mid D/c^2$.
A principal row of degree $c$ must hit every $p$-free degree.
On a $p$-free class $a$ of multiplicity $n$, set $g=\gcd(a,c)$,
$s=g\bmod p$ and $u=a/g$. Its normalized squared norm on this
class is an integer $Q\equiv ng^2\bmod p$ with
$Q\ge\max(n,ns(p-s)/(p-1))$ by Proposition~\ref{new:block-census}.
The least such integer costs $u^2Q$. Summing these costs and reserving
the vacuum gives the list of possible principal degrees.

Connect two degree classes when Theorem~\ref{rev:phase-overlap} forces
equality. A loop requires at least two labels in the class. Every
connected component with an edge carries just one primary twist.
A component meeting the vacuum or a class with $p$-unit codegree
is inactive. If $p\ge5$, a constant primary phase on any entire
degree class is fixed by all square units and is therefore trivial.
At any odd $p$, the active part of a remaining degree class of size
$n$ occupies at most $h\lfloor n/h\rfloor$ labels, $h=(p-1)/2$.
Each connected component of constant phase supplies at most one
distinct phase. Compare these upper bounds with the packet and mass
lower bounds in Proposition~\ref{blk:packets} and
Lemma~\ref{blk:activity}. The same argument counts individual
labels in odd dimension, with the additional pairing already proved
in Lemma~\ref{alg:odd-parity}.


\begin{lemma}\label{blk:admissible}
If $w$ preserves $\mathbb C v_z$, then $p\nmid f_z$. Row $z$ meets
an active column exactly when $p\mid f_z$. Activity of $z$ requires
$p\mid f_z$, and for odd $p$ its primary twist order divides $f_z$.
\end{lemma}
\begin{proof}
Use the local Weyl criterion of \cite[Lemma~3.4(1) and Theorem~3.6]{BDPcore}.
Its rotation identity (3.3) is \eqref{eq:rotation-entry} here, and the
prime divisors of the twist order are those of $D$ by the modular-data
axioms. Block orthogonality and positivity of the vacuum row therefore
give the stated equivalences. For odd $p$ the same conductor argument
applies by Lemma~\ref{conduct} and square-unit covariance.
\end{proof}

\begin{corollary}\label{blk:primepower}
Each label of $p$-power codegree belongs to the principal $p$-block.
In its presence, all $p$-free degrees belong to that block and have
$p$-power twists.
\end{corollary}
\begin{proof}
If $f_i$ is a $p$-power, $\{i\}$ is a block at every other divisor
of $D$, and factors outside $D$ act by scalars. Thus $w'$ preserves
$\mathbb C v_i$ and $sv_i$ lies in the $p$-block containing $i$.
Its nonzero vacuum entry makes this block principal. For $p\nmid d_j$,
$\chi_i(j)\equiv d_j\ne0$ above $p$; Lemma~\ref{blk:support} puts
$j$ in the same block, and Lemma~\ref{blk:admissible} excludes
activity at every other prime.
\end{proof}

\subsection{Masses of primary twists}
\begin{lemma}\label{blk:moments}
In a full $p$-block or the whole label set, fix any complementary
primary phases, $p$ odd. If $v_p(D)$ is even, all primitive $p$th
roots have equal mass. If it is odd, the mass at order $p$ is $(p-1)$
times the mass at phase one.
\end{lemma}
\begin{proof}
The projector onto the indicated fiber commutes with the local factor by Lemma~\ref{blk:local}. Apply the Bruhat identity to this projector as in \cite[Propositions~3.9 and~19.1]{BDPcore}. For a local unit $a$ the moment $D^{-1}\sum_{i\in F}d_i^2\theta_{i,p}^{a}$ is a fixed scalar times $(\frac{-a}{p})^{v_p(D)}$. Trace to $\QQ(\zeta_p)$ in even valuation, or to $\QQ$ in odd valuation. Higher-order roots have zero trace and the $p$th roots have sole rational relation $1+\sum_{j=1}^{p-1}\zeta_p^j=0$.
\end{proof}

Write $h_p=(p-1)/2$. At fixed degree and complementary phases the
square units act on a phase of order $p^j$ with orbit size
$p^{j-1}h_p$, also within each full block by Lemma~\ref{blk:local}.
We count in \emph{packets} of $h_p$ equal-degree labels. For the
categorical support cycles see \cite[Lemmas~4.6--4.7]{BGNPRW}.

\begin{proposition}\label{blk:packets}
For $A=\{t\mid q:t\text{ odd prime}\}$ put
\[
 P=\prod_{t\in A}(h_t+1),\qquad
 B(q)=\max\bigl(\{2^{\mathbf1_{2\mid q}}P\}\cup
      \{(h_t+2)P/(h_t+1):t\in A\}\bigr).
\]
Each nonempty fiber fixing odd-primary components outside $A$ contains
at least $B(q)$ labels. For any set $J$ of odd primes of even positive
valuation in $D$, at least $\sum_{p\in J}(p-1)$ nontrivial $J$-phase
tuples occur, accounting for at least $B(q)\sum_{p\in J}(p-1)$ labels.
\end{proposition}
\begin{proof}
For a fiber $F$ complementary to an odd $t\mid q$, the normalized slice gcd
satisfies \cite[Lemma~19.2]{BDPcore}
\begin{equation}\label{blk:fibergcd}
 t\mid\sum_{i\in F}(d_i/g)^2,\qquad g=\gcd_{i\in F}d_i .
\end{equation}
Apply the successive fiber and projected-tuple counts of
\cite[\S19.2 and Proposition~19.8]{BDPcore}. Their inputs here are
the slice identities of Lemma~\ref{blk:moments}, the two-label lower
bound at a ramified dyadic prime from Lemma~\ref{blk:admissible},
and the occurrence of each active coordinate from $\operatorname{rad}(N)=
\operatorname{rad}(D)$. These counts hold for each choice of $J$.
\end{proof}

\begin{lemma}\label{blk:activity}
Let $V_p=\sum_{\theta_{i,p}\ne1}d_i^2$. If the primary twist changes
along an edge $Y\to Z\subset XY$ of the fusion graph of a basis
element $X$ of degree $x$, then $2pxV_p\ge(p-1)D$. This applies to
every $p$-active $X$ and every based generator, in particular one of
least nonunit degree in the simple case.
\end{lemma}
\begin{proof}
Apply the weighted rotation-trace argument of \cite[Proposition~3.8]{BDPcore}. Its integer is $L_{N_{p'}}^0(X)$ by Lemma~\ref{alg:rotation-arithmetic}. The positive weights $N_{ij}^X d_i d_j/D$ sum to $x$. A changing edge gives a strict deficit, hence a deficit at least one. Averaging and using the Ramanujan means bounds the deficit above by $2pxV_p/((p-1)D)$. An active label has a changing edge at the unit, and for a generator this follows from connectedness.
\end{proof}

In a perfect ring and at odd $p$ of even valuation, a nonunit degree
$d$ is \emph{eligible} if $p\mid D/d^2$. Let $m_d$ be its multiplicity,
$n_d=\lfloor m_d/h_p\rfloor$. The difference $k_d$ of order-$p$
packet counts on the two square classes satisfies $|k_d|\le n_d$ and
$\sum k_dd^2=0$. If only two packets fit, activity forces one of each
class at the same degree. The packet optimization of \cite[\S19.4]{BDPcore} gives $V_p\le\max h_p\sum_d(n_d-\epsilon_d)d^2$ over the admissible
$k_d$, where $\epsilon_d=1$ when $n_d-|k_d|$ is odd and either less
than $p$ or $p^2\nmid D/d^2$, and $\epsilon_d=0$ otherwise.

\subsection{Simultaneous prime divisors of a twist}\label{sec:spherical}
Call the primes dividing $M$ but not $q$ \emph{allowed}, and set
$C=\prod_{t\mid q,\ t\ {\rm odd\ prime}}(t+1)/2$.
As before $b_i=M/d_i$.

\begin{theorem}\label{thm:spherical}
If a twist is active at every prime in a set $J$ of allowed primes,
at least $C\prod_{p\in J}p$ labels have $\prod_{p\in J}p\mid b_i$;
for odd $p\in J$, at least $C(p-1)\prod_{t\in J\setminus\{p\}}t$
such labels are active at $p$.
For any nonempty set $P$ of allowed primes at least $C\sum_{p\in P}p$
labels have denominator divisible by some $p\in P$. If $i$ is active
throughout $J$ and another allowed prime $\ell$ divides $b_i$, some
twist is active throughout $J\cup\{\ell\}$.
In odd dimension these lower bounds count duality orbits, so a
bound of $L$ gives at least $2L-1$ labels.
\end{theorem}
\begin{proof}
Decompose the cyclic subrepresentation generated by $v_0$ into
irreducibles for the product of local groups. The projection of $v_0$
onto every summand is nonzero. At odd $p$ its upper unipotents fix it,
and its torus character is $\varepsilon$: trivial at allowed primes
and quadratic at odd $t\mid q$. By Frobenius reciprocity each local
constituent occurs in the representation induced from this character
of the upper triangular subgroup. At level $p$ the trivial character
gives the constants and Steinberg, of dimension $p$ with $p-1$
active dimensions in the projective-line permutation model;
the quadratic character gives two constituents of dimension $(p+1)/2$
\cite[Lemma~4.3]{Nevins}; see also \cite[\S2.5]{PrasadSL}.
At conductor $p^a$, $a\ge2$, the dimension is at least
$p^{a-2}(p^2-1)/2$ \cite[Lemmas~4.2--4.3 and Proposition~4.4]{Nevins}.
Some upper-unipotent character has order $p^a$: otherwise its kernel
contains the conjugates of $I+p^{a-1}E_{12}$, which generate the last
congruence kernel since their reductions span
$\mathfrak{sl}_2(\mathbb F_p)$. The square units give at least
$p^{a-1}(p-1)/2\ge p-1$ active eigenvalues. Thus an active allowed
odd factor has dimension at least $p$ and active dimension at least
$p-1$; every ramified odd factor costs $(t+1)/2$. At two an active
constituent also contains the projected vacuum eigenvalue, so has
dimension at least two.

Every phase tuple occurs in some summand, since its spectral projector
has nonzero value on $sv_0$, whose entries are positive. Within this
tensor product each local eigenvalue can also be replaced by the
vacuum eigenvalue, so deleting any primary components of a twist
preserves occurrence. For allowed $p\nmid b_i$, $p\nmid f_i$ and
$\{i\}$ is a block. Its coordinate line is a local subrepresentation
with vacuum unipotent eigenvalue. An active local constituent has
dimension greater than one, hence is orthogonal to every such line
by equivariance of orthogonal projection. A summand witnessing $J$
is supported on the required labels. Tensoring its full local spaces,
or the active space at $p$, gives the two product bounds.
For the sum bound let $\mathcal F$ be the maximal active supports
in $P$ and choose a witnessing summand for each. Its active primes
in $P$ are exactly $F$, by the tensor-product property and maximality,
so distinct $F$ give orthogonal summands. Their dimensions add to at
least $C\sum_F\prod_{p\in F}p\ge C\sum_{p\in P}p$, since
$\mathcal F$ covers $P$ and a product of distinct primes dominates
their sum.

Finally the fiber fixing the indicated $J$-phases is locally invariant
at $\ell$. If no label there were $\ell$-active, its two generating
unipotents would both act by the vacuum scalar, so its Weyl element
would preserve $\mathbb C v_i$, contrary to $\ell\mid f_i$ and
Lemma~\ref{blk:admissible}.
Charge conjugation is central and fixes $v_0$, so the cyclic module
lies in its $+1$-eigenspace. In odd dimension this space has one
coordinate per duality orbit, with only the vacuum fixed
(Lemma~\ref{alg:odd-parity}), proving the last assertion.
\end{proof}

In square dimension with no twist active at two odd primes, an odd
$p$ either has at least $p-1$ active labels of $p$-power denominator,
or has at least $2(p-1)$ active labels with $2p\mid b_i$.
The first alternative holds when $p$ is isolated, by the support
assertion; otherwise apply the active bound to $J=\{2,p\}$.

\begin{corollary}[Ranks in simultaneous twist sectors]\label{new:sector-rank}
Fix a nonempty set $P$ of allowed primes in Theorem~\ref{thm:spherical}.
Let $\mathcal F$ be the inclusion-maximal sets of primes of $P$ that
occur together in a twist. For each $F\in\mathcal F$ and $K\subseteq F$
one can simultaneously reserve
\[
             C\prod_{p\in K}(p-1)
\]
labels with $\prod_{p\in F}p\mid b_i$ whose active primes in $P$ are
exactly $K$. Hence, writing $n_K$ for the size of this exact sector,
\[
 n_K\ \ge\ C\,\#\{F\in\mathcal F:K\subseteq F\}
                   \prod_{p\in K}(p-1),
 \qquad
 r\ \ge\ C\sum_{F\in\mathcal F}\prod_{p\in F}p .
\]
In odd dimension the reservations count duality orbits, giving
$r\ge 2C\sum_F\prod_{p\in F}p-1$.
\end{corollary}
\begin{proof}
Use the orthogonal witnesses for maximal supports constructed in
Theorem~\ref{thm:spherical}. Their local neutral dimensions are at
least one and their active dimensions at least $p-1$. Tensoring
these spaces with the ramified factors gives the bound for each
$K$, with the same denominator support. The resulting spaces for
different $(F,K)$ are orthogonal, so every union of coordinate
supports has size at least their total dimension. Hall's theorem
gives simultaneous reservations, also on duality orbits in odd
dimension. Summing over $K$ proves the rank bounds.
\end{proof}


\begin{lemma}[Keeping the vacuum out of the count]\label{lem:vac-reserve}
Consider pairwise orthogonal spherical summands as in
Theorem~\ref{thm:spherical}, each supported on a prescribed set of
eligible labels. If a label is ineligible for every summand, one may
reserve the indicated numbers of labels outside the vacuum.
More generally, if $I$ is ineligible for every summand, one may discard
any set $J$ of coordinates such that $s_{I,J}$ has column rank $|J|$.
\end{lemma}
\begin{proof}
Write $U$ for the direct sum. Any vector in $U$ supported on $J$ has
$s$-image zero on $I$, hence vanishes by the rank assumption. Projection
away from $J$ is therefore injective on $U$ and on every subsum.
Hall's theorem applies to the remaining coordinate supports.
For $J=\{0\}$ use $s_{i0}=d_i/\sqrt D>0$.
\end{proof}

\begin{lemma}[Deleting several prime-free columns]\label{new:many-reserve}
Let $I$ be the full class of $m$ rows of codegree $f=p^a$ and $J$ a
set of $p$-free labels. If $|J|(f-m)<f$, then $s_{I,J}$ has full
column rank. The same holds with equality if $J$ contains the vacuum
and a label of degree greater than one. Thus the last assertion of
Lemma~\ref{lem:vac-reserve}
can discard all of $J$ from spherical summands that avoid $I$.
\end{lemma}
\begin{proof}
The restricted Gram matrix is a positive contraction by
\eqref{eq:projection-master}. Each of its diagonal entries is at least $m/f$: the
squared character moduli in that column have Galois-invariant sum,
and each character value is a nonzero algebraic integer by
Corollary~\ref{blk:primepower}. The complementary Gram matrix has
trace at most $|J|(1-m/f)$. A zero eigenvalue of the first Gram
matrix would require an eigenvalue one in the complementary matrix.
At equality this would exhaust its trace, so the complementary
matrix would have rank one and every diagonal entry would equal
$1-m/f$. Thus the columns restricted to $I^c$ would be proportional,
all of the same norm. If $J$ contains a $p$-free label of degree
greater than one, then $0\notin I$: otherwise $f=D=p^a$ and every
$p$-free degree would be one. The vacuum coordinates of these
complementary columns are therefore $d_Y/\sqrt D$. Their unequal
moduli at $Y=0$ and at the label of degree greater than one contradict
proportionality with equal norms.
\end{proof}

For any $p$-free column $Y$, put
$F=\QQ(\chi_X(Y):X\in I)\subseteq K_Y$.
Proposition~\ref{r5:field-rank} gives the uniform bound
\[
 [K_Y:\QQ]\left(\frac mf+
       \frac{\mathbf1_{0\notin I}}{f_Y}\right)
       \le [F:\QQ].
\]
Indeed one appends the degree row exactly when it is absent from
$I$. Since $I$ is a full codegree class, $0\notin I$ is equivalent
to $f<D$. Under this condition, for $Y\ne0$ the bound gives
$[K_Y:F]<f/m$. Thus $2m\ge f$ forces
$K_Y=F\subseteq\QQ(\zeta_f)$, using Lemma~\ref{conduct}.
At odd $p$ a nonreal such field contains the unique subfield of
$\QQ(\zeta_p)$ of degree $2^{v_2(p-1)}$, since its stabilizer
in the cyclic Galois group has odd order \cite[Ch.~6]{Washington}.
In a non-pointed based-simple ring the condition $f<D$ holds for
every prime-power codegree, because $D$ has at least three prime
divisors.

In the non-pointed based-simple square case every denominator is at
least four (Section~\ref{sec:small}). For $K=\prod_{p\in J}p$,
\eqref{eq:recip} gives $1\le1/K+(r-K)/16$. Twist orders divisible by
$21$ or $35$ therefore require rank at least $37$ or $51$, respectively.

\paragraph{An integral norm for each phase.}
Let $J$ be any set of odd allowed primes, and let $F$ be the product
of primes dividing $M$ outside $J$ and not dividing $q$. There are
integers $a_i$ supported on twists trivial outside $J$ (and on
$P_p$ when $J=\{p\}$), with
\begin{equation}\label{eq:phase-column}
 \gcd(b_i,F)\mid a_i,\qquad
 \sum_{i\in E} d_i^2=q\sum_{i\in E} d_i^2a_i^2
\end{equation}
for each fiber $E$ of the $J$-phases. Indeed take
$z=w_Jv_0$, $w_J=\prod_{p\in J}W_p$. Covariance sends $w_J$ to
$h_{a,J}w_J=w_Jh_{a,J}^{-1}$; the local vacuum torus characters
are trivial at the allowed odd primes, so $z$ is rational.
Take $k>0$ equal to one on the $J$-part of $L$ and zero on
its complement. The Bruhat identity
$U(WUW^{-1})U=W^{-1}$ on the chosen factors gives
$|(sT^ks^{-1})_{0i}|=|z_i|$. Thus
$L_k^0(f_i e_i)=f_i|z_i|^2=q(b_i z_i)^2$ is an algebraic integer.
Since $q$ is square-free, $a_i=b_i z_i$ is integral. For
$\ell\mid\gcd(b_i,F)$ use \eqref{eq:rotation-entry} at $k$:
$(f_ie_i)^\ell=f_i^{\ell-1}(f_ie_i)$ gives
$\ell\mid q a_i^2$, hence $\ell\mid a_i$. Complementary unipotents
preserve the vacuum phase of $z$, and a single local Weyl preserves its
principal block. The complementary Weyl commutes with the projector
of $E$ and sends $z$ to $sv_0$, proving \eqref{eq:phase-column}.

In particular, if $g=\gcd_{i\in E}d_i$, $\ell\mid F$ and
$e=\max_{i\in E}v_\ell(b_i)>0$, then
\[
 \sum_{i\in E:\ v_\ell(b_i)=e}(d_i/g)^2\equiv0\pmod{\ell^2}.
\]
Indeed the other normalized degrees are multiples of $\ell$ and the
displayed indices have $\ell\mid a_i$. For odd $q$, if any denominator
in the fiber is even, the number of its denominators of largest dyadic
valuation is therefore a positive multiple of four.

\begin{proposition}[Moments and counts in a phase fiber]\label{phase:moments}\label{odd:edge-budget}
Write $D=qM^2$ with $q$ square-free, and let $p\mid D$ be an odd
prime not dividing $q$. Fix a fiber $F$ of complementary primary phases and
give each label weight $w_i=b_i^{-2}$. Let $J$ be its $p$-active
weight and $C$ its weight at any one primitive $p$th root.
On its $p$-inactive labels with $p\mid b_i$ there are integers
$a_i\in p\mathbb Z$ satisfying
\[
       \sum_i(1-a_i)w_i=C,\qquad
       \sum_i(a_i^2-1)w_i=J.
\]
In odd dimension, if the fiber excludes the vacuum, the same
identities hold on dual pairs with weight $b_i^{-2}$. If $k$ pairs
are active and $c$ have a chosen primitive phase, then
$k\equiv c\pmod2$ and $k>0$ implies $k\ge2(p-1)$.
At $p=3$, an active fiber contains at least eight pairs with
$3\mid b_i$.
\end{proposition}
\begin{proof}
Put $v=sv_0$, $z=W_p^{-1}v$ and $a_i=\sqrt q\,b_i z_i$.
The allowed local torus fixes $v$, since it fixes the vacuum.
Covariance gives $\sigma z=(\sigma\sqrt q/\sqrt q)z$, so $a_i$
is rational. The rotation entry $f_i|z_i|^2=a_i^2$ is integral;
hence $a_i\in\mathbb Z$. Its Frobenius congruence at the
complementary exponent, as in \eqref{eq:rotation-entry}, gives
$p\mid a_i$ when $p\mid b_i$. The column $z$ is supported on
$p$-inactive labels. If $p\nmid b_i$, the singleton local factor
has Weyl scalar one, by the Bruhat relation and its vacuum
unipotent character. Thus $a_i=1$ on those labels.

Let $E$ project onto the fiber. It commutes with $W_p$, so
$\|Ez\|^2=\|Ev\|^2$ gives the second identity after canceling
$1/q$. For the first, $Q=s^{-1}Es=sEs^{-1}$ commutes with $W_p$,
and $\langle z,Ev\rangle=(QW_p)_{00}$. The Bruhat identity in
Lemma~\ref{blk:moments} identifies this with the local Galois mean
of $q^{-1}\sum_{i\in F}\theta_{i,p}/b_i^2$. The root means at
orders one, $p$, and higher $p$-powers are $1,-1/(p-1),0$.
The primitive masses agree, giving the first identity. The terms
with $p\nmid b_i$ cancel in both.

In odd dimension only the vacuum is self-dual
(Lemma~\ref{alg:odd-parity}), and $z$ is charge-even; divide by two
to count dual pairs. Since
$w_i\equiv1\pmod8$ and $a_i^2\equiv a_i\pmod2$, the two
identities give $k\equiv c\pmod2$. If only primitive phases occur,
their counts agree modulo eight; their total $k$ is even, so each
count is positive and even. A higher phase supplies at least
$p(p-1)/2$ pairs. This suffices for $p\ge5$; at three the only
smaller higher orbit would have $k=3,c=0$, contradicting parity.
Finally, at three the second identity expresses $k$ modulo eight
as a sum of $0,3,7$. For $k=4,5,6,7$ this requires respectively
at least $4,3,2,1$ inactive eligible pairs.
\end{proof}

\begin{corollary}[Squares in an active phase fiber]\label{sq:fiber-squares}
Let $p\ge5$ be an allowed prime and let $\ell(n)$ be the least
number of rational squares summing to $n$. Whenever $p$ is active
in a fiber of complementary phases, at least $p-1+\ell(p+1)$
labels in that fiber have $p\mid b_i$. If some twist is also
active at a second odd allowed prime $t$, at least
\[
 [p-1+\ell(p+1)](t-1)+(p-1)
\]
labels have $pt\mid b_i$ and are active at $p$ or $t$.
\end{corollary}
\begin{proof}
Count individual labels. Adding twice the first phase moment to the
second gives
\[
       \sum_{\theta_{i,p}=1,\,p\mid b_i}
                \left(\frac{1-a_i}{b_i}\right)^2=J+2C .
\]
A phase of order at least $p^2$ contributes $p(p-1)/2$ active
labels. Otherwise the $p-1$ primitive phases all occur with equal
mass $C$. If each contains at least two labels, there are already
$2(p-1)$ active labels. Both counts suffice since $\ell(p+1)\le4$
by Lagrange's four-square theorem \cite[Theorem~369]{HardyWright}. In the remaining case a singleton
phase gives $C=b^{-2}$. The displayed sum represents
$(p+1)/b^2$, so the inactive terms alone require $\ell(p+1)$ labels.

At least $t-1$ occupied nontrivial $t$-phases meet $p$-activity,
by equality of primitive masses or the square-unit orbit of a
higher phase. Apply the first bound in each. Also at least $p-1$
nontrivial $p$-phases meet $t$-activity. The second $t$-moment
in each forces an additional $t$-inactive label with $t\mid b_i$.
These labels are disjoint from the preceding $t$-active ones.
\end{proof}

\begin{proposition}[Mixed primary twists]\label{round:mixed}
Let $p$ be an odd prime with $p\nmid q$. Suppose a twist is active
both at $p$ and at a different prime $t\mid M$, with $t\nmid q$.
There are at least $h(t)(p-1)$ labels active at $p$ whose denominators
are multiples of $pt$, where
\[
 h(2)=4,\qquad h(t)=3\ (t\equiv3\bmod4),\qquad
 h(t)=2\ (t\equiv1\bmod4).
\]
In the non-pointed based-simple square case, simultaneous activity at two and $p$ implies
$1\le (p-1)/p^2+(r-4(p-1))/16$.
\end{proposition}
\begin{proof}
Fix the nontrivial $t$-phase of the given label.
If a $p$-phase of order at least $p^2$ occurs in this fiber,
its square-unit orbit contains at least $p(p-1)/2$ phases.
Otherwise Lemma~\ref{blk:moments} gives all $p-1$ primitive roots
there. In either case at least $p-1$ distinct nontrivial $p$-phase
fibers contain a label with $t\mid b_i$, by
Lemma~\ref{blk:admissible}. Apply \eqref{eq:phase-column} with
$J=\{p\}$ to each such fiber. The labels whose denominator has
maximal $t$-valuation give a nonempty sum of unit squares zero
modulo $t^2$ after dividing their degrees by the fiber gcd.
At two this requires a multiple of four terms. At odd $t$ it
requires at least two, and at least three if $-1$ is a nonsquare.
Every label counted also has $p\mid b_i$. Summing their reciprocal
squares and bounding the rest by $1/16$ gives the last claim.
\end{proof}

Now specialize to square dimension and $J=\{p\}$.

In a fiber of size at most two, all $|a_i|=1$: a vanishing coefficient
in a two-label fiber would make $a^2-1$ a rational square for an
integer $|a|\ge2$, which is impossible. Thus every such label has
$p$-power codegree. Without a $p$-power codegree every fiber has at
least three labels; the vacuum and at least $p-1$ active phases
occur by Lemma~\ref{blk:moments} and the orbit sizes, so $r\ge3p$.

Let $W$ be the mass at each primitive $p$th phase and $H$ the mass
at orders $\ge p^2$. Higher square-unit orbits have zero root sum
by summing over lifts one level down, so $G=\sum_i d_i^2\theta_{i,p}
=D-pW-H$. The vacuum Bruhat identity gives $|G|=D|z_0|$, hence
$MF\mid G$, equivalently $MF/p\mid W+H/p$. For the underlying
categorical Gauss identity see \cite[Proposition~3.3]{NSWGauss}.

\begin{lemma}\label{blk:two}
If a local factor preserves the coordinate plane of $X,Y$ but $w$
moves $\mathbb C v_X$, then $p=2$, $d_X=d_Y$ and the primary twist
ratio has order four.
\end{lemma}
\begin{proof}
Apply the two-dimensional local Weyl calculation of
\cite[Lemma~19.5]{BDPcore}. It uses the relations $(wu)^3=w^2$,
unitarity, and the positive rational ratios of vacuum entries,
all part of Section~\ref{alg:modular}.
\end{proof}

When $2\mid q$ and $p\mid D$ is odd of even valuation,
$\sum_{d\ {\rm eligible}}(\lfloor m_d/h_p\rfloor+
\lfloor m_d/(2h_p)\rfloor)\ge6$. This is \cite[Lemma~19.6]{BDPcore}, applied to the phase fibers
of Lemma~\ref{blk:moments} and the two-label alternative of
Lemma~\ref{blk:two}.

\begin{proposition}[Isolation at two primes of valuation two]
\label{new:two-square-primes}
In the non-pointed based-simple rotation-admissible setting, suppose that exactly four distinct primes $p,q,u,v$ divide $D$, with
$v_u(D)=v_v(D)=2$. Then $u$ and $v$ are isolated in the twist graph
and $D$ is a square.
\end{proposition}
\begin{proof}
This is the argument of \cite[Proposition~9.3]{BDPcore}, expressed
in the present axioms. A $u$-active row vanishes on every $u$-divisible
degree by Lemma~\ref{blk:admissible}. Reducing its orthogonality to the
vacuum modulo $q$ gives a nonunit degree free of $u$ and $q$.
It has at least two prime factors, so it has the form $p^iv$ with
$i>0$. Interchanging $u,v$ also gives a degree $p^ju$.
Every $p$-free nonunit in the principal $q$-block meets both degrees
by principality. Coprime vanishing forces it to be divisible by $uv$.
For a $q$-active row, orthogonality to the vacuum on this block gives
a $p$-free nonunit hit, by reduction modulo $p$. The hit is divisible
by $uv$, so that row is inactive at both $u$ and $v$. Interchanging
$p,q$ separates the two pairs of primes. A row active at both $u,v$
would have only $p$- and $q$-divisible nonunit hits, again contradicting
vacuum orthogonality modulo $p$.
At a $u$-active label all complementary unipotents are scalar on its
nontrivial $u$-phase space; their Weyl conjugates have the same
property there. Thus each complementary local Weyl preserves the label
line, and Lemma~\ref{blk:admissible} makes its codegree a power of $u$.
Activity and $v_u(D)=2$ force it to be $u^2$, proving that $D$ is a
square.
\end{proof}

\subsection{Blocks of defect one}
\begin{lemma}\label{blk:defectone}
A full $p$-block containing a label of defect one ($v_p(f_i)=1$,
equivalently $p\mid q$ and $p\nmid b_i$) consists of two tame inertia
orbits, both of defect one and with the same complementary phases.
At two it is an equal-degree pair with dyadic twist ratio of order
four. At odd $p$, for the inactive and active orbit sizes $a,b$
and degrees $u,v$ respectively,
\[
 a,b\mid p-1,\qquad (p-1)/2\mid b,\qquad bv^2=(p-1)au^2 .
\]
At three the degrees are $(u,u,u)$ or $(u,u,2u)$ up to order and
rescaling, with denominators $(z,z,z)$ or $(z,2z,2z)$.
At five $(a,b,v/u)$ is $(1,4,1),(2,2,2)$ or $(4,4,2)$.
\end{lemma}
\begin{proof}
Pass to an unramified splitting extension $F/\QQ_p$ for the residue
characters. The full block is a local symmetric order over $\mathcal O_F$
by idempotent lifting. Its symmetrizing trace is
$\sum_j\Tr_{E_j/F}/f_j$ on the inertia coordinates. Thus
\cite[Proposition~2.13]{BDPcore} applies. At odd $p$ it gives two
tame defect-one orbits. By Lemmas~\ref{conduct} and~\ref{blk:admissible}
their sizes divide $p-1$ and their phases have order $1$ or $p$.
The mass identity of Lemma~\ref{blk:moments}, also within each fixed
complementary fiber, gives the common complementary phases, one neutral
and one active orbit, and $bv^2=(p-1)au^2$, as in
\cite[Corollary~3.10]{BDPcore}. Square-unit covariance gives
$(p-1)/2\mid b$. Substitution gives the two lists.

At two, the defect-one coordinate is unramified by
Lemma~\ref{conduct}, excluding the single ramified quadratic alternative
of the cited proposition. In the two-branch alternative, write
$e=2^k$ for the degree of the other branch. Its codegree has valuation
$k+1$ and its conductor exponent as an order branch is one
\cite[Lemma~2.12 and Proposition~2.13]{BDPcore}. If $k>0$,
Lemma~\ref{conduct} embeds that branch in $F(\zeta_{2^{k+1}})$, hence
identifies it with this extension by equality of degrees. Its different
exponent is $ke$, by the derivative of $X^e+1$. The order conductor
exponent would therefore be $e(k+1)-ke=e>1$, a contradiction. Both
orbits are singletons. Neither line is locally invariant by
Lemma~\ref{blk:admissible}. Splitting the block by complementary phases
must therefore keep both together, and Lemma~\ref{blk:two} finishes.
\end{proof}

\section{Small formal codegrees}\label{sec:small}
Fix a nonunit $X$ of degree $n$ and codegree $f$. Put $u_Y=\chi_X(Y)$,
$c_Y=\operatorname{mean}|u_Y|^2$ over Galois conjugates, and
$H=\{Y\ne0:u_Y\ne0\}$. The vacuum modular relation, row norm and lcm integrality give the
three inequalities of \cite[\S20, (20.1) and the next sentence]{BDPcore}:
\begin{equation}\label{st}
 \sqrt f\le\frac1n\sum_Yd_Y\sqrt{c_Y},\qquad
 \sum_Y c_Y=f,\qquad c_Y\ge(d_Y/\gcd(d_Y,n))^2\quad(Y\in H).
\end{equation}
Also $\sum_{H}d_Yu_Y=-1$, so the nonunit support has gcd one.
Under perfection it therefore contains at least two distinct degrees
greater than one. In the simple case each is at least six.

\subsection{Square-free and small codegrees}
\begin{theorem}[Square-free formal codegrees]\label{thm:squarefree-codegree}
For an odd prime $\ell$ let $c_\ell$ be the largest positive integer whose
square divides $2(\ell-1)$, and put $c_2=1$.
Suppose that $f_X$ is square-free. For every $Y$ with $s_{YX}\ne0$, both
$d_X/\gcd(d_X,d_Y)$ and $d_Y/\gcd(d_X,d_Y)$ divide
$\prod_{\ell\mid f_X}c_\ell$. In particular
\[
                         d_X\mid\prod_{\ell\mid f_X}c_\ell .
\]
Consequently square-free formal codegrees do not occur when $D>1$ is
square-full, in particular in a nontrivial perfect ring.
\end{theorem}
\begin{proof}
Write $s$ as the product of the commuting local Weyl operators.
Those at primes outside $f_X$ preserve the coordinate line of $X$
by Lemma~\ref{blk:admissible} and the scalar action outside $D$.
Apply the remaining operators $W_\ell$ in increasing order of $\ell$.
Whenever an input coordinate has defect one at $\ell$, local invariance
restricts its image to its full $\ell$-block (Lemma~\ref{blk:local}).
At two the degrees in that block coincide. At an odd $\ell$ they are
$u,v$ with $(v/u)^2=(\ell-1)a/b$, where $a\mid\ell-1$ and
$b$ equals $(\ell-1)/2$ or $\ell-1$ (Lemma~\ref{blk:defectone}).
Thus $v/u$ is a positive integer dividing $c_\ell$. No prime
$\ge\ell$ divides this ratio. It follows by induction that all current
coordinates still have defect one at every unprocessed prime.
Along any path supporting a nonzero matrix entry the degree changes
at each step by $k$ or $1/k$ with $k\mid c_\ell$. Taking the product along this path gives both divisibilities. The vacuum entry $s_{0X}$ is nonzero, giving the
assertion for $d_X$.

If $f_X>1$, its largest prime divisor $p$ divides none of the factors
$c_\ell$. Hence $p\nmid d_X$ and $v_p(D)=v_p(f_Xd_X^2)=1$.
If $f_X=1$, the divisibility gives $d_X=1$ and hence $D=1$.
Use Lemma~\ref{lem:basic-table} for perfection.
\end{proof}

For a spherical fusion category, Schopieray classified the
nonintegral formal codegrees of square-free field norm as the two
Fibonacci values $(5\pm\sqrt5)/2$ \cite[Theorem~4.13]{SchopFormal}.
Theorem~\ref{thm:squarefree-codegree} treats integral codegrees
using the modular representation.

\begin{lemma}\label{small}
In every nontrivial perfect rotation-admissible ring, each nonunit formal
codegree is composite and exceeds four.
\end{lemma}
\begin{proof}
The square-free case follows from Theorem~\ref{thm:squarefree-codegree}.

The vacuum and two nonunit hits give $f\ge3$. Suppose $f=4$.
By \eqref{st} the nonunit norm budget of three forces $d_Y\mid n$
on $H$. Two hits have coprime nonunit degrees, both proper divisors
of $n$. Cauchy--Schwarz in \eqref{st} gives
$2\le 1/n+\sqrt6/2<2$. For three hits all $c_Y=1$.
If all three degrees are proper, then $2\le1/n+3/2<2$, since $n$ has
a proper nonunit divisor. Otherwise precisely one degree equals $n$,
and the other two, say $a,b$, are coprime nonunit divisors. Now
$ab\mid n$ and $2\le1+(1+a+b)/n$. For coprime $2\le a<b$ this forces
$a=2$, $b=3$, $n=6$. The support then has degrees $1,2,3,6$.
By Corollary~\ref{blk:primepower} the row is principal at two and
therefore meets every odd degree. This gives exactly two odd labels,
whereas $4\mid D=\sum_Yd_Y^2$ requires a multiple of four.
\end{proof}

\subsection{Prime-square codegrees and the prime three}
\begin{lemma}[Prime-free trace support]\label{prime-square}
In a rotation-admissible ring suppose a nonunit label $X$ has
$f_X=p^2$ and $d_X=n$. The nonunit degrees at entries of nonzero
absolute field trace are prime to $p$, divide $n/p^{v_p(n)}$,
and have gcd one. This holds for every prime $p$ and every height
$v_p(n)$.
\end{lemma}
\begin{proof}
For $u_Y=\chi_X(Y)$ put $h=[K_X:\QQ]$, $e=[\QQ(u_Y):\QQ]$ and $t=\Tr_{\QQ(u_Y)/\QQ}u_Y$. The orbit trace identities \cite[Lemma~2.4, equations~(2.4)--(2.5)]{BDPcore} give
\[
 p^2=\sum_O(h/e)(|O|/e)t^2,\qquad 0=1+\sum_{O\ne\{0\}}d_Y(|O|/e)t.
\]
The vacuum contributes $h>0$ to the first identity. Thus $|t|<p$
off the vacuum, and the second identity gives gcd one on its
nonzero-trace support.

Corollary~\ref{blk:primepower} puts $X$ in $P_p$. If $p\mid d_Y$,
principality therefore gives $u_Y\equiv d_Y\equiv0$ at the place
above $p$. Its field lies in $\QQ(\zeta_{p^2})$ by Lemma~\ref{conduct},
so that place is unique. Every conjugate reduces to zero, and its
absolute trace is an integer divisible by $p$. Hence $t=0$.
Finally $d_Y\mid pn$, so the remaining degrees divide the prime-to-$p$
part of $n$.
\end{proof}

\begin{theorem}[Tame projection of a prime-power row]\label{thm:tame-orbit}
In a non-pointed based-simple rotation-admissible ring suppose
$f_X=p^a$ for an odd prime $p$, and write
$[K_X:\QQ]=p^s m$ with $p\nmid m$. Let $b$ be the number of
column orbits having nonzero mean in the Galois average of row $X$,
including the vacuum orbit. Then
\[
 4\le b\le p^{a-s-1}\left(1+\frac{p-1}{m}\right).
\]
In particular $K_X$ cannot contain $\QQ(\zeta_{p^a})^+$,
and the order of $\theta_X$ divides $p^{a-1}$.
\end{theorem}
\begin{proof}
Write $F=\QQ(\zeta_{p^a})$ and $F_0=\QQ(\zeta_p)$.
Lemma~\ref{conduct} gives $K_X\subseteq F$.
Since $\operatorname{Gal}(F/\QQ)$ is cyclic, the subgroup
$W=\operatorname{Gal}(F/F_0)$ induces an orbit of $p^s$ rows
among the $p^s m$ conjugates of $X$. Denote the averages of row
$X$ over $W$ and over the full Galois group by $v$ and $w$.
The conjugate rows are orthogonal of squared norm $p^a$, so
\[
 \|v\|^2=p^{a-s},\qquad \|w\|^2=p^{a-s}/m.
\]
Each coordinate $v_Y$ is in $\ZZ[\zeta_p]$: the average over $W$ of
a power of $\zeta_{p^a}$ fixes it if its order divides $p$ and is zero
otherwise. In particular $w_Y$ is the absolute Galois mean of $v_Y$.

Fix a column orbit $B$ of length $k$ and put $e=[\QQ(v_Y):\QQ]$
for $Y\in B$. Each conjugate of $v_Y$ occurs $k/e$ times on $B$.
The powers $\zeta_p^u$, $1\le u<p$, form an integral basis on
which Galois acts regularly. Hence the coefficients of $v_Y$ are
constant on $e$ cosets, each of size $(p-1)/e$. Denote their
integer values by $a_1,\ldots,a_e$. The absolute trace $t$ of $v_Y$
from its own field is $-\sum_j a_j$, so $w_Y=t/e$.
Using $\operatorname{Tr}_{F_0/\QQ}(\zeta_p^i)=-1$ for
$p\nmid i$ and $p-1$ otherwise gives
\[
 \operatorname{Tr}_{\QQ(v_Y)/\QQ}(v_Y\overline{v_Y})
 =p\sum_{j=1}^e a_j^2-\frac{p-1}{e}\,t^2 .
\]
Multiplying by $k/e$ and summing over all column orbits yields
\[
 \sum_B \frac{k}{e}\sum_{j=1}^e a_j^2
 =\frac{\|v\|^2+(p-1)\|w\|^2}{p}
 =p^{a-s-1}\left(1+\frac{p-1}{m}\right).
\]
Each orbit with $w_Y\ne0$ contributes a positive integer.

Orthogonality to the vacuum gives
$0=\sum_B d_Y(k/e)t$. The vacuum term is one, so the degrees
of the nonunit orbits with $w_Y\ne0$ have gcd one.
Every such label lies in the principal $p$-block by
Lemma~\ref{blk:support}, applied to a conjugate row meeting it.
Its nonunit degrees are pairwise noncoprime by
Lemma~\ref{blk:witness}; their gcd one therefore forces three
distinct degrees, proving $b\ge4$.

If $K_X$ contained the maximal real subfield of $F$, its degree
would be at least $p^{a-1}(p-1)/2$ and divisible by that number.
Thus $s=a-1$ and $m\ge(p-1)/2$, making the upper bound at most
three. The $p$-primary twist condition now follows by the
square-unit covariance used in Lemma~\ref{blk:admissible};
all other primary parts of $\theta_X$ are trivial by that lemma.
\end{proof}

\begin{lemma}\label{nine}
Under the hypotheses of Lemma~\ref{small}, a label of codegree nine and
degree $n$ forces $3\nmid n$. The principal $3$-block then consists
precisely of the labels of degree prime to three, and every twist of
order divisible by three has $3$-power order. In the non-pointed
simple case, nine is not a formal codegree.
\end{lemma}
\begin{proof}
Write $D=9n^2$ and suppose $I=\{X:d_X=n\}$ is nonempty. Every
$I$-row is supported in $P=P_3$, whose degrees divide $n$ by
Lemmas~\ref{blk:support} and~\ref{blk:height}.

For $3\mid n$, the wild-entry norm and dual-pair cancellation in
\cite[proof of Theorem~20.8, first paragraph]{BDPcore} give
$3\le 1/n+8/3$. Its inputs hold here by Lemmas~\ref{conduct},
\ref{prime-square}, principality and the vacuum modular relation.
This leaves $n=3$, hence $D=81$. Then every nonunit degree is
divisible by three, contrary to $D=1+\sum_{Y\ne0}d_Y^2$.

Now $3\nmid n$. By Corollary~\ref{blk:primepower}, each $X\in I$
belongs to $P$ and meets every $3$-free degree. Its support $P$ consists exactly of the $3$-free labels, as $d_Y\mid n$
there. By Lemma~\ref{blk:admissible}, the row sees only $3$-power
twists, and every $3$-active label has $3\mid f_Y$, hence $3\nmid d_Y$.
This proves the assertions under perfection. In the simple case the
last assertion contradicts Lemma~\ref{isolated}.
\end{proof}

\subsection{Weyl columns on isolated components}
Join two primes dividing $D$ if both divide the order of a twist.
Let $J$ be a nonempty union of components consisting of odd primes,
and write $W_J$ for the product of their local Weyl operators.
On any fiber of nontrivial $J$-phase, each complementary unipotent is
scalar with its vacuum eigenvalue. Its Weyl operator preserves the
fiber, so the conjugate unipotent has the same scalar there. The local
Bruhat relation then makes the complementary Weyl scalar as well.
Thus the codegree of each label in this fiber is supported on $J$
by Lemma~\ref{blk:admissible}. In the same way the codegrees of
labels active outside $J$ are prime to every prime in $J$.

\begin{proposition}[The local Weyl column]\label{new:component-column}
Suppose $D=M^2$. Write $M_J$ for the $J$-part of $M$,
$w=W_Jv_0$, and $u_i$ for the $J$-part of $b_i$.
Then $\epsilon=\gamma^3$ is a sign and, for every $i\ne0$,
\[
 \epsilon w_i=A_i/u_i,\qquad |A_i|<u_i,\qquad
 A_i(b_i/u_i)\equiv1\pmod{u_i}.
\]
In particular $w_i\ne0$ exactly when $u_i>1$. For $J=\{p\}$,
the principal $p$-block consists precisely of the labels with
$p\mid b_i$. If $b_i>1$ is supported on $J$, then $|w_i|=1/b_i$.

The positive number
\[
 g_J=1-w_0=
 \sum_{i:\,\theta_{i,J}\ne1}
 \frac{1}{b_i^2}\left(1-
 \frac{\mu(\operatorname{ord}(\theta_{i,J}))}
      {\varphi(\operatorname{ord}(\theta_{i,J}))}\right)
\]
has reduced denominator $M_J$. Its numerator satisfies
$(M/M_J)(M_Jg_J)\equiv-\epsilon\pmod{M_J}$. Moreover
\begin{equation}\label{eq:component-norm}
 \sum_{i\ne0}w_i^2=2g_J-g_J^2,\qquad
 \sum_{i\ne0}w_i/b_i=g_J(\epsilon+1/M).
\end{equation}
\end{proposition}
\begin{proof}
Set $W=W_J$, $V=W_{J'}$ and let $E$ be the $J$-local span of
the nontrivial $J$-phase fibers. The scalar action just proved gives
$V=\gamma^3 I$ on those fibers and hence on $E$. On $E^\perp$ all
the generating $J$-unipotents act trivially. Therefore
\begin{equation}\label{eq:ternary-split}
 s=VW=V+\gamma^3(W-I).
\end{equation}
Covariance sends a partial Weyl to itself times a local torus on the
right. The vacuum torus character is trivial in square dimension.
Thus $w,Vv_0$ are rational and $W^{-1}v_0=w$. On an active
$J$-coordinate $i$ the complementary operator preserves the vacuum
$J$-phases, so $(Vv_0)_i=0$ and $\gamma^3 w_i=1/b_i$. This gives
$\epsilon=\pm1$. If $b_i$ is supported on $J$, each complementary
local factor preserves $\mathbb C v_i$ by singleton blocks, and
$s=VW$ gives $|w_i|=1/b_i$.

Choose $k$ equal to one on the $J$-part of $L$ and zero on its
complement. The Bruhat relation identifies $W^{-1}$ with
$sT^ks^{-1}$ up to diagonal root-of-unity factors. Its vacuum entry
is unchanged since the $J$-unipotents fix $v_0$. The argument
of \eqref{eq:phase-column} gives $b_iw_i\in\mathbb Z$. In
\[
 (sT^ks^{-1}-I)_{i0}=\sum_y(\theta_y^k-1)\chi_y(i)/f_y
\]
all nonzero summands have denominator supported on $J$.
Interchanging $J,J'$ gives codegrees prime to $J$. The diagonal
factors are units, hence $w$ is integral away from $J$ and $Vv_0$
is integral at $J$. The vacuum entry of the Bruhat relation, averaged
over Galois, now gives the formula for $g_J$ by the Ramanujan sum
\cite[Theorem~8.6]{Apostol}. Its summands are positive and some
twist is active at each prime of $J$.

At the vacuum, \eqref{eq:ternary-split} reads
$1/M=(Vv_0)_0-\epsilon g_J$. It gives exactly the stated denominator
and congruence, and shows $w_0\ne0$. At other coordinates,
$A_i=\epsilon u_iw_i$ is integral by the two integrality assertions,
and multiplying \eqref{eq:ternary-split} by $b_i$ gives its
congruence. Unitarity and $w_0\ne0$ give $|A_i|<u_i$. For a single
prime, $w$ lies in the principal block by local invariance, and the
reverse inclusion follows from Lemma~\ref{blk:height}.
Finally $z=(W-I)v_0\in E$ satisfies
$\langle z,sv_0\rangle=\epsilon\langle z,Wv_0\rangle
=\epsilon g_J$. Subtract the vacuum contribution and use
$\|w\|=1$ to obtain \eqref{eq:component-norm}.
\end{proof}

For an isolated $p$ with $h$ active labels of denominator at least
$B$, the following entry bound, with active mass at most $h/B^2$
and its negative quadratic term discarded, gives
\begin{equation}\label{eq:cost-large-valuation}
 \sum_{i\ne0:\,1<u_i<b_i}u_i^{-2}
 <\frac{p+1}{p-1}\frac h{B^2}.
\end{equation}
The inequality is strict because $g_{\{p\}}>0$.
When denominator $p$ is absent, $B\ge p^2$; hence a nonunit
with $p\Vert b_i$ requires $h>p^2(p-1)/(p+1)$.

\begin{proposition}[A local entry bound from active mass]
\label{new:local-entry-mass}
Suppose $D=M^2$ and $p$ is an isolated odd prime. Let $E$ be the
reciprocal-square mass of the $p$-active labels, and suppose
$E\le E_0$. For any set $J$ of nonunit labels with mixed denominators
$b_i=u_i k_i$, where $u_i>1$ is a power of $p$, $k_i>1$ and $(k_i,p)=1$,
put
\[
 C_J=\sum_{i\in J}\frac{\|k_i^{-1}\|_{u_i}^{\,2}}{u_i^2},
 \qquad
 F_p(x)=\frac{p+1}{p-1}x-\frac{p^2}{(p-1)^2}x^2.
\]
Here $\|a\|_u=\min_{z\in\mathbb Z}|a-zu|$.
Then
\[
 C_J\le
 F_p\!\left(\min\left\{E_0,\frac{p^2-1}{2p^2}\right\}\right).
\]
In particular $C_J\le(p+1)^2/(4p^2)$, independently of rank.
In odd dimension, any occurring nonunit denominator $pk$, with
$(k,p)=1$ and $k>1$, therefore satisfies
$8\|k^{-1}\|_p^2\le(p+1)^2$.
If denominator $p$ is absent and $h$ labels are $p$-active, one
may take $E_0=h/p^4$. In odd dimension, one occurring mixed
denominator supplies a dual pair in $J$.
\end{proposition}
\begin{proof}
Let $g=g_{\{p\}}$ and $w=W_pv_0$. Proposition~\ref{new:component-column}
gives $|w_i|\ge\|k_i^{-1}\|_{u_i}/u_i$. The active labels are
pure and have $|w_i|=1/b_i$, so \eqref{eq:component-norm} gives
$C_J+E\le2g-g^2$. The Ramanujan means give
$0<g\le\alpha E$, where $\alpha=p/(p-1)$.
Consequently
\[
 C_J\le(2-1/\alpha)g-g^2,
 \qquad0<g\le\alpha E_0.
\]
Maximizing this concave quadratic proves the formula; no assumption
that $\alpha E_0\le1$ is necessary. When denominator $p$ is absent,
all active denominators are pure powers at least $p^2$.
\end{proof}

\paragraph{Valuation-one classes in odd dimension.}
In odd square dimension, let $p\ge5$. There is an exact version
when a complete denominator-$p$ class is
known and $p\Vert M$. With $m$ its size and $j=A/(p-1)$,
$g=j/p$. Subtracting its entire norm contribution $m/p^2$ gives
the proper norm budget $H=2pj-j^2-m$. Thus any proper dual pair
of denominator $pk$ obeys
\[
                 2\|k^{-1}\|_p^2+4\le H.
\]
The additional four comes from the other two distinct proper
degrees supplied by the gcd-one, pairwise-noncoprime proper
support in Lemma~\ref{new:weyl}; each has a dual pair and each
nonzero integral Weyl coefficient has magnitude at least one.
This is the same Weyl budget as \eqref{eq:weyl-cost}, now used
as a denominator floor before any phase enumeration.

\begin{proposition}[Neutral orbits joining twist components]\label{new:mixed-neutral-rank}\label{odd:neutral-orbits}
Suppose $D=M^2$ and the twist graph has $c\ge2$ components.
At least $c-1$ Galois orbits consist of nonunit twist-one labels
whose denominators meet two or more components. If $R$ is also
non-pointed and based-simple, these labels have at least three
distinct degrees, so there are at least $\max\{c-1,3\}$ such orbits.
In odd dimension each of these orbits contains at least two labels.
They are disjoint from every neutral orbit whose denominator is
supported on one component.
\end{proposition}
\begin{proof}
Normalize the local factors so that their upper unipotents fix the
vacuum. Odd factors already have this normalization. An odd
component exists, so Proposition~\ref{new:component-column} gives
$\epsilon=\gamma^3=\pm1$, hence $\gamma=\epsilon$. At two, if
$\epsilon=-1$, tensor the local representation with the sign
character of $\mathrm{SL}_2(\mathbb F_2)\simeq S_3$; it is $-1$
on the upper unipotent and the Weyl element, and $1$ on the torus.
A scalar factor outside the dimension support can be discarded.

For a component $J$ write $w_J=W_Jv_0$ for this normalized Weyl
column and $g_J=1-(w_J)_0$. The local span of its nontrivial phase
spaces is fixed pointwise by every complementary normalized factor;
its own factor is trivial on the orthogonal complement of that span.
Thus the differences $(W_J-I)v_0$ for different components are
orthogonal. They are nonzero: otherwise the local upper and lower
unipotents would fix $v_0$, then also $sv_0$, whose entries at the
active labels are nonzero. Since $W_J^{-1}v_0=W_Jv_0$, this gives
$g_J>0$ and
\begin{equation}\label{eq:component-cross}
 \sum_{i\ne0}(w_J)_i(w_K)_i=-g_Jg_K\qquad(J\ne K).
\end{equation}
Covariance makes $w_J$ rational. Every local torus fixes it:
on $J$ commute the torus across $W_J$ and invert it, and on the
complement use commutation. Hence its coordinates are constant on
Galois orbits. If $b_i$ is prime to the primes of $J$, its singleton
local blocks have trivial normalized action, so $(w_J)_i=0$.

Restrict the columns to labels whose denominators meet at least two
components. All cross terms removed from
\eqref{eq:component-cross} are zero. These labels are neutral,
because every active label has codegree supported on its own twist
component. Compress an orbit of size $h$ to its common coordinate
times $\sqrt h$, and call the restricted vectors $v_J$. Their Gram
matrix is
\[
 \operatorname{diag}\bigl(\|v_J\|^2+g_J^2\bigr)-gg^t,
 \qquad g=(g_J)_J.
\]
The diagonal matrix is positive definite, so subtracting a rank-one
matrix leaves rank at least $c-1$. There must be at least that many
orbit coordinates. If $R$ is non-pointed and based-simple,
Theorem~\ref{new:prime-cut}, applied to the twist components, gives
three distinct degrees on this same Galois-stable set, hence at
least three orbits. Lemma~\ref{alg:odd-parity} gives the
odd-dimensional count. The last assertion follows from the supports.
\end{proof}

In odd square dimension Proposition~\ref{new:weyl-parity} identifies the reductions of these columns modulo two with the field components at every activity level.

\begin{proposition}[Coupled residue lifts across components]
\label{new:coupled-weyl}
Suppose $D=M^2$, the twist graph has at least two components,
and normalize the local Weyl operators as in
Proposition~\ref{new:mixed-neutral-rank}. For a nonunit denominator
$b$, let $u_1,\ldots,u_t>1$ be its nontrivial parts on the twist
components, so that $b=\prod_j u_j$. Put
\[
 r_j=(b/u_j)^{-1}\bmod u_j\in\{1,\ldots,u_j-1\},\quad
 q_j=r_j/u_j,\quad K=\sum_j q_j-1/b\in\mathbb Z.
\]
The signed local coordinates $\epsilon w_j$ are $q_j$ or $q_j-1$,
with exactly $K$ negative choices. Moreover the sum of their positive
coordinates is strictly less than one. These are exactly the constraints
given by the congruences and strict absolute bounds for every union
of components at this coordinate.

Among all these choices, allowing prescribed signs, both the least
total squared norm and the least sum of positive coordinates select
the largest free $q_j$ for the remaining negative positions. Thus a
single sorted calculation either proves these entry constraints
inconsistent or gives their exact minimum norm. In odd dimension,
a prescribed $2$-adic parity of $w_j$ fixes its sign choice, since
$r_j$ and $r_j-u_j$ have opposite parity.
\end{proposition}
\begin{proof}
The mutually orthogonal local nontrivial spans used in
Proposition~\ref{new:mixed-neutral-rank} give
\[
 \prod_jW_j=I+\sum_j(W_j-I)=\epsilon s.
\]
Consequently at a nonunit coordinate the local entries sum to
$\epsilon/b$, and a union entry is the sum over that union.
Proposition~\ref{new:component-column} gives the congruence and
strict bound for unions of odd components. If a union contains
the component at two, use its complementary odd union and subtract
its entry from $\epsilon/b$. Write $K=J'$ and $g_K=A/M_K$. The complementary vacuum congruence
gives $M_JA\equiv-\epsilon\pmod{M_K}$, so
\[
 g_J=\frac{M_JM_K-M_JA-\epsilon}{M_JM_K}
\]
has reduced denominator $M_J>1$. At a nonunit coordinate write
$b=u_Ju_K$ and $\epsilon w_K=a/u_K$. Then
$\epsilon w_J=(1-au_J)/(u_Ju_K)=c/u_J$, where $c$ is integral and
$cu_K\equiv1\pmod{u_J}$. Since $g_J$ is nonintegral, the vacuum
entry is nonzero and unitarity gives $|c|<u_J$. The full union is
$\epsilon s$ and satisfies the same bound directly.

Thus $\epsilon w_j=q_j-e_j$ with $e_j\in\{0,1\}$.
Summing gives $\sum e_j=K$. For any subset the union coordinate
has modulus less than one. Taking all positive coordinates gives
their sum less than one; the absolute sum of negative coordinates
is smaller by $1/b$, so these two bounds in turn imply the strict
bound for every subset. The union congruences follow from the
Chinese remainder theorem. Finally changing $q_j$ to $q_j-1$
changes its squared norm by $1-2q_j$ and decreases the positive
sum by $q_j$. Both objectives therefore choose the largest
remaining $q_j$. The parity assertion follows because $u_j$ is odd.
\end{proof}

Writing $L_i$ for this minimum squared norm at each nonunit label,
the column identities also give the uniform necessary inequality
\[
 \sum_{i\ne0}L_i\le
 \sum_J(2g_J-g_J^2),\qquad
 \sum_Jg_J=1-\epsilon/M.
\]
When signs are prescribed by a field or inertia constraint, they
must be imposed before taking the minimum. Separate minima for
individual component entries are also valid lower bounds, but
need not attain the joint minimum simultaneously.

\begin{proposition}[Activity, height and the last denominator gap]
\label{new:activity-height-gap}\label{new:sparse-component-floor}
Suppose $D=M^2$, and let $p$ be an isolated odd prime.
Put $\delta=1$, or $\delta=2$ in odd dimension.
Let $p^a$ be the largest active denominator and let $p^b$ be
the next smaller active denominator, with $b=0$ if none exists.
Write $h_a$ for the number of active labels at $p^a$ and
$H=v_p(M)$. Then $H\le2a-1$. If $H\ge2b$, then
\[
                  h_a\ge\delta(p-1)p^{2a-1-H};
\]
if $H\le2b-1$, then
\[
                  h_a\ge\frac\delta2(p-1)p^{2(a-b)}.
\]
In particular, $h_a<\delta p(p-1)$ implies
\begin{equation}\label{eq:sparse-height}H=2a-1.\end{equation}
If $H$ is even, then $h_a\ge\delta p(p-1)$.
If $H\le2a-3$, then $h_a\ge\delta p^2(p-1)/2$.
These assertions make no restriction on the activity at smaller
denominators. If the total activity $h$ satisfies
$h<\delta p(p-1)$ and $B$ is the smallest active denominator,
the $p$-part of every nonunit denominator divisible by $p$
is at least $B$. Under this total-activity hypothesis, absence of denominator $p$
forces $p^2\mid b_i$ whenever $p\mid b_i$.
\end{proposition}
\begin{proof}
At denominator $p^a$, let $\delta u,\delta v$ be the primitive
phase multiplicities in the two square classes. Equal primitive
phase masses and the gap to the preceding active level give
\[
                      p^{2(a-b)}\mid u-v.
\]
If $u\ne v$, this already accounts for at least
$\delta(p-1)p^{2(a-b)}/2$ labels. Otherwise primitive phases
contribute $\delta u$ to the coefficient $c_a$ in the charge
$pc_a/p^{2a}$. Every higher-order square-unit packet contributes
an integer multiple of $\delta(p-1)/2$ to $c_a$; hence in this
balanced case $c_a$ is a positive integer divisible by $\delta$.
In both cases its contribution costs at least $(p-1)c_a$ labels.

More explicitly, if $h_{a,1}$ counts primitive active labels and
$h_{a,>1}$ counts higher-order ones, then
\[
 c_a=\frac{h_{a,1}}{p-1}+\frac{h_{a,>1}}p>0,
 \qquad h_a\ge(p-1)c_a.
\]
The coefficient is always integral at $p$, even without balance.
Every lower charge has $p$-denominator at most $p^{2b-1}$.
The reduced-denominator assertion for the local Weyl column gives
$H\le2a-1$. If $H\ge2b$, comparison of valuations forces
$v_p(c_a)=2a-1-H$; if $H\le2b-1$, it forces
$p^{2(a-b)}\mid c_a$.

In the balanced case, these divisibilities and $\delta\mid c_a$
give the stated bounds, with a factor two to spare in the second
one. In the unbalanced case use the primitive count already
obtained. For the first assertion, its exponent
$2a-1-H\le2(a-b)-1$, so the unbalanced bound is larger than
the claimed bound because $p/2>1$. The three final consequences
follow with $a-b\ge1$.


For the denominator floor suppose $h<\delta p(p-1)$. The same
largest-imbalance argument balances primitive multiplicities at every
denominator. Square-unit covariance and duality exclude phases above
order $p^2$ and allow at most one minimal order-$p^2$ packet. Thus
\[
 h=\delta\bigl((p-1)j+p(p-1)k/2\bigr),\qquad
 k\in\{0,1\},\quad
 j\le p-1\ (k=0),\quad j\le(p-1)/2\ (k=1).
\]
Let $E$ be the active reciprocal-square mass. Primitive phases
have coefficient $p/(p-1)$ in $g$, and order-$p^2$ phases have
coefficient one. Hence
\[
 2g-E\le\frac{\delta}{B^2}
       \bigl((p+1)j+p(p-1)k/2\bigr)
 \le\frac{\delta(p^2-1)}{B^2}.
\]
For $k=1$ the intermediate numerator is at most
$\delta(p-1)(2p+1)/2$, which is smaller. A nonunit denominator
with $p$-part $1<u<B$ is inactive and its Weyl coordinate costs
at least $u^{-2}$; its dual supplies the second copy when
$\delta=2$. Subtracting the active mass from
\eqref{eq:component-norm} would give
\[
 \delta p^2/B^2\le\delta/u^2\le2g-g^2-E
 <2g-E\le\delta(p^2-1)/B^2,
\]
a contradiction. This argument includes inactive pure denominators.
\end{proof}

\begin{proposition}[The leading local Gauss coefficient]\label{new:leading-gauss}
Let $D=M^2$ and let $p$ be an isolated odd prime. Write $p^a$ for the
largest active denominator, $p^b$ for the preceding one ($b=0$ if
there is none), $H=v_p(M)$ and $n=M/p^H$. On the largest level let
$h_1,h_>$ count labels whose primary twist has order $p$, respectively
strictly greater than $p$. Put $c=h_1/(p-1)+h_>/p$. If $H\ge 2b$,
then $u=p^{H+1-2a}c$ is a $p$-adic unit and
\[
                  nu\equiv-\epsilon
          \pmod {p^{\min(H,H+1-2b)}} .
\]
Set $\delta=2$ in odd dimension and $\delta=1$ otherwise.
If the total activity is less than $\delta p(p-1)$, then
$H=2a-1$, $c\in\delta\mathbb Z_{>0}$, $(p-1)c\le h_1+h_>$, and
\[
            n c\equiv-\epsilon
         \pmod {p^{\min(2a-1,2(a-b))}}.
\]
The sign is common to all such primes. In particular the square class
of $c$ modulo $p$ depends only on this sign and the parities of the
other prime heights.
\end{proposition}
\begin{proof}
The Ramanujan means express $g_p$ as a sum of charges $p c_j/p^{2j}$
over the active denominator levels. The $c_j$ are $p$-integral by
square-unit symmetry. The numerator $p^H g_p$ is a $p$-adic unit and
$n p^H g_p\equiv-\epsilon\pmod {p^H}$ by
Proposition~\ref{new:component-column}. Each lower-level contribution
to $p^H g_p$ is divisible by $p^{H+1-2b}$. This proves the first
assertion, including the case of a single level. Under the activity
bound, Proposition~\ref{new:activity-height-gap} gives $H=2a-1$;
its largest-level balance shows that the two primitive square classes
have the same multiplicity there. Their common multiplicity is
divisible by $\delta$, as is $h_>/p$ by the higher-order packet sizes.
The count follows directly from the definition of $c$.
\end{proof}

\begin{lemma}[An unweighted Gauss congruence]\label{new:dyadic-trace}
In odd dimension, let $t=\operatorname{Tr}(T)$ and
$p^+=\sum_i d_i^2\theta_i$. In the ring of cyclotomic integers one has
\[
 t\equiv p^+\pmod {16},\qquad
                  t\overline{t}\equiv r\pmod {16}.
\]
\end{lemma}
\begin{proof}
On each nonunit dual pair the degrees and twists are equal, and the
common odd degree $d$ satisfies $8\mid d^2-1$. Thus $p^+-t$ is divisible by sixteen. Apply Lemma~\ref{alg:odd-parity}.
Use $p^+\overline{p^+}=D$.
\end{proof}

For an exact odd square rank, $D\equiv r\pmod{16}$ fixes the possible
parities $e_q$ of all heights by
$(\prod_{e_q=1} q)^2\equiv r\pmod{16}$. At a sparse prime absent as
a denominator, $a\ge2$, so the leading coefficient is a positive
even integer below $2p$, and $nc\equiv-\epsilon\pmod{p^2}$.
For each parity vector and common sign, compute the possible
cofactor residues from the cyclic subgroups generated by $q^2$
modulo $p^2$, with fixed heights imposed directly. If the least
admissible coefficient is $c$, reserve $(p-1)c$ labels at the
highest active level. A dense branch instead reserves $2p(p-1)$
active labels. All these reservations are made on disjoint field
components. The same test modulo $p$ uses the prescribed multiplicity
at primes of height one. 

\begin{lemma}[The vacuum column at valuation one]\label{new:weyl}
Suppose $D=M^2$, $p\Vert M$ is odd and the ring is non-pointed and
based-simple. Assume that a denominator $p$ occurs, or more generally
that every $p$-active twist has $p$-power order. Put $n=M/p$.
Each primitive $p$th root occurs $j$ times, $1\le j<p$, and these
are all the $p$-active labels, of degree $n$. There are integers
$|a_i|<p$ ($i\ne0$) and a sign $\epsilon$ satisfying
\begin{equation}\label{new:local-column-sums}
 a_i\equiv d_i/n,\quad nj\equiv-\epsilon\pmod p,\qquad
 \sum_{i\ne0}a_i^2=2pj-j^2,\quad
 \sum_{i\ne0}(d_i/n)a_i=pj+\epsilon j/n .
\end{equation}
The nonzero coefficients occur precisely on the $p$-free nonunits.
Their degrees divide $n$ and share a prime in pairs, but have gcd one
even after deleting the copies of $n$. In particular at least three
distinct proper divisors of $n$ occur.

For $p\ge5$, let $m$ be the multiplicity of $n$, $z=pj-m$,
$H=j(p-j)$, and $k_i=n/d_i$ on proper $p$-free nonunits. All
coefficients at degree $n$ are one, and on the proper labels
\begin{equation}\label{eq:weyl-cost}
 \sum a_i^2=H+z,\quad \sum a_i/k_i=z+\epsilon j/n,\quad
 k_i a_i\equiv1\pmod p,\quad p\,|z|\le H+(p+1)j/n .
\end{equation}
\end{lemma}
\begin{proof}
Every active label has $p$-free degree. If a denominator-$p$ row occurs,
it is principal at $p$ and meets every such degree, forcing its twist
to be $p$-primary (Corollary~\ref{blk:primepower}). The preceding
splitting then makes each active codegree a $p$-power, necessarily
$p^2$ since it is a nontrivial square dividing $M^2$. Theorem~\ref{thm:tame-orbit}
gives order $p$. Lemma~\ref{blk:moments} gives common multiplicity $j>0$
and Proposition~\ref{sq:saturation} gives $(p-1)j\le p^2-2$, hence $j\le p$.

Apply Proposition~\ref{new:component-column} with $J=\{p\}$. Each
active denominator is $p$, so $g_J=j/p$. Its reduced denominator is $p$,
hence $j<p$. Set $a_i=p\epsilon w_i$. The component congruences,
using $d_i/n=p/b_i$, and \eqref{eq:component-norm} give
\eqref{new:local-column-sums}.

The bound and congruence identify the nonzero $a_i$. All $p$-free labels
meet a principal active row, hence lie in $P_p$. Their degrees divide $n$
and are pairwise noncoprime by principality and Lemma~\ref{lem:basic-table}
for nonunits. Lemma~\ref{prime-square} gives gcd one on a subset of
them. Deleting copies of $n$ preserves this gcd. For $p\ge5$ the $m$
coefficients at degree $n$ are $1\bmod p$. Their squared norm in
excess of $m$ is at most
$2pj-j^2-m\le(p+1)j-j^2\le(p+1)^2/4<(p-1)^2-1$, since $m\ge(p-1)j$.
All are one. Removing them gives the first two identities of
\eqref{eq:weyl-cost}. For $k_i>1$, $k_i a_i\ge p+1$ when positive
and $k_i|a_i|\ge p-1$ otherwise, whence
$-(H+z)/(p-1)\le z+\epsilon j/n\le(H+z)/(p+1)$, as claimed.
\end{proof}

\begin{lemma}[Neutral compression]\label{r8:steinberg-compression}
In a unitary $\mathrm{SL}_2(\mathbb F_p)$-representation with a unit
Borel-fixed vector $v_0$ generating a nontrivial module, suppose every nontrivial unipotent
eigenspace is one-dimensional and meets the subrepresentation
generated by $v_0$. If $W$ is Weyl and $Q$ projects onto the
unipotent-fixed space, put $g=1-\langle Wv_0,v_0\rangle$ and
$h=Q(W-I)v_0$. Then $g>0$ and $Q(W-I)Q=-hh^*/g$.
\end{lemma}
\begin{proof}
The vacuum generates a quotient of the permutation module
$\mathbf1\oplus\mathrm{St}$ (Theorem~\ref{thm:spherical}).
Steinberg exhausts the nontrivial eigenspaces, so its orthogonal
complement is trivial. The Weyl operator is self-adjoint and
$g>0$, since $Wv_0=v_0$ would make $v_0$ invariant. Its neutral
compression has image and coimage on the Steinberg neutral line;
the vacuum column $h$, with $h_0=-g$, determines the stated formula.
\end{proof}

\begin{theorem}[Prime denominators and their neutral orbits]
\label{new:one-packet}\label{new:phase-trace}\label{new:prime-den}\label{new:rational-prime-square}
In the non-pointed based-simple rotation-admissible setting, suppose
$D=M^2$ and a prime denominator $p$ occurs. Then $p\ge5$, $p\Vert M$,
and at least four primes divide $M$. Every $p$-active label has
denominator $p$ and primitive $p$-phase. Write $n=M/p$ and let $j$
be the common multiplicity of a primitive phase. Then $j$ is a proper
divisor of $p-1$. Exactly $pj$ labels have degree $n$, and their $j$
neutral members form a single Galois orbit. A denominator-$p$ row
is rational precisely when it is neutral and $j=1$.
For every primitive packet $A_k=\{X:\theta_{X,p}=\zeta_p^k\}$
and every neutral nonunit $Y$,
\begin{equation}\label{eq:packet-trace}
                \sum_{X\in A_k}\chi_X(Y)=-\epsilon a_Y .
\end{equation}
On the proper $p$-free nonunits, with $k_Y=n/d_Y>1$,
\[
 \sum a_Y^2=j(p-j),\qquad \sum a_Y/k_Y=\epsilon j/n,
 \qquad k_Y a_Y\equiv1\pmod p,\qquad\lcm_Y k_Y=n .
\]
Every $p$-inertia orbit of proper neutral labels has size $h<p-1$
and $j\mid h a_Y$. For each such orbit $O$ and each active $X$, $\sum_{Z\in O}\chi_X(Z)=-\epsilon h a_Y/j$. The local representation contains a unique Steinberg constituent.
\end{theorem}
\begin{proof}
Lemmas~\ref{small} and~\ref{nine} give $p\ge5$.
Theorem~\ref{rev:height} and Lemma~\ref{new:weyl} give $p\Vert M$,
isolation, local level $p$, and the integers $a_Y$.
Set $u_k=\sum_{X\in A_k}e_X$, and let $Q$ project onto the
$p$-neutral space. The phase-$k$ projection of $W_pe_0$ is
$\epsilon u_k/p$. Thus $u_k$ belongs to the nontrivial part of the
local vacuum cyclic module, a single Steinberg representation by
Theorem~\ref{thm:spherical}. Its neutral line is spanned by
$g=Q(W_p-I)e_0$, which is orthogonal to the trivial summand and
has $g_0=-j/p$. The complementary Weyl factor acts by $\epsilon$ on
$u_k$, so $Qsu_k$ lies on the same line. Its vacuum entry is $j/p$.
Consequently $Qsu_k=-g$. Since $g_Y=\epsilon a_Y/p$ at neutral
nonunits and $s_{YX}=\chi_X(Y)/p$ for $X\in A_k$, this proves
\eqref{eq:packet-trace}.

Let $O$ be a $p$-inertia orbit of neutral labels of degree $d$,
with $h=|O|\mid p-1$. For each active $X$ put
$b_X=\sum_{Y\in O}\chi_X(Y)$. These are ordinary integers, since
the row field at $X$ lies in $\QQ(\zeta_p)$ by local level and
Lemma~\ref{conduct}. They are congruent to $hd$ modulo $p$, and their
sum in every primitive packet is $-\epsilon h a_Y$. The Weyl
coefficient is constant on $O$. If $j\nmid h a_Y$, the $j$
integers in each packet are nonconstant. The identity
$j\sum b_X^2=(\sum b_X)^2+\sum_{X<X'}(b_X-b_{X'})^2$
then gives squared norm at least
$((h a_Y)^2+p^2(j-1))/j$ per packet: at least $j-1$ pairs differ.
Unitarity, applied to $\sum_{Y\in O}e_Y$ and retaining the
vacuum coordinate separately, gives
\[
       \sum_{X\ \mathrm{active}} b_X^2
             \le hp^2-h^2(d/n)^2<hp^2 .
\]
Thus $j\mid h a_Y$ whenever $h\le(p-1)/2$.

At a neutral denominator-$p$ label $a_Y=1$, $d=n$. The only
other orbit size to consider is $h=p-1$. Under $j\nmid h$, comparison
of the above budgets becomes
\[
 (p-1)^2+p^2(j-1)\le j(p^2-(p-1)),
 \qquad j(p-1)\le2p-1 .
\]
This forces $j\le2$, contradicting $j\nmid p-1$. Hence all neutral
orbit lengths at degree $n$ are divisible by $j$. Adding the
$j(p-1)$ active labels gives $j\mid m$, where $m$ is the size of
the degree class.

Put $z=pj-m$. Equation~\eqref{eq:weyl-cost} gives $|z|<2j$.
If $z=j$, its signed upper estimate gives
$1+\epsilon/n\le (p-j+1)/(p+1)$. If $z=-j$, its signed lower
estimate gives $1-\epsilon/n\le (p-j-1)/(p-1)$.
In either case $\epsilon=-z/j$ and $jn\le p+z/j$.
Since $jn\equiv-\epsilon\pmod p$ and $n>2$, equality holds.
Equality in the signed estimate makes $k_Ya_Y=(z/j)(p+z/j)$
at every proper label, hence $a_Y=-j\epsilon d_Y$ there.
Equation~\eqref{eq:packet-trace} would express $j d_Y$ as a sum
of $j$ entries, all of magnitude strictly less than $d_Y$.
There are at least three such proper degrees by Lemma~\ref{new:weyl}.
This excludes both nonzero multiples of $j$. Thus $m=pj$,
leaving exactly one neutral orbit, of size $j$, at degree $n$.
Its field is $p$-primary, so this is also its full Galois orbit.
Since $p\ge5$, square-unit covariance excludes a rational active row;
the neutral orbit proves the rationality assertion. The excluded
multiplicity $p^2-p$ in Lemma~\ref{lem:congruent-rows} makes $j$ a
proper divisor of $p-1$. Substitution in \eqref{eq:weyl-cost} and
the gcd-one property give the displayed identities. Every proper
$p$-free nonunit degree has at least two prime factors by
Lemma~\ref{lem:basic-table}. Their gcd one therefore forces at least
three primes other than $p$. Finally
Proposition~\ref{round:neutral-inertia} bounds the total neutral
inertia deficit by $2(j-1)\le p-3$, so no neutral orbit has size $p-1$.

For any neutral nonunit inertia orbit $O$ put
$c=-\epsilon h a_Y/j$, including $a_Y=0$ when $p\mid d_Y$.
We have proved $c\in\ZZ$ and $c\equiv h d_Y\pmod p$ since
$jn\equiv-\epsilon$. In each packet the integers $b_X$ have mean $c$.
A nonconstant packet costs at least $2p^2$ in squared deviations.
Inertia permutes the packets of each square class transitively preserving
their multisets of $b_X$. The resulting cost $(p-1)p^2$ exceeds the
column budget $hp^2$. Hence $b_X=c$ for every active $X$.
Thus projection onto the active coordinates after $W_p$ maps each
neutral orbit sum, and also $e_0$, into the line spanned by the
indicator of all active labels. These sums span the Borel-fixed space
by the unsigned torus action. Frobenius reciprocity puts this space
in the trivial and Steinberg isotypic parts. In each Steinberg copy the
indicated projection is injective on its Borel line: otherwise $W_p$
preserves that unipotent-neutral line, which then is invariant under
$\langle B,W_p\rangle=\mathrm{SL}_2(\mathbb F_p)$. Orthogonal
copies have orthogonal images. Hence the vacuum Steinberg copy is unique.
\end{proof}



\begin{proposition}[A pure denominator beyond the active labels]\label{new:pure-extra}
In the non-pointed based-simple square case, at least $p$ nonunit
labels have $p$-power denominator at every isolated odd prime $p$.
If $p$ itself is absent as a denominator and the dimension is odd,
there are at least $2p$ such labels. Suppose $p$ is absent. Write $s=1$ for the general
bound and $s=2$ for the paired bound in odd dimension.
If only $s(p-1)$ labels are active at $p$, they have a common
denominator $B=p^a$, $a\ge2$, and the inactive nonunit labels of
$p$-power denominator have mass greater than $s/(pB)^2$.
\end{proposition}
\begin{proof}
In the paired case every nonunit occurs with a distinct dual of
equal degree and twist (Lemma~\ref{alg:odd-parity}). A fiber containing
a higher phase and its square-unit translates therefore costs at least
$sp(p-1)/2>s(p-1)$ labels. With just $s(p-1)$ active labels, each
primitive $p$-phase occurs on one label or on one dual pair. Equality
of the masses gives their common denominator $B$.
Proposition~\ref{new:component-column} gives $g=g_{\{p\}}=sp/B^2$
and $M_p=B^2/p$.

Let $V$ be the mass of all nonunit pure $p$-power denominators and $E=B^2 V-s(p-1)$ the rescaled inactive part. On every pure denominator $\epsilon w_i=1/b_i$ by its absolute value and congruence. Proposition~\ref{new:activity-height-gap} gives $u\ge B$ for every nontrivial $p$-part of a denominator. On a mixed denominator $b_i=u k$, the congruence $A_i k\equiv1\pmod u$ and $k>1$ imply $|A_i|k\ge u-1$. Hence $|w_i|/b_i\le w_i^2/(B-1)$. Subtract the pure contributions from both identities in \eqref{eq:component-norm} to obtain the two-sided estimate
\begin{equation}\label{eq:packet-shell}
 \left|s-E+s\epsilon p/M\right|
 \le \frac{s(p+1)-E-(sp/B)^2}{B-1}.
\end{equation}
In particular $E/s\ge1-1/(p-1)-1/(6p)>1/p^2$: here $B\ge p^2$ and the active nonunit degree is at least six. In odd dimension the inactive pure labels occur in pairs.

More than the minimum number of active labels already gives the stated
count: in odd dimension their number is even.

If $p$ occurs among only $p-1$ pure labels, all are active with the
same denominator $p$. Now $g=1/p$ gives $p\Vert M$; Lemma~\ref{nine}
excludes $p=3$, and Theorem~\ref{new:one-packet} gives $p$ copies for $p\ge5$.

\end{proof}

There is also a bound for unequal active denominators. Suppose $p$ is isolated, the denominator $p$ is absent, and $h<p(p-1)/2$ labels are active. They all have primitive phase by square-unit covariance. Write $V$ for their mass and $P$ for the mass of inactive nonunit pure $p$-powers. Then $g=pV/(p-1)$, and Proposition~\ref{new:activity-height-gap} gives $u\ge p^2$ on every mixed denominator. Its congruence gives $|w_i|/b_i\le w_i^2/(u-1)$ as in \eqref{eq:packet-shell}. Subtracting the pure mass from both identities in \eqref{eq:component-norm} gives
\begin{equation}\label{eq:unequal-pure-shell}
 P\ \ge\ g-V-g/M-(2g-V)/(p^2-1)
   =\frac{V}{p-1}\left(1-\frac pM-\frac1{p-1}\right)>0.
\end{equation}
Here $M\ge6p^2$. In particular, for $p\ge5$, at most $2(p-1)$ pure nonunit positions force a single packet. Indeed $2(p-1)<p(p-1)/2$ and the primitive multiplicities agree at each denominator by the descent proving \eqref{eq:sparse-height}. There are at least $p-1$ active labels, fewer than $2(p-1)$ by the strict inequality above. Proposition~\ref{round:packet-idempotent} then gives exactly $p$ copies at their common denominator $B$, the smallest pure denominator, and $M_p=B^2/p$.

\begin{lemma}[An inactive prime-power denominator]\label{odd:inactive-pure}
Work in the non-pointed based-simple integral modular-data framework
with compatible rotations and odd square global dimension. Let $p$ be
an isolated odd prime, absent among the prime denominators. Suppose fewer than $p(p-1)$ labels are
$p$-active. Write $B$ for their smallest denominator and $K$ for an
upper bound on the number of inactive nonunit labels of $p$-power denominator.
If
\[
 K<2p^2\left(1-\frac1{p-1}-\frac1{6p}\right),
\]
there is an inactive dual pair of $p$-power denominator whose degree is
divisible by $p$. Every inactive $p$-power denominator is at least $B$.
\end{lemma}
\begin{proof}
A higher phase costs at least $p(p-1)$ labels, so all active phases
have order $p$. The proof of Proposition~\ref{new:activity-height-gap}
balances their multiplicities at every denominator and gives
$M_p=B_{\max}^2/p\ge pB$. That proposition also puts every
nontrivial $p$-part at least $B\ge p^2$, proving the last assertion.
Thus the proof of the shell estimate \eqref{eq:unequal-pure-shell} applies.
For active mass $V$ and inactive pure mass $P$ it gives
\[
 P\ge\frac{V}{p-1}\left(1-\frac pM-\frac1{p-1}\right)
 \ge\frac2{B^2}\left(1-\frac1{6p}-\frac1{p-1}\right),
\]
using $V\ge2(p-1)/B^2$ and $M\ge6p^2$.
If all inactive pure degrees were prime to $p$, their denominators
would equal $M_p\ge pB$, giving $P\le K/(pB)^2$, a contradiction.
\end{proof}



\begin{lemma}\label{isolated}
In the non-pointed based-simple case, if $D=(3n)^2$ and $3\nmid n$,
some twist has order divisible by three and another prime.
\end{lemma}
\begin{proof}
Otherwise apply Lemma~\ref{new:weyl} at three. Subtract its two sums:
\[
 \sum_{i\ne0}\bigl(a_i^2-(d_i/n)a_i\bigr)
     =j(3-j)-\epsilon j/n\le2+2/n .
\]
Every nonzero term is at least $1-d_i/n$. Three distinct proper
degrees have indices at least $2,4,5$, since $3\nmid n$; their
contribution is at least $41/20$. Each degree has at least two
prime divisors, so gcd one requires at least three primes in $n$.
Thus $n\ge2\cdot5\cdot7$, a contradiction.
\end{proof}

\section{The largest degree classes}\label{sec:large}
Write $P_p$ for the principal $p$-block. For a row principal at two, suppose its entries on a class of $a$ odd-degree
labels have modulus one. Their sum is a rational integer by Galois
symmetry, congruent to $a$ modulo two by principality. Its absolute
value is at most $a$, so it is a sum of $a$ signs. If it equals
$\pm a$, every entry equals that sign. Also, if $m$ orthogonal rows
of squared norm $f$ have common orthogonal projection $v$, Bessel's
inequality applied to the line of $v$ gives $m\|v\|^2\le f$.

\subsection{Multiplicities at prime-power codegrees}
We use one estimate for both the original rows and the averages below.
Suppose $m$ orthogonal rows $v$ of squared norm $f$, with vacuum
entry one, are orthogonal to the degree row. Fix a label $Y$ with
$d_Y>1$ and put $w=\sum_v v(Y)$, $A=\sum_v|v(Y)|^2$, $k=f-m$, $l=f-A$.
If $A\ge m$ then
\begin{equation}\label{eq:saturation-gram}
                       0<kl-|w|^2\le k^2.
\end{equation}
The $0,Y$ compression of \eqref{eq:projection-master} has
diagonal $k/f,l/f$ and off-diagonal modulus $|w|/f$.
Equality in Cauchy--Schwarz would give proportional vectors of norm
ratio $\sqrt{l/k}\le1$, whereas their components along the degree
row have ratio $d_Y>1$. Both components are nonzero.

\begin{lemma}[Congruent orthogonal rows]\label{lem:congruent-rows}
In a nontrivial perfect ring, let $v_1,\ldots,v_m$ be a Galois-stable family of mutually
orthogonal rows with cyclotomic integer entries, squared norm
$f=p^a$, vacuum entry one, and orthogonal to the degree vector.
Suppose there is a common coordinate set $F$ such that the rows
vanish off $F$ and $v_i(Y)\equiv d_Y$ above $p$ on $F$.
Then $m\le f-\lceil\sqrt p\rceil$, and $m\ne f-p$ if $a\ge2$.
\end{lemma}
\begin{proof}
If $m>0$, degree orthogonality gives a nonunit $p$-free $Y\in F$.
Put $w=\sum_i v_i(Y)$ and $A=\sum_i|v_i(Y)|^2$.
Galois stability makes both rational integers. All summands at $Y$
are nonzero, so the norm bound gives $A\ge m$.
For $k=f-m$ and $l=f-A$, the congruences give
$w\equiv-kd_Y$ and $l\equiv kd_Y^2\pmod p$.
Thus \eqref{eq:saturation-gram} gives
$0<p\mid kl-w^2\le k^2$, proving the bound.
If $k=p$, the same argument gives $l=p$ and $w=0$ at every
$p$-free nonunit. At all other coordinates the integral sum
$w_Y=\sum_i v_i(Y)$ is divisible by $p$.
Consequently $\sum_{Y\ne0}d_Yw_Y/p=-m/p$ is impossible modulo
$p$ when $a\ge2$.
\end{proof}

\begin{proposition}\label{sq:saturation}
In a nontrivial perfect rotation-admissible ring, a nonunit degree $n$
of codegree $f=p^a$ has multiplicity $m\le f-\lceil\sqrt p\rceil$,
with $m\ne f-p$ when $a\ge2$. In particular $m\le f-3$ for dyadic
$f\ge4$. For odd $n$ this improves to $m\le f-6$ at $f=8$ and
$m\le f-8$ at dyadic $f\ge16$.
\end{proposition}
\begin{proof}
Set $I=\{X:d_X=n\}$ and $k=f-m$.
The rows in $I$ satisfy Lemma~\ref{lem:congruent-rows}, by
Corollary~\ref{blk:primepower}. This proves the first two assertions,
including the dyadic bound $m\le f-3$.

For odd $n$ and dyadic $f$, assume $k\le5$ at $f=8$, or $k\le7$ at
$f\ge16$. Each $I$-row sees precisely $P_2$ and meets all odd labels.
An even label in $P_2$ costs $4m\ge f$ in its column by lcm integrality.
Hence $P_2$ comprises just the odd labels, whose degrees divide $n$.
The integrality of $D^{-1}\sum_{P_2}d_Y^2$ at two and $d_Y^2\equiv1\pmod8$
give $8\mid|P_2|$. The nonunit hit degrees have gcd one, hence at
least two distinct proper divisors of $n$ occur. Thus
$f-8<m+3\le|P_2|\le f$ and $|P_2|=f$. All entries there are roots of unity.

There are $h=k-1$ proper nonunit labels of degrees $a_i$, with $2\le h\le6$
($h\le4$ at $f=8$). Each degree-class sum is a sum of signs, as above.
Thus $1+nt+\sum_i\epsilon_i a_i=0$, with $t\equiv m\pmod2$ and
$\epsilon_i=\pm1$. Here $n\ge15$ and $3\le a_i\le n/3$ by gcd one.
For $h=2,4$, the proper signed sum has modulus at least $n-1$.
This excludes $h=2$; at $h=4$, mixed signs give at most $n-3$,
so all signs agree and their common value is fixed by the unsigned
sum $n\pm1$. For $h=3,5$, $t$ is even and
$1+5n/3<2n$, hence $t=0$. Two sign solutions differ at an even
number of positions. Four differences would leave a signed degree
$-1$; two only exchange equal degrees. Thus for $h=3,4,5$ every
proper degree $b$ has a fixed class sum $t_b$. If its multiplicity
is $s_b$, then $\sum|t_b|\ge3$: at $h=4$ this sum is four,
and otherwise it is odd and $\sum b t_b=-1$.
Projection onto the class-constant lines and the vacuum now gives
$m(1+\sum t_b^2/s_b)\le f$, hence $m(1+9/h)\le f$ by
Cauchy--Schwarz. This contradicts $m=f-1-h$ in the stated ranges.
At $h=6$, $t$ is odd; Bessel applied to the degree
vector on the vacuum and proper labels gives
$m n^2\le f(1+\sum a_i^2)$. With at most four degrees $n/3$ this gives
$m/f\le1/225+4/9+2/25<9/16\le(f-7)/f$. The remaining case has five
degrees $a=n/3$ and one coprime divisor of $3a$, necessarily three.
Its singleton entry is a sign; thus $a\mid1\pm3$, impossible for $a\ge5$.
\end{proof}

\begin{proposition}[Removing an equal-degree defect-one block]\label{new:effective}
Let $p\ne\ell$ be primes. In the non-pointed based-simple
rotation-admissible case,
let $B$ be a full equal-degree $\ell$-block of codegree $\ell p^a$.
It has $\ell$ labels, a unique $\ell$-neutral member $X$, and constant
complementary phases. Put $F=\{Y:\theta_{Y,\ell}=1\}$ and
\[
 v_B=\ell^{-1}\sum_{i\in B}\chi_i .
\]
Then $v_B$ equals $\chi_X$ on $F$ and is zero elsewhere. It is
supported on $P_p$, has integral entries, vacuum entry one, and
field $K_X\subseteq\QQ(\zeta_{p^a})$. Distinct such blocks give
orthogonal rows of squared norm $p^a$, orthogonal to the degree row;
Galois and duality permute them as they permute their neutral members.
Their modular relations have the form
\begin{equation}\label{eq:effective-gauss}
 \sum_Yd_Y\theta_Yv_B(Y)
   =\eta_B\sqrt{D/\ell}\,\theta_X^{-1},\qquad |\eta_B|=1,
\end{equation}
where $\eta_B$ is unchanged on a dual pair and, for odd $p$,
under Galois conjugation at $p$. For odd $\ell$ these factors
differ only by a sign. Thus the norm, divisibility and Gauss estimates
\eqref{st} apply with effective codegree $p^a$.
If $m$ blocks occur, then $m\le p^a-\lceil\sqrt p\rceil$;
also $m\ne p^a-p$ when $a\ge2$.
\end{proposition}
\begin{proof}
Let $e_i$ be the coordinate vectors and write $s=W_\ell W_pV$.
For odd $\ell$, Lemma~\ref{blk:defectone} gives one inactive and
$\ell-1$ active members; at two Lemma~\ref{blk:two} gives two
different phases. Put $c=\ell^{-1/2}\sum_{i\in B}e_i$.
The nonzero projection of $se_0$ on $B$ is proportional to $c$.
Since $W_\ell^{-1}se_0=W_pVe_0$ is $\ell$-neutral,
$W_\ell^{-1}c$ spans the unique neutral coordinate $e_X$.
The central operator $W_\ell^2$ preserves that line, so
$W_\ell c$ spans it too. Complementary factors in $V$ preserve
$e_X$ by Lemma~\ref{blk:admissible}. Local invariance therefore
puts $X$ in $P_p$ and $sc$ in its span.

The orthogonal complement of $c$ in $B$ maps under $s$ off $F$,
whereas $sc$ is supported on $F$. Hence
$\sqrt{p^a}\,sc=\ell^{-1}\sum_{i\in B}\chi_i$ agrees with
$\chi_X$ on $F$. This proves the norm, support, integrality and
divisibility assertions. Galois preserves the blocks and $F$.
An automorphism fixes $v_B$ exactly when it preserves $B$,
equivalently when it fixes $X$; thus the field is $K_X$.
Inertia at $\ell$ preserves the residue block and hence fixes $X$,
so Lemma~\ref{conduct} removes $\ell$ from its conductor.
Principality gives $v_B(Y)\equiv d_Y$ above $p$ on $F$,
and Lemma~\ref{lem:congruent-rows} proves both multiplicity bounds.

Average the vacuum modular relations over $B$. Its local inverse
twist sum is $1\pm i$ at two and a quadratic Gauss sum at odd
$\ell$, of modulus $\sqrt\ell$ at every embedding
\cite[Ch.~6]{Washington}. Divide out the common complementary
phase $\theta_X^{-1}$. This gives \eqref{eq:effective-gauss};
the local sum is unchanged by duality and by Galois multipliers
chosen to be one at $\ell$. The anomaly is fixed under the
$p$-part when $p$ is odd; at $p=2$, its conjugates still have
modulus one, which suffices for \eqref{st}.
\end{proof}

\begin{corollary}[Prime-square rows and prime support]
\label{new:prime-free-square-trace}
In the non-pointed based-simple setting, write $D=qM^2$.
Suppose either $q=1$ and denominator $p$ occurs, or
$q=\ell\ne p$ is prime and an equal-degree $\ell$-block has
denominator $p$. Then the nonunit degrees on the nonzero-trace
support of the actual or effective row are prime to $p$, have
gcd one, and are pairwise noncoprime. In particular at least three
distinct degrees and at least three primes other than $p$ occur:
\[
 \gcd\{d_Y:Y\ne0,\ p\nmid d_Y\}=1,\qquad \omega(M)\ge4.
\]
\end{corollary}
\begin{proof}
For an effective row, Proposition~\ref{new:effective} supplies all
the averaging identities of Lemma~\ref{prime-square}; moreover its
row equals $\chi_X$ on the neutral set, with $X\in P_p$, and is
zero elsewhere. The same direct congruence therefore gives prime-free
trace support. Both kinds of support lie in $P_p$. For two labels
$Y,Z$ in it, $\chi_Y(Z)\equiv d_Z\ne0$ above $p$; coprime
vanishing gives $\gcd(d_Y,d_Z)>1$. Gcd one then requires three
distinct degrees. A pairwise-intersecting family of prime supports
with empty total intersection uses at least three primes, all different
from $p$. Enlarging the trace support preserves gcd one.
\end{proof}

\begin{theorem}[Height of a prime-square row]\label{new:prime-square-height}\label{rev:height}
In the non-pointed based-simple rotation-admissible setting, let $p$ be odd. Suppose a nonunit degree $n$ has codegree $p^2$, or is
the common degree of a full equal-degree $\ell$-block of codegree
$\ell p^2$, where $\ell\ne p$ is prime. Then $p\nmid n$.
Equivalently the occurrence of such an actual or effective row forces
$p\Vert M$, in every rank.
\end{theorem}
\begin{proof}
Let $O$ be the Galois orbit of the actual or effective row, $h=|O|$
and $v=\sum_{X\in O}v_X$, where $v_X=\chi_X$ for an actual row.
All row fields lie in $\QQ(\zeta_{p^2})$, so the orbit is obtained
using only the $p$-part of Galois. Its sum is integral, with
$v_0=h$ and $\|v\|^2=hp^2$. By Lemma~\ref{prime-square} and
Corollary~\ref{new:prime-free-square-trace}, a nonzero $v_Y$ has
$p$-free degree dividing the $p$-free part of $n$ and has $p$-primary
twist: prime-power rows meet only $p$-primary phases by local admissibility,
and effective rows vanish on the phases active at $\ell$.
Theorem~\ref{thm:tame-orbit} gives twist order one or $p$ on $O$.
For effective rows the same proof applies to the integral orthogonal
translates of squared norm $p^2$ from Proposition~\ref{new:effective}.
Their nonzero-mean support has at least three pairwise noncoprime
nonunit degrees with gcd one by Corollary~\ref{new:prime-free-square-trace}.
Thus their field also omits $\QQ(\zeta_{p^2})^+$, and the neutral
member has no complementary twist. Write $\eta$ for their common Gauss
factor. For effective rows this common factor exists also when
$\ell=2$, since the chosen multipliers are one on the complementary
conductor, by Proposition~\ref{new:effective}. The Gauss equations give
\[
  \sum_Y d_Y\theta_Y v_Y=\eta pn\sum_{X\in O}\theta_X^{-1}.
\]
For a neutral orbit the right side has modulus $pnh$. If $p\mid n$,
the triangle inequality and $|v_Y|\le v_Y^2$ give
\[
  pnh\le h+(n/p)(hp^2-h^2)<pnh,
\]
since $n/p>1$ by based simplicity.

For an active orbit the square-unit action makes all its primitive
phases lie in one quadratic-residue class $e\in\{\pm1\}$. Apply
the Legendre-weighted average over the $p$-part of Galois used in
the quadratic Gauss sum formula \cite[Ch.~6]{Washington}. The integral vector $v$ and
$\eta$ are fixed. An order-$p$ twist $\zeta_p^k$ contributes
$(k/p)\tau/(p-1)$, with $\tau$ the quadratic Gauss sum; order one
gives zero and higher orders give zero by averaging first over
lifts modulo $p$. Every inverse row twist contributes
$(-1/p)e\tau/(p-1)$. Cancelling the common nonzero factor gives
\[
 \left|\sum_{\operatorname{ord}(\theta_Y)=p}
               (k_Y/p)d_Y v_Y\right|=pnh,
 \qquad \theta_Y=\zeta_p^{k_Y}.
\]
There is no vacuum contribution. If $p\mid n$, the left side is at
most $(n/p)(hp^2-h^2)<pnh$. This proves the theorem.
\end{proof}


\begin{proposition}[Codegrees of the form $4\ell$ and codegree eighteen]\label{new:small-effective}
In the non-pointed based-simple case let $\ell$ be an odd prime and
$h=(\ell-1)/2$. No equal-degree $\ell$-block has formal codegree
$4\ell$. The multiplicity of any degree with that codegree belongs to
$\{h,2h,\ldots,7h\}$. A degree of codegree eighteen is prime to three and has multiplicity at most ten. Consequently codegree eighteen requires at least four prime divisors of the dimension. More generally, the degree of an equal $\ell$-block of codegree $9\ell$, $\ell\ne3$ prime, is prime to three.
\end{proposition}
\begin{proof}
By Proposition~\ref{new:effective}, an equal $\ell$-block of
codegree $4\ell$ has an effective row satisfying \eqref{st}
with squared norm four. The proof of Lemma~\ref{small} forces
hit degrees $1,2,3,6$, contrary to Lemma~\ref{lem:basic-table}.

For any degree $n$ of codegree $4\ell$, an unequal block has
degrees in an integer ratio $k>1$, as in
Theorem~\ref{thm:squarefree-codegree}. The degree $n$ cannot
be smaller: the other codegree $4\ell/k^2$ would force $k=2$
and the forbidden square-free codegree $\ell$. Thus the $n$-labels
are active, in packets divisible by $h=(\ell-1)/2$. Their blocks
have total mass $m n^2\ell/(\ell-1)<D=4\ell n^2$, since they omit
the vacuum. This gives $m\in\{h,2h,\ldots,7h\}$.

For an equal $\ell$-block of codegree $9\ell$, $\ell\ne3$,
Theorem~\ref{rev:height} gives $3\nmid n$ for its common degree $n$.

At codegree eighteen every dyadic block is equal. Each effective
row meets the neutral member of every such block, and its nonunit
hit degrees share a prime in pairs but have gcd one. At least three
distinct proper divisors of $n$ are therefore hit. Their norm cost
gives $1+m_{18}/2+3\le9$, hence $m_{18}\le10$.

The four-prime conclusion at codegree eighteen is the effective-row
case of Corollary~\ref{new:prime-free-square-trace}.
\end{proof}

\begin{corollary}\label{sq:odd-small}
In the non-pointed simple case with $D=M^2$ and odd $n=M/4$, the
multiplicity $m$ of $n$ is at most six if $n$ has three prime divisors
or $r<38$.
\end{corollary}
\begin{proof}
Assume $m\ge5$. The preceding column argument ($4m>16$) makes $P_2$
exactly the odd labels, of degrees dividing $n$. They share a prime
in pairs by principality and coprime vanishing, but have gcd one
after the copies of $n$ are removed. Thus at least three distinct
proper degrees occur, and $m+4\le|P_2|\le16$, with $8\mid|P_2|$.
All sixteen entries are roots of unity, all twists dyadic. At least
three odd primes divide $n$. Four would require $2+4+6+10$ further
labels, so exactly three occur under either hypothesis.

Among the $15-m$ proper nonunits let $J_p$ be the $p$-free ones.
These three sets are nonempty and disjoint. For every degree-$n$ row
its class signs satisfy $\sum_{J_p}d_i\epsilon_i\equiv-1\pmod p$.
On a singleton or pair the signs are fixed: changing one alters the
residue, changing both negates it. They agree on repeated degrees,
so their actual entries are fixed. At $m\ge7$ two such fixed columns
and the vacuum would give $3m\le16$ by Bessel. With at most eight
positions across the three sets, the alternative requires a singleton
and a triple.
On a triple two solutions can differ only at two positions, exchangeable
after orienting their degree residues modulo sign. In each residue
class $C$ up to sign the oriented sum $t_C$ is therefore fixed in
all rows, also for actual entries since orientation depends only on
degree. Some $t_C\ne0$; the corresponding class lines contribute
$\sum t_C^2/|C|\ge1/3$. With the vacuum and singleton, Bessel now
gives $m(2+1/3)\le16$.

\end{proof}

If $D=5M^2$ and $9\mid M$, Proposition~\ref{new:small-effective}
therefore puts every denominator-three label in an unequal quinary
block. Lemma~\ref{blk:defectone} gives the same number of
denominator-six labels in that block, hence $m_3\le m_6$.

\begin{proposition}[A neutral duality orbit in a phase fiber]
\label{new:one-neutral-fiber}
Let $D=M^2$ in non-pointed based-simple rotation-admissible data.
For distinct primes $p,\ell$, fix a primitive $p$th phase $\eta$
and trivial phases away from $p,\ell$. In the resulting fiber
$F_\eta$ of labels with $\ell\mid b_i$, suppose the $\ell$-neutral
labels form one duality orbit of denominator $\ell p$. Then $p\Vert M$ and at least three other primes
divide $M$. If the full Galois orbit of these neutral labels has
square-free size and $\ell$ is odd, every active $\ell$-phase in
$F_\eta$ has order $\ell$.
\end{proposition}
\begin{proof}
Let $d\in\{1,2\}$ be the size of the neutral orbit, $y$ a member,
$H=\sum_{F_\eta}b_i^{-2}$, $c=H^{-1/2}\sum_{F_\eta}e_i/b_i$,
and $u$ the normalized sum of its neutral orbit. The charge-even
part of $F_\eta$ is $\ell$-invariant, since the deleted
$\ell\nmid b_i$ lines are singletons. Its only neutral line is
$\mathbb C u$. Projecting $W_\ell^{-1}s e_0$ onto this subspace
shows that $W_\ell^{-1}c$ spans $u$. So does $W_\ell c$, since
$W_\ell^2$ preserves the neutral line.

Put $a=b_y(W_\ell^{-1}s e_0)_y$. It is a nonzero integer
divisible by $\ell$: for odd $\ell$ this is the coefficient of
Proposition~\ref{phase:moments}, and for $\ell=2$ use
\eqref{eq:phase-column} on the complementary odd factors.
All neutral members have this coefficient by charge conjugation,
so projection gives $H=d a^2/b_y^2$. In particular $F_\eta$ has
active labels, since $|a|\ge\ell>1$. Write $k=|a|/\ell\ge1$.

As in Proposition~\ref{new:effective}, projection of $su$ onto
the full $\ell$-neutral space gives $\langle u,c\rangle sc$.
Consequently $v_\eta=sc/\sqrt H=H^{-1}\sum_{F_\eta}
b_i^{-2}\chi_i$ equals the mean of the $d$ neutral character rows on
that space and vanishes elsewhere. At $\ell\nmid b_Z$ only the
neutral rows contribute to its defining mean, giving
$v_\eta(Z)=v_\eta(Z)/a^2$. All terms vanish at $p\nmid b_Z$
by $p$-activity. Thus every hit degree divides $n=M/(\ell p)$.
As in Proposition~\ref{new:effective}, $(W_\ell^{-1}se_0)_y\ne0$
puts $y$ in $P_p$ since the factors outside $p,\ell$ preserve its
line. Up to a scalar $sc$ is $W_pu$, so its support lies in $P_p$
and has $p$-primary twists.

Vary $\eta$ through its square class. The square-unit action
reaches every phase in it, and complex conjugation interchanges
the members of each neutral pair. Thus the neutral labels form a
single Galois orbit of size $h$, with $d$ in each phase.
The sum $v=d\sum_\eta v_\eta$ is the field trace of $\chi_y$ on
the $\ell$-neutral coordinates and zero elsewhere. Its vacuum
entry is $h$ and its squared norm is $hp^2/k^2\le hp^2$ by disjointness. This norm is an integer strictly
greater than $h^2$, since vacuum orthogonality forces a nonunit hit.
Hence $k<p$ and $k^2\mid h$. When $h$ is square-free this gives
$k=1$. For odd $\ell$, the two phase moments then give $a=-\ell$
and $J=(\ell-1)C$, where $J$ is the total active weight and $C$
the weight at a primitive root. This proves the last assertion.
The orbit norm argument of Lemma~\ref{prime-square} gives
$|t|<p$ for each nonunit trace
$t=\operatorname{Tr}_{\mathbb Q(\chi_y(Z))/\mathbb Q}\chi_y(Z)$
on these coordinates. When $p\mid d_Z$, every conjugate vanishes
at the chosen place over $p$ since all row translates belong to
$P_p$. Thus $p\mid t$ and $t=0$. Vacuum orthogonality, grouped
by column orbits as in that lemma, gives gcd one on the nonunit
trace support. Its degrees are pairwise noncoprime by principality
and coprime vanishing, requiring at least three primes other than $p$.

Let $m=\langle c,T_\ell^{-1}c\rangle$ be the weighted mean of
the inverse $\ell$-phases. It is unchanged along the $p$-square
class: choose multipliers trivial away from $p$. Its absolute value
is $1/|a|$. Indeed $c$ is a unit scalar times $W_\ell u$, the
normalized upper unipotent $U$ acts on $u$ by a scalar, and
$W_\ell^{-1}U^{-1}W_\ell=UW_\ell U$ by Bruhat. Hence
$|m|=|\langle u,W_\ell u\rangle|=\sqrt d/(b_y\sqrt H)$.

The vacuum Gauss relation for $v_\eta$ has right side
$p^+m\eta^{-1}$. If $p\mid n$, every hit degree of $v$ divides
$n/p$. For odd $p$, the quadratic-character average of
Theorem~\ref{rev:height} fixes both the dyadic anomaly and $m$,
giving $pnh/k\le(n/p)(hp^2/k^2-h^2)<pnh/k$. For $p=2$ there
is one inverse phase, so the triangle inequality instead gives
$2nh/k\le h+(n/2)(4h/k^2-h^2)<2nh/k$, as $n>2$.
This contradiction proves $p\Vert M$.
\end{proof}

\Needspace{4\baselineskip}
\subsection{Codegree eight}
\begin{lemma}\label{lem:eight-support}
In the non-pointed based-simple case a largest degree $x$ of codegree
eight is even and occurs twice. Its two rows are conjugate over
$\QQ(\sqrt2)$, supported on these two labels and the four odd labels,
all of dyadic twist order. The three nonunit odd degrees are distinct
divisors of $x/2$, pairwise noncoprime and of gcd one. In particular
at least three odd primes divide $x$.
\end{lemma}
\begin{proof}
Here $D=8x^2$. A maximal row $u$ is principal at two, sees only
degrees dividing $x$, and meets every odd label
(Lemmas~\ref{blk:support}, \ref{blk:height}). The nonunit odd degrees
therefore share a prime in pairs. The number of odd labels is
divisible by four, by the dimension formula modulo four.

If $x$ is odd, the support is exactly the odd labels. Gcd one on its
nonunit degrees requires three distinct proper divisors after removing
all maximal labels, and $x\ge105$. The support has between five and
eight labels, hence eight; every entry there is an eighth root.
The multiplicity is $m\le2$ by Proposition~\ref{sq:saturation}.
At $m=1$, six proper degrees including three distinct ones give
$\sqrt8\le1/x+1+4/3+1/5+1/7$ by \eqref{st}. At $m=2$ the sum at
the two maximal labels is an even integer, necessarily zero since
the five proper degrees total at most $5x/3$. If the rows are conjugate,
their quadratic field is generated by nonreal roots of unity, so they
are dual and of equal twist. Cancellation of the two maximal terms
then gives $\sqrt8\,x\le1+5x/3$. Otherwise the rows are rational,
hence signs. Their symmetric maximal submatrix has zero row sums:
the rows disagree on both maximal columns and, by orthogonality, on
two proper columns. The disagreeing proper entries have opposite
signs within each row and equal degree $y$ by vacuum orthogonality.
The remaining three proper degrees have signed sum $-1$ and hence
are distinct by pairwise noncoprimality. For the half-sum $A$ and
half-difference $B$ of the two vacuum modular relations,
$|A|\le1+(1/3+1/5+1/7)x<7x/10$ and $|B|\le2x+2y\le8x/3$,
contrary to $|A|^2+|B|^2=8x^2$.

Now $x$ is even and $x\ge30$. The even-degree calculation in
\cite[proof of Theorem~20.7]{BDPcore} applies to this row. It uses
principality, $(x/d_Y)u_Y\equiv0$ above two, and the norm and vacuum
identities \eqref{st}. It excludes a rational row and gives
$u_Y=A_Y+B_Y\sqrt2$ with $A_Y,B_Y\in\ZZ$. Off the vacuum there are
three nonzero $A_Y$, all signs, and two opposite nonzero $B_Y$, both
signs, at a common degree $y$. The $A_Y$ identify the nonunit odd
degrees by principality; the congruence gives $d_Y\mid x/2$ there.
Gcd one and pairwise noncoprimality make these degrees distinct and
require at least three odd primes in $x$. Consequently $y$ is even.
If $y<x$, then $y\le x/2$ and the vacuum relation gives
$\sqrt8\,x\le1+(1+1/2+1/3)x/2+\sqrt2\,x$. Hence $y=x$.
All maximal rows have the same signs on the odd labels: two distinct
solutions of $1\pm a\pm b\pm c=0$ would yield a vanishing proper
subsum, impossible for four distinct positive weights with signed
total zero. Bessel on those four columns gives $4m\le8$. The dyadic
twists follow from Lemma~\ref{blk:admissible}.
\end{proof}

In particular, when $D=2M^2$, every odd prime denominator $p$ occurs
with $p\Vert M$. Indeed the dyadic defect-one block there has equal
degrees by Lemma~\ref{blk:defectone}, so Theorem~\ref{rev:height} applies.

\begin{lemma}[The norm-nine reduction with three odd primes]
\label{ext:eight-nine}
In the non-pointed based-simple rotation-admissible setting, suppose
the largest degree has codegree eight and exactly three odd primes
divide the dimension. Write $D=2M^2$.  If denominator three
occurs, then $3\Vert M$, its multiplicity is two, and its unique
effective row has nine nonzero entries, all signs.  Besides the
unit, the unique three-free odd nonunit $U$, and a neutral
denominator-three label, its support consists of six even proper
nonunits.  They are dyadic-neutral and inactive at the other two
odd primes $l<t$.  Four have denominator at least six; the other
two have denominators at least $3l$ and $3t$.
\end{lemma}
\begin{proof}
Proposition~\ref{new:small-effective} gives $3\Vert M$.
Put $n=M/3$ and let $d=d_U\le n/4$.  Write the multiplicity of
denominator three as $2k$, using the dyadic defect-one pairs.
Their effective rows have norm nine and common support on the
dyadic-neutral three-free labels.  They lie in $P_3$, whose degrees
divide $n$, and all their entries there are nonzero by principality.
Orthogonality to the mean of the two maximal rows fixes their entries
at $0,U$ as $1,-s_U$, where $s_U$ is the maximal mean's sign.
Their nonunit hit degrees have gcd one.  Besides $U$ and the $k$
neutral degree-$n$ labels, there are distinct even hits $A,B$ omitting
$l,t$, of degrees at most $n/l,n/t$: $d$ contains both $l,t$, and
no nonunit degree is a prime power.  All hits are inactive at $l,t$.

Bessel on $0,U$ gives $k\le4$.  The sum of the rows is integral,
since Galois permutes their blocks.  At $k=4$ it has squared norm
four off these two coordinates, but is nonzero modulo three at the
four neutral degree-$n$ labels and at $A,B$.  At $k=3$ its remaining
squared norm is nine and its entries are multiples of three.  Its
sole nonzero value there is $\pm3$, at degree $|1-s_Ud|$ by degree
orthogonality.  This is coprime to $d$, contrary to pairwise
noncoprimality in the principal three-block.

At $k=2$, if the two rows are rational, their difference is divisible
by three and has squared norm eighteen.  Its two nonzero entries are
$3,-3$ at equal degrees, by degree orthogonality.  Equality of the
row norms makes the entries there exchanged, costing at least five
in each row.  There are at least six hits, so precisely six occur,
with exchanged entries $1,-2$ up to sign.  Only the two degree-$n$
hits have equal degree.  Degree orthogonality gives
\[
 n\le1+d+n/l+n/t<n,
\]
a contradiction.  Otherwise the rows are conjugate in the unique
quadratic subfield of $\mathbb Q(\zeta_{18})$, hence dual.  Their
sum $z$ and $(v-\overline v)/\sqrt{-3}$ are integral, agree modulo
two, and have squared norms eighteen and six.  The latter has
opposite entries on dual pairs, hence exactly six signs.  At $0,U$
the entries of $z$ have magnitude two.  At the two neutral degree-$n$
labels they coincide and are nonzero modulo three.  Both have
magnitude one, since magnitude two would leave norm two for six
odd entries.  Off these four coordinates, $z$ has squared norm eight
with four odd entries, hence absolute sum six.  The two effective
Gauss relations have identical right sides, by duality and
Proposition~\ref{new:effective}, each of modulus $3n$.  All other
hit degrees divide $n$ properly, so
\[
 6n\le2+2d+2n+6n/2<6n,
\]
again a contradiction.

Thus $k=1$ and the row is rational.  It has at least five hits.
A double entry would force six hits.  At degree $n$, degree
orthogonality would imply
$2n\le1+d+n/l+n/t+n/2<2n$; elsewhere its Gauss relation would
imply $3n\le1+d+n+n/l+n/t+n<3n$.
All nine entries are therefore signs.  After $0,U$ and the neutral
degree-$n$ label, six even proper hits remain.  Their denominators
are at least six, with the two omitting $l,t$ at least $3l,3t$.
\end{proof}

\begin{lemma}[Odd prime denominators with four dimension primes]
\label{ext:eight-common-fiber}
In non-pointed based-simple rotation-admissible data with $D=2M^2$
and exactly four dimension primes, at most one odd prime occurs
as a denominator.
\end{lemma}
\begin{proof}
An odd prime denominator $p$ has formal codegree $2p^2$.
By Lemma~\ref{blk:defectone}, its full two-block is an equal-degree
pair, of common degree $M/p$. Theorem~\ref{rev:height} gives
$p\nmid M/p$, hence $v_p(D)=2$. Two such primes would force $D$
to be a square by Proposition~\ref{new:two-square-primes},
a contradiction.
\end{proof}

\begin{corollary}[Codegree eight through rank forty-five]\label{round:eight-rank}
In the non-pointed based-simple rotation-admissible case, a largest
degree of formal codegree eight requires rank at least $46$.
\end{corollary}
\begin{proof}
Suppose $r\le45$. Write $D=2M^2$. By Lemma~\ref{lem:eight-support}
there are two labels of denominator two and four odd-degree labels,
including the unit. All six have dyadic twists. For every odd prime $s$
reserve $2(s-1)$ further labels active at $s$. The union bounds of
Proposition~\ref{blk:packets} permit these reservations simultaneously
by Hall's theorem. At least three odd primes occur. Four would cost
at least $44$ places besides the six, and even the triple $5,7,11$
would cost $40$. Thus the odd primes are $3<l<t$.

The three odd nonunits have pairwise noncoprime degrees with gcd one,
so they omit different primes. Their degrees divide $M/4$, giving
denominator floors $12,4l,4t$ and vacuum floor $12lt$. Write $2k$ for
the multiplicity of denominator three. Lemma~\ref{ext:eight-nine}
gives $k\le1$. Among the four ternary reservations and the $2k$ threes,
one may therefore charge $2k$ places at three and $4-2k$ at six.
For $s=l,t$ the $2(s-1)$ reservations have floor $s$ if denominator
$s$ occurs, and $2s$ otherwise. Fill the remaining places to $45$ at
floor four. By Lemma~\ref{ext:eight-common-fiber}, at most one odd
prime denominator occurs. If $k=1$, both $l,t$ are absent. If $k=0$,
the largest bound allows $l$ to occur, since $(s-1)/s^2$ decreases
for $s>2$. In both cases every prime-dependent mass term decreases
as $l$ or $t$ increases, as does the number of remaining places at
floor four. Thus the maxima occur at $(l,t)=(5,7)$ and are
$7607/3920$ for $k=0$ and $109793/58800$ for $k=1$.
Both are below two, contrary to \eqref{eq:recip}.
\end{proof}

\begin{proposition}[A rational label detected by dyadic row mass]\label{new:dyadic-mass-witness}
Suppose $D$ is even in non-pointed based-simple rotation-admissible data. Let
$w=\sum_X1/f_X$ over the nonunit labels in the principal two-block
whose character fields have degree a power of two. Some nonunit
$Y$ of odd degree has an odd Galois orbit of size $h$, and every such
orbit satisfies $h(w+d_Y^2/D)\le1$. In particular $w\ge1/3$
forces a rational nonunit row of odd degree. If a nonunit row of
two-power codegree occurs, $Y$ can be chosen with twist $\pm1$.
Six denominator-four labels in square dimension suffice for both conclusions.
\end{proposition}
\begin{proof}
The dimension identity modulo two gives an odd number of nonunit labels
of odd degree. These labels form full Galois orbits. On an odd orbit
$O$ of size $h$ the value $\chi_X(Y)$ belongs to $K_X\cap K_Y=\QQ$,
by symmetry and $[K_Y:\QQ]=h$. It is constant on $O$ and an odd integer
by principality. The image under $s$ of $h^{-1/2}\sum_{Y\in O}e_Y$
has squared mass at least $hw$ on the selected rows and $h d_Y^2/D$
at the vacuum. Thus $hw<1$.

In the presence of a two-power row every odd-degree label has dyadic
twist by Corollary~\ref{blk:primepower}. Since $D$ is square-full,
$2\mid M$. The vacuum Gauss sum $M\sqrt q\,\xi$ is divisible by two.
Thus $\sum_{d_i\ {\rm odd}}\theta_i$ belongs to twice the integers
of a dyadic cyclotomic field. Its constant coefficient in a cyclotomic
power basis is the difference between the counts at $1$ and $-1$,
so their sum is even. Removing the vacuum leaves a Galois-stable
set of odd size at these two twists. Finally each denominator-four
row in square dimension contributes $1/16$ by Lemma~\ref{conduct}
and Corollary~\ref{blk:primepower}.
\end{proof}

\subsection{Large classes in square dimension}
\begin{lemma}[Support of denominator-four rows]\label{sq:heights}
In the non-pointed simple case with $D=M^2$, a denominator-four row
of even degree, or of degree occurring at least four times, is
supported on $4\mid b_Y$. Its nonunit hit denominators have lcm $M$;
for each $\ell^e\Vert M$, one is divisible by $\lcm(\ell^e,4)$.
If exactly three primes divide $M$, four denominators four force
$8\mid M$.
\end{lemma}
\begin{proof}
The support assertion is \cite[\S21.2 and Lemma~21.6]{BDPcore}; its
column-norm argument uses lcm integrality, principality above two
and the unit column norm. The nonunit hit degrees have gcd one by
\eqref{eq:foundorth}, giving the lcm assertions.
If $4\Vert M$, every hit degree is odd. With exactly three dimension
primes, each nonunit hit contains both odd primes by
Lemma~\ref{lem:basic-table}, contradicting gcd one.
\end{proof}

\subsection{Orbit projections and prime denominators}
\begin{proposition}[Neutral inertia orbits]\label{round:neutral-inertia}
Let $p$ be an odd prime of even positive valuation in $D$.
Suppose that every $p$-primary twist has order at most $p$.
Let $A$ count the $p$-active labels, and let $c$ be the number
of $p$-neutral labels minus the number of tame $p$-inertia orbits
among them. Then
\[
                 A\ \ge\ (p-1)(1+c/2).
\]
In particular at $A=p-1$ all neutral labels are individually fixed
by $p$-inertia.
\end{proposition}
\begin{proof}
The local group is $G=\mathrm{SL}_2(\mathbb F_p)$, with upper
unipotent $U$, diagonal torus $H$ and Borel $B=UH$.
On the coordinate space $V^U$ of the neutral labels, $H$ acts
by the inertia permutations with positive sign. Thus
$c=\dim V^U-\dim V^B$. Some copy of the Steinberg representation
occurs by the vacuum argument of Theorem~\ref{thm:spherical};
it contributes $p-1$ active dimensions and no nontrivial
$H$-character in $V^U$.

For $I_\alpha=\operatorname{Ind}_B^G\alpha$, the two Bruhat cells
give intertwining multiplicities
$\mathbf1_{\alpha=\beta}+\mathbf1_{\alpha=\beta^{-1}}$
\cite[Proposition~2.14]{PrasadSL}. The published decomposition
\cite[Lemma~4.3]{Nevins} makes $I_\alpha$ irreducible when
$\alpha\ne\alpha^{-1}$ and splits the nontrivial quadratic case
into two distinct constituents of dimension $(p+1)/2$.
The fixed point and free orbit on $G/B$ give
$I_\alpha|_U=\mathbf1\oplus\mathrm{Reg}_U$
\cite[\S15.9]{DigneMichel}. Thus the first case has two nontrivial
torus lines and $p-1$ active dimensions. In the quadratic case
reciprocity gives one such line in each constituent; their
$U$-fixed dimensions sum to two, so each has precisely one and
$(p-1)/2$ active dimensions. Reciprocity exhausts the nontrivial
torus types, none of which belongs to Steinberg. Summing gives
the bound.
\end{proof}

\begin{corollary}[Small prime packets]\label{new:five-reduction}
Under the hypotheses of Theorem~\ref{new:one-packet},
at denominator five there is one label per primitive phase, and four
proper five-free labels with Weyl coefficients of magnitude one. Their indices $k_i=n/d_i$ have lcm $n$, with $k_i\equiv a_i\pmod5$ and $\sum a_i/k_i=\epsilon/n$.
At denominator seven there is one or two per primitive phase. If
there are two, the proper seven-free labels consist of one inertia
pair with coefficients of magnitude one and two singleton labels
with coefficients of magnitude two, at three distinct degrees.
In odd dimension neither prime denominator occurs.
\end{corollary}
\begin{proof}
The neutral orbit at degree $n$ spends $j-1$ units of inertia
deficit, leaving at most $j-1$ on all other neutral labels.
Each proper orbit of size $h$ costs $h a_Y^2$ of the norm
$j(p-j)$ and satisfies $j\mid h a_Y$. At least three distinct
proper degrees are needed.

At five, $j=1$ or $2$. At $j=2$ there is room for at most one
moved pair. Since the norm is six, at most one further label can
have a nonzero even coefficient, whereas every fixed label has an
even coefficient. This gives at most two degrees. Thus $j=1$,
and norm four on at least three degrees gives four signs.

At seven the possible $j$ are $1,2,3$. At $j=3$ the remaining
deficit is two. A moved orbit of size two has coefficient divisible
by three and costs at least eighteen, exceeding the norm twelve.
An orbit of size three uses the entire remaining deficit and needs
two fixed labels to supply three degrees, costing at least
$3+2\cdot9>12$. With only fixed labels the contradiction is stronger.
At $j=2$ the norm ten cannot be spent on fixed labels alone, since
their coefficients are even. There is just one moved pair and the
three-degree condition requires two other labels. Over the norm
ten their absolute coefficients are forced to be $1,1,2,2$,
and the last two degrees are distinct singletons. In odd dimension
all nonunits have distinct duals of equal degree, excluding these
singletons and the sole neutral row at $j=1$.
\end{proof}


\begin{proposition}[The idempotent of a prime packet]\label{round:packet-idempotent}
In the non-pointed based-simple square case, let $p$ be an isolated odd prime with exactly $l(p-1)$ active labels of a common denominator $B$. If $B\ne p$, assume $l=1$ or that $D$ is odd and $l=2$, and suppose that the number of inactive nonunit pure $p$-power labels is less than
\[
 l p^2\bigl(1-1/(p-1)-1/(6p)\bigr).
\]
The denominator-$B$ class $I$ consists of $lp$ labels, with one neutral packet of size $l$. For $B=p$ this packet is a single Galois orbit. Put $M_p=B^2/p$ and let $w=W_pv_0$. For $H$ the nonunit labels outside $I$ with $p\mid b_Y$ one has
\[
 M_p e_I=l\,1-\sum_{Y\in H} M_p w_Y x_Y, \qquad M_pw_Y\in\ZZ.
\]
Every packet in $I$ has row sum $-M_pw_Y$ at $Y\in H$. At a member of the neutral packet the active packet sums are $-\epsilon M_p/B$ and the neutral sum is $\epsilon(p-1)M_p/B$. When $B=p$ this gives $p e_I=l\,1-\epsilon\sum_H a_Yx_Y$.
\end{proposition}
\begin{proof}
For $B=p$ the class size follows from Theorem~\ref{new:one-packet}. Otherwise a pure denominator smaller than $B$ alone exceeds the remaining Weyl norm $l(p+1)/B^2$, so denominator $p$ is absent and Proposition~\ref{new:pure-extra} gives $M_p=B^2/p$. All pure denominators are at least $B$. Without a neutral packet at $B$, each inactive pure label contributes at most $1/p^2$ to $E$, contradicting the lower bound following \eqref{eq:packet-shell}. The same equation gives $E/l\le1+1/(p-1)+1/(6p)<2$, proving uniqueness.

For $B=p$, Theorem~\ref{new:one-packet} puts each active packet sum in the vacuum Steinberg constituent. The neutral packet sum is Borel-fixed, so Steinberg uniqueness and Frobenius reciprocity place it in the sum of the trivial constituents and this vacuum copy. Work on that sum. For $B\ne p$, use the full or charge-even representation at minimal multiplicity. Its primitive phases exhaust the active packet lines and their vacuum Steinberg has trivial complement there. Write $u=\sum_Iv_X$, $A=B\epsilon\sum_H w_Yv_Y$ and $Q=lpv_0/B-\epsilon A$. The vacuum column gives $P=B\epsilon(W_p-I)v_0=u-\epsilon Q$ and $W_pP=-P$. Its neutral projection $g$ spans the neutral Steinberg line. It has $g_0=-l\epsilon p/B$, coefficient one on each neutral member of $I$, and coefficients $A_Y$ on $H$. Since $\|P\|^2=2lp$ and the active coordinates contribute $l(p-1)$, $\|g\|^2=l(p+1)$. Hence $Bv_0-Q$ is perpendicular to $g$ and fixed by $W_p$. It follows that $W_pQ=Q+\epsilon P$, and $W_pu=\epsilon Q$. The complementary Weyl product is $\epsilon$ on $I$ by local admissibility and Bruhat. Thus $su=Q$; expansion in the fusion basis gives the idempotent formula. Integrality follows from $u_Y\mid M_p$ in Proposition~\ref{new:component-column}. For an active packet $A_0$ the neutral projection of $s\sum_{A_0}v_X$ is proportional to $g$ and its vacuum entry is $l/B$. This gives the active row sums; subtract them from $su=Q$ for the neutral packet. At $l=1$ its singleton row is Galois-fixed and rational.
\end{proof}

\begin{lemma}[Two coprime denominator classes]\label{new:cross-weyl}
In the non-pointed based-simple rotation-admissible case, suppose $D=M^2$, with $p\Vert M$ odd, and denominator $p$ occurs. With $a_Y,j,\epsilon$ as in Lemma~\ref{new:weyl}, for any label $X$ of denominator $B>1$ prime to $p$,
\[
         \sum_{Y\text{ proper }p\text{-free}}a_Y\chi_X(Y)=\epsilon j.
\]
Here proper excludes the vacuum and denominator $p$.
\end{lemma}
\begin{proof}
In square dimension $s_{XY}\gcd(b_X,b_Y)$ is an algebraic integer by lcm integrality. If $\gcd(b_X,b_Y)=1$ and both labels are nonunit, this integer vanishes: all its Galois conjugates have modulus less than one by unitarity (the vacuum column never vanishes). Since $p\nmid f_X$, Lemma~\ref{blk:admissible} makes the two local generating unipotents act on $e_X$ by the same scalar, the vacuum eigenvalue; it is one since $p$ is odd and $\gamma$ is dyadic. Hence $W_p e_X=e_X$. As $W_p$ commutes with $s$, $\langle W_pe_0,se_X\rangle=1/B$. Insert $(W_p)_{00}=1-j/p$ and $(W_pe_0)_Y=a_Y/(p\epsilon)$, then multiply by $pB$. The denominator-$p$ terms vanish.
\end{proof}

More generally, let $I$ consist of $m$ rows of common codegree $f$ with denominators prime to $p$. Among the proper $p$-free labels take any set $Z$ containing all those hit by the rows in $I$. Then
\[
                \sum_{i\in Z}a_i^2\ge\frac{mj^2}{f-m}.
\]
Indeed, write $P$ for the orthogonal projection off the normalized rows in $I$, $v=\sum_{i\in Z}a_ie_i$, and $c=\epsilon j$. By Lemma~\ref{new:cross-weyl}, the projections of $v$ and $ce_0$ onto the rows agree. Thus $\|Pe_0\|^2=1-m/f$, $\langle Pe_0,Pv\rangle=-cm/f$ and $\|Pv\|^2=\|v\|^2-c^2m/f$. Cauchy--Schwarz gives the inequality. Equality requires $Pv=-cmPe_0/(f-m)$; taking the scalar product with the degree vector gives $\sum_{i\in Z}a_id_i=-cm/(f-m)$.

For twelve four-rows and five denominator-five labels take $Z$ to be
the proper five-free labels with $4\mid b$. At least three of the four
signs occur in $Z$. Exactly three would force the omitted signed degree
to be $4\epsilon$ by the equality case and \eqref{eq:weyl-cost}.
A nonunit degree has two prime factors, so all four satisfy $4\mid b$.

\begin{corollary}\label{new:seven-shape}
In the non-pointed based-simple rotation-admissible square case, suppose denominator seven occurs with $7\Vert M$ and denominator four has multiplicity twelve. Then there are seven denominator-seven labels and six proper seven-free nonunits. Their denominators are $7k$ with $k\equiv\pm1\pmod7$ and $k\ge6$.
The same conclusion about the six proper labels holds for eleven fours
if the seven-class has size seven.
\end{corollary}
\begin{proof}
The mass of the fours leaves room for at most twelve sevens. Theorem~\ref{new:one-packet} gives $7j=m_7\le12$, hence $j=1$ and $m_7=7$. The squared norm of the proper coefficients is six, on at least three distinct degrees, allowing six signs or the magnitudes $2,1,1$.

In the latter case all three columns are singleton degree classes, hence rational. Label the columns of coefficient magnitude one $A,B$. Since $\sum a_Yd_Y=\epsilon$, exactly one, say $A$, has odd degree. For $m=11$ or $12$ four-rows the entries at $A$ are odd integers by
principality. Their total square is strictly below sixteen by the
vacuum contribution, so all are signs. If their sum is $w$, the
Cauchy argument of Proposition~\ref{sq:saturation} gives
$|w|<16-m$. Since $w\equiv m\pmod2$, each sign occurs at least
$m-7\ge4$ times. At $B$ the entries are even integers, so at most
three are nonzero. In Lemma~\ref{new:cross-weyl} the remaining term
$\pm2\chi_X(C)$ has fixed residue modulo four by dyadic principality.
Consequently one of the signs at $A$ forces a nonzero entry at $B$
on at least four rows, a contradiction.

All six coefficients are signs, giving the congruence and lower bound on $k$.
\end{proof}

\begin{proposition}[The five- and seven-packets together]\label{prop:joint-five-seven}
In the non-pointed based-simple rotation-admissible case, suppose the dimension is square with four prime divisors $2,t,5,7$. Assume denominator five occurs with $5\Vert M$, seven occurs in one packet with $7\Vert M$, and at least eight labels have denominator four. Then exactly six proper nonunits have seven-free degree. Their denominators $7k$ satisfy $k\equiv\pm1\pmod7$ and $k\ge6$. Exactly one is also five-free, so the two packets require nine distinct proper neutral labels.
\end{proposition}
\begin{proof}
Write $H_p$ for the proper $p$-free labels and $a_{p,Y}$ for their Weyl coefficients. At five the four coefficients are signs; at seven the squared norm is six on at least three distinct degrees. By orthogonality of the component columns,
$\sum_{H_5\cap H_7}a_{5,Y}a_{7,Y}=-1$. The five-free witnesses omitting two and $t$ are distinct and outside $H_7$, as every nonunit degree has two prime divisors. Thus the intersection has at most two labels; when the seven coefficients are all signs its odd cardinality gives precisely one.

It remains to exclude coefficients $2,1,1$ in absolute value. The three proper seven-free labels then omit $2,t,5$, respectively, and each has singleton degree class. The common label $C$ omits five, so its coefficient is a sign by the cross identity. Let $Y$ carry the double. Its degree is even, since eight odd integer four-row entries and the seven-packet column mass $4/7$ would exceed the column norm. Thus the other sign lies at the odd label $A$. The five-packet has a unique odd proper label $U$: its signed degree sum is odd and the omissions of $7,t$ require two different even labels. Hence $U$ is also a singleton.

Select eight four-rows. For $V=A,Y,C$ orient its column by $\epsilon\operatorname{sgn}(a_{7,V})$, and orient $U$ using five. Write $\widetilde d_V$ for its signed degree. By Proposition~\ref{round:packet-idempotent} the oriented columns on the seven-rows at $A,Y,C$ have respective values $-1,-2,-1$, and $U$ vanishes there. On the five-rows $A,Y$ vanish and $U$ has value $-1$. All entries on the selected eight rows are integers by singleton rationality. At $A,U$ they are signs by principality and column norm; at $Y$ they are even and at most one can be nonzero, of magnitude two.

Write $r_i=(1,a_i,y_i,u_i)$ for the eight rows on columns
$0,A,Y,U$, ordering them so that $y_2=\cdots=y_8=0$. On these
columns the seven- and five-packets have rows $v=(1,-1,-2,0)$
and $w=(1,0,0,-1)$ respectively. Subtracting their contributions to the
column Gram matrix gives
\[
 G=560I-80v^tv-112w^tw-35\sum_{i=1}^8r_i^tr_i\succeq0.
\]
There are $12\binom{10}3=1440$ choices of $r_1$ and of the multiset of
the seven sign pairs $(a_i,u_i)$ for $i>1$. The fifteen principal
minors are nonnegative only for $r_1=(1,-1,2,-1)$ and multiplicities
$(0,1,2,4)$ in lexicographic sign order
(\supfile{arithmetic/gram_packets.py}{Gram calculation}).
Then $G(6,-1,-7,1)^t=0$. The vacuum row remains, giving
$6-\widetilde d_A-7\widetilde d_Y+\widetilde d_U=0$. The seven-packet
identity gives $\widetilde d_A+2\widetilde d_Y+\widetilde d_C=1$.
Eliminating $Y$ gives $5+5\widetilde d_A+7\widetilde d_C+2\widetilde d_U=0$,
whereas $t\ne5$ divides all three degrees. The index condition follows
from Theorem~\ref{new:one-packet}.
\end{proof}

\begin{lemma}[The Gauss sum on complete degree classes]\label{new:partial-gauss}
In square dimension put $T_b=\sum_{i:b_i=b}\theta_i$.
Remove complete degree classes of zero twist sum, then select further
classes indexed by a set $J$. Let $L$ be the lcm of their denominators
and $R$ the reciprocal-square mass outside all the removed and selected
classes. If $L^2(R+1/M)<1$, then
$\sum_{b\in J}(L/b)^2T_b=0$.
\end{lemma}
\begin{proof}
Every conjugate of this algebraic integer has modulus at most
$L^2(R+1/M)$ by the vacuum Gauss formula. A nonzero algebraic
integer has a conjugate of modulus at least one.
\end{proof}
For complete denominator classes in square dimension, write $T_b=\sum_{b_X=b}\theta_X$ and let $R$ be the reciprocal-square mass outside the classes displayed, after deleting known zero-sum classes. Lemma~\ref{blk:admissible} supplies the prime support of the twists. When $m_4$ is even, the following tests follow from Lemmas~\ref{r7:coupled-gap} and~\ref{r7:gauss-projection}:
\[
 \begin{array}{lll}
 m_6\le3,&36(R+1/M)<1,&T_4=T_6=0,\\
 m_6\le8,&72(R+1/M)<1,&T_4=T_6=0,\\
 m_6\le5,\ m_8\in\{0,2\},&96(R+1/M)<1,&T_4=0.
 \end{array}
\]
Here $R$ omits the classes $4,6$ in the first two tests and $4,6,8$ in the third; the vacuum Gauss relation is multiplied by $144$ and $576$.

\begin{corollary}\label{new:small-Gauss}
In the non-pointed based-simple setting of Theorem~\ref{alg:rank},
suppose $D=M^2$, $5\Vert M$, and there are five denominator-five labels.
With eleven fours there are at most three sixes. With twelve fours,
the four proper five-free nonunit denominators are divisible by
twenty and at least eighty, denominator six is absent, and there
are at most two eights. If there are two, their twist sum is zero
and denominator nine cannot occur exactly once.
\end{corollary}
\begin{proof}
The five twists are all the fifth roots by
Lemma~\ref{new:weyl}, hence $T_5=0$. Eleven fours and four sixes
would leave mass $1/720$, with $M\ge420$ since at least four primes
divide $M$. The lemma gives $9T_4+4T_6=0$, while at a place above
two $T_4\equiv11$ by Lemma~\ref{blk:admissible}.

For twelve fours the complementary mass is $1/20$. Since
$16(1/20+1/420)<1$, the preceding lemma gives $T_4=0$.
A single six would now give $1/M\ge2/36-1/20=1/180$ by the
triangle inequality. On the four proper five-free labels the
projection after Lemma~\ref{new:cross-weyl} gives $4\mid k_i$.
Since $k_i\equiv a_i\pmod5$, the only possibility below sixteen is
$k_i=4$ with sign $-1$. At most three positive terms, each at most
$1/16$, cannot offset $-1/4$ to $\pm1/n$ since $n>16$.
Thus every $5k_i$ is at least eighty.

Finally three eights leave mass $1/320$ beyond the fours and fives.
The lemma would give $T_8=0$, contrary to $T_8\equiv3$ above two.

With two eights, the remaining mass is $3/160$. Every conjugate of
$T_8$ has modulus at most $64(3/160+1/420)=142/105<\sqrt2$.
Its twists have two-power order by Lemma~\ref{blk:admissible}, so
Lemma~\ref{r6:dyadic-mean} forces $T_8=0$. A singleton nine would
then give $2/81\le3/160+1/M$ by the triangle inequality, impossible
for $M\ge420$.
\end{proof}

\begin{proposition}[Weighted prime-free row bound]\label{sq:joint-projection}
Select disjoint classes of $s_q$ nonunit rows of codegree $f_q$ a
power of a prime $q$, and a prime $p$ different from each $q$. Let
$U$ be a Galois-stable set of $p$-free nonunits containing every
such column of nonzero absolute mean in a selected row. Choose
positive integer orbit-constant weights $w_Y$ prime to $p$ with
$q\mid w_Yd_Y$ for each selected $q$. For $1\le a<p$ put
$c_Y=w_Y[ad_Y/w_Y]_p$ on $U$ and zero elsewhere, using least
absolute residues (the representative $1$ at $p=2$).
Let $0<u_q<pq$ satisfy $u_q\equiv-a\pmod p$, $q\mid u_q$.
Then for real $t$,
\[
 \|c\|^2\ge\sum_q\frac{s_q}{f_q}
       \operatorname{dist}(t+u_q,pq\ZZ)^2-t^2 .
\]
For one class write $q,f,s,u$ for these parameters, $l=u-pq$,
$h=\max(1,(p-1)/2)$ and
\[
 B_a=\frac1f\min_{0\le j\le s,\ j\in\ZZ}
  \left((s-j)l^2+ju^2+\frac{((s-j)l+ju)^2}{f-s}\right).
\]
Then
\[
 \sum_{Y\in U}w_Y^2\ge
 H_f(p;q,s):=\left\lceil
 \frac{\sum_{a=1}^h\lceil B_a\rceil}{\sum_{a=1}^h a^2}
 \right\rceil .
\]
If the selected rows and $U$ both consist of dual pairs without
fixed points, restrict $j$ to even integers and round both ceilings
up to even integers.

For one class of actual character rows there is a stronger bound
when $p\mid D$ is odd, $U$ is the set of all $p$-free nonunits, and
$w_Y=1$ (so $q\mid d_Y$ on $U$). For $1\le a\le(p-1)/2$, let
$\mathcal L_{p,a}$ be the continuous convex interpolation of $t^2$
on $-a+p\ZZ$; on $[v,v+p]$, $v\in-a+p\ZZ$, it equals
$(2v+p)t-v(v+p)$. With $l,u$ as above put
$S_j=(s-j)l+ju$ and $Q_j=(s-j)l^2+ju^2$. Then
\begin{equation}\label{eq:full-census}
 \sum_{Y\in U}[ad_Y]_p^2\ge
 \min_{0\le j\le s,\ j\in\ZZ}
 \frac{Q_j+(f-s)\mathcal L_{p,a}(-S_j/(f-s))}{f}.
\end{equation}
For paired rows and labels, restrict $j$ to even integers and round
the resulting norm bound up to an even integer.
\end{proposition}
\begin{proof}
The integers $z_X=\sum_Yc_Y\chi_X(Y)$ are rational by Galois
symmetry. Orthogonality to the vacuum gives $z_X\equiv-a\pmod p$:
omitted $p$-free orbits have row sum zero. Principality at $q$
(Corollary~\ref{blk:primepower}) and $q\mid w_Yd_Y$ give
$z_X\equiv0\pmod q$. Apply \eqref{eq:projection-master} to $c+te_0$. For one class, maximizing the resulting quadratic in $t$
gives $f\|c\|^2\ge\sum_X z_X^2+(\sum_Xz_X)^2/(f-s)$;
indeed $s<f$ by Bessel at $e_0$, including the vacuum row.
At fixed sum, bringing two residue levels closer decreases the
first term; if all levels have one sign, moving them towards zero
decreases both. The minimum therefore uses only $l,u$.
Each $Y$ contributes $w_Y^2\sum_a a^2$ after summing over $a$.
In the paired case the real row sums and the squared entries of $c$
occur in pairs, giving the indicated rounding.

For the stronger bound define $z_X=\sum_{Y\in U}[ad_Y]_p\chi_X(Y)$
on every character row. Galois invariance makes these integers.
Orthogonality gives $z_X\equiv-a\pmod p$ on nonunit rows, and
$p\mid D$ gives the same congruence at the vacuum. With weights
$1/f_X$ their mean is zero and their second moment is $\|c\|^2$.
The selected values are multiples of $q$ by principality. At fixed
sum $S$, replacing two of them by nearer lifts decreases their
squared sum, so a minimum uses consecutive lifts in $l+pq\ZZ$.
The remaining spectral mass $(f-s)/f$ costs at least
$(f-s)\mathcal L_{p,a}(-S/(f-s))/f$ by convexity.

Only the lifts $l,u$ are needed. Below $l$, raising an entry lowers
its square and the positive complementary mean, where
$\mathcal L_{p,a}$ is nondecreasing. Above $u$, lowering by $pq$
decreases its square by at least $pq(2u+pq)$; the complementary
mean stays negative, where the interpolant has slopes at most
$p-2a$, so its cost increases by at most $pq(p-2a)$.
This proves \eqref{eq:full-census}. In the paired case
$z_{X^*}=z_X$; apply the same exchanges to pairs. Both costs double,
and $j$ and $\|c\|^2$ are even.
\end{proof}
\begin{lemma}[Two pairs of prime-free degrees]\label{odd:two-free}
In the non-pointed based-simple case of odd dimension, suppose a
nonunit row has $p$-power codegree, with $p=5$ or $7$.
At least six nonunit labels have degree prime to $p$. If instead
there are exactly four nonunit labels of degree prime to eleven
and a row of three-power codegree, one of their degrees is prime
to three, provided eleven divides the dimension.
\end{lemma}
\begin{proof}
For $p\mid D$ and $0<a<p$, put $c_Y=[ad_Y]_p$ on all $p$-free
nonunits and zero elsewhere. As in the full-spectrum proof of
Proposition~\ref{sq:joint-projection}, its integral spectrum has mean
zero and second moment $\sum_Yc_Y^2$. Averaging
$(z+a)(z+a-p)\ge0$ gives $\sum_Yc_Y^2\ge a(p-a)$, with equality
only when every eigenvalue lies in $\{-a,p-a\}$.

At $p=5,7$, the dimension congruence excludes zero or one $p$-free
dual pair. With two pairs of degrees $A,B$, their absolute residues
are respectively $(1,1)$ and $(1,3)$. Taking $a=1$ and $2$ gives
squared norms $4$ and $10$, attaining the bound. The even vacuum
eigenvalue is therefore $4$ and $-2$, respectively, giving
$\pm A\pm B=2$ and $\pm2A\pm B=-1$. Since $A,B$ are odd,
both force $\gcd(A,B)=1$, contrary to
Corollary~\ref{blk:primepower} and Lemma~\ref{blk:witness}.

At eleven two pairs have absolute residues $1,2$, and $a=1$ gives
second moment ten. Neither possible eigenvalue $-1,10$ is
divisible by three. On a row of three-power codegree the spectrum
would reduce to zero above three if both degrees were multiples of
three, by principality.
\end{proof}

In fact, if six nonunit labels have degree prime to seven, the mass
of the union of all principal $t$-blocks with $t\ne7$ is at most
$3/7$, where the union ranges over primes dividing the global
dimension. The three dual pairs have absolute degree residues
$(1,1,1)$ or $(2,2,3)$ modulo seven: these are the solutions of
$1+2\sum r_i^2=0\pmod7$ with $1\le r_1\le r_2\le r_3\le3$.
Take $a=1$ or $3$, respectively. In each case
$2\sum_i[ar_i]_7^2=a(7-a)$. Thus every eigenvalue of the
least-residue sum is either $-a$ or $7-a$. Its augmentation is even,
hence equals $7-a$. Since the spectral mean is zero, this value
has total mass $a/7$. Every row of a principal $t$-block for
$t\ne7$ has the augmentation value modulo $t$; the two possible
eigenvalues differ by seven, so on the union it equals $7-a$.
In particular, at least eight nonunit labels have seven-free degree
as soon as there is a seven-power codegree and the rows of other
prime-power codegrees have mass greater than $3/7$. Either
$35$ nines, or $20$ nines and $22$ elevens, suffice, since
$20/81+22/121-3/7=1/6237$.

\subsection{Prime support in small rank}

\begin{theorem}[Three prime divisors in square dimension]\label{sq:three-prime-rank}
In the non-pointed based-simple case, if $D=M^2$ and exactly three
primes divide $M$, the rank is at least $45$.
\end{theorem}
\begin{proof}
Suppose $r\le44$. Lemma~\ref{blk:witness} gives $p^2\mid M$ for
every $p\mid M$, and Theorem~\ref{new:one-packet} excludes prime
denominators. If $M$ were odd, every denominator
would be at least nine, contradicting \eqref{eq:recip}. Thus the
support is $\{2,p,t\}$, $p<t$, and $M\ge900$. Write $m\le13$ for
the multiplicity of four (Proposition~\ref{sq:saturation}). Other nonunit denominators have floor six
(eight if $p>3$). We use the following position counts to bound the
sum of reciprocal squares, padding to $44$ places.

Put $A=4$ if $v_2(M)=2$, and $A=8$ otherwise. The first case has
$m\le3$ by Lemma~\ref{sq:heights}. Odd degrees occupy $3\bmod4$ nonunit places; these include
the $m$ fours when $A=4$, and one further denominator divisible by
$At$ from $P_t$. The other odd places have floor $4p$ if $A=4$,
and $8$ otherwise. For $s=p,t$ let $n_s$ count the $s$-free
nonunits. At least $k_s=H_{16}(s;2,m)$ have denominators divisible
by $4s^{v_s(M)}$ if $m\ge4$; put $k_s=0$ otherwise. Use
$n_s\ge\max(2,k_s)$, $n_3\equiv2\pmod3$, $n_5\ge3$, and the
stronger bound \eqref{eq:full-census} at five and seven when $m\ge4$.
The three omission sets are disjoint. When $m<4$ one $p$-free
denominator is divisible by $p^{v_p(M)}t$ and one $t$-free denominator
by $t^{v_t(M)}p$, by principal witnesses.

If $t$ mixes with two or $p$, reserve respectively $4(t-1)$ or
$p(t-1)$ active multiples of $2t$ or $pt$. These avoid the four-row
hits and, when $m>0$, the odd degrees, whose twists are dyadic. For
the odd interaction with $m\le3$ simply reserve the entire $pt$
places of the spherical summand, allowing the vacuum among them.
If $t$ is isolated, first test $t(t-1)$ pure places of floor $t^2$,
allowing them to supply $t$-free places other than four-row hits or
the principal-$p$ witness. Below this count sparse height gives
$M_t=t^{2c-1}$, $c\ge2$, and every nontrivial $t$-part has size
at least $t^2$. Nonminimal activity then reserves $2(t-1)$ pure
$t$-divisible degrees, by primitive balance or a higher orbit.
With fewer than $t(t-1)$ pure labels, minimal activity leaves fewer
than $(t-1)^2$ inactive pure labels. Thus Proposition~\ref{round:packet-idempotent} gives $t$
such degrees at the common denominator $t^c$.
Adding the reciprocal squares of these floors leaves only
$(p,t)=(3,5)$ with minimal activity, $A=8$, and
$m\in\{4,5,7,8,9,10,11,12,13\}$.

Write $M=2^a3^b5^{2c-1}$. The vacuum Weyl congruence gives
$2^a3^b\equiv\pm1\pmod{125}$. Truncate $a,b,c$ at the first
exponents with respective powers exceeding $K=65\,536$, the
last value representing its whole tail. Reserve the vacuum, the
$m$ fours, the five places at $5^c$ and the minimum omission
counts above. For each remaining position enumerate the permitted
divisors $d\le K$ of $M$ with their exact reciprocal squares;
a nonunit degree omits at most one prime. A larger denominator of
floor $F$ contributes in $[0,\max(K+1,F)^{-2}]$, also allowing
padding. The vacuum has positive mass at most $(2^a3^b5^{2c-1})^{-2}$ at the chosen lower endpoints.
If a valuation was truncated, omit its congruence test when necessary
and use the lower endpoint zero for a truncated packet mass.
Enumerate in order within each group, discarding intervals that
cannot contain one. The nine profiles leave no intervals after
$18\,095$ nodes (\supfile{arithmetic/three_prime_cutoff.py}{exact recursion}).
\end{proof}

The single-class bound also applies with any number of dimension primes.
For $m\ge7$ four-rows the degree $x=M/4$ is even if $m\ge9$ by
Proposition~\ref{sq:saturation}, or if $M$ has exactly four prime divisors
by Corollary~\ref{sq:odd-small}: otherwise $x$ has exactly three odd
prime divisors. Whenever $x$ is even, every label hit by these rows has $4\mid b$ by Lemma~\ref{sq:heights}. All odd degrees are hit by principality,
and every hit is inactive at the odd primes. For an odd prime $p\mid M$,
let $e_p,o_p$ count the even and odd nonunit hits of $p$-free degree,
outside the four-rows. Then
\[
                     e_p+4o_p\ge H_{16}(p;2,m).
\]
This is Proposition~\ref{sq:joint-projection} with weights one on even
degrees and two on odd degrees.

\begin{lemma}[Two inactive entries]\label{next:two-inactive}
In a complementary phase fiber at an odd allowed prime $p$, suppose every
$p$-active label has exact $p$-primary order $p$, and each
primitive $p$-phase contains exactly one label, of common reciprocal
square mass $C=1/B^2$, and exactly two eligible labels are inactive at
$p$. Their total reciprocal square mass $W$ satisfies
\[
                         W\le C/(p+1).
\]
\end{lemma}
\begin{proof}
Write their moment coefficients as $a_i\in p\mathbb Z$ and their
weights as $w_i=1/b_i^2$. The phase moments are
$\sum a_iw_i=W-C$ and $\sum a_i^2w_i=W+(p-1)C$.
If $a_1=0$, the second coefficient is $a_2=pk\ne0$, and solving
the two equations gives $w_1/C=(k+1)/k$. A positive rational square
of this form would make the coprime positive integers $|k|,|k+1|$
consecutive squares, which is impossible; $k=-1$ gives zero.
Thus both coefficients are nonzero, and
$p^2W\le W+(p-1)C$ proves the claim.
\end{proof}

\begin{corollary}\label{cor:odd-isolation}
In the non-pointed based-simple even square case through rank $44$,
no twist is active at two odd primes. Every label active at an odd
prime $p\ge7$ has prime-power denominator. An odd prime that is not
isolated has at least $4(p-1)$ active denominators divisible by $2p$.
\end{corollary}
\begin{proof}
The three-prime theorem supplies at least three odd dimension primes.
The existing spherical reservations exclude an odd edge of product
$K\ge33$, and for $K=21$ leave only the third prime five. They also
exclude a fourth odd prime: even granting every unreserved position
denominator four, the largest bound is
$11/16+1/15+1/7+1/11<1$. A fourth-prime spherical summand cannot
coincide with the first one because its dimension would exceed the
rank. A four must occur in the two remaining cases, so the vacuum
is reserved separately as before.

For $K=21$ a second odd edge, or dyadic interaction at five, gives
insufficient mass: for edges $15,21$ use $12/225+18/441+14/16<1$;
for $21$ and dyadic five use $16/100+18/441+10/16<1$.
At $K=21$ reserve its $22$ mixed labels.
Absence of five leaves mass below one even with thirteen fours.
Its presence reserves five labels and four proper witnesses; at
most eleven fours have insufficient mass, and thirteen overshoot.
At twelve all positions are reserved. The four proper five-free
denominators are at least $80$ by Corollary~\ref{new:small-Gauss}.
If any of the $22$ multiples of $21$ exceeds $21$, then
\[
 1\le12/16+1/5+21/21^2+1/42^2+4/80^2+1/420^2<1.
\]
Thus all equal $21$. At equality in Corollary~\ref{sq:fiber-squares}
there are exactly two nontrivial ternary fibers meeting seven-activity,
with eight eligible labels apiece. A higher ternary phase would require
three such fibers, a higher seven-phase $21$ labels in one fiber.
Each therefore has six singleton primitive seven-phases and two eligible
seven-neutral labels, all of the same weight. The fiber moments give
$(1-a_1)^2+(1-a_2)^2=8$ with $a_1,a_2\in7\ZZ$, impossible.

At $K=15$, denominator five is absent, for its presence would isolate five. The third odd prime $u$ is isolated: a second odd edge is excluded as above, and a dyadic reservation at $u\ge7$
together with the eighteen multiples of $15$ leaves insufficient mass. Thus
$u<44$ by the $u$ pure places. Reserve the vacuum, $m$ fours, these
eighteen places, and either $u$ at $u^2$ or $uj$ at $u$. In the latter
case partition $j(u-j)$ into at least three positive squares $a_i^2$.
The proper witness floors are $uk_i$, with $k_i>1$ least satisfying
$a_ik_i\equiv\pm1\pmod u$; use Corollary~\ref{new:seven-shape} at seven.
For $m\ge7$, Corollary~\ref{sq:odd-small} gives $8\mid M$, since exactly
three odd primes divide $M$. Reserve three odd degrees, imposing
$8\mid k_i$ on their overlap with the proper witnesses. The proper
witnesses include an odd degree by gcd one and an even degree since
every odd $u$-free degree contains three and five.
At $j=1$ the prime twist sum vanishes, so odd $m$ requires
$16(1-m/16-1/u+1/(30u))\ge1$. Every remaining upper mass is below one
(\supfile{arithmetic/odd_edge_mass.py}{exact mass bounds}).

Finally, a two--seven interaction reserves twenty-four labels of
floor fourteen; larger primes cost more. If five is present, its
packet, four proper witnesses, and two ternary active labels give
mass below one. Otherwise reserve six active positions at the other
odd primes. At most ten fours give the upper bound
\[
                  \frac{10}{16}+\frac{24}{196}
                       +\frac9{36}+\frac1{210^2}
                         =\frac{87977}{88200}<1.
\]
Eleven or more fours require at least three additional odd nonunits,
and hence at least $1+24+6+11+3=45$ labels. This excludes every
two--$p$ interaction for $p\ge7$. The remaining assertions follow
from the existing spherical and mixed-prime results.
\end{proof}

\begin{lemma}[Simultaneous proper supports]\label{new:proper-support-union}
A pairwise noncoprime family of integers greater than one with
collective gcd one has at least three primes dividing two or more
members. Consequently, in the non-pointed based-simple square case,
if the dimension primes are $2,3,5,7,11$ and minimal packets occur
at $5,7,11$, their proper free-degree supports have union of size
at least six.
\end{lemma}
\begin{proof}
Primes occurring in only one member cannot account for a pairwise
gcd. A pairwise intersecting family of nonempty subsets of a
one- or two-element set has nonempty intersection, proving the
first assertion. Write $H_p$ for the proper $p$-free labels.
Their squared Weyl norms are $4,6,10$, and each contains at least
three distinct degrees (Lemma~\ref{new:weyl}). If their union had
at most five labels, the only possible squared-norm partitions
would give $|H_5|=4$, $|H_7|=3$, $|H_{11}|=4$, with coefficient
magnitudes respectively $(1,1,1,1)$, $(2,1,1)$, $(2,2,1,1)$.
At most one $H_7$ degree is then divisible by five, and at most one
by eleven. None is divisible by seven, leaving only two primes
that can divide two $H_7$ degrees, a contradiction.
\end{proof}

\begin{proposition}\label{sq:many-primes}
Suppose no twist is active at two distinct odd primes. In the
non-pointed based-simple even square case, at least five dimension
primes require rank at least $45$, and at least six require rank at least $48$.
\end{proposition}
\begin{proof}
At an isolated odd prime reserve $p$ pure labels, and at a prime
mixed with two reserve $4(p-1)$ labels
(Propositions~\ref{new:pure-extra} and~\ref{round:mixed}). Five odd
primes cost at least $3+5+7+11+13=39$ nonunit positions. Through
rank $47$ only the prime three can mix with two. If all are isolated the
mass is at most
\[
 1/27+1/5+1/7+1/11+1/13+7/16+1/2310^2<1;
\]
if three mixes with two replace $1/27+7/16$ by $8/36+2/16$.
Larger primes or more primes only lower these bounds. This proves
the assertion for six dimension primes.

Now suppose $r\le44$ and exactly four odd primes occur. Reserve at each $p$ one of $(p,p^2)$,
$(4(p-1),2p)$, or $(p,p)$ according as $p$ is absent and isolated,
mixed, or present. A present denominator contributes at least three
proper witnesses at $2p$, or four at $20$ if $p=5$. For $m\ge9$
reserve three odd labels of floor eight, combining their floors with
the witnesses by Lemma~\ref{new:reservation-calculus}.
Use $M\ge2310$ and default floor six. Enumeration gives $49$
profiles, of which the $38$ with only prime-power twists are excluded by the
recursion of Lemma~\ref{round:deficits}. Here at least one prime denominator
is present; subtract its $j/p$ and $(m-\langle T_4\rangle)/16$,
and charge the vacuum, fours, prime classes and largest proper
witness quota before enumerating the remaining digits.

In the other eleven profiles eight reserved denominators are multiples
of six supported on two and three. Reserve also $p$ pure powers
for each absent isolated prime and $p$ copies of each present prime.
At most eight positions remain. A present prime has one packet:
this holds at five, and another packet at $p\ge7$ would leave fewer
than the three required proper degrees. On the remaining positions
its Weyl coefficients satisfy
\[
 \sum_i a_{pi}^2=p-1,\qquad \sum_i a_{pi}/b_i=\epsilon/M,\qquad
 \sum_i a_{pi}a_{ti}=-1\quad(p\ne t).
\]
Indeed $(W_p-I)e_0$ and $(W_t-I)e_0$ are orthogonal since no summand
is nontrivial at both primes; deleting the vacuum gives the last
identity. If $p\nmid b_i$ then $a_{pi}=0$; otherwise it is the least
absolute inverse of $b_i/p$ modulo $p$.

Enumerate the coefficient magnitudes with the indicated norms, at
least three nonzero entries per row and odd pairwise dot products.
Within each resulting column type try the compatible denominators
through $256$ in increasing order, allowing repetitions. A larger
denominator contributes mass in $[0,257^{-2}]$ and at most
$|a_{pi}|/257$ in modulus to the signed sum. The zero endpoints
allow padding to rank $44$. For the vacuum use
$M\ge4(\prod_{p\mid M,\,p>2}p)\prod_{t}t$, where $t$ runs over
the absent isolated primes. The mass test leaves eighteen intervals. The signed sums exclude all eighteen (\supfile{arithmetic/support_five_moments.py}{signed-moment calculation}).
\end{proof}

\begin{lemma}[Primary Gauss deficits]\label{round:deficits}
Suppose $D=M^2$, $\xi=\pm1$, and all active twists have prime-power
order. For an odd prime $p$, its contribution to
$1-\xi/M=\sum_i(1-\langle\theta_i\rangle)/b_i^2$ is a sum of
terms $p^{1-2k}$, $k\ge1$, with $p^k\mid M$. Each term accounts for
at least $p-1$ active labels. At a prime denominator $p$ with
$p\Vert M$ the contribution is $j/p$ in the notation of
Lemma~\ref{new:weyl}. At two each active label of denominator
$2^k$ contributes $2^{-2k}$ or $2^{1-2k}$.
\end{lemma}
\begin{proof}
An active denominator is a power of its twist prime by
Theorem~\ref{thm:spherical}. The primitive $p$-phases have equal mass
$W$ by Lemma~\ref{blk:moments}; their total deficit is $pW$.
Compute $W$ on the least numerous fiber. Higher phases have mean
zero and occur in same-degree square-unit orbits of size divisible
by $p$. Group their equal contributions $p$ at a time. The dyadic
assertion is the root trace formula.
\end{proof}

\begin{proposition}[High multiplicity at denominator four]
\label{sq:thirteen-34}\label{next:high-four44}
In non-pointed based-simple rotation-admissible data of even square
dimension through rank $44$, let $m\ge9$ labels have denominator
four. Exactly four primes divide the dimension, and the only
nontrivial twist component is $\{2,3\}$. There are thirty-nine
necessary mixed omission signatures: respectively $12,11,9,6,1$
at $m=9,10,11,12,13$.

Through rank $41$ one has $m\le12$. Through rank $40$ the support
is $\{2,3,5,7\}$, with five and seven isolated. In this range,
absence of denominator five forces seven at $m=9,10,11$; at $m=11$
they cannot both occur, and at $m=12$ seven is absent.
\end{proposition}
\begin{proof}
Theorem~\ref{sq:three-prime-rank}, Corollary~\ref{cor:odd-isolation}
and Proposition~\ref{sq:many-primes} leave exactly four dimension
primes. Proposition~\ref{sq:saturation} makes $M/4$ even and
$m\le13$. The vacuum, the fours and three odd nonunit hits are
disjoint from odd-prime activity, so
$\sum_{p\text{ odd}}(p-1)\le r-m-4$.

\emph{Common reservations.}
At an isolated odd prime $p$ absent as a denominator reserve $p$
pure labels of floor $p^2$; at a mixed prime reserve $4(p-1)$ active
labels of floor $2p$. A present prime supplies at least $p$ copies,
and exactly five at $p=5$. These are
Propositions~\ref{new:pure-extra}, \ref{round:mixed} and
Theorem~\ref{new:one-packet}, with Corollary~\ref{new:five-reduction}.
For isolated $p\ge7$ absent as a denominator, the remaining positions are fewer than
$p(p-1)$, so Proposition~\ref{new:activity-height-gap} gives
$p^3\mid M$ and puts every nontrivial denominator $p$-part at least $p^2$.

At isolated five without denominator five, twenty active labels give mass at most
$20/625+10/16+3/64+7/49+3/36+1/210^2<1$.
With eighteen, the three five-free nonunits are inactive of floor
$125$. Eleven fours still give mass below one; twelve force support
$2,3,5,7$, three pure ternary labels, seven pure seven-labels, and
three remaining positions, all five-free and odd. Absence of seven
gives mass below one, while its presence requires an odd proper
seven-free witness of prime-power degree. Thus at most sixteen
labels are five-active, and sparse height applies here too.
Twenty-three is excluded by
$9/16+1/27+1/5+1/23+3/64+1/2310^2<1$; larger primes exceed the rank.

Set $P_p=p$ for present or mixed primes, $P_3=9$ for isolated three,
and $P_p=p^3$ for isolated $p\ge5$ absent as a denominator.
A label omitting the odd primes in $J$ has denominator divisible
by $c\prod_{p\in J}P_p$, where $c=8,4,1$ for odd degrees, even
four-row hits and other even degrees. An odd degree omits at most
one odd prime, and an even degree at most two. Apply
Proposition~\ref{sq:joint-projection} on the hits with weights $1,2$
as above, then with unit weights and both possible common parities
of the row sums. Proper free degrees at a prime denominator omit
each other dimension prime. Exactly one of the four proper
five-free labels is odd: the signed sum is odd, and the other two
odd-prime omissions require distinct even labels. At least
$\lceil m/(16-m)\rceil$ are hits by the projection after
Lemma~\ref{new:cross-weyl}. At twelve fours use
Corollaries~\ref{new:small-Gauss} and~\ref{new:seven-shape}.
Thirteen fours exclude five by mass and seven by
Lemma~\ref{new:partial-gauss}: one packet leaves mass $5/112$,
forcing $T_4=0$ despite $T_4\equiv1$ above two; two packets exceed
the mass. The \supfile{arithmetic/large_four.py}{omission calculation}
uses these floors and $M\ge8\prod P_p$.

\emph{Mixed twists.}
The \supfile{arithmetic/high_four44.py}{rank-$44$ omission check}
gives forty-one signatures. Only two involve five: ten or eleven
fours, support $2,3,5,7$, isolated three and denominator seven.
Their five and one omission profiles have maximum masses
$77237/75600$ and $233293/226800$. Three odd neutral labels omit
respectively three, five and seven; two even four-row hits omit
$\{3,7\}$ and $\{5,7\}$, with an additional $\{3,5\}$ hit at
eleven fours.

An extra active pair at five loses $2(1/36-1/100)$, so exactly
sixteen occur. Each primitive five-phase has four labels of a common
positive dyadic height; a higher phase costs forty labels. Unless
all sixteen denominators are ten, equal phase masses lose at least
$4(1/100-1/400)=3/100$, exceeding both excesses. The three pure
ternary positions similarly have denominator nine, since the next
common level loses $8/243$. Their shape is a higher ternary orbit,
or a primitive pair and the neutral label of
Proposition~\ref{round:packet-idempotent}. At ten fours at most one
additional pure ternary label fits, again at nine; at eleven none
fits, and $11/16+16/100+3/81+1/7>1$ already excludes that case.

At ten fours, four nines leave one free position: a six overshoots
and a denominator at least eight undershoots. Thus there are three
nines, with deficit $1/27$ and $M_3=27$. The two free denominators
include exactly one six and one at least twelve. The five neutral
witnesses and vacuum must fill a gap of at least $29/75600$.
Write $M=2^a27\cdot5^c7$. If $c\ge2$, their floors
$216,200,56,756,700$ give mass at most $43949/119070000$.
If $c=1$, the seven-packet congruence gives
$135\cdot2^a\equiv\pm1\pmod7$, hence $a\equiv2\pmod3$ and $a\ge5$.
The floors $864,160,224,756,140$ give at most $2873/25401600$.
Both are below the gap; the
\supfile{arithmetic/two_five44.py}{exact rational check} records all
profiles and inequalities. The remaining thirty-nine signatures
are all mixed at two and three.

\emph{Prime-power twists.}
Let $t=\langle T_4\rangle$. Subtract the positive deficits of
Lemma~\ref{round:deficits} from $1-(m-t)/16-\sum j_p/p$, with final
error $\epsilon/M$. Reserve the vacuum, fours, prime classes and
proper witnesses. When five and another prime denominator occur,
at least six proper witnesses are needed: two at the other prime
omit two and the remaining odd prime, and neither is five-free.
Take the maximum of this quota and the three neutral labels supplied
by Proposition~\ref{new:mixed-neutral-rank} for the four singleton
components, allowing overlap.

If no prime denominator occurs, a zero intermediate remainder is
still impossible: at an isolated $p\ge7$, the terms summing to
$1/M$ could neither omit $p$ nor include a $p$-term, which is at
least $1/M_p>1/M$. A nonzero remainder $A$ with $s$ terms left
and largest term $1/u$ satisfies $|A|\le s/u+1/\sqrt u$.
This bounds the next prime power in the
\supfile{arithmetic/phase_deficits.py}{deficit recursion}, carrying
the required divisors of $M$.
The \supfile{census/preflight.py}{structural preparation} gives
$46$ primary profiles. Its
\supfile{census/arithmetic/int_deficits.py}{integer deficit recursion}
leaves only $M=136080,15482880$, both with nine fours, five fives
and seven sevens: two seven-packets would give
$9/16+1/5+12/49>1$.

The signed proper-degree unions number $13,39$. Since $T_5=T_7=0$,
the remaining mass gives $|\langle T_4\rangle|\le1$.
Enumerating ternary deficits $3/b^2$ at powers $b\ge9$ dividing
$M$, each costing two labels, and binary splittings of the remaining
dyadic deficit gives respectively zero and $218$ profiles.
For ternary deficit $G$, $h$ remaining positions and $v$ nines,
\[
                    2G/3\le v/81+(h-v)/729.
\]
After reserving the nines, all remaining degrees are multiples of
$35$ dividing $M$. Dividing the residual dimension by $35^2$ leaves
$5092$ nonnegative integer targets with at most five slots.
The \supfile{census/arithmetic/gauss_divisors.py}{divisor square-sum test}
excludes them all. Zero padding and the at-most-$44$ budget cover
every smaller rank.

Finally, the same preparation and the
\supfile{arithmetic/frontier_cases.py}{lower-rank omission check}
regenerate the tables through ranks $40,41,42,43$ and compare them
with the native rules. The ten mixed profiles through rank $40$
give the stated denominator restrictions; the rank-$41$ table
excludes thirteen fours.
\end{proof}

The classification above is used as a prefix restriction through rank
$44$. The complete enumeration of Section~\ref{sec:finish} covers
both mixed and all-primary branches at that rank.

\section{Odd dimension}\label{sec:odd-new}
\begin{theorem}\label{odd:rank}
A non-pointed based-simple rotation-admissible ring of odd dimension
has rank at least $217$, and at least $219$ if its dimension is nonsquare.
\end{theorem}
Throughout this section assume the hypotheses of this theorem.
Write $D=qM^2$ with $q$ square-free, $b_i=M/d_i$ and $r=|I|$.
Lemma~\ref{alg:odd-parity} gives odd degrees and denominators,
distinct nonunit dual pairs, and $r\equiv D\equiv q\pmod8$. Every nonunit degree has at least two
prime divisors and at least three primes divide $D$
(Lemma~\ref{lem:basic-table}).

\begin{proposition}[The diagonal determines the twist]\label{odd:diagonal-twist}
In odd global dimension,
\[
 \chi_i(i^*)\equiv\theta_i^2\pmod{2\mathcal O_{\QQ(\zeta_N)}},
 \qquad \QQ(\theta_i)\subseteq\QQ(\chi_i(i^*))\subseteq K_i .
\]
If $\theta_i$ has order $m$, its Galois orbit has cardinality divisible
by $\varphi(m)$. At each phase represented on this orbit the number of
labels is divisible by $2^{\omega(m)}$.
\end{proposition}
\begin{proof}
Write the unnormalized modular relation as $STS=p^+T^{-1}ST^{-1}$,
where $p^+=\sum_j d_j^2\theta_j$. At the $(i,i^*)$ entry, the two terms
of each nonunit dual pair are equal. The vacuum contributes $d_i^2$.
The same pairing gives $p^+\equiv1\pmod2$ and $d_i$ is odd. Thus
$S_{ii^*}=d_i\chi_i(i^*)\equiv\theta_i^2\pmod2$. If an automorphism
fixes this entry, the congruence shows that it fixes the reduction of
$\theta_i^2$. Roots of odd order inject under reduction above two.
Hence it fixes $\theta_i$. The orbit length equals $[K_i:\QQ]$.
The square-unit action on phases has image of size
$\varphi(m)/2^{\omega(m)}$ with equal fiber sizes on this orbit.
\end{proof}

\begin{lemma}[Simultaneous residue norms]\label{rev:two-residue}
Let $u\equiv1\pmod2$ be an algebraic integer of odd cyclotomic
conductor. For an odd prime $p$ let $h_p$ be the order of a subgroup
of tame $p$-inertia fixing $u$, and put $e_p=(p-1)/h_p$.
If $u$ reduces to $s\in\mathbb F_p$ at every place over $p$, write
$b\equiv s\pmod p$ for its even lift with $|b|<p$ and put
$\delta_p=\operatorname{dist}(b,h_p(e_p+2\ZZ))$. Then
\[
 (p-1) C(u)\ge p(p-1+\delta_p^2)-b^2.
\]
If $u$ vanishes at every place above every $p$ in a finite set $P$, then
\[
 C(u)\ge\prod_{p\in P}p\left(1+
 \frac{\mathbf1_{\{e_p\ {\rm odd}\}}h_p^2}{p-1}\right).
\]
Both bounds are sharp for the prescribed inertia subgroups.
\end{lemma}
\begin{proof}
Fix $p$, abbreviate $h=h_p$, $e=e_p$, and first work in
$K(\zeta_p)$, with $K$ cyclotomic of odd conductor prime to $p$.
As in Lemma~\ref{nt:residue-norm}, write
$u=\sum_{i=0}^{p-1}c_i\zeta_p^i$ with $c_i\in\mathcal O_K$ and
$\sum_i c_i=b$. Unramifiedness of $K$ at $p$ gives this unique
normalization. In the integral power basis modulo two,
$c_i-c_{p-1}\equiv\delta_{i0}$ for $i<p-1$, hence $c_0\equiv0$
and $c_i\equiv1$ for $i>0$. The latter coefficients have common
values $v_1,\ldots,v_e$ on the $e$ exponent orbits. Put $z=\sum_jv_j$.
The trace identity is
\[
 (p-1)C(u)=p\{C(b-hz)+h\sum_j C(v_j)\}-b^2.
\]
Take the nearest $t\equiv e\pmod2$ to $b/h$ and write $z=t+2x$.
Here $x$ is integral, $|b-ht|=\delta_p\le h$, and
$C(x)\ge|\Re\langle x\rangle|$ by Cauchy--Schwarz and the algebraic
norm bound for nonzero integers. Expansion gives $C(b-hz)\ge\delta_p^2$.
Every $v_j$ is nonzero. Equality is attained in the fixed field of
the indicated group by taking $v_j=\pm1$ with sum $t$, possible since
$|t|\le e$. In particular for $h=1$ equality in $\QQ(\zeta_p)$
means precisely $c_0=0$, $c_i=\pm1$ for $i>0$, $\sum c_i=b$.

For the product bound use induction, starting with $C(u)\ge1$.
Take $p\in P$, so $b=0$. If $q\in P\setminus\{p\}$, the radical of
$q\mathcal O_K$ extends to that of $K(\zeta_p)$ by unramifiedness.
The power basis puts all $c_i-c_{p-1}$ in the radical; $\sum c_i=0$
then does the same for all $c_i$ since $q\ne p$. The other inertia
subgroups fix every $c_i$ by uniqueness. Apply induction to each
$v_j$, and also to $z$ when $e$ is odd: it then reduces to one at two.
The trace identity gives the asserted factor, since $c_0=-hz$.
The product of the prime-level extremizers attains equality by linear
disjointness.

At higher $p$-power level apply the inertia-equivariant random
specialization in Lemma~\ref{nt:residue-norm}. In the integral power
basis over $K(\zeta_p)$, the constant coefficient is one modulo
two and all others vanish modulo two. At each other required prime
all coefficients are in its radical by unramifiedness. Thus every
specialization has the required reductions, including the one at $p$
since the substituted roots all reduce to one there. It retains the
other prescribed inertia symmetries because the other inertia groups
fix every $p$-power root. Its expected mean squared modulus is $C(u)$,
so the level-$p$ argument applies for arbitrary conductor of $K$.
\end{proof}

\begin{corollary}[Diagonal orbit cost]\label{rev:diagonal-cost}
In odd dimension let $i$ be nonunit, $h_i=[K_i:\QQ]$, and
$P_i=\{p:p\mid d_i,\ i\in P_p\}$. Then
\[
 \sum_{Y\in G i^*}|\chi_i(Y)|^2\ \ge\ h_i\prod_{p\in P_i}p,
 \qquad f_i\ge1+h_i\prod_{p\in P_i}p.
\]
More generally, if $i\in P_p$ and $\theta_i$ has $p$-power order,
the same orbit contributes at least
$h_i\bigl(p-b^2/(p-1)\bigr)$, with $b$ the even representative of
$d_i\bmod p$ of absolute value less than $p$.
\end{corollary}
\begin{proof}
Proposition~\ref{odd:diagonal-twist} gives
$u=\chi_i(i^*)/\theta_i^2\equiv1\pmod2$. At every place above
$p\in P_i$ the numerator reduces to $d_i\equiv0$.
In the more general assertion $u$ reduces to $d_i$, since the
twist reduces to one. Symmetric Galois covariance makes the displayed
orbit sum $h_i C(u)$, as division by a root of unity preserves each
squared modulus. The vacuum contributes one outside this orbit.
\end{proof}
One may use the full product in Lemma~\ref{rev:two-residue} with the
stabilizers of the normalized diagonal entry. In particular a prime
$p\in P_i$ whose tame inertia acts trivially on $K_i$ contributes $p^2$.

\subsection{Dual pairs and primary phases}
\paragraph{Counting pairs at fixed phases.}
In a fiber excluding the vacuum, dual pairs inherit square-unit symmetry and the mass identities of Lemma~\ref{blk:moments} (divide all masses by two). For a set $J$ of odd primes $p$ dividing $D$ with even valuation, the proof of Proposition~\ref{blk:packets} gives at least $\sum_{p\in J}(p-1)$ distinct nontrivial tuples. If every such fiber has at least $l$ labels, the same bound for every subset of $J$ and Hall's theorem select disjoint sets of $l(p-1)$ labels active at each $p$.

For each prime $t\mid q$ a nonempty fiber fixing any or all of the other
odd-primary phases has at least
\[
 C(t)=\begin{cases}t,&t\not\equiv1\pmod{16},\\
 (t-1)/2+2,&t\equiv1\pmod{16}\end{cases}
\]
labels, and at least $C(t)$ \emph{pairs} when it excludes the unit.
Indeed, order $t^2$ or higher gives at least $t(t-1)/2$ by squared
Galois symmetry. Otherwise the mass identity forces both trivial and
primitive phases. Write $h=(t-1)/2$; the number of primitive phases
counted with multiplicities is $hk$. Since every squared degree is
one modulo eight, $hk\equiv2h n_0\pmod8$, where $n_0>0$ counts
objects at phase one. If $8\nmid h$, then $k$ is even and the fiber
has size at least $2h+1$. If $8\mid h$, the bound $h+2$ follows
from Proposition~\ref{blk:packets}, whose proof works identically for
pairs. At every such $t$ at least $h+1$ different phases occur in
the fiber by the same proposition.

\paragraph{The trace on a dual pair.}
For a nonunit $X$ in odd dimension, choose one $Y$ in each nonunit dual
pair and put $z_Y=\chi_X(Y)+\overline{\chi_X(Y)}$. Then
\[
 \sum_Y z_Y^2=f_X-2,\qquad
 \sum_Y d_Yz_Y=-1,\qquad
 \sqrt D\le1+\sum_Y d_Y|z_Y|.
\]
Indeed $X$ and $X^*$ give orthogonal rows, each of norm squared $f_X$.
Their sum has vacuum entry two and takes the same value on $Y,Y^*$.
The second identity is row orthogonality to the degree row. For the
last inequality add the vacuum modular relations for $X,X^*$: both
right sides are $p^+\theta_X^{-1}$ and $|p^+|=\sqrt D$. Galois acts
on pairs, and on each orbit all conjugates of $z_Y$ occur equally often.

\begin{lemma}\label{odd:27}
If exactly three primes divide the odd dimension, no label has
formal codegree $27$.
\end{lemma}
\begin{proof}
A label $X$ of codegree $27$ gives $D=3M^2$, $n=d_X=M/3$.
Write $3,p,l$ for the three primes. Lemma~\ref{blk:witness}
gives $s^2\mid M$ for each of them.
The row $X$ belongs to and is supported on the principal $3$-block,
so all hit denominators are multiples of three
(Lemmas~\ref{blk:support}, \ref{blk:height} and
Corollary~\ref{blk:primepower}).

The nonreal row field is contained in $\QQ(\zeta_{27})$. Its
degree cannot be $18$ by Theorem~\ref{thm:tame-orbit} (the upper
bound on the number of nonzero-mean orbits would be two).
It therefore lies in $\QQ(\zeta_9)$, by cyclicity.
For $3\mid d_Y$ the algebraic integer $z_Y$ vanishes at the place over three
by principality. Each orbit of such pairs on which $z_Y\ne0$
costs at least nine in $\sum z_Y^2$. A nonzero rational value
is divisible by three. Otherwise $z_Y$ generates the real cubic field of ninth roots.
All three conjugates reduce to zero above three. Its nonzero integral
norm and the trace of its square are divisible by three. By
arithmetic--geometric mean this trace is at least $3\cdot3^{2/3}>6$
and hence at least nine, as in Lemma~\ref{nine}.

The identity $\sum d_Yz_Y=-1$ requires a $p$-free and an $l$-free
degree with nonzero $z_Y$. These are distinct Galois orbits,
both at degrees divisible by three. They spend at least eighteen
of the available $25$, leaving no further nonzero orbit at a
three-divisible degree. On any union of nonzero orbits the number
of pairs is at most its squared norm, by the arithmetic--geometric
mean applied to the conjugates of each nonzero algebraic integer
squared. Cauchy--Schwarz bounds the sum of absolute values by the
same squared norm. The three-free pairs thus contribute at most
seven to $\sum |z_Y|$, and all pairs at most $25$. Three-free hit
degrees are at most $n/3$, and $s$-free hit degrees for $s=p,l$
are at most $n/s^2$, since $3s^2\mid b_Y$. The trace inequality gives
$3\sqrt3\,n\le 1+7n/3+25n/\min(p,l)^2$, a contradiction.
\end{proof}

\paragraph{Odd multiplicities at defect one.}
At an odd ramified prime $t$, both degrees $u,v$ of a defect-one
block are odd. In Lemma~\ref{blk:defectone} the square $(v/u)^2$
is either $a$ or $2a$. The latter has positive $2$-adic valuation,
so only the former occurs, with $b=t-1$ and $a$ an odd square.
In particular, when $t-1$ has no nontrivial odd square divisor
each block consists of $t$ labels of equal degree, exactly one
$t$-inactive. This count persists at fixed complementary phases. Thus if every
$t\mid q$ has this property and $\gcd(b_i,q)=1$, successive restriction
to inactive phases divides the multiplicity by $q$. The remaining
labels have trivial inertia at every $t\mid q$, so their row fields
lie in $\QQ(\zeta_{b_i^2})$ by Lemma~\ref{conduct}. This excludes
$b_i=1$ by nonreality. If $b_i$ is a power of a prime $p$, the
stabilizer in this cyclic Galois group has odd order, since its unique
involution is complex conjugation. The multiplicity of the degree is
therefore divisible by $q\,2^{v_2(p-1)}$.

\paragraph{The first active denominator.}
Suppose $q=t$ as in the preceding paragraph and $p\ne t$ is an
unramified odd prime of $D$ (that is, $v_p(D)$ is even).
In a degree of denominator $b$ divisible by $p$, activity at order
$p^j$ occupies multiples of $p^{j-1}v$ labels on each square class,
where $v=p-1$ if $t\mid b$, and $v=t(p-1)$ otherwise. Indeed
square covariance supplies $p^{j-1}(p-1)/2$ phases, duality
pairs the labels, and the equal-degree $t$-blocks supply the factor
$t$ at each fixed phase when $t\nmid b$. If $b=p$ and $r<10v$,
the two primitive square classes have equal multiplicity at this degree.
Otherwise their normalized mass difference is at least $v/p^2$,
and Lemma~\ref{blk:moments} requires the same discrepancy among the
remaining $p$-multiple denominators, all at least $3p$.
This costs at least $9v$ additional labels. Consequently the activity
in $m$ copies of denominator $p$ is at most the largest integer
$2v h+pv k\le m$ with $h,k\ge0$ in this range.

\subsection{Odd nonsquare dimension}
Suppose $D=tM^2$ is odd, $t$ prime, and $t-1$ has no nontrivial
odd square divisor. At an allowed prime denominator $p$, all
$t$-blocks are equal. Theorem~\ref{rev:height} gives $p\Vert M$ and its proof gives twist order one or $p$ on each effective row. Proposition~\ref{new:effective} gives invariance of the Gauss factors on $p$-inertia orbits and dual pairs.

Suppose $p\Vert M$ and all $p$-active labels have denominator $p$.
Write $m$ for the number of defect-one blocks of this denominator,
$n=M/p$, and $j$ for the number of such blocks per primitive $p$-phase.
These counts agree by Lemma~\ref{blk:moments}, since all active
labels have the same degree and order $p$ at $p$. Duality makes $j$
even and the effective multiplicity bound gives $0<j<p$.
Then $m\ne p^2-p$ and
\begin{equation}\label{eq:effective-weyl-bound}
           p|pj-m|\le j(p-j)+(p+1)j/n .
\end{equation}
The exceptional multiplicity is excluded by
Lemma~\ref{lem:congruent-rows} on the effective rows.

Here are the details for the cost. Set $w=W_pe_0$ and let
$L_Y$ be the sum of the effective rows of the active blocks at $Y$.
For $k\ne0\bmod p$, choose an exponent inducing $k$ at the local
$p$-factor and zero on its complement. The vacuum Bruhat identity,
as in Lemma~\ref{blk:moments}, and the vacuum rotation entry give
\[
  \theta_{Y,p}^{-1/k} w_Y-\delta_{Y0}
     =p^{-2}\sum_{B\ {\rm active}}(\theta_{B,p}^k-1)v_B(Y).
\]
In fact each of the $t$ summands of a block has codegree $tp^2$.
The local torus fixes the vacuum since $v_p(D)$ is even.
Averaging over $k$ gives
\[
 p(w_Y-\delta_{Y0})=
 \begin{cases}
   -L_Y/(p-1),&\theta_{Y,p}=1,\\
   L_Y,&\theta_{Y,p}\ne1 .
 \end{cases}
\]
Thus $A_Y=pw_Y$ vanishes outside the $t$-neutral coordinates,
$A_0=p-j$, and $A_Y\equiv-jd_Y\pmod p$ on those coordinates.
Indeed $w$ is rational by \eqref{eq:phase-column}, the unaveraged
identity makes it integral away from $p$, and $b_Yw_Y\in\ZZ$ with
$p^2\nmid b_Y$. For the unique $t$-neutral member $X$ of every
denominator-$p$ block one has $|w_X|=1/p$: the projection of $se_0$
onto its block has that norm and $W_t^{-1}$ sends its line to
$\mathbb C e_X$; all remaining complementary Weyls preserve this
coordinate line. Thus $A_X=\delta$, one common sign by the congruence.

Write $a_Y=\delta A_Y$, and on the proper nonunit support use
$k_Y=n/d_Y>1$, since $w$ is supported on the principal $p$-block.
The congruence gives $k_Ya_Y\equiv1\pmod p$.
As $\sum d_YL_Y=0$, the averaging identity gives
$\sum d_Y A_Y=p+pnj\delta$:
in the active part precisely $(p-1)j$ neutral singletons contribute
$\delta n$ each. Removing the vacuum and the $m$ singletons from
this identity and $\sum A_Y^2=p^2$ gives exactly
\eqref{eq:weyl-cost} with $z=pj-m$, and hence \eqref{eq:effective-weyl-bound}.

\begin{lemma}[Two neutral rows]\label{odd:two-neutral}
Suppose $p\ge5$, $p\Vert M$, and $j=2$, either in odd square
dimension under Lemma~\ref{new:weyl} or in the effective rows above.
Then $m=2p$, $p\equiv3\pmod4$ and the proper Weyl coefficients,
one per dual pair, have squared sum $p-2$ with exactly one odd term.
In particular $p\ge11$, and at eleven the magnitudes are $1,2,2$.
\end{lemma}
\begin{proof}
In the square case $m=2p$ is Theorem~\ref{new:one-packet}.
For effective rows, use the $t$-neutral part of the charge-even space.
It is $p$-invariant and satisfies Lemma~\ref{r8:steinberg-compression};
its normalized rows are unit scalar multiples of $W_pv_X$, equal
on dual blocks (Proposition~\ref{new:effective}). Thus the same
active-column comparison and signed estimates, with
\eqref{eq:effective-weyl-bound}, exclude $z=\pm2$ and give $m=2p$.

Write $e=\epsilon$ or $\delta$ and let $g$ be the neutral projection
of $pe(W_p-I)v_0$. The lemma gives compression $-gg^*/(2p)$,
where $g_0=-2e$ and $g_Y=\sqrt2a_Y$ on dual pairs. On the neutral
denominator-$p$ pair $a_Y=1$; its normalized row has vacuum entry
$\sqrt2/p$, so its scalar multiplying $W_p$ is $e$.
Consequently the neutral row $X$ has pair trace
$\chi_X(Y)+\overline{\chi_X(Y)}=-e a_Y$ on every proper pair,
$-e$ on the active pairs and $e(p-1)$ on its own pair.
For the active entries use \eqref{eq:packet-trace} and symmetry
at equal degree (in the effective case use the $t$-neutral members).
Off its support the effective row is zero.

The neutral pair has nonreal quadratic field in $\QQ(\zeta_{p^2})$,
hence $p\equiv3\pmod4$ and the field is $\QQ(\sqrt{-p})$
\cite[Ch.~6]{Washington}. Write
$\chi_X(Y)=(z_Y+w_Y\sqrt{-p})/2$ on the support.
The integers $z_Y,w_Y$ have the same parity. Orthogonality of the
two dual rows gives $\sum_Y w_Y^2=p$, over nonunit pairs.
The $p-1$ active pairs have odd trace. On the proper pairs
$\sum_Y a_Y^2=p-2$ is odd, so exactly one more odd trace occurs.
At least three distinct proper degrees occur: their gcd is one,
their degrees divide $n$, and any two share a prime. Thus
$p-2\ge1+4+4$, with equality precisely at the indicated triple.
\end{proof}

\begin{proposition}\label{odd:ternary-multiplicity}
In odd dimension $D=3M^2$, every nonquadratic row of denominator
three has twist of order nine. If either $27\mid M$, or $9\Vert M$ with exactly four prime
divisors of $M$ and at most two labels of denominator nine, then the
three cube roots occur equally often among the remaining rows of
this denominator.
In particular its multiplicity is divisible by six.
\end{proposition}
\begin{proof}
Put $n=M/3$. The row field is nonreal and contained in $\QQ(\zeta_9)$,
by the codegree-$27$ argument of Lemma~\ref{odd:27};
Theorem~\ref{thm:tame-orbit} gives $\theta_X^9=1$.
For a nonquadratic row take the three conjugates of the trace row
$z_Y=\chi_X(Y)+\overline{\chi_X(Y)}$ under the cubic subgroup.
Their mean $w$ is integral: at each coordinate this is the normalized
trace on $\ZZ[\zeta_9]^+$, which is integral by averaging powers of
$\zeta_9$ over $\zeta_9\mapsto\zeta_3^j\zeta_9$. On nonunit dual
pairs their squared norms are $25$ and mutual inner products $-2$
by row orthogonality. Hence $\sum w_Y^2=7$ and $\sum d_Yw_Y=-1$.
If $\theta_X^3=1$, the three rows have the same twist by square-unit
covariance, so $3\sqrt3\,n\le1+\sum d_Y|w_Y|$.
Every hit degree divides $n$ by principal-block height. Put
$a=\sum_{d_Y=n}|w_Y|$. The right side is at most
$1+n(a+(7-a)/3)$, forcing $a\ge5$ (a nonunit degree has at least
two distinct odd prime divisors). The proper terms in orthogonality,
including the vacuum, now have total magnitude at most $1+2n/3<n$.
Thus $\sum_{d_Y=n}w_Y=0$ and $a=6$. All six contributions have
magnitude one, since any larger integer spends at least two extra
units of squared norm. The sole remaining proper term would give
$1\pm d_Y=0$, impossible.

For a quadratic denominator-three row, the traces on nonunit pairs
are integers with squared sum $25$. A nonzero trace at a
three-divisible degree is a multiple of three, hence equals
$\pm3$; there are at most two. When $27\mid M$, every
three-free degree is at most $n/9$, and a proper three-divisible
degree is at most $n/3$.

Suppose instead that $9\Vert M$, exactly four primes divide $M$,
and $m_9\le2$. At most one three-free pair has degree $n/3$,
and every other three-free degree is at most $n/15$; a proper
three-divisible hit has degree at most $n/5$. If the free squared
norm is $F$, its absolute degree-weighted trace sum is at most
\[
 \frac n{15}\max_{a\in\ZZ_{\ge0},\ a^2\le F}(5a+F-a^2).
\]
For $F=25,16,7$ the maxima are respectively $31,22,13$.
In either valuation case, unless two traces occur at degree $n$,
the Gauss bound is at most $1+43n/9<3\sqrt3\,n$.
Thus the two traces occur at $n$. The remaining weighted trace
sum is at most $13n/15<n-1$, so their signs are opposite.
Each pair is individually Galois-stable, being uniquely specified
by its trace sign, and hence is quadratic. The error in the cube-phase comparison below is at most
$(1+13n/15)/(3n)<1/3$, since $n\ge15$. Quadratic rows have
cube-root twists: the multiplier four fixes the row and raises
its twist to the sixteenth power.

Write $\theta_+,\theta_-$ for the twists at the positive and negative
traces. The vacuum modular relation, divided by $3n$, puts
$\theta_+-\theta_-$ within $1/3$ of $\pm\sqrt{-3}\,\theta_X^{-1}$:
indeed $p^+=\pm M\sqrt{-3}$ since $p^+/(M\sqrt{-3})$ is a
two-power root of unity in the odd-conductor twist field.
The differences of cube roots are zero and six points of modulus
$\sqrt3$, separated by at least $\sqrt3$. Thus equality holds.
A quadratic pair of phase $\zeta_3^a$ has precisely one trace neighbor
of exponent $1-a$ and one of exponent $2-a$ modulo three.
This incidence is symmetric by Drinfeld duality at equal degrees.
Counting incidences from exponent zero to each of the other two
therefore gives equal cardinalities in all three phase classes.
\end{proof}

\begin{lemma}[Constant pair traces]\label{odd:constant-pair-trace}
Let $I$ be $m$ nonunit rows of common formal codegree $f$ in integral
modular data. Let $(Y_j,Y_j^*)$, $1\le j\le t$, be disjoint nonunit
dual pairs, with $Y_j\ne Y_j^*$, such that
$\chi_X(Y_j)+\overline{\chi_X(Y_j)}=z_j\in\mathbb Z$
for every $X\in I$. Then
\[
             m\left(1+\frac12\sum_jz_j^2\right)\le f,
\]
strictly if $1+\sum_jd_{Y_j}z_j\ne0$. In particular the inequality
is strict for one pair with $d_Y>1$ and $z\ne0$.
\end{lemma}
\begin{proof}
Restrict the orthonormal columns $se_0$ and
$s(e_{Y_j}+e_{Y_j^*})/\sqrt2$ to $I$. Their Gram matrix is
$(m/f)vv^t$, where $v=(1,z_1/\sqrt2,\ldots,z_t/\sqrt2)^t$.
Its largest eigenvalue $(m/f)\|v\|^2$ is at most one.
Equality makes the corresponding column combination vanish outside
$I$, whereas its vacuum coordinate is
$(1+\sum_jd_{Y_j}z_j)/\sqrt D$.
\end{proof}

\begin{corollary}[A ternary class at minimal valuation]\label{odd:minimal-ternary}
Let $D=qM^2$ be odd with exactly four prime divisors, where
$q\in\{1,3\}$. Suppose $3^a\Vert M$ and $f=q3^{2a}\ge27$.
In non-pointed based-simple rotation-admissible data, at most $f-11$
labels have denominator $3^a$. In particular the bound is sixteen
when $D=3M^2$ and $3\Vert M$, and seventy when $D=M^2$ and $9\Vert M$.
\end{corollary}
\begin{proof}
Put $n=M/3^a$ and suppose the multiplicity $m$ is at least $f-9$.
The selected rows lie in the principal three-block $P$. Every label
$Y\in P$ meets all of them, since $\chi_Y(X)$ reduces to the three-free
integer $n$. If $3\mid d_Y$, lcm integrality makes
$\chi_X(Y)/3$ integral; Galois averaging column unitarity, with its
positive vacuum term, gives $9m<f$, a contradiction. Conversely every
three-free label is hit by a selected row, whose
support lies in $P$ by Lemma~\ref{blk:support}. Thus $P$ consists
exactly of the three-free labels, and $F=|P|\le f$ by
Galois averaging. The integer $F$ is odd and divisible by three.

Removing the copies of $n$ from the nonunit degrees in $P$ preserves
their gcd one, since $n$ contains every prime other than three.
Each proper degree contains at least two of those three primes.
Hence at least three distinct proper degrees, or six labels, occur,
and $F\ge m+7$. As $m\ge f-9$, this forces $F=f$.

Norm saturation makes every selected row entry on $P$ a root of unity.
There are at most four proper dual pairs. Each of the other three
primes is omitted by some such pair, and each pair omits at most one.
Some prime $t\ne3$ is therefore omitted by a unique pair $Y,Y^*$,
which is Galois-stable and has a quadratic row field. By symmetry its
selected-row entries lie in
$\mathbb Q(\zeta_f)\cap K_Y\subseteq\mathbb Q(\sqrt{-3})$, so its traces
are $2\epsilon$ or $-\epsilon$, with common sign $\epsilon$
fixed by reduction above three. Vacuum orthogonality modulo $t$
fixes the same trace on every selected row, since the two choices
differ by $3\epsilon$. Lemma~\ref{odd:constant-pair-trace} gives
$3m/2<f$, contradicting $m\ge f-9\ge2f/3$.
Evenness of $m$ now gives $m\le f-11$.
\end{proof}

\paragraph{Large quinary classes.}
If a five-power codegree $f$ occurs with multiplicity $m>f/2$
in odd dimension, every five-free nonunit character field contains
$\QQ(\zeta_5)$, by the observation after
Lemma~\ref{new:many-reserve}. Its Galois orbit therefore has length
divisible by four. Deletion requires at least two such orbits.
There cannot be exactly eight five-free nonunits: their two
four-element orbits would give $1+4a^2+4b^2\equiv0\pmod5$,
which has no solution with $5\nmid ab$. Thus there are at least
twelve five-free nonunits.

\begin{lemma}[Omission witnesses at a real quadratic prime]\label{odd:omission-four}
In odd-dimensional non-pointed based-simple rotation-admissible data,
let $p\equiv1\pmod4$, and let $I$ consist of $m$ rows of common
$p$-power formal codegree $f$, with $mp^2>f$. Fix another prime $q$.
Suppose every $q$-free nonunit label that meets a row in $I$ has degree
and denominator divisible by $p$. Then at least four such labels occur.
\end{lemma}
\begin{proof}
Each row in $I$ has a nonunit support of degree gcd one, so it meets
a $q$-free nonunit. The union $J$ of these hits is Galois-stable and
stable under duality. If it had only two labels, they would form one
Galois-stable dual pair $Y,Y^*$. Its character field is imaginary
quadratic, whereas the unique quadratic subfield of
$\mathbb Q(\zeta_f)$ is real. Hence every $\chi_X(Y)$, $X\in I$,
is rational. The rows are principal at $p$, so these values are
integers divisible by $p$. Each is nonzero, since each row must meet
$J$. Column norm then gives $mp^2/f\le1$, a contradiction.
\end{proof}

\begin{proposition}[A ramified cost for an allowed edge]\label{r8:ramified-edge}
Suppose $D=3M^2$ is odd, and let $p,t\ge5$ be distinct primes.
If some twist is active at both $p$ and $t$, then at least
$12(p-1)(t-1)$ labels are active at both. In addition, at least
$2(t-1)$ labels are active at $t$ and inactive at $p$, and at least
$2(p-1)$ are active at $p$ and inactive at $t$.
\end{proposition}
\begin{proof}
Apply the allowed-prime two-moment identities in each fiber of
complementary phases. If all $p$-active phases in that fiber are
primitive, the counts of its primitive phases agree modulo eight,
since their masses agree and odd inverse squares equal one modulo
eight. Their total count in dual pairs is even. The moment parity
$k\equiv c\pmod2$ therefore makes each primitive-phase count even.

First suppose the jointly active phases are primitive at both primes.
Equality of primitive masses supplies all $(p-1)(t-1)$ phase pairs.
At each pair, the count in each fixed three-phase is even, by the
preceding paragraph. If a higher three-phase occurs, its square-unit
orbit has at least three phases, and hence at least six dual pairs.
Otherwise let $n_0,n_1$ count the neutral and primitive three-phase
pairs. Both are positive and even, and the ramified mass identity
implies $n_1\equiv2n_0\pmod8$. Consequently $n_0+n_1\ge6$.
Every phase pair therefore costs at least twelve labels.

If a jointly active phase has $p$-order at least $p^2$, its
square-unit orbit has at least $p(p-1)/2$ phases. At each of them
there are at least $t-1$ nontrivial $t$-phases: primitive phases have
equal positive mass, and a higher orbit has at least $t(t-1)/2$
phases. A nonempty fiber of these allowed phases costs at least
six labels by the ramified three-mass identity and square-unit
orbits. The resulting bound $3p(p-1)(t-1)$ exceeds the claimed
bound. Interchange $p,t$ for a higher $t$-phase.

Finally, in each occupied nontrivial $t$-phase fiber, the second
$p$-moment has positive active mass. It forces a $p$-inactive dual
pair with denominator divisible by $p$. There are at least $t-1$
such phases. The analogous argument at $t$ gives the other solely
active family, and both are disjoint from the jointly active labels.
\end{proof}

\paragraph{Prime support.}
Call a prime $p\mid D$ unramified when $p\nmid q$. Suppose $q>1$ and $r\le218$. Then
Theorem~\ref{thm:squarefree-codegree} gives $9q<r$ and $q\le23$. For $q=3s=15,21$ an
unramified $c$ gives $r\ge9q+4s(c-1)(1-9/c^2)>218$
($c\ge7,5$, respectively). At $q=19,23$ use
$r\ge9q+2C(q)(p-1)(1-9/p^2)>218$ for an unramified $p\ge5$.
Thus $q=t\le17$ is prime; all $t$-blocks at an unramified prime
denominator have equal degrees.

On three-prime support Lemma~\ref{blk:witness} gives $p^2\mid M$
for every $p\mid M$, excluding unramified prime denominators by the height theorem. For $t=3$ Lemma~\ref{odd:27}
now gives $r\ge81t$, and for $t\ge7$ one gets $r\ge49t$.
For $t=5$ we have $r\le213$, so $m_5>62$ by the mass bound
$62/25+151/81<5$. Reserve twelve five-free labels at floor $25$.
At each other prime $p$ a $p$-free hit of a five-row has five-divisible
degree and denominator divisible by $5p^2$; the preceding omission bound
reserves four such labels. These sets and the active sectors are disjoint.
Reserving twenty active labels at nine and sixty at twenty-one gives
$112/25+20/81+60/441+12/625+4/45^2+4/245^2+11025^{-2}<5$.

With at least four primes, whenever $j=2$ one has $2t(p-1)\le r$ and $n\ge3\cdot5\cdot7$, so Lemma~\ref{odd:two-neutral} applies. The inequality $C(t)(1+2\sum_{p\ne t}(p-1))\le r$
leaves $t=3,5$. For five the possible supports are $3,5,7,11$ and
$3,5,7,13$, with disjoint unramified activities (a $3,7$ edge costs
$120$, apart from one hundred at the last prime). The effective
multiplicity at denominator three is even, at most seven and unequal to
six by Lemma~\ref{lem:congruent-rows}, giving $m_3\le20$.
Concentration at three forces $m=4,j=2$ in \eqref{eq:effective-weyl-bound},
hence $6\le2+8/n$. Concentration at seven excludes $j=2$ by
Lemma~\ref{odd:two-neutral}; $j\ge4$ costs $120$ labels beside
twenty and one hundred. At three an imbalance costs one hundred
further labels besides the other $160$; at seven use the first-active
denominator bound. Removing the balanced equal quinary blocks therefore
requires the full $20,60$ residual activity at floors nine and
twenty-one. The mass is at most
$20/9+12/25+20/81+60/441+100/121+1155^{-2}<5$.

Hence $D=3M^2$, $r\le211$. Four unramified primes require precisely
$P=\{5,7,11,13\}$. The edge bounds above make their activities disjoint,
costing at least $A_p=6(p-1)$ apiece. At $p=5,7,13$, $j=2$ is impossible
by Lemma~\ref{odd:two-neutral}, while $j\ge4$ plus the other costs and the
vacuum gives $217$ labels. An imbalance at $p$ already costs
$30(p-1)+\sum_{s\ne p}A_s>211$. Removing balanced ternary blocks and
using the ramified mass identity therefore reserves $A_p$ labels at $3p$
at these three primes. The upper mass is
$18/9+24/15^2+36/21^2+60/11^2+72/39^2+15015^{-2}<3$.
There remain three unramified primes, with $3+\sum A_p\le211$ (fourteen
supports). Their activities are again disjoint and their primitive phase
classes balanced at prime denominators. If activity persists after
deleting the ternary blocks at $p$, it costs $A_p$ labels at floor $3p$;
otherwise use \eqref{eq:effective-weyl-bound} and Lemma~\ref{odd:two-neutral}.

\paragraph{The first two prime powers.}
Write $u=m_3$ and $n=M/3$. When $9\mid M$ and $u>0$, let $F$
be the number of three-free labels including the vacuum.
It is odd, divisible by three and at least nine by deletion and the
dimension identity. Every three-free column costs at least one in
the Galois-averaged norm of a three-row, so $F\le27$. Equality
forces every three-row to vanish on all other columns. Bessel applied
to the vector $(d_Y\theta_Y)$ on these columns and the modular
relation then gives $u/9\le\sum_{3\nmid d_Y}b_Y^{-2}\le27/81$.
Thus $F\le21$ if $u>3$.

Suppose $9\Vert M$, $u>3$, exactly three other primes divide $M$,
and $m_9=2k\ge4$. We claim every denominator-three row is
nonquadratic. Otherwise use its integral traces and put
$f=(F-1)/2\in\{4,7,10\}$. The $k$ free pairs at degree $n/3$
have traces congruent to the same nonzero residue modulo three;
all other free degrees are at most $n/15$. A three-divisible
proper hit has degree at most $n/5$ by height. If none of the
$j\le2$ nonzero three-divisible traces, all of magnitude three,
occurs at $n$, their absolute degree-weighted sum together with
the free terms and vacuum is at most
\[
 1+n\left(3j/5+\lfloor\sqrt{f(25-9j)}\rfloor/3\right)
       <3\sqrt3\,n .
\]
Here $j=2$ requires $f\le7$. Thus a trace of magnitude three
occurs at $n$. With just one such hit, the free squared norm $16$
consists of $e=(16-f)/3$ doubles and $f-e$ singles: all are
nonzero modulo three and none can have magnitude four. If $c$
doubles occur at degree $n/3$, its signed trace sum has magnitude
$|k-3c|$. For $k\le7$ the free terms have total magnitude at most
\[
 \frac n{15}\bigl(5|k-3c|+f-k+e-c\bigr)\le40n/15 .
\]
Indeed the parenthesis has respective maxima $40,34,40$ for
$f=4,7,10$ over $2\le k\le\min(f,7)$ and
$\max(0,k-f+e)\le c\le\min(e,k)$. This cannot cancel $3n$
in vacuum orthogonality. If $k\ge8$, at most two free degrees
besides $n/3$ occur. They each omit at most one allowed prime,
contrary to $\sum d_Yz_Y=-1$.

With two three-divisible hits the free squared norm is seven.
A second proper hit cannot cancel the first, since
$1+(3/5+7/3)n<3n$. The two hits at $n$ must have opposite signs.
At $f=7$ all free traces are singles, and $kn/3$ exceeds
$1+(7-k)n/15$. At $f=4$ there is one double. Unless $k=3$ and
the double is at $n/3$, the signed trace sum there is nonzero;
its weight exceeds the vacuum and remaining free terms, of
total magnitude at most $1+3n/15$. In the exception just one
free pair remains, giving $1\pm d_Y=0$, impossible.
This proves the claim. All twists at denominator three therefore
have order nine by Proposition~\ref{odd:ternary-multiplicity}.

\begin{proposition}[The four-prime ternary bound]\label{new:ternary-cap}
Let $D=3M^2$ be odd with precisely four prime divisors. In non-pointed
based-simple rotation-admissible data, the multiplicity $m_3$ of
denominator three is at most sixteen, and at most twelve if $9\mid M$. If $9\mid M$ and either
$27\mid M$, $m_9\le2$, or $m_3>3$, then $m_3>0$ forces, for every
prime $p\ne3$ dividing $M$, a nonunit dual pair of degrees
coprime to $3p$.
\end{proposition}
\begin{proof}
If $3\nmid M$, denominator three is absent. Corollary~\ref{odd:minimal-ternary} covers $3\Vert M$. Take
$9\mid M$, put $n=M/3$, and write $F$ for the three-free labels,
including the vacuum. Its cardinality is an odd multiple of three
and is at least nine, by deletion and the dimension identity.

First take a nonquadratic denominator-three row $X$. Its field is
$K=\mathbb Q(\zeta_9)$ and its twist has order nine by
Proposition~\ref{odd:ternary-multiplicity}. Write
$u=\chi_X(X^*)/\theta_X^2$. Proposition~\ref{odd:diagonal-twist} gives
$u\equiv1\pmod{2\mathcal O_K}$, and principality gives $u\equiv0$
at the prime above three because $3\mid n$. Thus
$z=(u-1)/2$ and $z+1$ are both nonzero algebraic integers.
Let $C$ denote the Galois mean of squared absolute values.
The norm inequality gives
\[
 C(u)=2C(z)+2C(z+1)-1\ge3.
\]
Equality forces $z$ and $z+1$ to be roots of unity, so
$u=\pm\sqrt{-3}$. At the six conjugate labels of $X^*$ the entries
are Galois conjugates: apply symmetry at equal degree and row
Galois covariance. Their total mean squared cost is at least eighteen.
Each entry on $F$ is nonzero by principality, and costs at least one.
These sets are disjoint and the row has norm $27$. Consequently
$|F|=9$, all its entries on $F$ are roots of unity, $u=\pm\sqrt{-3}$,
and all remaining entries vanish.

Average the row over $\operatorname{Gal}(K/\mathbb Q(\zeta_3))$.
Every entry on the orbit of $X^*$ is the product of $\pm\sqrt{-3}$
and a primitive ninth root, and has average zero. The three rows are
orthogonal of squared norm $27$. Their mean has squared norm nine on
the nine entries of $F$. Since each has absolute value at most one,
the original roots on $F$ are already fixed by this cubic group.
The mean of the dual orbit is orthogonal to the first mean. Thus
$\sum_{Y\in F}\chi_X(Y)^2=0$. For the four free dual pairs, a pair
of real units contributes two and a pair of nonreal sixth roots
contributes minus one. Exactly one of the four pairs has real entries.
Its pair trace has magnitude two and the other three have magnitude
one. Vacuum orthogonality of the averaged row gives
$1+\sum_i d_i z_i=0$ on these four free pairs, where $z_i$ is their
pair trace.

For a quadratic row in the range $27\mid M$ or $m_9\le2$, the
quadratic calculation in Proposition~\ref{odd:ternary-multiplicity}
gives two opposite traces of magnitude three at degree $n$ and
no other nonzero traces outside $F$. The free traces have squared
sum seven and satisfy $1+\sum_i d_i z_i=0$. In the remaining
range $9\Vert M$, $m_9\ge4$ and $m_3>3$, the preceding argument
excludes quadratic rows. In all the stated ranges, reduction of
the free trace equation modulo each $p\ne3$ gives the asserted pair.

Suppose $m_3>12$. If $|F|>9$, every row is quadratic and there
are at least seven free pairs. Seven units of squared trace force
exactly seven pairs with traces of magnitude one. The sign on each
fixed pair is determined by reduction modulo three.
When $|F|=9$, every row instead has one double and three single
traces, in the quadratic case by its free squared trace sum seven.
The $i$th trace is $2\epsilon_i$ or $-\epsilon_i$, with
$\epsilon_i=\pm1$ fixed by principality. If pair $i$ is doubled,
vacuum orthogonality gives
\[
             3\epsilon_i d_i=\sum_j\epsilon_j d_j-1 .
\]
The four nonunit degrees have gcd one by the free trace equation.
Each contains at least two of the three allowed primes, so at least
three degrees are distinct. At most two indices can attain the fixed
value of $\epsilon_i d_i$ required for a double.
Thus two pair traces are constant signs throughout the family.
Lemma~\ref{odd:constant-pair-trace} gives $2m_3<27$: its vacuum
coordinate is nonzero because $1\pm d_{Y_1}\pm d_{Y_2}$ is odd.
Evenness forces $m_3\le12$.
\end{proof}

\begin{lemma}[A minimal ramified fiber]\label{r7:six-fiber}
In odd dimension $D=3M^2$, let $F$ consist of six nonunit labels
at fixed complementary phases. Suppose its coordinate space is
invariant under the local three-factor and under three-inertia and
duality. Then all six degrees agree, and their three-phases are
$1,1,\omega,\omega,\omega,\omega$, or their conjugates.
\end{lemma}
\begin{proof}
On the three charge-even dual-pair lines the local inertia
multiplier $-1$ acts as $-P$, where $P$ is a permutation involution.
A negative-central representation of the local three-group has even
dimension: quaternion elements $A,B$ with
$A^2=B^2=-I$ and $AB=-BA$ interchange the $i$- and $-i$-eigenspaces
of $A$. Thus $P$ is a transposition. Its exchanged pair lines have
equal degree and the same phase. This excludes a primitive ninth-root
orbit and three distinct cube phases. The ramified mass identity
says that the nonneutral squared-degree mass is twice the neutral
mass. If the exchanged pair lines were neutral, the third degree
would be twice theirs, impossible in odd dimension. They are
therefore at the same primitive cube root, and the mass identity
makes the remaining, neutral degree equal to theirs.
For the quaternion elements take
$A=\left(\begin{smallmatrix}0&1\\-1&0\end{smallmatrix}\right)$ and
$B=\left(\begin{smallmatrix}1&b\\b&-1\end{smallmatrix}\right)$,
where $b^2=-2$ modulo the local three-power by Hensel lifting
\cite{SerreLocal}.
\end{proof}

\paragraph{Multiplicity and mass.}
At most one unramified prime $p$ is a denominator, with $p\Vert M$ and
the other two heights at least two (Proposition~\ref{new:two-square-primes}).
Bessel on its effective rows, subtracting the vacuum row from its modular
image, gives
\[
 m_p/(3p^2)\le\sum_{\theta_{Y,3}=1,\ \theta_{Y,p}\ne1}
                     |\theta_{Y,p}-1|^2/b_Y^2.
\]
The averaged squared factor is at most $2p/(p-1)$ and at most a third
of the mass at a nontrivial $p$-phase is three-neutral. With no activity
at $p$ itself this requires $9(p-1)m_p/(2p)$ active labels at $3p$ or
above. For general $t$ replace both factors of three by $t$.

Two joint restrictions on $m_3,m_5$ are needed. Choose a three-free
nonunit hit $Y$ in an effective five-row; it is also five-free and its
orbit meets every row in both families. At $(m_3,m_5)\ge(8,36)$ their
masses sum to more than $3/4$. Their fields are disjoint, so
Proposition~\ref{r5:field-rank} forces one family rational on $Y$ and
$|G Y|\le\max(27/8,25/12)<4$. On this pair the five-values are nonzero
integers, and every ternary pair trace is nonzero modulo three; the
symmetric pair column costs $24/25+m_3/54>1$. At $(m_3,m_5)\ge(6,48)$
the same test forces the three-family rational since the five-family
alone exceeds half. Now $|G Y|\le27/6$ and the five-family forces
$\QQ(\zeta_5)\subset K_Y$ by the pair column bound, so equality of the
fields holds. Every five-trace in this quartic field is $4d_Y\ne0$
modulo five. Projection on the constant orbit line costs at least
$4m_3/27+m_5/300>1$.

The \supfile{odd_nonsquare/census.py}{multiplicity enumeration} checks $m_3,m_9$ and the three prime multiplicities
($m_5\le48$ by the effective bound, divisibility by twelve and the
exclusion of effective multiplicity twenty). A remaining nonvacuum
denominator has floor $15$. After the preceding inequalities every
survivor has a pure three-power row; all three-free labels are then
inactive at the other primes. The deletion bound requires eight
three-free nonunits. A prime of height one forces a dual pair omitting
it and each other prime by principal-block height and gcd one;
different omitted pairs are disjoint. The free trace equation of
Proposition~\ref{new:ternary-cap} also forces the omissions $\{3,p\}$ in its
range. A pair omitting $p,s\in P$ is inactive at every unramified
prime by Theorem~\ref{thm:spherical}; if its denominator is coprime to
three it belongs to six equal-degree labels, so fewer than four spare slots force a two-label witness to have denominator divisible by three. Use the full known powers
of the omitted primes in these floors and the vacuum bound. At
$9\Vert M$ the unrecorded three-free labels have floor $45$; if also
$p\Vert M$, a residual $p$-active label has denominator exactly $3p$
(its support lies in $\{3,p\}$ and at $9p$ it would be three-free).
If $3\Vert M$, a residual active denominator at $p$ is a pure
$p$-power $\ge p^2$, since $3\mid b$ makes its degree three-free.
This excludes residual activity for $p\Vert M$. In this case count
the recorded three-free labels and their reservations against the
bound $|F|\le27$; once exhausted, all further denominators are
prime to three and at least $25$. Credit each position once.

For the initial multiplicity count, put
\[
 s=210-m_3-m_9-\sum_{p\in P}m_p,\qquad
 E=3-\frac{m_3}{9}-\frac{m_9}{81}-\sum_{p\in P}\frac{m_p}{p^2},
 \qquad M_{\min}=3\prod_{p\in P}p .
\]
Besides the preceding multiplicity restrictions, the
\supfile{odd_nonsquare/receipt.json}{census} imposes
\begin{equation}\label{eq:odd-nonsquare-initial-mass}
 s\ge0,\qquad 0<E\le M_{\min}^{-2}+\frac{s}{225}.
\end{equation}
Indeed $M_{\min}\mid M$, and every remaining nonunit denominator
is at least $15$. Counting on the fourteen supports gives
$4\,189\,686$ profiles before the upper mass bound in
\eqref{eq:odd-nonsquare-initial-mass}, and $685\,040$ after it.
The subsequent valuation and reservation tests leave four branches.
Each has $P=\{5,7,11\}$, $M_0=27\cdot5\cdot49\cdot121\mid M$ and
$(m_3,m_5,m_7,m_9,m_{11})=(12,24,0,k,0)$, $k=36,38,40,42$.
Six disjoint unrecorded pairs outside the activity have denominators
divisible respectively by $135,1323,3267,27,245,605$, from the three
omissions $\{3,p\}$, a fourth three-free pair and the principal five-block.
The three activity reservations cost $6(p-1)$ each and leave at most
$42-k\le6$ spare labels. Higher $p$-order costs at least $3p(p-1)$, so
all phases are primitive. Their counts after deleting the five-blocks
are equal: the two square classes agree modulo eight in each phase
and an imbalance needs $4(p-1)$ extra labels. A common increase needs
$2(p-1)$ by duality. Hence every primitive phase has six labels.

Each six-label fiber is stable under three-inertia and duality.
The removed denominator-five labels form complete equal-degree
three-blocks, so their complement remains invariant under the local
three-factor; the complementary phase projectors commute with it.
Lemma~\ref{r7:six-fiber} therefore makes the six degrees equal.
Their denominators are supported on $\{3,p\}$ by
Theorem~\ref{thm:spherical}. Equality of primitive $p$-phase masses
then gives one denominator $B_p=3^i p^j\ge3p$ at each $p$.

Thus the activity mass is $\sum 6(p-1)/B_p^2$. Check $B_p\le500$ and
bound each larger reciprocal square in $[0,501^{-2}]$, reserving the
six pairs, $M_0^{-2}$ for the unit and $1/225$ per spare position.
Only $k=36$, $(B_5,B_7,B_{11})=(15,21,33)$ remains, forcing $r=211$.
There the residual identity is
$5023/266805=M^{-2}+\sum_{i=1}^{9}2/b_i^2$. For residual $E$ with
$h$ pairs left the next sorted denominator satisfies
$\sqrt{2/E}<b\le\sqrt{(2h+1)/E}$. Try all supported $b$ in this interval
with $v_5(b)\le1$, all possible assignments of the six divisors, and at
a leaf require that $M$ be a multiple of $M_0$ and all $b$.
The $72$ nodes have four leaves, none admissible. This excludes nonsquare odd dimension at every rank below $219$.

\subsection{Odd square dimension}

\begin{lemma}\label{odd:nine-support}
Suppose $D=M^2$ is odd and $m$ labels have denominator nine. If
$27\mid M$ or $m\ge9$, their rows are supported on $9\mid b_Y$.
\end{lemma}
\begin{proof}
At a hit, $3\mid b_Y$ by Lemmas~\ref{blk:support} and~\ref{blk:height}.
If $b_Y=3k$ with $3\nmid k$, write
$u=\chi_X(Y)/3=\chi_Y(X)/k$. It is integral by lcm integrality.
If $9\Vert M$, the degree $M/9$ is prime to three, so principality
of $Y$ makes $u$ nonzero in every denominator-nine row.
The sum of their squared moduli is strictly below $81$ for the
unscaled entries, by column unitarity and the vacuum term;
Galois averaging gives $9m<81$.

Now suppose $27\mid M$. Principality gives $u\equiv0$ at three
since $3\mid d_X$, and the two row conductors place $u$ in
$\QQ(\zeta_{81})\cap\QQ(\zeta_{9k^2})=\QQ(\zeta_9)$.
By total ramification, the positive integer
$T=\Tr_{\QQ(u)/\QQ}(u\bar u)$ is divisible by three and
$T\ge h\,3^{2/h}$ for $h=[\QQ(u):\QQ]$.
The conjugates in the different denominator-nine rows cost at
least $T$ in the budget $\sum_X|\chi_X(Y)/3|^2<9$. Thus $h=2$,
$T=6$, and the orbit of $X$ has size two: otherwise it repeats
these conjugates. In particular $K_X=\QQ(\sqrt{-3})$ and $u$ has
norm three. Write $\chi_X(Z)=(z_Z+c_Z\sqrt{-3})/2$ in this
quadratic field. Orthogonality of $X,X^*$ applied to their sum
and difference gives, over the nonunit dual pairs,
$\sum z_Z^2=79$ and $\sum c_Z^2=27$. But the Eisenstein integers
of norm three give $|z_Y|=9$ or $|c_Y|=6$ for $\chi_X(Y)=3u$,
a contradiction.
\end{proof}

\begin{lemma}[Localization of a row field]\label{odd:field-component}
In square dimension the conductor of every row field is supported
on a single component of the twist graph. In odd square dimension let $p$
be a prime denominator, let $j$ be its primitive multiplicity and put
$d=2^{v_2(p-1)}$. Then $d\mid j$. If $j=d$, at least two distinct
proper $p$-free degrees have a label fixed by $p$-inertia.
For every proper $p$-free label $Y$ fixed by that inertia,
\[
      2pj\|d_Y\|_p^2<p^2,\qquad
      |a_Y|\ge \|j d_Y\|_p .
\]
\end{lemma}
\begin{proof}
Let $J$ be a component of odd primes and let $E$ be the $J$-local
span of the nontrivial $J$-phase spaces. As in
Proposition~\ref{new:component-column}, every complementary upper
unipotent and its lower conjugate are scalar on each of these spaces.
The complementary tori are therefore scalar there. Their actions
are unsigned permutations in square dimension, so each scalar equals
one on the coordinate vectors of a nonempty phase space. These tori
fix $E$ pointwise by commutation, and the $J$-groups act trivially
on $E^\perp$. For torus permutations $P,Q$ from the two sides it follows
that $(P-I)(Q-I)=0$. The torus operators are unsigned in square
dimension, and factors outside the dimension support act by scalars,
trivial on the vacuum for these torus elements. If $P,Q$ both move
a basis vector $v_Y$, the $Y$-coordinate of
$(P-I)(Q-I)v_Y$ is at least one. At least one of any two distinct
components consists of odd primes. Hence inertia can move $Y$ in
at most one component, as asserted.

Theorem~\ref{new:one-packet} gives local level $p$, a class of
$pj$ labels and a single neutral orbit of size $j$. By localization,
every moved label has its full Galois orbit under $p$-inertia.
Its field is a nonreal subfield of $\QQ(\zeta_p)$, so its stabilizer
in the cyclic inertia group has odd order. Every such orbit length,
including $j$, is therefore divisible by $d$.
If $j=d$, the neutral class orbit spends $d-1$ of the inertia
budget $2d-2$ from Proposition~\ref{round:neutral-inertia}.
Each further moved degree spends at least $d-1$. Of the three distinct
proper degrees supplied by Lemma~\ref{new:weyl}, at least two
therefore have inertia-fixed labels.

For an inertia-fixed proper label $Y$, every denominator-$p$ row
entry at $Y$ belongs both to its field and to $\QQ(\zeta_p)$,
so is a rational integer congruent to $d_Y$ modulo $p$.
Bessel on the normalized dual-pair sum gives the strict displayed
bound, since the vacuum uses positive mass. The coefficient residue
is \eqref{eq:weyl-cost}.
\end{proof}
\begin{corollary}[Odd part of a prime denominator minus one]
\label{new:odd-prime-profile}\label{odd:four-five-edge}
Let $p$ be a prime denominator in odd square dimension, and write
$p-1=dv$ with $d$ a power of two and $v$ odd. Then $v\ge5$.
If $v=5$ then $j=d$. The proper labels form three distinct degree
classes: one $p$-inertia orbit of size $d$ with coefficient magnitude
one, and two $p$-fixed dual pairs with coefficient magnitude $d$.
In particular thirteen is excluded and eleven has this shape with $j=2$.
\end{corollary}
\begin{proof}
Lemma~\ref{odd:field-component} and Theorem~\ref{new:one-packet}
give $d\mid j\mid p-1$ and $j<p-1$. This excludes $v=1$ and forces
$j=d$ when $v=3$ or $5$. Every further moved orbit likewise has
size divisible by $d$ by Lemma~\ref{odd:field-component}.
At $j=d$ the remaining inertia deficit $d-1$ therefore permits
at most one further moved orbit, necessarily of size $d$.
All $p$-fixed labels have coefficients divisible by $d$ and occur
in dual pairs. Three distinct proper degrees with only fixed labels
cost at least $6d^2$. With a moved orbit they cost at least
$d+4d^2$. Compare to the available norm $d((v-1)d+1)$.
This excludes $v=3$ and gives equality and the asserted shape at $v=5$.
\end{proof}


\begin{proposition}[Inertia support at an isolated prime]
\label{odd:neutral-parity}\label{odd:inertia-eight}
Suppose $D=M^2$ is odd and $p$ is isolated in the twist graph.
Put $d=2^{v_2(p-1)}$. Let $A$ count the $p$-active labels, $Z$ the
$p$-neutral labels moved by $p$-inertia, and $N=A+Z$. Then
\[
 (p-1)\mid A,\qquad d\mid Z,\qquad
 \frac{A}{p-1}\equiv\frac Zd\pmod2,\qquad
                  \operatorname{lcm}(8,2d)\mid N.
\]
The moved sets at different isolated primes are disjoint. If all
$p$-primary twists have order at most $p$, then
$2(p-1)\mid A$ and $2d\mid Z$. Under this local-level hypothesis,
an inactive nonunit label of $p$-power denominator forces
\[
 A\ge d(p-1),\qquad Z\ge2d,
 \qquad\text{at least $d$ inactive pure labels.}
\]
Under the same local-level and inactive-pure hypotheses, at
$A=2(p-1)$ the moved neutral labels are exactly two dual pairs,
both with field $\mathbb Q(\sqrt{-p})$, and $p\equiv3\pmod4$.
\end{proposition}
\begin{proof}
Localization puts each moved row field in a $p$-power cyclotomic
field. It is nonreal, so its stabilizer has odd order and its
orbit length is an odd multiple of $d$. An active stabilizer has
square fixing a nontrivial $p$-power twist, hence has trivial image
modulo $p$; its orbit length is $(p-1)$ times a power of $p$.
The fields of precisely these orbits contain
$K_p=\mathbb Q(\sqrt{(-1)^{(p-1)/2}p})$.
Quadratic incidence makes their number even. This proves the
first three assertions and $2d\mid N$; localization proves
disjointness.

Only $p\equiv3\pmod4$ needs the additional divisibility by eight.
If $N>0$, write its $2k$ orbit lengths as $h_i=2p^{s_i}u_i$,
where $u_i\mid(p-1)/2$ is odd. Lemma~\ref{odd:quadratic-incidence}, with
$K_p=\QQ(\sqrt{-p})$ and $u=M^2/p$, makes
$\Delta=(-1)^k\prod h_i$ a norm from $\QQ(\sqrt p)$.
For a rational norm $p^ev=x^2-py^2$, with $p\nmid v$, the two
summands have valuations of opposite parity; their leading units
give $(v/p)=(-1)^e$. Quadratic reciprocity gives
$(u_i/p)=(-1)^{(u_i-1)/2}$, since $p\equiv1\pmod{u_i}$.
Applying these facts to $\Delta$ yields
\[
 k+\sum_i s_i+\sum_i (u_i-1)/2\equiv0\pmod2.
\]
But $N=2\sum_i p^{s_i}u_i$ is congruent modulo eight to four
times this integer. Thus $8\mid N$, with no restriction on the
local $p$-power level.

At local level $p$ every active orbit has size $p-1$.
Write $2a,2b$ for the label counts per phase in the two square
classes. Equal primitive-phase masses and odd squares modulo eight
give $2a\equiv2b\pmod8$. Thus $A/(p-1)=a+b$ is even, and $2d\mid Z$.
An inactive pure label has a nonreal $p$-field, and its entire
orbit is pure and has at least $d$ labels. The positive number of
moved neutral orbits is even, so $Z\ge2d$. Their inertia deficit
is at least $2(d-1)$; Proposition~\ref{round:neutral-inertia}
gives $A\ge d(p-1)$. At $A=2(p-1)$ the same deficit is at most
two, forcing exactly two quadratic neutral orbits and $d=2$.
\end{proof}

\begin{corollary}[Pure support and neutral rounding]
\label{odd:isolated-dyadic-cost}\label{odd:fourfold-neutral}
In odd square dimension, let $p$ be isolated and absent as a
denominator. With the notation above,
\[
 A\ge d(p-1),\qquad N\ge d(p+1),\qquad
 \#\{i\ne0:b_i\text{ is a $p$-power}\}\ge dp.
\]
If $A<p(p-1)$, there are at least $2d$ moved neutral labels,
and their denominators are divisible by $p^2$.
For every isolated $p\equiv3\pmod4$, whether or not denominator
$p$ occurs, one has $Z\equiv-A\pmod8$. At local level $p$, if
$A/(p-1)\equiv0\pmod4$ and an inactive pure label exists, then $Z\ge8$.
\end{corollary}
\begin{proof}
If $A<p(p-1)$, square covariance and duality exclude higher
$p$-power phases. The argument of Lemma~\ref{odd:inactive-pure},
before its extra hypothesis on $K$, gives an inactive pure label
and excludes every mixed denominator with $p$-part $p$.
Apply the proposition; all active labels are pure and every moved
row has $p$ in its denominator. This gives all three bounds,
including $A+d\ge dp$ pure labels.
Otherwise $A\ge p(p-1)\ge dp$. It also exceeds $d(p-1)$;
as $N$ is a multiple of $2d$ and $(p-1)/d$ is odd,
$N\ge d(p(p-1)/d+1)\ge d(p+1)$. The last assertions are the
proposition's congruence rounded upward from a positive $Z$.
\end{proof}

\begin{corollary}[Inertia at minimal activity]\label{r4:sparse-inertia}
Let $p$ be an isolated prime in odd square dimension. At least
$2(p-1)$ labels are active at $p$, and at least $4(p-1)$ if
$p\equiv1\pmod4$. If precisely $2(p-1)$ labels are active, then
$p\equiv3\pmod4$. They have a common denominator $B=p^a$, with
$M_p=B^2/p$, where $M_p$ is the $p$-part of $M$. Exactly two
$p$-neutral dual pairs are moved by $p$-inertia; each has row
field $\QQ(\sqrt{-p})$.
\end{corollary}
\begin{proof}
Put $d=2^{v_2(p-1)}$. If denominator $p$ occurs, Theorem~\ref{new:one-packet}
and Lemma~\ref{odd:field-component} give $A=(p-1)j\ge d(p-1)$;
otherwise Corollary~\ref{odd:isolated-dyadic-cost} gives the same
bound. This proves both activity lower bounds.

At $A=2(p-1)$ one has $d=2$. Higher-order phases would require
$p(p-1)$ labels, so every primitive phase is one dual pair.
Equality of their masses gives a common denominator $B=p^a$,
and \eqref{eq:sparse-height} gives $M_p=B^2/p$.
An inactive pure label exists by Theorem~\ref{new:one-packet} if
$p$ occurs as a denominator, and Proposition~\ref{new:pure-extra}
otherwise. Proposition~\ref{odd:neutral-parity} now identifies the
two moved neutral pairs and their quadratic fields.
\end{proof}

\begin{proposition}[Partial Weyl columns modulo two]\label{new:weyl-parity}
Let $D=M^2$ be odd. For a set $J$ of dimension primes, let $W_J$
be the partial Weyl operator and $C_J=W_J^2$ its local sign involution.
At every place above two,
\[
              (W_J)_{i0}^2\equiv (C_J)_{i,i^*}.
\]
If $J$ is a twist component, a nonunit coordinate of $W_Jv_0$ is
odd precisely when the conductor of its field belongs to $J$.
Consequently the population of every twist component, counted on
nonunit labels by their fields, is divisible by eight. At a fixed
denominator $b$, specifying the component of its field now fixes
all signs in Proposition~\ref{new:coupled-weyl}.
\end{proposition}
\begin{proof}
The matrices are integral at two: $M$ is odd, the unipotents have
unit diagonal entries, and $s,s^{-1}$ are integral there. Each partial
Weyl matrix is symmetric. Indeed its inverse is the product of the
local upper unipotent, its lower Weyl conjugate and the upper
unipotent again. The middle matrix has form $sUs^{-1}$ with $U$
diagonal commuting with $C=s^2$, hence is symmetric.
All three operators commute with $C$.
In $(W_J^2)_{i,i^*}$ the two terms indexed by $Y,Y^*$ coincide,
and the vacuum term is $(W_J)_{i0}^2$. The local torus at $-1$
acts as the unsigned Galois permutation in square dimension.
By Lemma~\ref{odd:field-component} each nonunit has nonreal field
on just one component, so $C_Ji=i^*$ exactly for the fields on $J$.
The entries of the component vacuum column are rational with odd
denominators. On dual pairs their squared sum is $g_J-g_J^2/2$
by \eqref{eq:component-norm}. Its undivided form gives $g_J=2t$
with $t\in\ZZ_{(2)}$. The sum equals $2t(1-t)$, zero modulo four.
This gives the divisibility by eight. The two possible lifts of a
local coordinate have opposite parity by Proposition~\ref{new:coupled-weyl}.
Finally distinct component vectors on nonunit dual pairs have even
inner product by \eqref{eq:component-cross}; their supports modulo two
are in fact disjoint.
\end{proof}

\begin{corollary}[Compatibility of minimal components]
\label{new:minimal-component-tournament}
For any collection of isolated minimal-activity primes, let $Y_p$
denote the mixed moved neutral pair at $p$, if it exists.  There is
at most one such pair for each $p$.  Set
\[
 c_{pq}=w_q(Y_p)/2\pmod2
\]
and set $c_{pq}=0$ if $Y_p$ is absent.  Then
$c_{pq}+c_{qp}=1$ for distinct $p,q$.
Thus these coefficients orient a tournament; an arrow $p\to q$
requires $q$ in the denominator of $Y_p$.  At least $t-1$ distinct
mixed moved neutral pairs occur for $t$ such primes.
\end{corollary}
\begin{proof}
The indicator formula makes the displayed quotient integral at two.
In the cross sum modulo four only the pairs $Y_p,Y_q$ can contribute.
Their contributions are respectively $2c_{pq},2c_{qp}$, while the
right side is two.  At most one vertex of a tournament has no outgoing
edge.  Field localization makes the pairs attached to different
vertices distinct.
\end{proof}

\begin{proposition}[Residues on the proper support]
\label{new:minimal-proper-residue}
Suppose $p\ge5$ is isolated, $p\Vert M$, and precisely $2(p-1)$
labels are active at $p$.  On the proper $p$-free nonunit dual pairs
let $a_Y$ be the coefficients of Lemma~\ref{new:weyl}.  Then
\[
 \sum_Ya_Y^2=p-2,\qquad
 a_Y\equiv-2\epsilon d_Y\pmod p.
\]
Exactly one coefficient is odd.  Its pair has field
$\mathbb Q(\sqrt{-p})$, and its square is at most $p-10$.
Every other proper pair is fixed by $p$-inertia.  For each such pair
there is an integer $\nu_Y\equiv d_Y\pmod p$ with
$a_Y=-2\epsilon\nu_Y$, and
\[
                   4\|d_Y\|_p^2\le p-7,
\]
where $\|\cdot\|_p$ is the least absolute residue.
The proper degrees divide $M/p$, have gcd one, and share a prime
in pairs.
\end{proposition}
\begin{proof}
Theorem~\ref{new:one-packet} gives $m=2p$ and the square budget.
Corollary~\ref{r4:sparse-inertia} and Proposition~\ref{new:weyl-parity}
give exactly one moved proper pair, with odd coefficient; the
other moved neutral pair is the pure denominator-$p$ pair $X$.
The congruence follows from $a_Y\equiv d_Y/n$ and
$2n\equiv-\epsilon\pmod p$, where $n=M/p$.
The neutral compression formula in the proof of
Lemma~\ref{odd:two-neutral} gives
$\chi_X(Y)+\overline{\chi_X(Y)}=-\epsilon a_Y$ on proper pairs.
For an inertia-fixed $Y$, this character value is a rational
integer $\nu_Y\equiv d_Y\pmod p$ by principality, proving
$a_Y=-2\epsilon\nu_Y$.

Lemma~\ref{new:weyl} gives the stated support properties and at
least three distinct proper degrees. Besides a selected even pair,
an odd pair and another even pair cost at least $1+4$ in the
square budget. Besides the odd pair, two even pairs cost at least
eight. These are precisely the two asserted bounds.
\end{proof}

\begin{proposition}[Character values at minimal activity]\label{rev:minimal-packet}
Suppose $D=M^2$ is odd, $p\ge7$ is isolated, $p\Vert M$, and
precisely $2(p-1)$ labels are active at $p$. Use $n=M/p$, the sign
$\epsilon$ and the proper-pair coefficients $a_Y$ of
Proposition~\ref{new:minimal-proper-residue}. For any active $X$,
\[
 \chi_X(Y)=-\epsilon a_Y/2\quad(a_Y\text{ even}),\qquad
 \chi_X(Y)=(-\epsilon a_Y\pm\sqrt{-p})/2\quad(a_Y\text{ odd}).
\]
The two signs occur on the members of the unique odd pair. Moreover
$\chi_X(X^*)/\theta_X^2=\sum_{i=1}^{p-1}c_i\zeta_p^i$, where
$c_i\in\{1,-1\}$ and $\sum_i c_i$ is the even representative of
$n\bmod p$ of magnitude $(p+1)/2$. On every other Galois orbit of
degree $n$, each character value $\chi_X(Y)$ has
$C(\chi_X(Y))=(p+1)/4$.
\end{proposition}
\begin{proof}
Corollary~\ref{r4:sparse-inertia} and Theorem~\ref{new:one-packet}
give $p\equiv3\pmod4$, denominator $p$ on the active labels and
$2p$ labels of degree $n$. Local level $p$, field localization and
Proposition~\ref{odd:diagonal-twist} give $K_X=\QQ(\zeta_p)$,
so $G X^*$ has $p-1$ labels. Since $2n\equiv-\epsilon\pmod p$,
the even representative $b$ of $n$ has magnitude $(p+1)/2$.
Corollary~\ref{rev:diagonal-cost} charges at least
$p(p-1)-(p+1)^2/4$ to this orbit. Every value on the other $p+1$
degree-$n$ labels reduces to $n$ above $p$. Summing
Lemma~\ref{nt:residue-norm} on these orbits gives at least
$(p+1)^2/4$. Subtracting both bounds and the vacuum from $f_X=p^2$
leaves at most $(p-1)/2$ for the proper $p$-free nonunit pairs,
counting one squared modulus per pair.

The trace identity \eqref{eq:packet-trace} gives
$2\Re\chi_X(Y)=-\epsilon a_Y$. On an inertia-fixed pair the value
is rational by field localization and symmetric Galois covariance.
On the unique moved proper pair it belongs to $\QQ(\sqrt{-p})$.
Its odd trace implies a nonzero odd coefficient of $\sqrt{-p}/2$.
The proper pairs therefore cost at least
$(\sum_Ya_Y^2+p)/4=(p-1)/2$ by
Proposition~\ref{new:minimal-proper-residue}.
Equality is forced throughout. This identifies the imaginary
coefficient, saturates each of the preceding orbit bounds and
supplies the signed expansion by the equality statement in
Lemma~\ref{rev:two-residue}.
\end{proof}

\begin{corollary}[A quadratic identity in the fusion algebra]\label{new:packet-square}
In Proposition~\ref{rev:minimal-packet}, let $Y,Y^*$ be the unique
proper pair moved by $p$-inertia. Put $I=\{X:d_X=M/p\}$,
$e_I=\sum_{X\in I}e_X$, and $u=[Y]-[Y^*]$. Then
$u^2=-p e_I$ and $u^3=-pu$.
\end{corollary}
\begin{proof}
On each of the $2p$ rows $X$ in $I$, the pair has odd trace:
use \eqref{eq:packet-trace} on the active rows and the proof of
Lemma~\ref{odd:two-neutral} on the neutral pair.
Its entry lies in $\QQ(\sqrt{-p})$, so
$|\chi_X(Y)-\chi_X(Y^*)|^2\ge p$. Character orthogonality gives
$\sum_X|\chi_X(Y)-\chi_X(Y^*)|^2/(2f_X)=1$. The $2p$ indicated
rows already exhaust this sum since $f_X=p^2$. Thus $u$ takes
values $\pm\sqrt{-p}$ on $I$ and zero elsewhere. Evaluation on all
characters gives both identities.
\end{proof}



\begin{proposition}[Rational values on a dual pair]\label{rev:rational-pair}
Let $Y\ne Y^*$ be a nonunit and $I$ a set of nonunit rows in the
principal $p$-block, for an odd prime $p$. Suppose $\chi_X(Y)$ is
rational for every $X\in I$. If $a,b\in\ZZ$, $b\ne0$, $2b\nmid a$
and $a+2b d_Y\equiv0\pmod p$, then
\[
                \sum_{X\in I} f_X^{-1}<\frac{a^2+2b^2}{p^2}.
\]
In particular, in odd square dimension, if $m_9\ge28$, every
nonunit three-free $Y$ has $\QQ(\sqrt{-3})\subseteq K_Y$. In particular
local inversion at three moves $Y$, and its row field lies in the
twist component containing three.
\end{proposition}
\begin{proof}
Column orthogonality gives $a^2+2b^2$ for the total squared norm
of the character evaluations of $a1+b(x_Y+x_{Y^*})$, weighted by
$1/f_X$. At $X\in I$ its value is $a+2b\chi_X(Y)$, a nonzero
integer divisible by $p$ by principality and $2b\nmid a$. The
vacuum also contributes positively. For denominator nine use $p=3$,
$b=1$ and $a=\pm1$ with $a+2d_Y\equiv0\pmod3$. The nine-rows are principal at three and have fields in $\QQ(\zeta_{81})$.
If every $t_X=\chi_X(Y)$ there were real, then $a+2t_X$ would be
one modulo two, zero above three and fixed by the full tame inertia
at three (complex conjugation). Lemma~\ref{rev:two-residue} with
$h_3=2$ gives $C(a+2t_X)\ge9$. Summing over the Galois-stable
class again gives $m_9<27$. Thus a nonreal $t_X$ occurs. Its field
in the cyclic ternary cyclotomic field has even degree and contains
the unique quadratic subfield. Symmetry puts it in $K_Y$. Apply
Lemma~\ref{odd:field-component}.
\end{proof}

\paragraph{A uniform finite residue relaxation.}
First suppose $p$ is isolated, $p\Vert M$, and exactly $2(p-1)$
labels are active at $p$, as in
Proposition~\ref{new:minimal-proper-residue}. Put $P=\operatorname{supp}(M)$.
For a required proper $p$-free degree avoiding $Q$ and having field
in another component, enumerate supports
$T\subseteq P\setminus(Q\cup\{p\})$ with $|T|\ge2$ and residues
$\prod_{q\in T}q^{e_q}\bmod p$. A known valuation upper bound
restricts $1\le e_q\le v_q(M)$; otherwise use the full cyclic
subgroup generated by $q$. Every actual degree is included, and
some residue $r$ must satisfy $4\|r\|_p^2\le p-7$.
Failure after imposing a trial cap $e_q\le h$ proves $v_q(M)>h$.
If $m_9\ge28$, Proposition~\ref{rev:rational-pair} supplies this test
with $Q=\{3\}$ at every such $p$ outside the three-component,
since the proper degrees have gcd one.

For the whole budget, select one odd pair and any even pairs whose
supports intersect pairwise and have empty total intersection.
Their residue costs are respectively $\|2r\|_p^2$, allowed only
when that least residue is odd, and $4\|r\|_p^2$.
The minimum total is at most $p-2$. Choosing supports and residues
independently only weakens the necessary condition.

More generally, at any set $P_1$ of isolated valuation-one primes,
the primitive multiplicities satisfy
\[
 j_p\prod_{q\ne p}q^{v_q(M)}\equiv-\epsilon\pmod p
 \qquad(p\in P_1),\qquad \epsilon\in\{1,-1\},
\]
with one common sign by Lemma~\ref{new:weyl} and
Proposition~\ref{new:component-column}. Remove factors of known
valuation in $\prod_{p\in P_1}(\mathbb Z/p\mathbb Z)^\times$.
Each unknown exponent ranges over the cyclic subgroup generated
by $(q\bmod p)_{p\in P_1}$. At least one sign target must belong
to the product of the allowed sets. Trial exponent intervals give
valuation lower bounds in the same way. Resource limits retain
undecided cases.

\begin{lemma}[Two columns in the Gauss sum]\label{odd:centered-column}
Let $O$ be a Galois orbit of $h$ nonunit labels of common codegree $f$,
and put $u_X=h^{-1}\sum_{Y\in O}\chi_X(Y)$. Choose a proper set $F$
of rows. On its complement put
\[
 m=\sum_{X\notin F}\frac1{f_X},\quad
 \bar u=\frac1m\sum_{X\notin F}\frac{u_X}{f_X},\quad
 V=\sum_{X\notin F}\frac{(u_X-\bar u)^2}{f_X},\quad
 L=\sum_{X\notin F}\frac{\theta_X}{f_X}.
\]
At every cyclotomic embedding one has
\[
 \left|\sum_{X\in F}\frac{u_X\theta_X}{f_X}+\bar u L\right|
 \le f^{-1/2}+\sqrt{V\left(m-|L|^2/m\right)} .
\]
\end{lemma}
\begin{proof}
The orbit sums are rational by symmetry and Galois covariance;
column orthogonality gives $\sum_Xu_X/f_X=0$ and
$\sum_Xu_X^2/f_X=1/h$. The vacuum modular relation expresses
$\sum_Xu_X\theta_X/f_X$ as the mean of $h$ roots of unity
times a scalar of modulus $f^{-1/2}$ at every embedding.
On the complement of $F$ subtract the weighted means from $u_X$
and $\theta_X$. Cauchy--Schwarz bounds the resulting covariance
by $\sqrt{V(m-|L|^2/m)}$, as asserted.
\end{proof}

\begin{proposition}\label{odd:eleven-nine}
Under the hypotheses of Theorem~\ref{alg:rank}, suppose $D=M^2$
is odd, $11\Vert M$, and eleven is isolated. Suppose the complete
class of denominator eleven has $22$ labels, twenty active, and
contains all eleven-activity. Then at most $62$ labels have denominator nine.
\end{proposition}
\begin{proof}
Write $s=m_9\ge64$. Corollary~\ref{new:odd-prime-profile} gives three proper pairs of distinct degrees and coefficient magnitudes $2,2,1$. Their degrees have gcd one, and hence one pair $Y$ is three-free. Its quadratic field is $\QQ(\sqrt{-3})$ by Proposition~\ref{rev:rational-pair}. The eleven-row entries at $Y$ are rational, so $|a_Y|=2$ and their common sign is denoted by $\nu$. The pair $Z$ with $|a_Z|=1$ has field $\QQ(\sqrt{-11})$.

The nine-row entries at $Y$ are Eisenstein integers of the same
nonzero residue $\eta=\pm1$ at three. All are units: a nonunit
in this residue class costs at least four in squared norm and
occurs in a dual pair, whereas $1-s/81-2/11<6/81$. Let $t$
count the real entries. The other entries have real part $-\eta/2$.
Projection on the antisymmetric and symmetric unit vectors of the
$Y$-pair and the residual Gram determinant on the columns $0,Y$ give
\[
 \frac{3(s-t)}{162}\le1,\qquad
 \frac{s+3t}{162}+\frac4{11}<1,\qquad
 \left|\frac{(3t-s)\eta}{162}+\frac{2\nu}{11}\right|
 <1-\frac{s}{81}-\frac2{11}.
\]
The vacuum makes the second inequality strict. In the last,
equality would make the two residual columns proportional with equal
norms, impossible at their vacuum coordinates since $d_Y>1$.
As $s,t$ are even, the only solution with $s\ge64$ is
$(s,t)=(64,12)$ with $\eta=\nu$. Multiply column $Y$ by this sign
and orient column $Z$ by a sign so its eleven-row trace is $-1$.

Its eleven-row entries now have the form $(-1+c\sqrt{-11})/2$
with $c$ odd. On the antisymmetric unit vector of the $Z$-pair
they already use norm squared at least $(2/11)(11/2)=1$.
Thus $c=\pm1$ in every entry. The degree of $Z$ is divisible by
three: otherwise the same field argument would give
$K_Z=\QQ(\sqrt{-3})$. Its nine-row entries are therefore integers
divisible by three, by field intersection and principality. Write $3k$ for their value on each dual pair,
$E=\sum k^2$ and $K=\sum k$. After removing both classes,
the Gram matrix $G$ of $0,Y,Z$ is real, with
\[
 \begin{gathered}
 G_{00}=G_{YY}=\frac{25}{891},\qquad G_{0Y}=-\frac8{891},\qquad
 G_{ZZ}=\frac5{11}-\frac{2E}9,\\
 G_{0Z}=\frac1{11}-\frac{2K}{27}.
 \end{gathered}
\]
Hence $E\le2$. At $E=2$, $K$ is even of absolute value at most
two, and the $0,Z$ minor is negative. At $E=0$, $G_{YZ}=1/11$
and the determinant on $e_0+e_Y,e_Z$ is $-154/9801$.
At $E=1$ the $0,Z$ minor forces $K=1$.
Thus exactly one nine-pair has $Z$-entry three and the rest vanish.
In particular $9\mid b_Z$ by Lemma~\ref{odd:nine-support}, so
$99\mid b_Z$.

Apply Lemma~\ref{odd:centered-column} to the $Z$-pair with $F$
the nine- and eleven-rows, using the chosen sign for its real column.
It takes the value $3$ on one nine-pair, zero on all other nine-rows
and $-1/2$ on the eleven-rows. Thus
\[
 m=25/891,\qquad m\bar u=1/11-2/27=5/297,\qquad
 V=1/2-2/9-1/22-m\bar u^2=2/9.
\]
The eleven twists sum to zero. The nine twists sum to twice a
nonzero cyclotomic integer: the $32$ dual pairs have three-power
twists, so the half-sum reduces to $32$ above three.
At some embedding this integer has modulus at least one.
The vacuum Gauss relation therefore gives
$2/81-1/M\le |L|\le m$ at that embedding. If $\theta$ is the
twist on the distinguished pair, the known contribution of the
oriented column is $2\theta/27$. Since $b_Z\ge99$ and the odd nonunit degree $d_Z$ has at least
two prime factors, $M=b_Zd_Z\ge99\cdot15>1155$. Thus
the lemma would imply
\[
 \frac{14}{297}=\frac2{27}-\frac35m-\frac1{99}
 \le\sqrt{\frac29\left(m-\frac{(2/81-1/1155)^2}{m}\right)} .
\]
The left square exceeds the right square by
$1303828/2701400625>0$. This excludes $s=64$.
\end{proof}


\paragraph{The twist graph.}
We use $209$ as a common rank ceiling. We also record which graph restrictions extend to the next possible rank, $217$, since $r\equiv1\pmod8$. Write $P$ for the prime support of $M$. Join two primes if some twist
is active at both. A twist active at three primes forces rank at
least $385$. Indeed, choose three such primes $p,t,u$, taking
$p=3$ when possible.
Equality of primitive phase masses, or the square-unit orbit at a
higher phase, gives at least $(t-1)(u-1)$ disjoint complementary
fibers meeting $p$-activity; the local variations are independent.
If $p=3$, Proposition~\ref{phase:moments} puts eight dual pairs
in each fiber, giving at least $16\cdot4\cdot6=384$ nonunits.
If all three primes are at least five, the same proposition gives
at least $4(p-1)(t-1)(u-1)\ge960$ nonunits. Add the vacuum.
Thus no twist in the range considered here has three prime divisors.
Corollary~\ref{new:sector-rank} now gives
\[
 \sum_{pt\ {\rm edge}}pt+\sum_{p\ {\rm isolated}}p\le (r+1)/2.
\]
For each edge $pt$ it reserves $(p-1)(t-1)$ pairs active at both
ends, $p-1$ active at $p$ alone and $t-1$ active at $t$ alone, all of
denominator divisible by $pt$. An isolated vertex supplies $p-1$
active pairs of pure denominator by Theorem~\ref{thm:spherical}.

The prime classes are specified by Theorem~\ref{new:one-packet}
and Corollary~\ref{new:odd-prime-profile}. At denominator $p$,
\[
 m=pj,\qquad a=(p-1)j,\qquad
 d\mid j\mid p-1,\quad j<p-1,\quad (p-1)/d\ge5,
 \qquad d=2^{v_2(p-1)}.
\]
The prime is isolated, $p\Vert M$, its local level is $p$, and its
$j$ neutral labels form one Galois orbit. In particular all prime
denominators below eleven, and thirteen, are excluded.

We give the details of the finite check to specify its completeness.
Take the largest rank $R$ below the bound with $R\equiv1\pmod8$.
Choose $P$, with $|P|\ge3$ and $2\sum_{p\in P}p-1\le R$,
the prime-denominator classes just described, their graph, and the
even multiplicity $s\le76$ at denominator nine ($78$ is the
exceptional value $81-3$ in Proposition~\ref{sq:saturation}).
Reserve the vacuum and the full classes. Every other denominator is
at least $15$ and is supported on $P$. When bounding the mass of
a reservation, use the least such permitted denominator, or any
smaller lower bound. In the three-prime case the two non-three primes
$p$ also need distinct, mutually disjoint, $p$-free supports of
the nine-row sums, all of three-primary twist. For $s\ge9$ their
denominators are multiples of $9p^2$ by Lemmas~\ref{odd:nine-support}
and~\ref{blk:witness}. Proposition~\ref{sq:joint-projection} gives their sizes with unit weights:
every such degree is divisible by three. Subtract these positions
separately from the active reservations away from three.
At a valuation-one class the proper $p$-free labels in
Lemma~\ref{new:weyl} give at least three dual pairs of denominator
at least $3p$. With four support primes these may be chosen to
omit one distinct other prime $t$ each, so their denominators
are at least $p t^{v_t(M)}$. An isolated $t$ with no denominator
$t$ has $t^2\mid M$, since some $t$-active denominator is a pure
$t$-power. All proper pairs are neutral, since their twists are
$p$-primary and every $p$-active label has denominator $p$. Except where
disjointness has been shown, reservations are tested separately.

If $|P|=4$ and $s\ge72$, the $f=81$ case of
Corollary~\ref{odd:minimal-ternary} excludes $9\Vert M$, so $27\mid M$. Choose the vacuum and four complete nonunit three-free
dual pairs, using the deletion bound and the degree sum modulo three.
These nine columns are disjoint from the nine-class, and
$9(81-s)\le81$. Lemma~\ref{new:many-reserve} allows all nine to
be discarded from spherical summands avoiding that class. Every
edge summand and every isolated-prime summand other than the one
for three avoids it. In the charge-conjugation fixed space this gives
\[
 r\ \ge\ s+9+
 2\left(\sum_{pt\ {\rm edge}}pt+
             \sum_{p\ne3\ {\rm isolated}}p\right).
\]



In odd square dimension, Proposition~\ref{odd:edge-budget} gives $4(p-1)(t-1)$ jointly active labels on an edge $pt$:
there are at least $t-1$ occupied nontrivial $t$-phases, by
Lemma~\ref{blk:moments} or the size of a higher orbit. For a leaf
$t$ adjacent to three, at least $16(t-1)$ mixed-denominator labels
are $t$-active. Distinct leaves use disjoint such labels.

We describe the remaining finite phase test. For an edge $3t$,
$t=5,7$, take $B=3t$ and measure weights in units of $B^{-2}$.
Every mixed denominator above $B$ is at least $3B$. In a phase at
one endpoint let $p$ be the other, $k,h$ the joint and solely active
pair counts, and $j,v$ their respective counts at $B$. Take $c,u$ to
be the count and the number at $B$ in one primitive $p$-phase in this
fiber. Then $u\le C\le u+(c-u)/9$ and $J\le j+(k-j)/9$. The $v$
small solely active coefficients are multiples of $p$. Their sum $A$
and squared sum $H$ satisfy $H\le v+(h-v)/9+J$. The remaining $h-v$
coefficients must complete their two sums to $h-c,h+k$ modulo eight.
Cauchy--Schwarz on those same $h-v$ positions gives
\[
 (v-A-C)^2\le\frac{h-v}{9}\bigl(J-v+2A+2C-H\bigr).
\]
Indeed their linear sum with weights is $V+v-A-C$ and their weighted
squared sum is $V+v+J-H$, where their weight $V\le(h-v)/9$.
\begin{lemma}[Parity in a primitive fiber]\label{odd:parity-lattice}
In the notation of Proposition~\ref{odd:edge-budget}, suppose all active
phases have order $p$, and write $a_i=pk_i$. Let $E$ be the number
of even $k_i$, and let $W_{\rm e},W_{\rm o}$ be the total weights
on even and odd $k_i$, respectively. Then $4\mid E$ and
\[
 2p|V-C|\le V+J+p^2W_{\rm o},\qquad V=W_{\rm e}+W_{\rm o}.
\]
If $E=0$, all $a_i=-p$ and $C=(p+1)(W_{\rm e}+W_{\rm o})$.
\end{lemma}
\begin{proof}
Here $J=(p-1)C$. Eliminating $C$ from the two moments gives
$\sum_i(pk_i-1)(k_i+1)w_i=0$. Since every $w_i$ is one modulo
four, and the integer coefficient is $-1$ modulo four for even
$k_i$ and zero for odd $k_i$, one has $4\mid E$.
On odd integers the coefficient is nonnegative and vanishes only
at $k_i=-1$, proving the assertion when $E=0$.
For either sign, $(a_i\pm p)^2\ge p^2$ for even $k_i$,
whereas this square is nonnegative for odd $k_i$. Summing and
substituting the two moments proves both sides of the inequality.
\end{proof}

\begin{proposition}[Local signs in a complementary phase]\label{new:sign-fiber}
Assume odd square dimension. Fix a complementary phase fiber active at a prime $t\ne p$.
Count labels in dual pairs. Write $k$ for the number of pairs active
at $p$. Among the $p$-inactive pairs with $p\mid b_i$ let $E$ have
even coefficient $a_i$ in Proposition~\ref{phase:moments}. Put
$d=2^{v_2(p-1)}$. At each degree and fixed coefficient, the even
coefficients occur in multiples of $d$ pairs. At each degree and
tuple of phases, jointly $p,t$-active labels occur in multiples of
four (two pairs). If $k>0$ and all active $p$-phases in the fiber have
order $p$, then
\[
             k\ \ge\ (p-1)\bigl(1+E(1-1/d)/2\bigr).
\]
\end{proposition}
\begin{proof}
Let $K=C_{p'}$, the product of central signs away from $p$. It preserves this fiber by square-unit covariance. If $Ki$ belonged to the $p$-inertia orbit of $i$, some product of $K$ with a $p$-inertia element would fix $i$. This contradicts Proposition~\ref{odd:diagonal-twist}: such an element inverts the nontrivial $t$-part of $\theta_i$. At a joint label the two local signs are therefore independent.
As $W_p^{-1}sv_0=W_{p'}v_0$ up to the normalized scalar, the
coordinate test of Proposition~\ref{new:weyl-parity} says that $a_i$ is odd exactly when $C_pi=i$.
An even coefficient has stabilizer of odd order in cyclic
$p$-inertia, and therefore an orbit of size divisible by $d$.
The coefficients and phases are constant on this orbit for
$p$-inactive labels; the degree is constant as well. The disjoint
translate under $K$ doubles its size.

In the $K$-fixed subspace of the fiber use normalized $K$-pair sums.
Its active dimension is $k$. The vector obtained by projecting
$W_p^{-1}sv_0$ to the fiber is fixed by $K$ and the local Borel;
its Weyl transform has a nonzero coordinate at every active phase. The local torus fixes $se_0$ and conjugation by $W_p$ inverts the torus, so the projected vector is torus-fixed.
With local unipotent order $p$, the representation on this fiber
factors through $\mathrm{SL}_2(\mathbb F_p)$: the upper unipotents
at level $p$ and their conjugates generate the principal congruence
kernel (successively use the conjugates of $E_{12}$ in each graded
copy of $\mathfrak{sl}_2(\mathbb F_p)$). The Borel-fixed vector thus
supplies a Steinberg summand. On the neutral coordinates at least
$E(1-1/d)$ dimensions have nontrivial torus character, by the
orbit bound just proved. The principal-series bound in
Proposition~\ref{round:neutral-inertia}, applied to this representation,
gives the inequality.
\end{proof}
When $p\equiv1\pmod4$ and only primitive $p$-phases occur, in fact
$8\mid E$. In the eliminated moment
$\sum_i(pk_i-1)(k_i+1)w_i=0$ with $a_i=pk_i$, each odd $k_i$
contributes zero modulo eight. Each even $k_i$ contributes $-1$
modulo four, and equal coefficients at equal weights occur in
even-sized sets. Their total is therefore $-E$ modulo eight.

The finite phase test retains this parity information together with
its two moments. Suppose $v$ of the $h$ solely-active pairs have
smallest denominator $B$. Their coefficient sum, squared sum and
even count are $A,H,e$. For each $E\in4\ZZ\cap[e,e+h-v]$, the
large pairs satisfy
\[
 \begin{aligned}
 9(p+1)u&\le (p+1)^2(h-v)-p^2(E-e)
       +9\bigl((2p+1)v-H-2pA\bigr),\\
 0&\le(3p-1)(8u+c)+9\bigl(2pA-H-(2p-1)v\bigr)
          +(p-1)^2(h-v-E+e).
 \end{aligned}
\]
Apply the two square inequalities to these pairs, subtracting the
exact small moments. Use their weight bound $1/9$ and, in the
reflected inequality, discard the negative even-weight term.
When $E=0$, instead retain the exact identity $C=(p+1)W$;
thus $c\equiv(p+1)h\pmod8$ and the intervals
$[u,u+(c-u)/9]$ and $(p+1)[v,v+(h-v)/9]$ must meet. A lower
endpoint is strict whenever the corresponding set has a larger
denominator. In particular this recovers the rigidity for $h<4$.


A higher-order orbit costs at least $p(p-1)/2$ pairs in addition to
the primitive phases. If $c>0$, every primitive phase count is
positive and congruent to $c\bmod8$, so those phases together
require at least $c+(p-2)(1+((c-1)\bmod8))$ pairs.



There is also a height test. The squared degrees at the largest
$p$-height sum to zero modulo $p^2$ after normalization, by
\eqref{eq:phase-column}. If fewer than three denominators exceed $B$
for $p\equiv3\pmod4$, or fewer than two otherwise, no higher
$p$-height can occur. The small terms and inverse even powers of
the other prime must then sum to zero modulo $p^2$.
For $p=3$, fewer than three larger pairs therefore make every
larger relative denominator at least five; use $1/25$ in place of
$1/9$ in the corresponding moment bounds.

\begin{lemma}[The cost of an odd edge]\label{odd:edge-reservation}
Suppose the dimension is odd and square and no twist is active at
three distinct primes. For an edge $pt$ with $p,t\ge5$, at least
\[
 \min\{8(p-1)(t-1),\,
 4(p-1)(t-1)+2(t-1)\mu_p+2(p-1)\mu_t\}
\]
labels have denominator divisible by $pt$, where
$\mu_p=1,3,5$ according as $p\equiv1\pmod8$, $p\equiv5\pmod8$,
or $p\equiv3\pmod4$. In particular, an edge $5$--$7$ requires
at least $172$ such labels.
\end{lemma}
\begin{proof}
First suppose every joint phase has order $pt$. The primitive phase
weights are equal in both directions, so all $(p-1)(t-1)$ phase
counts, in dual pairs, are congruent modulo eight. They are positive
and even by Proposition~\ref{odd:edge-budget}. Galois symmetry makes
each count constant on its square orbit, of size
$(p-1)(t-1)/4$. Unless every count is two, there are at least
$8(p-1)(t-1)$ joint labels: increasing the common residue from two
to four doubles the count, and increasing one square orbit by eight
has the same cost. If a joint phase has order divisible by $p^2$,
there are at least $p(p-1)/2$ complementary $p$-phases, each
containing at least $2(t-1)$ joint pairs. This already gives
$2p(p-1)(t-1)\ge8(p-1)(t-1)$ labels; similarly at $t^2$.

It remains to treat two pairs at every primitive joint phase.
In a fixed primitive $t$-phase, the $p$-fiber has
$k=2(p-1)$ active pairs and primitive count $c=2$. Let $h$ be
its inactive pairs with denominator divisible by $p$.
If $h<4$, Lemma~\ref{odd:parity-lattice} gives $E=0$, so every
coefficient is $-p$ and $2\equiv(p+1)h\pmod8$. This gives $h\ge1$ for
$p\equiv1\pmod8$, $h\ge3$ for $p\equiv5\pmod8$, and excludes
$h<4$ for $p\equiv3\pmod4$. In the last case $h=4$ is also
impossible: the two moments modulo eight give
$\sum a_i\equiv2$ and $\sum a_i^2\equiv0$.
Four squares summing to zero modulo eight must all be even;
then $\sum a_i^2\equiv2\sum a_i\equiv4$, a contradiction.
Thus $h\ge\mu_p$. These pairs, the corresponding pairs with
$p,t$ interchanged, and the joint labels are disjoint and all have
denominator divisible by $pt$. Their sum is the asserted bound.
\end{proof}

\begin{proposition}[A ternary edge]\label{odd:edge35-reservation}
Suppose the dimension is odd and square and no twist is active at
three distinct primes. Let $t\ge5$ be prime and suppose $3$--$t$
is an edge. At least $18(t-1)$ labels are active at $t$ and have
denominator divisible by $3t$. Moreover at least
\[
 C_t=18(t-1)+4\eta_t,\qquad
 \eta_t=\begin{cases}
 1,&t\equiv1,3\pmod8,\\
 2,&t\equiv5\pmod8,\\
 5,&t\equiv7\pmod8
 \end{cases}
\]
labels are active at three or $t$ and have denominator divisible by $3t$. In particular,
$C_5=80$, $C_7=128$, $C_{11}=184$ and $C_{13}=224$.
The first reservations are disjoint for different neighbors of three.
\end{proposition}
\begin{proof}
Use the primitive fiber congruences established above. If every
joint primitive phase has $c=2$ pairs, a $t$-fiber requires at least
$\mu_t$ solely-three pairs, with $\mu_t$ as in
Lemma~\ref{odd:edge-reservation}; a three-fiber requires at least
five solely-$t$ pairs. Thus the joint, solely-three and solely-$t$
counts are at least
\[
       \bigl(8(t-1),\ 4\mu_t,\ 10(t-1)\bigr).
\]
For $c=4$ they are at least
\[
       \bigl(16(t-1),\ 4\nu_t,\ 2(t-1)\bigr),
 \qquad
 \nu_t=\begin{cases}
 2,&t\equiv1\pmod4,\\
 1,&t\equiv3\pmod8,\\
 5,&t\equiv7\pmod8.
 \end{cases}
\]
Indeed, below four solely-active pairs the coefficients are all
$-t$, so $4\equiv(t+1)h\pmod8$. At $t\equiv7\pmod8$ this
also excludes four odd coefficients; four even coefficients would
give $\sum a_i^2\equiv2\sum a_i\pmod8$, whereas the fiber
congruences give respectively four and zero.

The primitive phase counts are positive even integers, congruent
modulo eight, and constant on square orbits of size $(t-1)/2$.
Below $24(t-1)$ joint labels the only additional possibility is
residue two with one square orbit raised by eight. Its total joint
count is $16(t-1)$ and the first set of solely-active lower bounds
still applies. Thus all primitive cases give the two assertions,
since $\eta_t=\min(\mu_t,\nu_t)$.

A single minimal order-$9t$ family has $12(t-1)$ joint labels.
Its three ninth-root phases require $6\mu_t$ solely-three labels.
At each of the $t-1$ other phases the three-fiber has $k=6,c=0$.
Its congruences $\sum a_i\equiv h$ and
$\sum a_i^2\equiv h+6\pmod8$ exclude $h=0,1,2$.
Hence there are at least $6(t-1)$ solely-$t$ labels. Any larger
single family or two ninth-root families already use at least
$24(t-1)$ joint labels. A primitive family together with a ninth-root
family uses at least $20(t-1)$ joint labels and has at least one
solely-three pair at each of its five occupied three-phases.
These counts imply the assertions also when $t\equiv7\pmod8$.
An order at least $27$ at three requires at least $36(t-1)$ joint
labels. An order at least $t^2$ at $t$ requires at least
$4t(t-1)$, already sufficient for both assertions.

Finally, every counted label active at $t$ is inactive at any
other neighbor of three, since triple activity is excluded.
\end{proof}

\begin{proposition}\label{odd:large-ternary-edge}
In the odd square case, three has no neighbor other than five in
the twist graph through rank $217$.
\end{proposition}
\begin{proof}
There is no triple activity below rank $385$.
Proposition~\ref{odd:edge35-reservation} excludes neighbors at least
thirteen. For a neighbor eleven, its $184$ mixed positions and the
absence of denominators five and seven give
$105^{-2}+184/33^2+32/9^2<1$ as an upper bound for the entire mass.
Here denominator five is excluded by
Corollary~\ref{odd:four-five-edge}, and seven by the prime-class
bounds preceding that corollary.

Suppose seven is a neighbor. An additional edge $5$--$7$ costs
$172$ positions disjoint from at least $68$ three-active positions
in the $3$--$7$ sector. These reservations are taken in the
witnessing sectors with all complementary phases trivial.
The count $68$ is $48+20$ in the minimal primitive family; all
other joint families already have at least $72$ labels.
An edge $3$--$11$ costs $108+180$ disjoint leaf-active positions;
larger additional edges also exceed the budget. For the remaining
pair of edges $3$--$5$, $3$--$7$, reserve $72$ five-active,
$108$ seven-active and four centrally active positions, with floors
$15,21,15$. The mass is at most
$105^{-2}+72/225+108/441+4/225+32/81=3881/3969<1$.
Thus only the single edge $3$--$7$ remains.

An outside isolated prime $q$ requires
$A_q=4(q-1)$ active labels if $q\equiv1\pmod4$, and $A_q=2(q-1)$
otherwise. Without denominator $q$ their floor is $q^2$; a present
prime class has at least $4(q-1)$ or $2q$ labels, respectively.
The cheapest losses relative to floor nine for outside primes
$5$, $11$, and $q\ge13$ are therefore
\[
 c_5=16(1/81-1/625),\quad c_{11}=22(1/81-1/121),\quad
 c_{13}=48(1/81-1/169).
\]
Compare with $E=105^{-2}+128/441+88/81-1$.
The sums $c_5+c_{13}$, $c_{11}+c_{13}$ and
$c_5+20(1/81-11^{-4})$ exceed $E$. Consequently the support is
$\{3,7,q\}$ or $\{3,5,7,11\}$ with a present eleven-class.

For three-prime support, put $s=m_9$ and let $w_q(s)$ be the paired
nine-row omission bound. For $s\ge9$ reserve its $q$-free witnesses
at $9q^2$, disjoint from the $128$ mixed positions and the $A_q$
isolated active positions. The eight three-free labels and the
$w_7(s)$ seven-free witnesses may overlap the mixed positions;
the minimum losses per label, compared with the already assigned
floors, are
$\delta_3=21^{-2}-189^{-2}$ and
$\delta_7=21^{-2}-441^{-2}$. With
$k=216-s-128-A_q-w_q(s)\ge0$, the upper mass is
\[
 105^{-2}+s/81+128/441+A_q/q^4+w_q(s)/(9q^2)^2
       +k/225-8\delta_3-w_7(s)\delta_7.
\]
For $s<9$ omit the witnesses and their losses.

For four-prime support let $m\ge22$ be the eleven-class size.
Forty five-active positions would give mass at most
$1155^{-2}+128/441+22/121+40/625+26/81<1$.
Thus the sparse height identity makes all five-free labels neutral.
There are at least six by Lemma~\ref{odd:two-free}. The three
proper eleven-Weyl pairs omit respectively $3,7,5$, so only the
last pair can be five-free. Reserve four further five-free labels
at $25$. The proper pairs have floors $g,77,275$, where $g=99$
for $s>0$ and $g=33$ otherwise. With
$k=216-s-128-m-16-6-4\ge0$, the upper mass is
\[
 1155^{-2}+s/81+128/441+m/121+16/625
       +2(g^{-2}+77^{-2}+275^{-2})+4/625+k/225.
\]
Both formulas are strictly below one. Their respective maxima over
the even $0\le s\le76$ and the permitted $q,m$ are
$215990653/218791125$ and $33299377/33350625$.
The exact rational checks are recorded in
\supfile{verification/odd_theory/check_edge37.py}{the supplement}; no denominator search is involved.
\end{proof}


Before generating a phase profile, let $U$ be the booked ordinary
mass and let $k$ positions remain. If $E$ of these positions have
floor $3B$, while all others have floor at least fifteen, then
\[
 E\le\left\lfloor
 \frac{U+k/225-1}{1/225-1/(3B)^2}\right\rfloor.
\]
Generating only profiles within this necessary cap avoids enumerating
prefixes that ordinary mass would immediately exclude.

At order $3t$ the four square classes have a common phase weight in
both directions. At order $9t$, each three-primary square class gives
two classes with a common weight. For each such family enumerate
its positive counts per phase, congruent modulo eight, and its small
counts; a count $n$ with $u$ small pairs has weight between $u$ and
$u+(n-u)/9$, strictly above $u$ when $n>u$. Test the two moments in
each occupied phase at both endpoints. Higher orders can be bounded
separately by selecting $t(t-1)/2$ fibers in one higher $t$-orbit,
or nine in a higher three-orbit, and testing just this set of fibers.
For a central three with two leaves it suffices to retain the
$t$-phase bounds at each leaf. A fiber with $n$ mixed pairs needs at
least three above $B$ if $3\nmid n$, and one if $9\nmid n$, by the
height test. We record the neutral reservations used with these phase bounds.
At a sparse isolated prime $p$ absent among prime denominators,
\eqref{eq:sparse-height} makes all active degrees divisible by $p$.
Deletion and Corollary~\ref{blk:primepower} give at least four nonunit
$p$-free neutral labels, all with $p^3\mid b_i$. With four primes
and a valuation-one class, only one of its three selected proper
pairs omits $p$. With two valuation-one classes at least ten proper
labels can be reserved; their two triples of pairs can overlap only
on a pair omitting both class primes. With four primes $3,t,p,q$,
one edge $3t$ and sparse isolated $p,q$ absent as denominators,
there are at least eight neutral nonunits outside the nine-class.
Indeed, in each of $P_p,P_q$ the neutral nonunit degrees still have
gcd one and intersect pairwise in their prime supports. Each needs
three pairs. If their union had only three, each would belong to
both blocks. At least two would omit $p$ and at least two $q$,
and one would omit each of $3,t$. Each degree omits at most two
primes. Three two-element supports with pairwise intersections and
empty total intersection form a triangle, inconsistent with the
resulting omission multiplicities $2,2,1,1$.

With support $3,t,q$ and $m_9\ge9$ one has $27\mid M$ since
at $9\Vert M$ all three-block degrees seen by the nine-rows would
divide $M/9$ and contain both other primes. Book at least eight
nonunit three-free labels, as well as the Gram witnesses omitting
$t$ or $q$. These three sets are pairwise disjoint, since every
nonunit degree has at least two prime divisors. Every selected label
is seen by a nine-row and has three-primary twist. The $q$-free Gram
witnesses are neutral when $q$ is isolated. In the $t$-free set,
possible overlaps with solely-three-active mixed labels are charged
only once, using the larger denominator floor.



\subsection{The signed Gauss sum}
Average $\sum\theta_X/b_X^2=\xi/M$ ($\xi=\pm1$) over the cyclotomic
group. A label of twist order $e$ contributes to the deficit
$1-\xi/M$ the nonnegative rational weight
$(1-\mu(e)/\varphi(e))/b_X^2$ after averaging its whole orbit. Call
this weight its charge. At an occupied pure power $b=p^a$, the charge
is $pc/b^2$ with $c=h_1/(p-1)+h_2/p$, where $h_1,h_2$ count labels of
primitive and higher $p$-order. Thus $c$ is integral and costs between
$(p-1)c$ and $pc$ labels. If $a_p$ labels of denominator $p$ are active, the class contributes $a_p/[p(p-1)]$. The class of denominator nine, of size $s=m_9$, contributes $c/27$ with $0\le c\le s/2$. When fewer than
$p^2(p-1)$ positions remain for the current and all higher levels, the height and balance bounds of Proposition~\ref{new:activity-height-gap} give even $c$. Every subsequent coefficient is below $p^2$, so the top charge has a unique largest pole at $p$. Its reduced denominator, and therefore all lower-level charge denominators, divide $M$.

Subtract the charges of complete specified groups from one, leaving $U$.
Let $L\mid M$ contain all known divisors. If the next denominator is $b$
and at most $k$ slots remain, then
\begin{equation}\label{eq:odd-deficit-tail}
 U\ne0,\qquad -1/\lcm(L,b)\le U\le1/\lcm(L,b)+3k/(2b^2).
\end{equation}
Indeed every remaining charge rate is at most $3/2$. If $U=0$, the
positive remaining charges would sum to $1/M$. Group them by their
component $J$. Their denominators are supported on $J$, up to a factor
two for the joint orbit of order fifteen. Comparing at the odd primes
of $J$ makes each positive group a multiple of $1/(2M_J)$, where $M_J$
is the $J$-part of $M$. At least two components occur, so $M/M_J\ge3$.
For an empty tail require $U=\pm1/M$. This argument gives finite
branching even before $M$ is known, since $\lcm(L,b)\ge b$.

\paragraph{The empty graph.}

Reserve six neutral nonunits overall (Proposition~\ref{odd:neutral-orbits})
or the stronger proper and omission counts above. Three-row omission
witnesses are neutral since three-activity has pure denominator.
The isolated field populations round up independently to multiples of
$\operatorname{lcm}(8,2d)$. With $s\ge9$ on three-prime support, three-height at least three
reserves eight three-free labels beyond the nine class, which may also
be active. For a mixed denominator $b=uk$ with $u$ its full $p$-part,
a dual pair spends at least $2\|k^{-1}\|_u^2/u^2$ in the local Weyl
norm (here $\|\cdot\|_u$ is the least absolute residue) (Proposition~\ref{new:local-entry-mass}); at a prime class use
its proper norm after reserving the other proper degrees. This bounds
the smallest possible free denominator, with the full tail retained
above the numerical cutoff.

Visit pure powers in increasing order. At $p$ let $t_p$ be the shortfall
in its activity quota and $B_p$ its next possible power. Include $B_p$
in $L$ when $t_p>0$, or $B_p^2/p$ when fewer than $2p(p-1)$ active
slots remain in that field. Indeed in the latter case the largest
active denominator $B$ satisfies $M_p=B^2/p$. The positive charge
still available is at most
$\sum_p t_p\alpha_p+(k-\sum_p t_p)\max_p\alpha_p$, where
$\alpha_p=p/((p-1)B_p^2)$; allow error $1/L$ for the vacuum Gauss term.
Use the same reservation in ordinary mass. Enumerating $c$ at each
level, subtract its minimum label cost and credit its maximum activity,
rounding the field populations after each choice. \eqref{eq:odd-deficit-tail}
bounds every next level, and a terminal requires $L\mid M$.

\paragraph{The edge $3,5$.}
Let $z,x,y$ count respectively the joint, solely-three and solely-five
labels of mixed denominator. At primitive order fifteen, the eight
positive joint-phase counts in pairs are even and congruent modulo
eight, with a common count on each two-phase square orbit. An increase
of eight costs $32$ labels. A ninth-root square class costs at least
$48$ joint labels and has size divisible by $24$. Orders $25,27$ require
ten or nine phases containing at least four or eight joint pairs,
respectively. A trial with a higher order need only reserve one square orbit of
its phases, costing a multiple of $20$ or $18$ joint labels. Among the
surviving reservations the respective joint counts are $32,48,64,96,128$;
no projection to a higher-order orbit therefore survives. The remaining phase counts are exact. A solely-active mixed label shares its phase with
joint activity, since $J=C=0$ in the two moments of Proposition~\ref{phase:moments} would give
$\sum(a_i-1)^2w_i=0$.

\begin{corollary}[The mass of a minimal joint family]\label{new:rigid-family}
Suppose a $3,5$ twist component has either $32$ primitive joint
labels, or $48$ joint labels of order $45$ in one ninth-root square
class, and these comprise all joint activity. All joint denominators
have a common value $B$. If $z$ is the joint count, their mass is
$J=z/B^2$. The mixed labels active solely at three have mass exactly
$J/24$. Their primitive three-phase imbalance is zero. Together
these sectors contribute ordinary mass $25J/24$ and Gauss deficit
$30/B^2$ or $50/B^2$, respectively. The number $x$ of mixed labels
active solely at three satisfies
\[
   x\equiv12\pmod{16}\quad(z=32),\qquad
   x\equiv18\pmod{48}\quad(z=48).
\]
\end{corollary}
\begin{proof}
Each occupied joint phase contains four labels: for order fifteen
there are eight phases, and for the stated order-forty-five family
there are twelve. Each phase forms one local-sign orbit and has a
single denominator by Proposition~\ref{new:sign-fiber}. Equality
of the fifth-root phase masses at fixed three-phase, and equality
of the primitive third-root masses or transitivity of the ninth-root
square class, identify all these denominators.

Fix an occupied three-phase. It contains eight joint pairs, with
two at each primitive fifth root. Proposition~\ref{new:sign-fiber}
gives $E<4$; Lemma~\ref{odd:parity-lattice} gives $E=0$. Thus
every neutral mixed coefficient equals $-5$ and the solely-three
mass in pair units equals $1/(3B^2)$. Scaling by $B^2$, every pair
weight is the square of an odd $2$-adic unit. The number $h$ of
such pairs in this three-phase satisfies $h\equiv1/3\equiv3\pmod8$.
There are two primitive third roots, or three conjugate ninth roots
with equal counts. Doubling the pair counts gives the asserted
congruences. Equality of the two primitive masses in the first
case gives zero imbalance; the second case uses ninth roots only.
Summing the fiber identities gives the ordinary mass, and the
Ramanujan means give the two Gauss deficits.
\end{proof}

When $B=15$ no solely-three mixed label can have denominator fifteen,
since one pair would already exceed its entire fiber mass. If such
a fiber has exactly three pairs, all three denominators are $45$:
the next possible mixed denominator after $15$ is $45$, and three
weights $45^{-2}$ attain the required sum.

\begin{lemma}[Mass and saturation in a primitive fiber]\label{odd:fiber-lower-mass}
In Proposition~\ref{odd:edge-budget}, suppose all active phases have
order $p$. Put $V=\sum_iw_i$, so $J=(p-1)C$. Then
\[
 V\ge\frac{C}{p+1}=\frac{J}{p^2-1},\qquad
 \sum_i a_i(a_i-1)w_i=J+C.
\]
If positions of total mass $Z$ have coefficient zero, then
\[
 C^2+ZJ\le V(2C+J-Z),\qquad
 V+J-p(2k+1)(V-C)+p^2k(k+1)(V-Z)\ge0
 \quad(k\in\mathbb Z).
\]
If the coefficients lie in a sublattice $d\ZZ$, replace $p$ by $d$
in the second inequality. A position of weight $w$ has coefficient zero when
$J+C<p(p-1)w$. At equality, a nonzero coefficient there is $p$,
all other coefficients vanish, and $V=C+pw$.
\end{lemma}
\begin{proof}
The coefficient moments are $V-C$ and $V+J$. Their difference is
$J+C$, and $a(a-1)\ge p(p-1)$ for nonzero $a\in p\mathbb Z$,
with equality only at $a=p$. This proves the saturation assertions.
Cauchy--Schwarz outside the zero set gives
$(V-C)^2\le(V-Z)(V+J)$, including the first bound when $Z=0$.
Finally sum $(a_i-pk)(a_i-p(k+1))w_i\ge0$ outside that set.
\end{proof}

Moreover $p|V-C|\le V+J$ bounds $V$ above. These facts
permit exact treatment of fibers where a position of denominator
fifteen would saturate the squared norm; otherwise use the moment
intervals above, with next denominator at least $45$. In particular
in a primitive $32$-label fifth-root fiber a nonzero coefficient at
fifteen forces $V=5/225$; if there are fewer than thirteen solely-five
pairs all five must then have denominator fifteen.

\paragraph{Neutral positions.}
In the three-prime case keep disjoint three-free pairs, five-free
nine-row witnesses and $q$-free nine-row witnesses. At sparse isolated
$q$ their floors are $27,225,9q^3$; the first two can meet mixed
activity only in the solely-three sector, at floors $135,225$.
All $q$-free nonunit degrees are divisible by three, since each degree
contains at least two of the three primes. Their labels are neutral:
a pure $q$-row makes their twists $q$-primary, whereas $q$-active
degrees are divisible by $q$. To bound their number enumerate counts $v_t$ of pair residues
$t=1,\ldots,(q-1)/2$. Besides $1+2\sum v_tt^2\equiv0\pmod q$,
the quantity $2\sum v_t\|at\|_q^2$ must reach the projection bound
\eqref{eq:full-census} for each $a$, using the two consecutive
multiples of three in $-a+q\ZZ$ and even selected-row multiplicities.
Additional $q$-free labels have floor $q^3$.

For four primes with prime denominator $p$ and other isolated prime
$q$, proper pairs omit $3,5,q$ at floors $9p,5p,pq^3$ in the sparse
case. Only the last may be $q$-free. The seven-free count is at least
six and at least eight when the mass of the nine and marked classes
exceeds $3/7$, by the seven-free mass bound after Lemma~\ref{odd:two-free}; for
$q\ne7$ reserve an additional pair. With two prime classes $p,q$, the five distinct pair
floors $9p,5p,9q,5q,pq$ suffice. Combine these reservations by Lemma~\ref{new:reservation-calculus}.

\begin{lemma}[Gram bound for a denominator class]\label{odd:class-gram}
Let $F$ be a complete class of $s$ nonunit labels of denominator $b$ in odd square
dimension, and let $I$ contain $F$ and the supports of its rows.
Put $H=\sum_{Y\in I}b_Y^{-2}$, $E=H-s/b^2$, and $C_k=\sum_{X\in F}\theta_X^k/b^2$.
Then
\[
  (|C_1|-E)_+^2\le HE,
  \qquad |C_2-\overline{C_2}|\le2E .
\]
\end{lemma}
\begin{proof}
Write $v$ for the restriction of $se_0$ to $I$ and $P$ for orthogonal
projection on the corresponding rows $se_X$, $X\in F$. Since the
unrestricted vectors have the same projections, $Pv=0$, and the
coordinates of $P T^{\pm1}v$ in this basis are
$\epsilon\theta_X^{\mp1}/b$, by the modular relation and duality.
Both $(1-P)T^{\pm1}v$ therefore have squared norm $E$.
Pairing $(1-P)Tv$ first with $v$ and then with $(1-P)T^{-1}v$,
Cauchy--Schwarz gives the assertions after dropping the tail
$I\setminus F$ at cost at most $E$ in each inner product.
Here $T$ denotes the unnormalized twist diagonal.
\end{proof}

At $b=9$ let $c/27$ be the charge and $i/81$ the absolute primitive imbalance. Then
$81|C_1|=\sqrt{(s-3c)^2+3i^2}$ and
$81|C_2-\bar C_2|=2\sqrt3 i$ after averaging higher-order phases
over their wild Galois subgroup when necessary. Indeed at these
complete levels their primitive-root sums vanish also before averaging.

With $P,v$ as in the proof of Lemma~\ref{odd:class-gram}, put $A=(T+T^{-1})/2$. The squared residual norm of $Av$ is the
weighted sum of squared real parts on $I\setminus F$, since the
$F$-contribution cancels the projection. The Galois means of a real
part and of its square are $t$ and $(1+t)/2$, $-1/2\le t\le1$, by
odd order. If $e_0,u_0,h_0$ are the specified tail mass and its first
two real moments, with unspecified mass at most $E'$, Cauchy and
averaging give
\[
 2((s-3c)/81+u_0+u)^2\le(s/81+e_0+E')(2h_0+E'+u),
 \qquad -E'/2\le u\le E'.
\]
At a nine-multiple $b$, $a$ solely-three labels with $i'$ primitive
packets add $(a,-i',(a-i')/2)/b^2$ to the three specified moments.
Every supporting label has three-primary twist and is at a multiple
of nine by Lemma~\ref{odd:nine-support} in the cases tested; use the
least common multiple with nine of each omission divisor.

\paragraph{The mixed recursion.}
Enumerate the occupied denominators in increasing order, assigning
complete square-unit packets and their duals. The charge per label is
$7/(8b^2)$ at order fifteen, $1/b^2$ at $45$, $p/((p-1)b^2)$ at order
$p$ and $1/b^2$ at higher prime powers. For $p=3,5$, if $i_p$ primitive
packets occur, update the absolute difference $d_p$ of the two
primitive phase masses per phase at trivial complement to
$|d_p+j/b^2|$, $|j|\le i_p$, $j\equiv i_p\pmod2$.
Start with $d_3=i/81$, $0\le i\le c$, $i\equiv c\pmod2$, and $d_5=0$;
both vanish at completion. Subtract the minimum label cost and test
all residual ordinary, charge, phase and Gram intervals. In
\eqref{eq:odd-deficit-tail} one may use $M\ge\kappa b$, with $\kappa$
the least of $L/(L_3L_5)$ and $pL/L_p$ at sparse isolated primes
($b^2/p\mid M$ for a pure active $b\ge p^2$). With next level $b>15$
a free label contributes at most $1/225$ to ordinary mass plus
$2(b^2-225)/675$ times its charge. For remaining solely-three,
joint and pure-three masses $A,J,P'$ with $P'\le P$, booked
$A_0,J_0$ and $A+J+P'\le T$, the fiber bound
$16(A+A_0)\le9(J+J_0)$ and $2d_3\le A+P'$ imply
$50d_3\le9T+9J_0-16A_0+16P$. At order $45$, $2d_3\le P$.
At minimum isolated activity reserve the neutral pair at the unique
active denominator by Proposition~\ref{round:packet-idempotent};
otherwise use Lemma~\ref{odd:inactive-pure} for its floor.

A terminal determines $M$ by $U=\pm1/M$. Check divisibility by every
specified denominator and the prime heights. The residual ordinary
squared-degree sum is twice a sum of squares of admissible divisors
$d$ of $M$ at floor $M/d\ge15$, with the number of pairs congruent to
that half-sum modulo eight. Enumerating multiplicities of the six
largest available squares and bounding the rest by their floors gives
a necessary interval test. Bounds on activity and residual mass always
use the ceiling $209$; any count known only to be below the ceiling
retains the whole range.

\paragraph{Small mixed fibers.}
We record the integer consequences of the moment bound used at the last
levels. Scale by $15^2$, set $u_i=a_i-1$ and let $k_i=b_i/15$.
A fiber of $j$ joint pairs satisfies
$\sum u_i/k_i^2=-C$, $\sum u_i^2/k_i^2=J+2C$, $E\equiv-j\pmod8$,
the last by summing squares modulo eight in the even coefficient groups.
At any of the first cofactors $1,3,5,9$, enumerate
$u_i^2\le k_i^2(J+2C)$; on the unassigned tail use the moment and parity
intervals above. Modulo $32$ each equal-cofactor group uses one of
$1,9,17,25$ as common weight. If its mass is fixed to $v>0$, the next
cofactor is at most $\sqrt{h/v}$ for $h$ remaining pairs.

At $z=64$ primitive, divide each phase into sign orbits of four labels.
The positive numbers of such orbits per phase are equal modulo eight,
constant on two-phase orbits, and total sixteen, hence all two. If both
groups in a phase have the smallest denominator, the same holds in
every phase; if exactly one does in each phase, their second denominators
agree. A phase can omit the smallest $b$ only if the next $b'$ has
$b'^2<2b^2$. For $z=96,128$ with all joint denominators fifteen, and for each feasible decomposition into $15$- and $45$-order phases,
$z=16a+24(b+c)$ with $a,b,c$ even, the local $(J,C)$ are $(4t,t)$
at the three-phase of type $t=a,b,c$ and $(2a+3b+3c,a)$ at each fifth
root. These follow by counting the full phase pairs, with $a$ at each
primitive phase and $b,c$ on the ninth-root square classes.

In each coefficient group the nonnegative moment
$\sum a_i(a_i-1)/k_i^2=J+C$ charges at least $p(p-1)/k_i^2$ per
odd coefficient, and at least $2p(2p-1)/k_i^2$ per nonzero even one.
If the odd coefficients exhaust this budget at their minimum, they
all equal $p$, the even coefficients vanish, and the first moment fixes
the mass. For groups assigned cofactors below $75$ use their exact
weights and bound each remaining weight by $1/75^2$.

Below $84$ active labels at seven, the possible counts are
$12u+42v$, $v=0,1$, by the square classes of primitive and higher twists.
Reserve the groups at $49$, then $343$, bounding the rest from $2401$.
A lone final group has charge $14/B^2$ or $42/B^2$, with $B^2/7\mid M$;
in the sparse case the total count is $12,24,36$. A pure isolated row
on three-prime support requires separate three-free and five-free neutral
pairs at floors $3^{H_3}q^2,5^{H_5}q^2$, where $q$ is the isolated
prime and $H_i$ the known height, by principal support and gcd one.
They can coincide with the corresponding omissions already counted.
Initially a free label of floors $15$ neutral or $25,27$ active has
charge at most $1/486$ and ordinary mass plus $64/45$ times its charge
at most $1/225$.

To check remaining slots modulo eight, each primitive or higher packet
at three adds respectively (charge numerator, count) $(3,2),(6,6)$;
at five the pairs are $(5,4),(20,20)$. Odd reciprocal squares reduce
to one modulo eight. Keep the charge residue and the two primitive
imbalance parities, admitting every odd residue of $\pm1/M$; shortest
path length in label count gives a lower bound. At the largest sparse
seven-power $B$, clearing denominators gives $M_7=B^2/7$ and
$cM/M_7\equiv-\xi\pmod{49}$, $c=2,4,6$. A prime class at eleven gives
$2M/11\equiv-\xi\pmod{11}$ and hence $\xi=-1$ on the support
$3,5,7,11$ when $M_7=B^2/7$: the seven-height is odd, $3,5$ are
quadratic residues and $2,7$ nonresidues modulo eleven. Testing
multipliers of $L$ up to $10^5$ gives a lower bound for $M$; the
entire larger tail is retained.

\paragraph{The rank ceiling.}
The \supfile{odd_square/profiles.py}{enumeration} runs once with $209$
as an upper bound throughout. At every reservation it keeps unused slots
available to later steps. Thus the same computation covers every smaller
rank. From $16\,512$ prime profiles, the ordinary masses, phase sectors and
nine-class bounds leave $1\,151$ edgeless reservations and $310$ with
one edge $3$--$5$. The
\supfile{odd_square/check.py}{two recursions}
visit $6\,721\,348$ and $2\,753\,944$ nodes, respectively. Every
reservation is excluded. The \supfile{odd_square/census.json}{census}
and \supfile{odd_square/certificate.json}{certificate} record the input,
outcome and node count in the same order for an exact replay.

\begin{proof}[Proof of Theorem~\ref{odd:rank}]
The nonsquare census excludes all ranks below $219$. In square dimension
the preceding calculation excludes every rank through $209$, and
$r\equiv1\pmod8$ makes $217$ the next possible rank.
\end{proof}

At the first open odd rank, $217$, the
\supfile{verification/frontier217/check.py}{profile and graph screens}
retain $1\,133$ edgeless and $495$ single-edge assignments. The
\supfile{verification/frontier217/graphs.json}{indexed totals}
record the intermediate counts.
The $3$--$5$ edge and the empty graph are exhaustive at this stage:
the ternary separation applies through $217$ and a $5$--$7$ edge
would give mass at most $172/35^2+45/9^2<1$.
These assignments have yet to pass the terminal tests.


\section{The finite verification}\label{sec:verify}

Throughout this section the ring is non-pointed and based-simple and
has the data of Theorem~\ref{alg:rank}. Let $n$ be
the largest degree, $m$ its multiplicity, $B=M/n$ and $f=qB^2$.
By \eqref{eq:recip}, $f<r$. Equality $f=r-1$ would give
$\sum_{i\ne0}(n^2-d_i^2)=1$, although each positive summand is at
least $2n-1$. Thus $f\le r-2$. Theorems and lemmas above give $B\ge2$,
$f\ge8$ and $f\ne9$
(Theorem~\ref{thm:squarefree-codegree}, Lemmas~\ref{small} and~\ref{nine}).

For every odd $p$ with $v_p(f)=1$, the maximal-class count of
\cite[Proposition~19.9]{BDPcore} applies by Lemmas~\ref{blk:moments}
and~\ref{blk:height}. Its one-prime consequence used by the covering
enumeration is
\begin{equation}\label{eq:startmult}
 m\le f-1-(p-1)/2 .
\end{equation}

Every denominator strictly above $B$ is at least $B+1$. More generally let
$p\mid M$ be an odd prime outside $qB$, put $A=B(q)(p-1)$ and
$C=p\lceil(B+1)/p\rceil$. By Proposition~\ref{blk:packets} at
least $A$ labels are $p$-active; their denominators are multiples of $p$
by Lemma~\ref{blk:admissible}, hence at least $C$. Summing squared degrees gives
\begin{equation}\label{eq:startgap}
 f\le n^{-2}+m+(r-1-m-A)\left(\frac{B}{B+1}\right)^2
                     +A\left(\frac{B}{C}\right)^2 .
\end{equation}
One may take $A=0$. The generator ranges over $2\le B$ and square-free
$q$ with $8\le qB^2\le r-2$, $qB^2\ne9$.

It then considers the possible supports of $M$ and all multiplicities $m$ satisfying these bounds. Corollary~\ref{round:eight-rank} excludes $f=8$ throughout the computation. Empty ranges need no search. The odd-dimensional theorem permits the further requirement $2\mid M$ throughout our enumeration.

\begin{proposition}[Separated primary twists]\label{new:separated-primary}
Suppose $q\in\{3,5\}$. Put $F=(q+3)/2$ and let $p\ge5$ be a prime
of $D$ other than $q$. If fewer than $2F(p-1)$ labels are active
at $p$, then for every such label $X$ and every prime $t\mid D$
outside $\{p,q\}$ one has
\[
              t\nmid f_X,\qquad \theta_{X,t}=1 .
\]
There are at least $F(p-1)$ active labels. Writing $h=(p-1)/2$,
their multiplicity at each degree is divisible by $h$. When $2\mid D$
and $q\nmid b_X$, it is divisible by $qh$.
\end{proposition}
\begin{proof}
Every nonempty fiber fixing phases outside $q$ has at least $F$
members. Indeed the proof of Proposition~\ref{blk:packets} at $q$
also applies when a dyadic phase is fixed, since the complementary
unipotents commute with the $q$-factor. An order-$p^j$ phase with
$j\ge2$ would therefore require at least $Fph\ge2F(p-1)$ active
labels by the square-unit action. All active $p$-phases have order $p$.

The subspace with all complementary phases trivial is invariant under
the $p$-factor and contains the vacuum. It contains a $p$-active label:
otherwise the upper unipotent and its Weyl conjugate are scalar
there, so the local Weyl fixes the vacuum line, contradicting
Lemma~\ref{blk:admissible}. On this subspace the primitive $p$-phases
have equal mass by Lemma~\ref{blk:moments}, since $v_p(D)$ is even.
Thus all $p-1$ primitive roots occur with every complementary phase
trivial. Here and below the moment identity may be applied with
the complementary phase projectors alone.

If a $p$-active $X$ has $t\mid f_X$, its fiber fixing all phases
outside $t$ contains two distinct $t$-phases. Otherwise the local
unipotent acts there by a scalar, as do its Weyl conjugate and hence
the local Weyl, again contradicting Lemma~\ref{blk:admissible}.
In particular some label is simultaneously active at $p,t$.
Fix that nontrivial $t$-phase. Equality of primitive $p$-masses gives
all $p-1$ primitive roots in this phase also. For each primitive root
both choices of $t$-phase have at least $F$ labels. This contradicts
the upper bound. The assertion about $\theta_{X,t}$ follows from
local admissibility, and the lower bound from the $p-1$ primitive
roots. Square-unit symmetry gives divisibility by $h$.

If $2\mid D$, all active denominators are odd. A label with $q\nmid b_X$
has defect one at $q$, and its whole block has the same $p$-phase by
Lemma~\ref{blk:defectone}. Every unequal-degree block at three or five
contains denominators $z,2z$, impossible among these odd denominators.
The equal-degree blocks have $q$ labels each. Thus the multiplicity
at each fixed primitive root is divisible by $q$. Galois translation
across its orbit of $h$ primitive roots proves the last assertion.
\end{proof}

\subsection{Denominator and degree tests}
The \supfile{census/run.py}{generators} use \eqref{eq:interval}
and its terminal divisor tests. With two places left, the prime-to-$k$
part $U$ of $L$ gives $k^2\le(1+2U^2)/E$.
All fixed-width limits transfer the unfinished branch to exact arithmetic.

In square dimension the final two places cannot introduce a prime,
whose principal block would need three distinct nonunit denominators
(Lemma~\ref{blk:witness}).
With three places left and $k$ even, the equation
$Ek^2=1+d^2+e^2+g^2$ modulo four and eight gives $d,e,g$ odd,
$v_2(k)=1$ and $E$ odd. In nonsquare dimension the support $S$
contains the primes of $qB$, has at least three primes (four at
codegree eighteen), and satisfies
\[
 B(q)\Big(1+\sum_{p\in S,\ p\nmid2q}(p-1)\Big)\le r
\]
by Proposition~\ref{blk:packets}. This bounds the new primes;
the subsequent denominators are $S$-smooth.

For $p\mid q$, tile the $p$-free denominators by
Lemma~\ref{blk:defectone}. At two their multiplicities are even.
At three or five, remove the copies at $z$ already required by
$z/2$. For the remaining multiplicity $m$, reserve $2b$ copies
at $2z$, with the least $b\ge0$ satisfying
$b\equiv m\pmod3$ or $2b\equiv m\pmod5$, $2b\le m$,
respectively. It suffices to use the least choice: replacing $p$
mixed blocks by equal blocks at both denominators preserves the
tiling. The eight-label quinary block counts as two mixed blocks.

\begin{corollary}[Prime-cut reservation at a denominator prefix]
\label{new:prime-cut-prefix}
Let $S$ be a finite set containing the final prime support of $M$,
and partition $S=A\sqcup B$ so that each part contains a prime
already known to divide $M$. Consider a prefix of nonunit
denominators and let $C$ be its set of distinct denominators
divisible by a prime in each part. Put $k=\max\{0,3-|C|\}$.
Suppose $h$ nonunit positions remain, all with denominator at least
$b$, and put $\delta=1$, or $\delta=2$ in odd dimension.
Then $h\ge\delta k$. If $c_1<\cdots<c_k$ are the least integers
at least $b$, outside $C$, divisible by some $pq$ with $p\in A$
and $q\in B$, the remaining nonunit mass is at most
\begin{equation}\label{eq:prime-cut-prefix-mass}
 \frac{h-\delta k}{b^2}+\delta\sum_{j=1}^k\frac1{c_j^2}.
\end{equation}
For a residual which includes the vacuum, add its separately
certified upper bound. Each cut is tested separately.
\end{corollary}
\begin{proof}
The final support meets both parts, so Theorem~\ref{new:prime-cut}
forces $k$ previously unseen crossing degree classes. Reserve one
label from each, or a dual pair in odd dimension. When their distinct
denominators are ordered, the $j$th is at least $c_j$. These classes
are disjoint and occupy $\delta k$ future positions. Bound each
remaining position by $b^{-2}$. Allowing the integers $c_j$ to have
additional prime divisors only weakens this upper bound. Repeats of
an open prefix class do not supply a new class. Finally, the same
future class can cross several cuts, so their reservations must not
be added without a further disjointness argument.
\end{proof}

The integers $c_j$ require no smooth-number search: merge the
arithmetic progressions of multiples of $pq$, skipping duplicates
and members of $C$. A resource-limited implementation may omit cuts;
it may not infer exclusion from a truncated candidate search.

\paragraph{Disjoint reservations.}
The activity bound holds on every union of unramified odd-prime
sectors, so Hall's theorem reserves $B(q)(p-1)$ distinct labels
at each $p$ (Proposition~\ref{blk:packets}). A closed degree class
of size $s$ supplies at most $h_p\lfloor s/h_p\rfloor$ of them,
$h_p=(p-1)/2$; an open class may supply all $s$.
At a present denominator five, the other odd sectors use only
five-free denominators. Every unfilled reservation has the least
eligible multiple of its prime as a denominator floor; all other
positions have the next allowed denominator as floor. At the initial
class at most $m$ maxima can meet these reservations, and none
can serve a prime outside $B$.

These floors also give the next-denominator estimate of
\cite[\S21.4, following Lemma~21.6]{BDPcore}. If $v$
reserved suffix positions have floors $a_j$, put
$W=\sum_{j=1}^v\lceil L^2/a_j^2\rceil$. With $h$ positions
unfilled, $E>W+1$ implies
\[
                 b^2\le\frac{(h-v)L^2}{E-W-1}.
\]
The term one bounds the vacuum. If a floor depends on an undecided
prime valuation, take the weakest bound over its possible values.

\paragraph{Successive Gauss sums.}
Apply Lemma~\ref{new:partial-gauss} successively to complete
classes, deleting every forced zero sum. Such a vanishing is
impossible if the selected sum has nonzero reduction at a prime $p$.
This occurs when its unique class of maximal $p$-adic denominator
height is a singleton, or has $p$-power denominator and multiplicity
prime to $p$.

For one class the criterion is $b^2(R+1/M)<1$, with $R$ the
complementary mass. It also precludes a further copy of $b$ in a
prefix. A singleton cannot vanish, nor can an odd sum at a dyadic
denominator, by reduction above two. Round all mass bounds outward.

Lemma~\ref{r7:gauss-projection} supplies two-class zero sums.
With ten fours, five fives and six
sixes, $T_5=0$, the complementary mass is $1/120$, and the even
four-row degree gives $8\mid M$ and $M\ge840$. Thus
$72(1/120+1/840)<1$ forces $T_4=T_6=0$.
The remaining support prime $p\ge7$ is isolated. A denominator
$p$ would have valuation one, but its $p-1$ active labels already
exceed mass $1/120$, since $p-1\le41$. Sparse height therefore gives
$M\ge8\cdot3\cdot5\cdot7^3$. Every further denominator satisfies
$2/b^2\le1/120+1/M$: either its class has at least two labels,
or the singleton house bound gives the same inequality. Denominator
$15$ is excluded by the proper five-free residue condition.

Seven seven-rows give $T_7=0$. Apply the union test to
$9T_4+4T_6$, adding one selected twelfth-denominator twist when
$m_4$ is even. If its house is below one, reduction above two is
nonzero when $m_4$ is odd, or when the additional twist was selected.

\paragraph{Square prefixes.}
In even square dimension through rank $43$, $M$ has four prime divisors and
Corollary~\ref{cor:odd-isolation} separates the allowed odd
activities. If $p$ is isolated, at least $p$
nonunit positions have pure $p$-power denominators by
Proposition~\ref{new:pure-extra}. This count is taken separately
from the activity and omission bounds because an inactive position
can contribute to an omission. In the pure alternative at $p$, some class $p^a$ has
at least $p-1$ labels. Indeed an order-$p^2$ or higher orbit suffices;
otherwise multiply the common primitive-phase masses by the square
of the greatest active denominator and reduce modulo $p^2$.
At an isolated $p\ge5$ the greatest active exponent $a$ itself
satisfies $v_p(M)=2a-1$ by \eqref{eq:sparse-height}.
For an isolated $p\ge7$, the vacuum and the other two odd-prime activities
leave at most $36$ labels active at $p$. This is below both
$p(p-1)$ and $p^2(p-1)/(p+1)$. At isolated five, an occurrence of
five gives exactly four active labels. Otherwise there are at most
sixteen: for at least nine fours this was proved in
Proposition~\ref{sq:thirteen-34}; for at most eight fours, at least eighteen
active five-labels would give mass at most
\[
 18/625+6/49+2/36+8/16+8/36+1/210^2
 =1024277/1102500<1.
\]
The active count at five is even by square-unit covariance.
Thus sparse height applies in every case, and
\eqref{eq:cost-large-valuation} excludes $p\Vert b_i$ for every
isolated $p\ge5$ without denominator $p$.

Given the current lcm $L$ and last denominator $b$, eligible
greatest powers have $a\ge\lceil(v_p(L)+1)/2\rceil$. A class below
$b$ must already have $p-1$ members; only the class at $b$ can
gain members from the suffix. Taking the least still possible
$p^a$ gives the lower divisor $p^{2a-1}$ for every $p$-free
denominator. When $p$ is absent these witnesses are disjoint
from all odd-active labels. For a present prime denominator
$p$, its proper free degrees have gcd one. For each other
dimension prime $t$ they require a label omitting both $p,t$,
whose denominator is divisible by the full powers of $p,t$ in $M$.
Different unordered pairs require different labels, since a
nonunit degree contains at least two of the four primes.
Each remaining free-degree count charges a distinct omitted prime.
A nonunit can pay at most two charges, so their union uses at least
the largest individual count and at least half their total, rounded up.
Combine overlapping floors by Lemma~\ref{new:reservation-calculus}.

In the mixed $2,3$ alternative the required ternary phase fibers
each contain four denominators $2^i3^j$ of a common positive
two-adic height, by \eqref{eq:phase-column}. The difference of
the two primitive ternary phase masses is zero. Initial eligible
denominators of lcm $W$ contribute an integer multiple of $W^{-2}$
to this difference; when all remaining eligible mass, including the
unfilled positions at floor $b$, is smaller than $W^{-2}$, the
initial part therefore balances. Two successive such cuts give a
balanced segment. In a segment, a single label at its greatest
two- or three-adic denominator height must have coefficient zero
by reduction modulo that prime. Its closed degree class is a
singleton, so it cannot carry order nine or higher either, by
square-unit covariance. It may be removed and the argument repeated.
The other eligible labels can be assigned optimistically to quartets
in their height classes, using full groups first and then the
largest remainders. The same quartet count without the segment
test applies in a mixed $2,5$ alternative.

For four-prime even square support, fix the remaining odd prime
before enumerating its supported denominators. At an isolated $p$,
put $E_p=\sum_{b_i=p^a}b_i^{-2}$ and let $C$ be the sum of the
least squared Weyl coordinates at the observed mixed multiples of
$p$, as in Proposition~\ref{new:local-entry-mass}. The pure entries
use exactly $E_p$ in \eqref{eq:component-norm}, so
\begin{equation}\label{eq:prefix-component-cost}
 C<\frac{p+1}{p-1}E_p,
\end{equation}
since $0<g_p\le pE_p/(p-1)$.

There is an integer refinement when fewer than $p^2$ labels are
active. At each complete pure class $b$, equal primitive masses
have unique base-$p^2$ digits; an order-$p^2$ square-unit packet
has size $p(p-1)/2$, and higher orders exceed this activity bound.
Thus that class contributes $pc_b/b^2$ to $g_p$, with
$c_b\in\ZZ_{\ge0}$ and $(p-1)c_b\le n_b$. Subtracting its full
pure mass from $2g_p-g_p^2$ gives
\begin{equation}\label{eq:prefix-closed-component-cost}
 C\le\sum_{b\ {\rm closed}}
       \frac{2p\lfloor n_b/(p-1)\rfloor-n_b}{b^2}
       +\frac{(p+1)\lfloor H/(p-1)\rfloor}{B^2}.
\end{equation}
Here $H$ bounds all remaining pure positions, including an open
class, and $B$ is their least possible denominator. Each further
packet uses at least $p-1$ of its $2p$ norm units on its own pure
mass, which explains the last numerator. The activity bounds above
justify this refinement for $p\ge7$ and for isolated five with at
least nine fours. Both estimates use the observed mixed terms and
an upper bound for the unfilled pure mass.

Suppose five is present, $m\ge9$ fours occur, and a twist is active
at two and three. Reserve the vacuum, the fours, five fives, four
proper five-free labels, two further odd labels and eight ternary-active
positions. At most $H=r-m-20$ pure positions remain at $p\ge7$.
If $p$ is absent as a denominator and $H\le2(p-1)$,
\eqref{eq:unequal-pure-shell} forces $p$ copies at the first pure
level $B=p^c$, $c\ge2$. Then $M_p=B^2/p$ and
$M/M_p\equiv-\epsilon\pmod{M_p}$ by Proposition~\ref{new:component-column}.
Reserve the remaining copies and test this congruence; larger pure
levels are allowed after the class is complete.

Suppose $r\le43$, denominator seven occurs, five does not, and
$m_4+m_6+27>r$. Seven is isolated, and its class has seven members
or at least fourteen by Theorem~\ref{new:one-packet}. A single packet also
has at least three proper seven-free labels. All have trivial twist:
their denominators are divisible by the isolated seven, and the
seven-active labels already belong to its packet. A mixed $2,5$
phase would require sixteen additional five-active labels, disjoint
from the fours, sixes and the above reservations. Including the
vacuum gives at least $1+m_4+m_6+7+3+16=m_4+m_6+27>r$ positions. With two or more
packets the fourteen sevens alone suffice. Thus five is isolated
whenever it occurs in the dimension support.

At denominator seven the multiplicity is $7j$ with $j=1,2,3$.
For each permitted $j$ the remaining squared-coefficient budget is
$C=j(7-j)-\sum a^2$, and the scaled partial signed sum is
$R=7\sum aL/b$, with $L$ the current lcm. The total budget is at most
twelve, so every proper coefficient is the unique $a\in\{\pm1,\pm2,\pm3\}$
congruent to $(b/7)^{-1}$ modulo seven. A negative $C$ is impossible.
For next-denominator floor $b_*$ let $d_\pm$ be the least $|a|k$ with
$1\le|a|\le3$, $k\ge\lceil b_*/7\rceil$ and $k|a|\equiv\pm1\pmod7$.
Since $|a|/k=a^2/(|a|k)$, the remaining positive and negative
contributions are at most $CL/d_+$ and $CL/d_-$. If $R>0$ then
$Rd_-\le CL+\lceil7jLd_-/M_{\min}\rceil$ for a lower bound $M_{\min}$
on $M$, with the analogous inequality at $R<0$. The remaining cost
and slots must also permit at least three distinct proper degrees in
total. The covering enumeration admits the weaker budgets in
\eqref{eq:weyl-cost}, which include all the cases just described.

For a single packet at $p=5,7$, write $b_i$ for the proper $p$-free
denominators and $a_i$ for their Weyl coefficients. Their lcm is $M$,
$\sum_i a_i/b_i=\epsilon/M$ and $\sum_i a_i^2=p-1$ by
Theorem~\ref{new:one-packet}. For a nonempty proper initial segment
put $W=\lcm_{i\le s}b_i$ and $S=\sum_{i\le s}a_iW/b_i$.
Since $(b_i/p)a_i\equiv1\pmod p$,
$S\equiv (W/p)\sum_{i\le s}a_i^2\ne0\pmod p$.
If the remaining absolute coefficients sum to $C$, then
$b_{s+1}\le(C+1/6)W/|S|$ as $M\ge6b_{s+1}$.
For the last witness $z$, take $d=W/\gcd(W,z)$ and $t=z/\gcd(W,z)$.
Then $M=Wt$, $z=Wt/d$ and $St+a_zd=\pm1$, so testing divisors
$d\mid W$ lists the possibilities. At fixed support the candidates are products of powers of the four primes, respecting the exact valuations at occupied prime denominators. The test applies to the four signs at five and either shape at seven. The least admissible first witness also bounds the floor of every outstanding witness. Order constraints, checked integer widths and a complete list are required for exclusion; an unfinished list retains the prefix.

When both classes have a single packet, there is exactly one common
proper label. Indeed each family also has separate labels omitting
one of the other two primes, by its gcd condition and four-prime
support. Thus at most two proper labels coincide, at most one in
the three-label seven case. For isolated five and seven the cross
identity following Proposition~\ref{new:component-column} gives
$\sum_{35\mid b_i,\ i\ne0}a_{5,i}a_{7,i}=-1$, since $g_p=1/p$
and the Weyl sign $\gamma^3$ is common. If both families consist of
signs, parity leaves one common label; in the three-label case the
sum likewise forces one with opposite signs. Writing its denominator
as $35k$, the two inverses have opposite signs, hence $k\equiv\pm17\pmod{35}$.
For a root fixed by the four five-free witnesses, the shared label has $a_7=\pm1$ and fixes its seven-coefficient by the opposite-sign rule. All other proper seven-free degrees have the form $5d$ with $d\mid M/35$. Their coefficients are five signs, or one sign and one of magnitude two. Testing the signed sum $\sum_Ya_{7,Y}d_Y=\epsilon$ uses the same sign as at five: $M/p\equiv-\epsilon\pmod p$ at both primes.

For the general tensor-square estimate let $x=d_X>1$ and $p\mid x$.
Set $a_p=\min_{p\mid y}v_p(y)$ over the nonunit degrees and let $c_p$
be the least sum of nonunit degrees prime to $p$ congruent to $-1$
modulo $p^{a_p}$. Reduction of $XX^*-1$ gives a contribution at least
$c_p$ from these degrees. Thus $x^2-1\ge\sum_p w_pc_p$ whenever
$w_p\ge0$ and $\sum_{p\nmid y}w_p\le1$ for every nonunit degree $y$;
a missing residue separately excludes the list.

At closure apply Lemmas~\ref{lem:basic-table} and~\ref{lem:deletion},
the dimension equation modulo four and five, and the dimension
partitions of products \cite[\S6.4]{ABPP}. In particular at least
three nonunit degrees omit two, and at least three omit five.

For distinct primes $t,p\mid M$, Lemma~\ref{blk:witness}
requires a $t$-free nonunit at denominator divisible by
$t^{v_t(M)}p$. With three support primes the same lemma gives
$p^2\mid M$ for every $p\mid M$, and different omitted primes
have disjoint nonunit witnesses.

For three-prime support, let $L$ be the prefix lcm and
$a_t=\max(t^2,t^{v_t(L)})$. If $t^{v_t(M)}=a_t$, each prime $p\ne t$ missing from the witnesses at this height requires a suffix multiple of $pa_t$; if the height rises, all $t$-free witnesses must come from the suffix. Use the weaker of the two reservations. At least two nonunit labels are $t$-free, at least three for $t=2,5$; their number is $3\bmod4$ at two and $2\bmod3$ at three by the dimension equation and Lemma~\ref{lem:deletion}. The reservations for different $t$ are disjoint. One for each $t$ must reach the largest outstanding $pa_t$, the others at least $a_t$. A $t$-power codegree also separates these witnesses from activity at other odd primes (Corollary~\ref{blk:primepower}). Under Proposition~\ref{new:separated-primary}, once the eligible prefix labels plus the remaining positions fall below its threshold, use odd multiples of $p$ in whole packets of $h$, or $qh$ if $q\nmid b$.

With exactly three primes, every prime omitted by a nonunit degree is
unique. If $d\le e\le a$ are the three least nonunit degrees counted with
multiplicity, the decomposition of $XX^*$ for $d_X=d$ must contain
a constituent omitting each prime of $d$, by reduction of $d^2-1$
at that prime. Two such constituents are distinct and different from
$X$; hence $e+a\le d^2-1$.

At an odd $p\mid M$ with $p\nmid q$, a phase of order at least
$p^2$ needs $p(p-1)/2$ distinct nontrivial $p$-phases by square
Galois symmetry. Each has at least $B(q)$ labels, as does the
trivial phase. Consequently
\[
 r<B(q)\bigl(1+p(p-1)/2\bigr)
 \quad\Longrightarrow\quad \operatorname{ord}(\theta_{i,p})\le p
 \quad\text{for every } i .
\]


For $q=3$ the fiber with all unramified odd phases trivial has at least
four labels. Indeed its intersection with $P_3$ contains the unit.
A twist of $3$-part at least nine would supply three other labels.
Otherwise the order-three mass equals twice the inactive mass in
this intersection. If there were at most three labels, the unit
could not be the sole inactive one, since the requisite active mass
would be two. The only remaining possibility is an inactive nonunit
degree $a$ and an active degree $b$ with $b^2=2(1+a^2)$.
Modulo four and three this makes $a$ odd and both degrees prime to
three, and the equation gives $\gcd(a,b)=1$. Both labels belong to
$P_3$, so their $S$-entry is nonzero modulo three after dividing by
the row degree, contrary to Lemma~\ref{lem:basic-table}.

For $q=3$ and $J=\{p\mid M:p\ne 2,3\}$, the support and higher-order bounds improve to $r\ge4+3\sum_{p\in J}(p-1)$ and $r\ge4+3p(p-1)/2$, respectively.

\paragraph{Testing the principal block.}
A candidate principal block contains the vacuum. Its nonunit degrees
have gcd one by Lemma~\ref{blk:witness}. They must satisfy the height
condition of Lemma~\ref{blk:height} and the character bounds of
Proposition~\ref{new:block-census}. If a $p$-power codegree occurs,
Corollary~\ref{blk:primepower} forces every label of such codegree
and every $p$-free label into the block. For each other eligible degree
of multiplicity $m$, allowing any count from zero to $m$ gives a
bounded subset-sum test for the block mass. For a prime-power row we also test the simultaneous class traces and
squared norms from Section~\ref{sec:dimensions}; failure at any degree or prime
excludes the type. Omitting a test retains the type.

\begin{lemma}[Partial Gauss deficits]\label{r4:prefix}
Suppose $D=M^2$. Let $J$ be a nonempty union of odd components of the twist
graph. If $M/M_J>2^{|J|}$, then for every union $S$ of complete nonunit
degree classes,
\[
 R_S=1-\sum_{i\in S}(1-\langle\theta_i\rangle)/b_i^2\ne0 .
\]
In particular the assertion holds at any isolated odd prime of a
non-pointed based-simple datum, in every rank.
\end{lemma}
\begin{proof}
Here $\langle\ \rangle$ denotes the rational Galois mean.
The twist sum on a degree class is an algebraic integer fixed by all
square multipliers: the Galois permutations preserve degree and
$\gamma^{a^2-1}=1$. Its mean has denominator a power of two,
since the quotient of the Galois group by its squares is a $2$-group.
On classes whose nontrivial twists are supported on $J$, that
denominator divides $2^{|J|}$: the units modulo each odd prime
power are cyclic. Any class active on $J$ has denominator supported
on $J$, and a class active outside $J$ has denominator prime to
$M_J$, by the component splitting in Section~\ref{sec:small}.
This also makes the two kinds of active classes disjoint.

If $R_S=0$, the remaining deficits are nonnegative and sum to
$-\xi/M$, hence to $1/M$. Write this sum as $P+Q$, with $P$ the
part active on $J$. The denominator of $2^{|J|}P$ is supported on
$J$, while $Q$ is integral at each prime of $J$. Thus the denominator
of $2^{|J|}P$ divides $M_J$. If $P>0$, then
$P\ge1/(2^{|J|}M_J)>1/M$. If $P=0$, the denominator of
$Q=1/M$ gives the contradiction. For a single isolated prime the
cofactor exceeds two by the three-prime bound.
\end{proof}

Use \eqref{eq:prefix-component-cost} and
\eqref{eq:prefix-closed-component-cost} for the mixed-entry and
pure-tail bounds. As degree classes close, apply
Lemma~\ref{r4:prefix}, retaining all possible deficits; below
$p^2$ active labels use the integral class coefficient $pc_b/b^2$
proved above. At other classes use the rational root means, the
known class traces and $[0,2/b^2]$ per unassigned label. Twice the
remaining reciprocal-square mass bounds the unassigned deficit,
with terminal error $\xi/M$.

\begin{lemma}[Equal-mass packet rigidity]\label{new:packet-rigidity}
Let $p$ be an odd prime of even positive valuation in $D$, and put
$h=(p-1)/2$. List in increasing order the nonunit degrees
$d_1<\cdots<d_s$ of multiplicities $m_j\ge h$ for which
$p\mid D/d_j^2$. Suppose that
\[
 m_j<ph\quad\text{whenever }p^2\mid D/d_j^2,
 \qquad
 d_j^2>\sum_{i<j}\lfloor m_i/h\rfloor d_i^2.
\]
Then every primitive $p$th-root phase contains at most
$\sum_j\lfloor m_j/(p-1)\rfloor$ labels. In particular this sum
is at least $B(q)$, where $q$ is the square-free part of $D$.
\end{lemma}
\begin{proof}
A $p$-primary twist of order greater than $p$ would require
at least $ph$ equal-degree labels by square-unit covariance,
contrary to the first hypothesis. Write $a_j,b_j$ for the
numbers of degree-$d_j$ labels at a phase in the two square classes
of primitive $p$th roots. Covariance makes these numbers independent
of the phase within its class, while the equal-mass identity gives
$\sum_j a_jd_j^2=\sum_j b_jd_j^2$ and
$a_j+b_j\le\lfloor m_j/h\rfloor$.
At the largest index with $a_j\ne b_j$, the second hypothesis
contradicts this equality. Thus $a_j=b_j\le\lfloor m_j/(p-1)\rfloor$.
Some nontrivial $p$-phase occurs, and Proposition~\ref{blk:packets}
gives the final inequality.
\end{proof}

\subsection{Verification of the rank bounds}\label{sec:finish}


\paragraph{Excluding a complete list.}
Apply five tests in order. First, if $d$ is the least nonunit degree and
$E_p=\sum_{i\ne0,\,p\mid f_i}d_i^2$, Lemma~\ref{blk:activity} requires
$2pdE_p\ge(p-1)D$ for every $p\mid D$.

Second, at a prime denominator $p\nmid q$, put $n=M/p$. For $q=1$ this
gives a row of codegree $p^2$. For $q=2$, or for $q=3,5$ when neither
$2n$ nor $n/2$ occurs, the defect-one block of the ramified prime has
equal degrees and gives an effective row of codegree $p^2$.
Corollary~\ref{new:prime-free-square-trace} then requires gcd one
among the nonunit degrees prime to $p$.
Third, tile the $p$-free denominator classes at each $p\mid q$
using Lemma~\ref{blk:defectone}, as in the prefix search.

Fourth, let $A_p$ consist of the nonunit degrees $y$ with $p\mid f_y$ such
that, for every $p$-free nonunit degree $x$ of multiplicity $m$,
\[
 g=(x,y)>1,\qquad 2mg>p,\qquad m\,\lcm(x,y)^2<D .
\]
For a principal row, Proposition~\ref{new:block-census} gives
$g>1$, $|A_x|<mg$, $A_x\equiv mg\pmod p$ and $Q_x\ge m$.
The trace conditions require $2mg>p$, and the norm identity gives
$m(x/g)^2\le f_y-1<f_y$. Thus $A_p$ contains every nonunit
degree of $P_p$, and Lemma~\ref{blk:witness} gives $\gcd A_p=1$.

The remaining lists have $q=1$ and a denominator five. By
Corollary~\ref{new:five-reduction} there are four proper five-free
nonunit labels, i.e. those of degree $d\ne M/5$ prime to five.
On these an active five-row has signs $s_d\equiv d\pmod5$
by \eqref{eq:packet-trace} and principality. It vanishes on
five-divisible degrees; its own degree class lies outside $P_3$
by principal height. Thus some subset of the four labels with
$d\in A_3$ must satisfy $1+\sum s_d d\equiv0\pmod{M_3}$
by \eqref{eq:principal-trace}.

The \supfile{census/reference-census.zip}{covering archive} contains
the generation sources and outputs of $12\,273$ completed tasks,
including $41\,915$ arbitrary-precision continuations. The
\supfile{census/audit.py}{auditor} checks the task partition, completion
records, arithmetic identities and unique ownership of the lists.
The \supfile{census/five_tests.py}{replay}
applies the five tests to every list and verifies both the
\supfile{census/rank-counts.json}{per-rank totals} and the
\supfile{census/rank-dispositions.jsonl.gz}{individual exclusions}.
The first two columns of Table~\ref{tab:counts} refer to this fixed
covering census. A generator applying further necessary conditions
may emit fewer lists. Such an early rejection does not change the
covering census or its terminal totals.

The stronger conditions on degree slices, common support and fusion
multiplicities also give a separate exclusion of the same lists.
In the order slices, overlap, fusion, prime-square structure, block
entries, simultaneous traces and tiling, the successive survivor counts
are $393,284,215,143,124,12,0$
(\supfile{census/recheck.py}{seven-test comparison}).

\paragraph{Completeness and verification.}
The coverage argument has three steps. First, the inequalities at the
start of this section give finite ranges for $q$, $B$, the maximal
multiplicity and the dimension-prime support. Each possible datum
belongs to one of these ranges. Second, its ordered nonunit
denominators satisfy \eqref{eq:interval} at every prefix. The exact
update of $(L,E)$ preserves the residual dimension equation, and the
terminal divisor equations recover all possible last degrees. Third,
the restrictions in Sections~\ref{sec:small}--\ref{sec:verify} discard
a prefix only when an upper bound for its possible completions fails
a necessary condition. Unknown valuations, overlapping reservations
and unassigned phases are allowed every possibility in taking these
bounds. Thus a possible datum reaches a complete list. An arithmetic
width limit transfers the exact unfinished state to the
arbitrary-precision recursion; a work limit on an optional test leaves
the state undecided. Neither operation is an exclusion.

The \supfile{proof-map.md}{implementation map} associates each pruning
predicate with the statement or displayed inequality proving its
necessity, including its extra hypotheses. The
\supfile{verification/native_boundary/check.py}{arithmetic-boundary check}
compares bounded continuations with exact arithmetic across the native
handoff. These checks concern the implementation of the enumeration.
The completion records identify the source snapshot recorded for the covering run;
checking them does not itself repeat that run.

A separate \supfile{verification/independent/check.py}{independent verifier}
reconstructs $D=qM^2$, every denominator and each of the five tests
directly from the archived degree lists, without importing the
enumeration or terminal-test modules. It writes an
\supfile{verification/independent/terminal-witnesses.jsonl.gz}{exact
first-obstruction witness for each of the $18\,199$ lists} and
reproduces Table~\ref{tab:counts}. It also independently counts the
initial odd profiles and excludes the last nonsquare nine-pair equation
in $172$ memoized states. Its
\supfile{verification/independent/certificate.json}{certificate}
identifies the inputs, checked claims and outputs. The odd square
recursions have their own exact replay certificates in
Section~\ref{sec:odd-new}. Terminal witnesses certify the listed
exclusions; the mathematical pruning implications and the source-bound
enumeration supply the separate coverage step.

\begin{proof}[Proof of Theorem~\ref{alg:rank}]
The preceding coverage argument puts every putative degree multiset
through rank $44$ in the covering census. Each of its lists violates
one of the five necessary conditions proved above, as recorded in
Table~\ref{tab:counts} and the individual witnesses. Hence no datum in
the stated scope has rank at most $44$. The odd-dimensional assertion is
Theorem~\ref{odd:rank}.
\end{proof}

\paragraph{A partial cover at rank $45$.}
A start is $(q,B,m,S)$, where $S$ is the prime support of $M$;
adjoining the maximal twist supports $\mathcal F$ of
Corollary~\ref{new:sector-rank} gives a refined start.
The \supfile{verification/sector_frontier/sector_prefix.py}{prefix replay}
closes two further denominator classes and retains $204$ starts and
$462$ refinements. Of these, $167$ starts and $417$ refinements have
square dimension and least denominator four. The
\supfile{verification/sector_frontier/sector-prefix45-depth2.json}{certificate}
records the tuples and intermediate counts. A retained start is
undecided; this calculation supplies no exclusion of rank $45$.

The initial ranges are $8\le qB^2\le43$, the multiplicity bounds
above, and $p\le45$ for $p\in S$: allowed primes satisfy
Theorem~\ref{thm:spherical}, while primes dividing $q$ satisfy the
first bound.
Apply Corollary~\ref{round:eight-rank}, the prime-denominator and
defect-one restrictions, and
\[
 C\sum_{\substack{p\in S\\p\nmid q}}p\le45,\qquad
 C\prod_{p\in F}p\le45\quad(F\in\mathcal F).
\]
Lemma~\ref{new:reservation-calculus} combines the floors for the
vacuum, the $m$ maxima, phase sectors and prime cuts. At an isolated
odd prime in square dimension the $p$ pure nonunit labels are
disjoint from all crossing witnesses and the vacuum.
Prescribed equal positions must also meet their floors.

For $t\in S$, let $k_t$ be the least $k\ge2$ for which $-1$ is a
sum of $k$ unit squares modulo $t^2$. Deletion and the dimension identity
require $k_t$ nonunit $t$-free degrees. By Lemma~\ref{blk:witness},
one has denominator divisible by
$t^{v_t(M)}\max_{p\in S\setminus\{t\}}p$.
A nonunit degree contains at least two primes, so it can pay at most
$|S|-2$ omission charges; apply the charged-floor rule of
Lemma~\ref{new:reservation-calculus}. When $|S|=3$,
Lemma~\ref{blk:witness} gives $v_t(M)\ge2$ for every $t\in S$.

For unramified odd primes $J$, the active union has size at least
\[
 \max\left\{
 B(q)\sum_{p\in J}(p-1),\
 C\sum_{F\in\mathcal F}
   \left(\prod_{t\in F}t-\prod_{t\in F\setminus J}t\right)
 \right\}.
\]
The second term sums the exact sectors of
Corollary~\ref{new:sector-rank}. The eligible active union is
Galois-stable, and square-unit covariance makes each of its orbits
divisible by $(p-1)/2$ for some eligible $p$.
A closed class of denominator $b$
and size $m$ receives the semigroup credit
$Q(m;((p-1)/2)_{p\in J,\ p\mid b})$ from
Lemma~\ref{new:reservation-calculus}; any unfilled demand has
denominator divisible by some $p\in J$.
At odd rank and even dimension both parity classes of nonunit
degrees have odd size. A Galois $2$-Sylow therefore fixes a label
of each parity. These labels are self-dual and inactive at every
$p\equiv1\pmod4$, whose nontrivial square-unit orbit has even size.
When every $p\in J$ is $1\bmod4$, reserve these two labels before
crediting each class's remaining capacity. Their denominators
have respectively the full and a smaller dyadic part of $M$.
All placements in the prefix or suffix are allowed, and their
bounds are combined by coordinatewise minima.

For several twist components, perform the same accounting at each
suffix floor $d$. Within each component use the simultaneous
reservations of all nonempty exact sectors, and separately the
disjoint odd-prime activity quotas and each mixed-pair bound in
Proposition~\ref{round:mixed} and
Corollary~\ref{sq:fiber-squares}; also use the pure-label quota at
an isolated odd prime in square dimension. In each reservation
retain only groups whose divisibility forces denominator at least
$d$, and subtract their eligible closed-class semigroup credits.
For nonempty exact active support $K$, independent square-unit
multipliers give orbit size divisible by
$\prod_{p\in K,\ p\ {\rm odd}}\bigl((p-1)/2\bigr)$; for activity
at either member of a mixed pair, either single-prime size is allowed.
Take the maximum residual demand within each component, then add
across components, since their nonempty active
supports are disjoint and an isolated pure label cannot be active
elsewhere. The vacuum is separate. These threshold counts bound
every suffix order statistic.

The replay uses \eqref{eq:interval}, all multiplicities of each next
class, and the least eligible $S$-smooth multiples as floors. It
tiles ramified classes and applies the prime-denominator and
sparse-height bounds. At an isolated $p$, the other facets reserve
outside the vacuum by Lemma~\ref{lem:vac-reserve}, since a pure
$p$-active label is ineligible for them. Thus at most
$44-\sum_{F\ne\{p\}}\prod_{t\in F}t$ labels are $p$-active.
If $q=\ell$ is prime and an odd prime denominator $p\ne\ell$ has known
$v_p(M)\ge2$, a full equal-degree $\ell$-block would give a prime-square
row of $p$-divisible degree, contrary to Theorem~\ref{rev:height}.
With four support primes in nonsquare dimension, at most one
unramified prime has valuation one; each possible exception is kept
fixed along a branch. Finally, for each closed pure-codegree class
write $M=Lt$ and test the integral slice bound at every residue of
$t$ modulo the principal prime, excluding zero only at a known
exact valuation. Each refined start is counted once, at its first
admissible prefix.

\section{Finite rotation conditions}\label{formal:lean}\label{formal:finite}

The finite inputs are the based ring, degree character, Drinfeld map
and congruence representation of Sections~\ref{alg:framework}--\ref{alg:modular}.
For the rank bounds, the rotation certificates are used through the
following finite consequences.

\begin{proposition}[Finite rotation criterion]\label{formal:criterion}
In Theorem~\ref{alg:rank} one may replace rotation certificates by the
following tests. Write $\nu_j=L_j^0$. At each $p\mid D$ fix a place
$\mathfrak p\mid p$ and a positive $k_p$ congruent to zero on the
$p$-part of $L$ and to one on its complement.
\begin{enumerate}[label=\textup{(F\arabic*)},leftmargin=3em]
\item For each full residue block $B$ at $p$,
$\sum_{z\in B}s_{iz}\overline{s_{jz}}=0$ if
$\theta_{i,p}\ne\theta_{j,p}$.
\item For each $x_c$, the numbers $\nu_d(x_c)$ are algebraic integers
for $d\mid N$, and
$\nu_{k_p}(x_c)^p\equiv\nu_{k_p/p}(x_c^p)\pmod{\mathfrak p}$.
For every set $J$ of allowed odd primes fix $k_J>0$ equal to one
on the $J$-part of $L$ and zero on its complement. Require the same
Frobenius congruence at $k_J$ for each prime $\ell\mid M$ outside
$J$ with $\ell\nmid q$ (using the fixed place above $\ell$).
\item For each $p,c$, the number $\nu_{N_{p'}}(x_c)$ is rational.
\end{enumerate}
\end{proposition}
\begin{proof}
Lemmas~\ref{blk:local} and~\ref{alg:rotation-arithmetic} give (F1)--(F3)
from rotations. Assume (F1)--(F3). Condition (F1) gives local invariance.
For $d=\gcd(j,N)$
choose a unit $b$ modulo $N$ with $j\equiv bd\pmod N$. The rational
weights $s_{0u}s_{0v}=d_ud_v/D$ in \eqref{alg:rotation} give
$\nu_j(x_c)=\sigma_b\nu_d(x_c)$, an algebraic integer by (F2).
Linearity and Frobenius extend the congruence to integral algebraic
coefficients. The identity \eqref{eq:rotation-entry}, (F1) and this integrality
give Lemma~\ref{blk:admissible}, Proposition~\ref{new:component-column}
and the effective-row Weyl identities. Condition (F3) supplies the
rational integers in Lemma~\ref{blk:activity}.

The same Frobenius congruences give the divisibilities in
\eqref{eq:phase-column} and Proposition~\ref{phase:moments}.
These are all the rotation consequences used in the exclusions, so
the remaining hypotheses of Theorem~\ref{alg:rank}, together with
(F1)--(F3), imply its rank conclusions.
\end{proof}

\section{Applications to modular fusion categories}\label{ext:classification}

\begin{proposition}[Categorical realization]\label{alg:realization}\label{cat:realization-section}
For an integral modular fusion category $\C$ over $\mathbb C$, the
positive spherical structure makes $K_0(\C)$ rotation-admissible. Its
symmetric Drinfeld map pairs $[X_i]/d_i$ with $[X_j]/d_j$ by
$S_{ij}/(d_i d_j)$. Perfection and simplicity of $\C$ correspond to
perfection and based-simplicity of $K_0(\C)$. A nontrivial perfect
$\C$ has at least three prime divisors of its dimension.
\end{proposition}
\begin{proof}
Use pseudo-unitarity \cite[Propositions~8.23--8.24]{ENOannals} and
Verlinde duality for the stated pairing
\cite[\S8.14]{EGNO}, \cite[\S10]{BPBurnside}. The exponent is $N$ and has
the prime divisors of $D$ \cite[Theorem~7.7]{NSfs},
\cite[Theorem~3.9]{BNRWFinite}. Since $\xi^8=1$, take $\gamma=\xi^3$.
The congruence and Galois theorems supply $\rho$ and covariance on $W$,
and $\sigma_a(T/\gamma)=(T/\gamma)^a$ gives covariance on $U$
\cite[Theorem~II, Remark~4.5 and Proposition~6.7]{DLNCongruence}.

Under $Z(\C)\simeq\C\boxtimes\C^{\rm rev}$ the central insertion
$(X_i,X_i^*)$ has twist one. Its rotation on
$\operatorname{Hom}(X_i\otimes X_i^*,V^{\otimes m})$ has $m$th power
one, and its $t$th power for $t\mid m$ has the one-step trace for
$V^{\otimes t}$ on $m/t$ strands \cite[Lemmas~2.5, 2.7 and Proposition~2.8]{NSind}.
The indicator formula has entry $s_{ia}\overline{s_{i^*b}}=s_{ia}s_{ib}$
and twist $\theta_a/\theta_b$, giving precisely \eqref{alg:rotation}
\cite[Corollary~5.6 and Propositions~5.7, 6.1]{NSind}.
Finally, dimension one detects invertibles and based subrings detect
fusion subcategories. A nontrivial category with at most two prime
divisors is solvable and has a nonunit invertible
\cite[Theorem~1.6 and Proposition~4.5(iv)]{ENOWGT}, \cite[Proposition~5.2]{ABPP}.
\end{proof}

This proves Corollary~\ref{main:rank}. Every nontrivial perfect modular
category of odd integral dimension is a Deligne product of non-pointed
simple modular categories \cite[Corollary~4.2 and Remark~4.3]{CPodd}.
Since rank is multiplicative under Deligne products, its rank is at least
$217$, and at least $217^2$ if it is nonsimple. The perfect category
$\mathcal Z(\mathrm{Vec}_{A_5})$ illustrates the role of simplicity in
the arithmetic restrictions. Its electric objects of degrees three
and four have $S_{ij}=12$, and codegree nine and denominator five both
occur although $25\nmid60$.

\begin{corollary}\label{ext:perfect}
Up to braided equivalence, the nontrivial perfect weakly integral
modular categories over $\mathbb C$ of rank at most $44$ are the following.
Write $\Gamma=\mathrm{SL}_2(5)$ and let $\mathcal P$ be the pointed
modular category on the even-weight vectors of $\mathbb F_2^5$ with
twist $q(v)=(-1)^{|v|/2}$, where $|v|$ is the Hamming weight.
\[
\begin{array}{r r l r}
\toprule
\text{Rank}&\text{Dimension}&\text{Category}&\text{Classes}\\
\midrule
22&3600&\mathcal Z(\mathrm{Vec}_{A_5}^{\omega})&18\\
32&28\,224&\mathcal Z(\mathrm{Vec}_{\mathrm{PSL}_2(7)}^{\omega})&12\\
37&7200&\text{semion centralizer in }\mathcal Z(\mathrm{Vec}_\Gamma^{\omega})&36\\
44&129\,600&\mathcal Z(\mathrm{Vec}_{A_6}^{\omega})&24\\
44&57\,600&\text{$A_5$-gauging of }\mathcal P&18\\
\bottomrule
\end{array}
\]
In the third row $[\omega]|_{Z(\Gamma)}\ne0$. The last row uses the
action permuting the five coordinates. Its categories are weakly
group-theoretical and are not group-theoretical. All other rows are
group-theoretical.
\end{corollary}

Here $\mathrm{Vec}_G^\omega$ is the category of $G$-graded vector spaces
with associator $\omega$. Thus $22$ is the least rank of a nontrivial perfect weakly integral
modular category, answering the rank part of \cite[Question~11.8]{BPBurnside}. The dimension-$7200$
family contradicts the dimension classification in \cite{XuDong}.
The issue in \cite[Lemma~3.1]{XuDong} is that a fixed invertible may
have a nontrivial projective stabilizer multiplier and hence no
invertible lift \cite[Corollary~2.13]{BN}. For example, the perfect category
$\operatorname{Rep}(\mathrm{SL}_2(5))$ contains $\operatorname{Rep}(A_5)$
and de-equivariantizes to $\operatorname{Rep}(C_2)$. Its nonunit label
is fixed by $A_5$, but has lifts of dimensions $2,2,4,6$
\cite[Table~2]{KacTodorov}. This disproves the lemma under its stated
hypotheses, which do not require nondegeneracy. The nonunit semion
label in the perfect rank-$37$ family realizes the same obstruction.
\medskip

\subsection{Rank in the equivariantization}\label{ext:equiv-rank}
A perfect weakly integral modular category is integral
\cite[Proposition~10.11]{BPBurnside}.
By Corollary~\ref{main:rank}, every nontrivial perfect category
of rank $r\le44$ is nonsimple. We classify these categories by
counting the simple objects in an equivariantization.
Let $\mathcal E$ be a minimal nontrivial fusion subcategory. It cannot
be nondegenerate: it would be non-pointed and simple, contrary to
Corollary~\ref{main:rank}. Its nontrivial M\"uger center is a fusion
subcategory of $\mathcal E$, so minimality makes $\mathcal E$ symmetric.
Write $\mathcal E=\operatorname{Rep}(H,z)$ in the super-Tannakian
classification \cite[\S9.9]{EGNO}. If $z\ne1$, then
$\operatorname{Rep}(H/\langle z\rangle)$ is a proper Tannakian
subcategory; minimality forces it to be trivial, leaving
$\mathcal E=\mathrm{sVec}$, contrary to perfection. Thus
$\mathcal E=\operatorname{Rep}(G)$ is Tannakian. Minimality and
perfection make $G$ nonabelian simple, as in
\cite[Corollary~5.4]{LPR}. Its de-equivariantization
$\mathcal F=\C_G$ is integral and faithfully $G$-graded, with modular
neutral component $\mathcal B$, $\mathcal F^G\simeq\C$ and
$\FPdim(\C)=|G|^2\FPdim(\mathcal B)$
\cite[Propositions~2.21, 4.26, 4.56 and Corollary~4.27]{DGNO}.
Indeed, a simple $Y$ of $\mathcal F$ with stabilizer $G_Y$ has a lift
of dimension $[G:G_Y]\dim(\pi)\FPdim(Y)$ for each irreducible
projective stabilizer representation $\pi$ \cite[Corollary~2.13]{BN}.
Thus $\FPdim(Y)$ is a rational algebraic integer. The neutral rank can increase: for the center of $\mathrm{Vec}_{C_7\rtimes C_2}$
with inversion action, the quotient onto $C_2$ gives ranks $28$ and
$49$ respectively \cite[Corollary~3.7]{GNN}.

Kirillov's formula gives
$n_g=\operatorname{rk}(\mathcal F_g)=|\Irr(\mathcal B)^g|$
\cite[Lemma~10.7]{Kirillov}, \cite[Proposition~6]{JMNR}.
Let $r_{[g]}$ count simples over the conjugacy class of $g$,
and $t,u$ count nonunit neutral orbits and their simple lifts.

\begin{lemma}\label{ext:sector}
For any twist-trivial $\operatorname{Rep}(G)$ in a modular category,
each $C_G(g)$-orbit in $\Irr(\mathcal F_g)$ contributes at least
$m=\operatorname{ord}(g)$ lifts, and exactly $|H|$ if its stabilizer
$H$ has cyclic $H/\langle g\rangle$. It contributes at least $2m$
lifts if $|H|/m$ is not a square. Consequently
\[
 r_{[g]}\ge\max\{m,m^2 n_g/|C_G(g)|\},\qquad
 r_{[g]}=mn_g\quad\text{if } C_G(g)=\langle g\rangle ,
\]
and, writing $k(G)$ for the number of conjugacy classes,
\begin{equation}\label{ext:rankcontrol}
 r\ge c(G)+u,\qquad c(G)=k(G)+\sum_{[g]\ne[1]}\operatorname{ord}(g),
 \qquad u\ge t\ge\left\lceil(\operatorname{rk}\mathcal B-1)/|G|\right\rceil .
\end{equation}
\end{lemma}
\begin{proof}
The crossed twist identifies $Y\in\Irr(\mathcal F_g)$ with $gY$, so
$\langle g\rangle\subset H\subset C_G(g)$. The lifts correspond to
simple modules of a twisted group algebra of $H$ \cite[Corollary~2.13]{BN}.
On each multiplicity space its basis element $u_g$ acts, up to a
scalar, by the twist of the equivariant simple \cite[Theorem~3.9]{Kirillov}.
Thus $u_g$ is central; normalized by $u_g^m=1$, its powers generate
$\mathbb C[C_m]$ in the center. Its $m$ eigenspace ideals each have
dimension $|H|/m$, as multiplication by $u_g$ cycles the group basis
in blocks of length $m$. A simple complex algebra has square
dimension, giving the nonsquare assertion. For cyclic quotient,
$u_g$ and one generator lift generate a commutative algebra.
Each orbit has size at most $|C_G(g)|/m$. In grade one the unit has
$k(G)$ lifts and each remaining orbit has at least one.
\end{proof}

\begin{lemma}[Crossed traces]\label{ext:crossed-trace}
In the above de-equivariantization, let $g,h\in G$ be commuting
involutions. For a neutral simple $X$ fixed by both, let $\alpha_X$
be its projective stabilizer multiplier. Then
\begin{equation}\label{ext:crossed-trace-formula}
 |\Irr(\mathcal F_g)^h|
 =\sum_{X\in\Irr(\mathcal B)^{g,h}}
       \frac{\alpha_X(h,g)}{\alpha_X(g,h)}.
\end{equation}
If $P=\langle g,h\rangle\simeq C_2^2$, a summand is $+1$ or $-1$
according as $[\alpha_X]|_P$ is trivial or nontrivial. Its sign is
independent of the ordered basis $(g,h)$ of $P$.
\end{lemma}
\begin{proof}
The crossed $S$-operator is $G$-equivariant and identifies the
span of the identity maps of simples in grade $g$ with the sum of
the $g$-twisted lines on the neutral component
\cite[Theorem~6.3, \S9 and proof of Lemma~10.7]{Kirillov}.
On the first space $h$ permutes the identity basis, so its trace
counts fixed grade labels. On the second, permuted lines have trace
zero. For a fixed line choose intertwiners $v_a:aX\to X$ satisfying
$v_a\,a(v_b)=\alpha_X(a,b)v_{ab}$, with the action coherence
understood. Conjugation by the lift of $h$ gives
\[
 v_h\,h(v_g)\,g(v_h^{-1})
 =\frac{\alpha_X(h,g)}{\alpha_X(g,h)}v_g.
\]
Reversing the line orientation inverts this quotient; it is a sign
for commuting involutions. This proves the formula. On $C_2^2$ the
commutator of a multiplier is either trivial or the unique nontrivial
alternating bicharacter, proving the final assertions.
\end{proof}

\begin{lemma}\label{ext:groups}
If $G$ is nonabelian simple and $c(G)<46$, then $G=A_5,\mathrm{PSL}_2(7)$
or $A_6$, with $c(G)=20,29,29$ respectively. In the last two cases $r\ge32+u$.
\end{lemma}
\begin{proof}
For $k=k(G)$, Cauchy's theorem and Burnside's theorem give
\begin{equation}\label{ext:cauchycensus}
 c(G)\ge3k-2+\sum_{\substack{p\mid|G|\\p\ \mathrm{odd}}}(p-2)
 \ge3k+2,
\end{equation}
since at least two odd primes divide $|G|$. Hence $k\le14$.
The complete census is recorded in Table~\ref{tab:simplecensus}.
Its classification inputs are \cite[Tables~I--II]{Landauer},
\cite[\S3]{AK}, \cite{VeraLopez85,VeraLopez86,VeraSang};
we use the explicit combined tables in \cite[\S4]{HM}.
In particular, there are no nonabelian simple groups with thirteen
classes. For $k\le10$, the displayed $c(G)$ are obtained by adding
$k-1$ nonidentity class orders to $k$, using the indicated pages of
\cite{AtlasGroups}; for $k\ge11$, \eqref{ext:cauchycensus}
gives the displayed lower bounds from the group orders on those pages.
Only $A_5,L_2(7),A_6$ have $c(G)<46$.

\begin{table}[ht]
\centering\small
\begin{tabular}{lrrlc}
\toprule
$G$&$k(G)$&$c(G)$&Census entry \cite[\S4]{HM}&ATLAS page\\
\midrule
$A_5$&5&20&$60/5$&2\\
$L_2(7)$&6&29&$168/42$&3\\
$A_6$&7&29&$360/118$&4\\
$L_2(11)$&8&51&$660/13$&7\\
$L_2(8)$&9&62&$504/156$&6\\
$L_2(13)$&9&67&$1092/25$&8\\
$A_7$&9&46&$A_7$&10\\
$M_{11}$&10&68&$M_{11}$&18\\
$L_3(4)$&10&51&$\mathrm{PSL}_3(4)$&23\\
$L_2(17)$&11&$\ge47$&$\mathrm{PSL}_2(17)$&9\\
$\mathrm{Sz}(8)$&11&$\ge50$&$\mathrm{Sz}(8)$&28\\
$L_2(19)$&12&$\ge55$&$\mathrm{PSL}_2(19)$&11\\
$L_3(3)$&12&$\ge46$&$\mathrm{PSL}_3(3)$&13\\
$M_{22}$&12&$\ge52$&$M_{22}$&39\\
$U_3(3)$&14&$\ge46$&$\mathrm{PSU}_3(3)$&14\\
$L_2(23)$&14&$\ge71$&$\mathrm{PSL}_2(23)$&15\\
$A_8$&14&$\ge49$&$A_8$&22\\
$U_3(5)$&14&$\ge49$&$\mathrm{PSU}_3(5)$&34\\
\bottomrule
\end{tabular}
\caption{All nonabelian simple groups with at most fourteen conjugacy
classes. Here $L_d(q)=\mathrm{PSL}_d(q)$ and
$U_d(q)=\mathrm{PSU}_d(q)$; $n/i$ denotes the identifier in the Small
Groups Library used in \cite[\S4]{HM}, pp.~6--10 of the linked author manuscript. The ATLAS page locates both the
group order and its conjugacy classes in \cite{AtlasGroups}.}
\label{tab:simplecensus}
\end{table}

In the last two groups an element $a$ of order four has centralizers
$C_G(a)=C_4$ and $C_G(a^2)=D_8$ \cite{AtlasGroups}. If $n_a\ge2$,
its sector adds at least four to $29+u$. Otherwise $n_{a^2}\equiv n_a=1
\pmod2$ by the action of $a$ on the neutral labels fixed by $a^2$.
The $2$-group $D_8$ fixes a simple in grade $a^2$. For that stabilizer
normalize lifts of $a$ and a reflection by $U^4=V^2=1$.
In $VUV^{-1}=\lambda U^{-1}$, centrality of $U^2$ gives $\lambda^2=1$.
Rescaling $U$ by a fourth root with square $\lambda^{-1}$ gives the
ordinary dihedral relations, hence five lifts in place of two.
\end{proof}

\subsection{The neutral component}\label{ext:neutral}
Put $B=\FPdim(\mathcal B)$. We prove the exhaustive alternatives in
Table~\ref{tab:neutral-cases} using the following orbit form of the
spherical bound. The
\supfile{arithmetic/perfect_counts.py}{arithmetic check} reproduces only
the finite prime-power lists and orbit budgets; the categorical
arguments are given here.
If a finite group acts on the integral modular category $\mathcal B$ by braided autoequivalences
with $t$ nonunit orbits, then
\begin{equation}\label{ext:primeorbits}
 1+t\ \ge\sum_{p\mid B} c_p,\qquad c_p=
 \begin{cases}(p+1)/2& p\text{ odd and }v_p(B)\text{ odd},\\
 p&\text{otherwise}.\end{cases}
\end{equation}
Indeed the permutation action commutes with the modular representation
and fixes the vacuum. Its invariant subspace contains the cyclic
subrepresentation in the proof of Theorem~\ref{thm:spherical}.
For each $p\mid B$, choose a summand witnessing $p$-activity there;
its local factor has dimension at least $c_p$ by that proof.
Within each chosen summand multiply the costs of its assigned primes;
their product dominates their sum. Distinct summands are orthogonal.
Combining with \eqref{ext:rankcontrol} gives $r\ge c(G)-1+\sum_{p\mid B}c_p$. In particular, if $B$ is square the bound is $1+t\ge\sum_{p\mid B}p$.

Every further de-equivariantization below takes place inside $\C$ and
is integral by the same lift-dimension argument. A $G$-invariant
Tannakian subcategory of $\mathcal B$ of dimension $m$ gives a larger
$\operatorname{Rep}(H)\subset\C$, with $|H|=m|G|$ and neutral
dimension $B/m^2$ \cite[Corollary~4.31(iii)]{DGNO}. Such an $H$ is
perfect, since a nontrivial linear character would be an invertible
of $\C$. These facts apply at each successive elementary-kernel step.

\begin{lemma}[Transitive binary subgroups]\label{ext:binary-reductions}
If a finite group acts linearly and transitively on
$V\setminus0$, where $V=\mathbb F_2^d$ and $d\ge3$, then its dual
action is transitive off zero and every invariant quadratic form
$q:V\to\mathbb C^\times$ is identically one. A perfect finite group acting by braided
autoequivalences on a nontrivial integral modular category of square
dimension has at least two nonunit orbits.
\end{lemma}
\begin{proof}
The dual action has the same fixed-space dimensions and hence the
same number of orbits, by Burnside's formula. An invariant quadratic
form takes one value $c$ off zero. For independent $x,y,z$,
$b(x,y)=b(x,z)=b(x,y+z)=c^{-1}$, so multiplicativity gives $c=1$.

If there were a single nonunit orbit, \eqref{ext:primeorbits} would
make the dimension a power of two. Solvability supplies nonzero
$2$-torsion invertibles, which exhaust that orbit. The category is
pointed on $\mathbb F_2^d$. Square dimension makes $d$ even. For
$d\ge4$ the first assertion contradicts nondegeneracy; for $d=2$
transitivity would give a nontrivial image of a perfect group in
$\mathrm{GL}_2(2)\simeq S_3$.
\end{proof}

\begin{table}[htbp]
\centering\footnotesize\setlength{\tabcolsep}{4pt}
\begin{tabular}{@{}p{.07\linewidth}p{.40\linewidth}p{.47\linewidth}@{}}
\toprule
Case & Hypotheses after the preceding reductions & Conclusion or terminal obstruction\\
\midrule
\hyperref[ext:neutral-N]{N} & $G=\mathrm{PSL}_2(7)$ or $A_6$ &
$\mathcal B=\mathrm{Vec}$; the nontrivial prime-power modules give
$r\ge45$, or centralizer ideals of nonsquare dimension $96,192$.\\
\hyperref[ext:neutral-A1]{A1} & $G=A_5$, $n_5>1$; a singleton and free orbits &
$\mathcal B$ is semion.\\
\hyperref[ext:neutral-A2]{A2} & A singleton, a size-$30$ orbit and free orbits &
A Tannakian extension has $c(H)\ge45$.\\
\hyperref[ext:neutral-A3]{A3} & A singleton and a size-$15$ orbit &
A Tannakian subgroup of order $16$ would require $256\mid32$.\\
\hyperref[ext:neutral-A4]{A4} & A size-$6$ orbit, with the permitted size-$15$, size-$30$ and free orbits &
Prime-power congruences, or $r\ge38+10=48$.\\
\hyperref[ext:neutral-B1]{B1} & $n_5=1$; at least three primes divide $B$ &
The orbit budget, \eqref{ext:crossed-trace-formula}, or
$n_3\ge4$ gives a contradiction.\\
\hyperref[ext:neutral-B2]{B2} & $n_5=1$, $B>1$ with at most two primes;
$|\operatorname{Inv}(\mathcal B)[p]|=81$ or $256$ &
Respectively $r\ge52$, or a negative value in
\eqref{ext:crossed-trace-formula}.\\
\hyperref[ext:neutral-B3]{B3} & The elementary subgroup has order $16$ and is transitive off zero &
For its neutral quotient $E$, odd $E$ gives $r\ge50$; even $E$
gives $n_g>8$ or a second Tannakian extension with $r\ge45$.\\
\hyperref[ext:neutral-B4]{B4} & The order-$16$ subgroup has orbits $5+10$ &
Degenerate form: $c(H)\ge47$. Nondegenerate form:
$\mathcal B=\mathcal P$ and $r=44$.\\
\bottomrule
\end{tabular}
\caption{Neutral-component cases for $r\le44$. In rows A1--A4,
$n_5>1$; the complete seven-pattern list appears before A1.
For $G=A_5$, the alternative $B=1$ gives $\mathcal B=\mathrm{Vec}$.}
\label{tab:neutral-cases}
\end{table}

\paragraph{Case N: the groups $\mathrm{PSL}_2(7)$ and $A_6$.}\label{ext:neutral-N}
We show that $\mathcal B=\mathrm{Vec}$ in either case for $r\le44$.
Write $\ell=7,5$, respectively. The two classes of Sylow $\ell$-generators
have cyclic centralizer. The gain of at least three in the order-two and
order-four sectors in Lemma~\ref{ext:groups} is independent of these
classes. Thus $r\ge32+u+2\ell(n_\ell-1)$, giving $t\le12$.
If $n_\ell>1$, a nonunit stabilizer $L$ contains a Sylow
$\ell$-subgroup. Its generators meet both $G$-classes. Each
has cyclic centralizer in $L$, so its centralizer acts trivially by
twisted conjugation on the corresponding basis line of any twisted
group algebra of $L$. Each of the two classes therefore contributes
a central class sum, in addition to the identity. This orbit supplies
at least three lifts, whence $r\ge32+3+2\ell>44$.
Hence $n_\ell=1$. The unique simple in grade $\ell$ has squared dimension
$B$ by faithfulness of the grading, and counting equal-degree cycles
of neutral labels gives $B\equiv1\pmod\ell$
\cite[Theorem~2.3(1)]{SchopierayFPF}. Three prime divisors would cost at least
$2+3+5=10$ invariant dimensions in \eqref{ext:primeorbits}. Since
$t\le12$, the only possibilities are $\{2,3,5\}$ for
$\mathrm{PSL}_2(7)$ and $\{2,3,7\}$ for $A_6$.
For $p=2,3$ a Sylow $p$-subgroup fixes a nonunit label, since otherwise
$B\equiv1\pmod p$. In the first group this gives $n_4,n_3\ge2$
and $r\ge29+9+4+3$. In the second,
$t\ge11$ and a stabilizer containing a Sylow two-subgroup has
order of $2$-valuation three. It cannot have just one projective
irreducible, since its order is nonsquare. Thus $u\ge t+1$ and
$n_4\ge2$ gives $r\ge29+12+4=45$. Thus $B$ has at most two prime divisors.

If $B>1$, solvability gives $V=\operatorname{Inv}(\mathcal B)[p]
\simeq\mathbb F_p^d\ne0$ for some $p$. The $\ell$-element acts freely
off zero. The $G$-action is therefore faithful, and $d\ge3$:
$|\mathrm{GL}_2(2)|=6$, and in odd characteristic $G$ would lie
in $\mathrm{SL}_2$, which has a unique involution. Enumerating prime
powers with $d\ge3$, $p^d\le1+12|G|$ and $p^d\equiv1\pmod\ell$ gives
\[
\begin{array}{c|l}
\mathrm{PSL}_2(7)&2^3,2^6,2^9,3^6,11^3\\
A_6&2^4,2^8,2^{12},3^4,7^4,11^3 .
\end{array}
\]
In characteristic two and $d\ge6$, Jordan form gives
$n_4\ge 2^{\lceil d/4\rceil}\ge4$ and
$n_2\ge 2^{\lceil d/2\rceil}\ge8$.
The four- and two-sectors therefore gain at least twelve and two
over $c(G)$, respectively. Since $2^d-1$ does not divide $|G|$
for any of the listed cases with $d\ge6$, $u\ge2$ and $r\ge45$.
For $3^6$ an order-three element fixes at least nine vectors, adding
$24$ to the baseline. In odd characteristic a faithful
four-dimensional representation has a two-dimensional involution
fixed space: determinant one excludes odd multiplicities of $-1$,
and neither scalar action is faithful. Thus $7^4$ gives an
involution sector of rank at least $49/2$ by Lemma~\ref{ext:sector}.
For $3^4$ the involution sector gains at least three, while both
order-three classes have at least nine fixed vectors and gain six
each, giving $r\ge29+u+3+12>44$.
For $11^3$ an order-four element has a conjugate pair of primitive
fourth-root eigenvalues of product one, and hence fixes a line.
Its sector already has rank at least $44$.

In the remaining cases $|V|=8,16$, respectively. All nonzero orbit
lengths are multiples of $\ell$. Hence $G$ is transitive off zero:
for $A_6$ a decomposition of fifteen would require an index-five
subgroup. Lemma~\ref{ext:binary-reductions} makes $V$ Tannakian and its dual
$W$ transitive off zero. Equivariantization gives
$\operatorname{Rep}(H)\subset\C$ with $1\to W\to H\to G\to1$.
Its stabilizers have order $24$. Choose an element $g$ of order three
fixing a nonzero $w\in W$, and an order-three lift $a$ (by
Schur--Zassenhaus). Then $a,wa$ have orders three and six and give
distinct classes above $[g]$. For $x$ of order four, $C_G(x)=\langle
x\rangle$. For any lift $b$ of $x$, the classes meeting its coset are
indexed by $W/(1-x)W$: conjugation by $W$ gives these cosets, and
$b$ acts trivially on the quotient. There are at least two, each
of order at least four. The extra kernel class and these two sector
gains give
\[
c(H)\ge c(G)+(1+2)+(1+6)+(1+4)=44,
\]
so equality in Lemma~\ref{ext:sector} forces the neutral component
for $\operatorname{Rep}(H)$ to be trivial. In the resulting twisted
double consider the flux of $0\ne w\in W$. It has order two
and centralizer of order $24|W|$ by transitivity off zero.
Since $12|W|$ is $96$ or $192$, Lemma~\ref{ext:sector} gives at least
four lifts in this sector instead of two, contradicting $r=44$.

\begin{lemma}[The deleted permutation action]\label{ext:deleted-action}
Let $V=\{(v_1,\ldots,v_5)\in\mathbb F_2^5:\sum_i v_i=0\}$ with
its natural $A_5$-action. If a pointed modular category on $V$ has
an $A_5$-invariant twist, every braided categorical action inducing
this permutation action is equivalent to the canonical action on the
centralizer of the diagonal semion in $\mathrm{Semion}^{\boxtimes5}$.
Moreover $H^1(A_5,V)=H^2(A_5,V)=0$, and the five-point and ten-point
orbits have respectively four and three simple lifts.
\end{lemma}
\begin{proof}
Write the polarization additively over $\mathbb F_2$, and put
$e_{ij}=e_i+e_j$. Its diagonal vanishes: invariance and
$e_{12}=e_{13}+e_{23}$ give $b(e_{12},e_{12})=0$.
Invariance makes $b(e_{ij},e_{kl})$ constant for intersecting
supports, say $a$, whereas
$b(e_{23},e_{14})=b(e_{12},e_{14})+b(e_{13},e_{14})=0$.
Nondegeneracy forces $a=1$. Thus $b$ is the permutation dot product.
The invariant quadratic refinement is unique, since the ratio of two
refinements would be an invariant character, and $V$ has no invariant
linear forms. It is $q(v)=(-1)^{|v|/2}$, precisely the restriction of
the twist $i^{|v|}$ on five semions.

The splitting $\mathbb F_2^5=\mathbb F_2\oplus V$ and Shapiro's
lemma identify
\[
 H^i(A_5,\mathbb F_2)\oplus H^i(A_5,V)
 \simeq H^i(A_4,\mathbb F_2),
\]
with the constants inducing restriction
\cite[Proposition~3.3 and equation~(3.9)]{BrownLectures}.
In degree one both constant-coefficient groups vanish, since the
abelianizations are $1$ and $C_3$. In degree two both have order two
\cite[\S10]{KraftSchwarz}, and restriction is an isomorphism by
restriction--transfer at the odd index five
\cite[equation~(3.12) and Remark~3.5]{BrownLectures}.

Categorical actions lifting the given permutation action form a torsor
under $H^2(A_5,V)$, using the nondegenerate pairing to identify
$V$ with its character group \cite[\S2.1 and \S4]{GannonJonesPermutation}.
Thus the canonical factor-permutation action is the only one.
In the explicit tuple model of the factor-permutation action, its
composition maps are identities \cite[Proposition~4.1]{GannonJonesPermutation}.
A stabilizer fixes the simple tuple literally, so identity isomorphisms
give trivial multipliers; the even-weight subcategory inherits these
maps. The stabilizers are
$A_4$ and $S_3$, with four and three ordinary irreducible representations.
\end{proof}

\paragraph{The group $A_5$.}
For $A_5$ let $f_j(L)$ count the points fixed by an element of order $j$
on a nonunit neutral orbit $G/L$, and $h(L)$ be the indicated lower
bound on its lifts. Put $w=h-1+f_2+3f_3+10f_5$.
\[
\begin{array}{c|rrrrrrrrr}
 |L|&1&2&3&4&5&6&10&12&60\\\hline
 h&1&2&3&1&5&3&4&3&4\\
 f_2&0&2&0&3&0&2&2&1&1\\
 f_3&0&0&2&0&0&1&0&2&1\\
 f_5&0&0&0&0&2&0&1&0&1\\
 w&0&3&8&3&24&7&15&9&17
\end{array}
\]
Indeed $f_j=|C_G(g)|\,|[g]\cap L|/|L|$ from the subgroup list of
$A_5$ \cite{AtlasGroups}. Projective degrees have square sum $|L|$.
Cyclic multipliers vanish. Neither $6$ nor $12$ is a sum of at most
two squares, and $60$ is not a sum of three squares (first reduce
modulo four, then divide by four and reduce modulo eight). At order
ten some degree is one, so the multiplier is trivial and there are
four lifts.

Write $n_j=1+\sum f_j(L)$ over nonunit orbits. In the involution
sector singleton and two-element orbits contribute four and two lifts
respectively, by Lemma~\ref{ext:sector}. The other centralizers are
cyclic, so with $\epsilon=n_2\bmod2$,
\begin{equation}\label{ext:A5rank}
 r\ge5+u+n_2+3n_3+10n_5+3\epsilon
 \ge19+t+\sum w(L)+3\epsilon .
\end{equation}
Suppose $n_5>1$. The weights in \eqref{ext:A5rank} give the following
nonfree orbit sizes and upper bounds on the number $j$ of free orbits:
\[
\begin{array}{c|rrrrrrr}
\text{nonfree sizes}&1&1,30&1,15&6&6,15&6,30&6,15,30\\\hline
j\le&7&3&0&6&5&2&1 .
\end{array}
\]
These are obtained by summing the displayed weights with
$t+3(n_2\bmod2)\le25-\sum w(L)$. In the three singleton cases $n_5=2$,
so the two simple degrees in a five-grade express $B$ as a sum of
two integral squares. In particular $v_7(B)$ is even.

\paragraph{Case A1: a singleton and free orbits.}\label{ext:neutral-A1}
For a singleton and $j$ free orbits, $B=1+d^2+60\sum e_i^2$ is prime to three.
The even rank $2+60j$ forces $B$ even, since in odd dimension every
simple degree is odd. In \eqref{ext:primeorbits} the only way to fit
three primes into $1+t\le9$ is $2,5,7$ with $v_7(B)$ odd. Thus $B$
has at most two prime divisors. Solvability supplies a nontrivial
elementary subgroup of invertibles. Its order is a prime power of the
form $1+a+60b$, $0\le a\le1$, $0\le b\le7$, hence one of
$2,61,121,181,241,361,421$. A faithful $A_5$-module has dimension at
least three, leaving only the singleton in this subgroup. Hence $d=1$
and $B\equiv2\pmod4$. Its order-two invertible generates a
nondegenerate subcategory: if symmetric, it lies in its centralizer
of dimension $B/2$, forcing $4\mid B$. Thus $\mathcal B$ splits
equivariantly as a semion factor and its centralizer
\cite[Theorem~3.13]{DGNO}. Every nonidentity element fixes just the
unit in this centralizer, since it fixes exactly the two semion labels
in $\mathcal B$. Its $30j$ nonunit labels therefore form at most three
free orbits. A nontrivial elementary subgroup of its invertibles would
have order $61,121$ or $181$, none admitting a faithful action of $A_5$.
Hence $\mathcal B$ is semion.

\paragraph{Case A2: a singleton and a size-thirty orbit.}\label{ext:neutral-A2}
For orbit sizes one and thirty with $j\le3$, $B$ is again even and prime
to three. By \eqref{ext:primeorbits}, $1+t\le6$, and $v_7(B)$ is even,
so only two and five can divide $B$. The order of a nonzero elementary
subgroup of invertibles is $1+a+30b+60c$ with $a,b\in\{0,1\}$ and
$c\le3$. This excludes five-torsion and leaves the two-power orders
$2,32$. If its order is two, then the singleton is invertible,
$B\equiv2\pmod5$ and $B$ is a power of two. Reducing
$B=2+30e^2+60\sum f_i^2$ modulo four forces $e$ odd, hence $e=1$.
The invertible group then has order $32$, the only power of two among
$32,92,152,212$. Its two-torsion cannot have order two, since a cyclic
$2$-group admits no nontrivial $A_5$-action. Thus in either case
there is a binary space $U$ of order $32$ with nonzero orbits of sizes
one and thirty. Translation by its fixed vector $z\ne0$ preserves
the other orbit, so $b(z,x)=q(z)^{-1}$ there. Multiplicativity forces
$q(z)=1$. The induced quadratic form on $U/\langle z\rangle$ is
constant off zero, hence trivial by Lemma~\ref{ext:binary-reductions}.
Thus $U$ is Tannakian. Equivariantization gives an extension $H$ of
$A_5$ by $U^*$. A five-element fixes a nonzero vector in $U^*$ and has
an order-five lift; multiplying by that vector gives order ten.
Each of the two five-classes adds eleven to $c(A_5)$, and a nonzero
kernel class adds three, so $c(H)\ge45$.

\paragraph{Case A3: a singleton and a size-fifteen orbit.}\label{ext:neutral-A3}
For nonfree orbit sizes one and fifteen with $j=0$,
\eqref{ext:primeorbits} has $1+t=3$, so $B$ is a prime power with
prime $2,3$ or $5$. A nonzero elementary subgroup of invertibles has
order $2,16$ or $17$, forcing $B$ to be a power of two. If only the
singleton is in that subgroup, then $B=2+15e^2$; an even $e$ gives
$B\equiv2\pmod4$ and an odd $e$ gives $e=1$ and $B=17$.
Otherwise the size-fifteen orbit consists of invertibles and
$B=16+d^2$ for $d$ a power of two, forcing $d=4$ and $B=32$.
Lemma~\ref{ext:binary-reductions} makes the order-sixteen subgroup
Tannakian, so its squared order must divide $B$.
This is impossible.

\paragraph{Case A4: a size-six orbit.}\label{ext:neutral-A4}
For a size-six orbit, $n_3=1$, so $B$ is a square prime to three.
By \eqref{ext:primeorbits}, $1+t\le8$ permits only powers of two,
five or seven, or products of two and five. A seven-power requires
$t\ge6$. Its nonzero seven-torsion is a faithful $A_5$-module of
order at least $343$. An orbit of size fifteen leaves at most
$1+6+15+5\cdot60=322$ labels. Without that orbit the only pattern
with $t\ge6$ is six and free orbits, whose unions cannot contain
$342$ nonzero vectors. The full neutral rank is at most $367$.
A five-power is impossible, since the order of nonzero five-torsion
would be one or two modulo five. A two-power would give
$4^k\equiv1+6\cdot4^a\equiv0$ or $2\pmod5$, contradicting
$4^k\equiv\pm1\pmod5$.

It remains that $2,5\mid B$, in which case $t\ge6$. Besides the
size-six orbit there are only free orbits and possibly one of size
fifteen. Solvability supplies two-torsion invertibles, since
five-torsion was excluded. Their order $2^d$ cannot include the
size-six orbit, since $2^d\not\equiv7\pmod{15}$. The only
possibilities are $16$ and $256$. An involution would fix at least
sixteen vectors at $256$, whereas $n_2\le6$. At $16$ it is transitive off zero, hence Tannakian by
Lemma~\ref{ext:binary-reductions}.
Its equivariantization gives an extension $H$ with $c(H)\ge38$ by the
involution calculation in \hyperref[ext:neutral-B3]{Case B3}.
The two five-classes have centralizer of order five, since a
five-element fixes no nonzero vector in this kernel. The neutral
dimension for $\operatorname{Rep}(H)$ is $B/16^2$, divisible by five.
Hence each five-element fixes a nonunit neutral label, since otherwise
counting squared degrees modulo five gives $B/16^2\equiv1\pmod5$.
The two sectors add at least ten to $c(H)$, giving $r\ge48$.

\paragraph{Case B1: prime support when $n_5=1$.}\label{ext:neutral-B1}
Hence take $n_5=1$, so $B$ is a square prime to five. We claim that
$B$ has at most two prime divisors. Three odd primes require
$t\ge3+7+11-1=20$ by \eqref{ext:primeorbits}; any nonfree orbit
costs at least three, and all-free orbits give $r\ge22+t$.
Thus $t\le22$ and all nonunit orbit lengths are sixty, thirty or fifteen
(the latter two have weight three). This forces $B\equiv1\pmod3$,
raising the cost of three odd primes to $7+11+13$.
If $2\mid B$ but $3\nmid B$, two odd primes require $t\ge19$.
Divisibility by four requires an odd nonunit orbit. Its length must
be fifteen, with stabilizer $P=C_2^2$. Two such orbits already cost
six in weights and three in parity. Thus there is exactly one;
all remaining orbits are free except possibly one of length thirty. Choose generators $g,h$ of $P$. The length-thirty orbit has cyclic
stabilizer, whereas the length-fifteen orbit has three points fixed
by $P$, with conjugate multipliers. Lemma~\ref{ext:crossed-trace}
therefore gives
\[
 |\Irr(\mathcal F_g)^h|=1+3\delta,\qquad \delta\in\{1,-1\}.
\]
The left side is a nonnegative integer, so $\delta=1$ and four
grade-$g$ labels are fixed by $h$. They are fixed by $g$ as well,
by the crossed twist. Each consequently has stabilizer $P$ and
contributes four lifts by Lemma~\ref{ext:sector}. This adds twelve
to the involution contribution in \eqref{ext:A5rank}, a contradiction.

Finally suppose $2,3,p\mid B$ with $p\ge7$. Then $t\ge11$ and
$n_3\ge3$ since $1$ and $1+d^2$ never vanish modulo three.
In fact $n_3\ne3$. Put $M=\sqrt B$. Each grade-three simple of
degree $y$ lifts to objects of dimension $20y$, since its centralizer
is cyclic of order three. Modular divisibility gives $b=3M/y\in\mathbb Z$;
the grade has squared degree sum $B$, hence $\sum 1/b^2=1/9$.
With three denominators in increasing order the first two must be
$(4,5),(4,6)$ or $(5,5)$ by the successive remaining-mass bounds.
The corresponding values of the remaining square $b_3^2$ are
$3600/31$, $48$ and $225/7$, respectively, so none is an integer
square. Hence $n_3\ge4$, but supplying three
nonunit fixed points in the table costs at least $8+7=15$ in weights,
more than $44-19-11$. This proves
the claim.

\paragraph{Case B2: the elementary invertible subgroups.}\label{ext:neutral-B2}
If $B>1$, solvability again supplies $V=\operatorname{Inv}(\mathcal B)[p]
\simeq\mathbb F_p^d\ne0$ for some prime $p$. As before $d\ge3$,
$5\mid p^d-1$, and $p^d\le1+60t\le1321$. This leaves
$2^4,2^8,3^4$. At $3^4$, Jordan form gives $n_3\ge9$. An involution
is nonscalar with determinant one, so fixes a plane. Hence
$r\ge5+1+9+27+10>44$.
At $2^8$ an involution fixes at least sixteen vectors. Burnside's
formula gives at least $(255+15\cdot15)/60=8$ nonzero orbits.
The rank bound forces $n_3=1$ and $u\le10$. An involution
fixes exactly sixteen vectors of $V$, since the next possible fixed
space size is thirty-two. Every nonunit stabilizer in $\mathcal B$
is $1,C_2$ or $C_2^2$.
If $a,b,c$ count the corresponding nonzero orbits in $V$, then
$60a+30b+15c=255$, $2b+3c=15$, and $a+2b+c\le10$.
Thus $(a,b,c)=(3,0,5)$. These orbits use at least eight lifts.
Every four-group multiplier must be nontrivial: a trivial one on
$V$ adds three lifts and a new orbit with trivial multiplier costs
four. For two independent involutions, each four-group orbit contributes
$-3$ to \eqref{ext:crossed-trace-formula}, and cyclic stabilizers
contribute zero. Its nonnegative left side would therefore be at
most $1-3\cdot5<0$, a contradiction.

\paragraph{Case B3: the transitive binary module.}\label{ext:neutral-B3}
Suppose $|V|=16$ and $V\setminus0$ is one orbit.
Lemma~\ref{ext:binary-reductions} gives a Tannakian extension
$1\to W=V^*\to H\to A_5\to1$ with $W\setminus0$ transitive. Each stabilizer is a Sylow
four-group, and an involution $s$ lies in a unique such group,
$C_G(s)$. Counting incidences shows that $s$ fixes three nonzero
vectors. Hence $C_G(s)$ acts trivially on $W^s=(1+s)W$.
For a lift $a$ of $s$, the squares $(wa)^2=(1+s)w+a^2$ run over
$W^s$ and distinguish at least four classes over $[s]$, three of
order four. Consequently $k(H)\ge9$,
$c(H)\ge38$. The four square fibers above a chosen $s$ have four
elements each, and each is one $W$-conjugacy class since $(1+s)W=W^s$.
Their squares are fixed by $P=C_G(s)$, so each lift has centralizer
of order $16$. We show that one order-four lift $g$ has centralizer
$C_4\times C_4$, with $g$ generating a direct factor. Put
$W_0=W^s$ and choose $t\in P\setminus\langle s\rangle$. The
three nonidentity $v\in P$ fix $W_0$ pointwise and induce
isomorphisms $f_v=1+v:W/W_0\to W_0$. Hence
$f_{st}=f_s+f_t$, so $J=f_tf_s^{-1}$ permutes the three nonzero
points of $W_0$ cyclically. Take an involution lift $a$ of $s$.
Its square fiber is one $W$-class preserved by $t$; thus some lift
$b$ of $t$ centralizes $a$. Write $z=b^2\in W_0$. For $x,y\in W$
with $f_s(y)=f_t(x)$, the lifts $g=xa$, $h=yb$ still commute.
Their squares are $w=f_s(x)$ and $z+J^2w$. If $z=0$ choose any
$w\ne0$; otherwise choose $w=z$, so $z+J^2w=Jw$. The squares
are independent. Thus $g,h$ generate all of $C_H(g)\simeq C_4^2$.

Let $\mathcal E$ be the neutral component for $\operatorname{Rep}(H)$,
of square dimension $E=B/16^2$. The available rank is
\begin{equation}\label{ext:kernel-budget}
 u_{\mathcal E}\le44-c(H)\le6.
\end{equation}
An orbit in grade $g$ has size $1,2$ or $4$ and contributes
respectively $16,8$ or $4$ lifts by Lemma~\ref{ext:sector}, since
$C_H(g)/\langle g\rangle$ is cyclic. If $E$ is odd, the neutral
rank and $n_g$ are odd (all degrees are odd and $g$ has order four).
The singleton in grade $g$ then gives $r\ge38+16-4$.
For even $E$, Lemma~\ref{ext:binary-reductions} gives at least two
nonunit lifts. There are consequently at most eight lifts in grade $g$;
the three possible orbit sizes therefore give
\begin{equation}\label{ext:kernel-fixed-bound}
 n_g\in\{2,4,8\}.
\end{equation}
The two order-five classes of $H$ have cyclic centralizers of order
five since a five-element acts freely off zero on $W$. A fixed
nonunit would add ten to $c(H)$, so $E\equiv1\pmod5$.
By \eqref{ext:primeorbits}, the prime divisors of the even square $E$
sum to at most seven. They are therefore two and possibly three. Solvability supplies a nonzero
$V'=\operatorname{Inv}(\mathcal E)[p]$ for $p=2$ or $3$.
In particular, $5\mid |V'|-1$.

If $p=2$, $U=(V')^W$ is nonzero, since a $2$-group acts on a
nonzero binary space. Its action factors through $A_5$ and
$|U|\le1+6\cdot60$. Since $5\mid |U|-1$, its order is $16$ or $256$.
At $256$, its binary involution fixed space has at least sixteen
vectors, contrary to $n_g\le8$. Thus $|U|=16$.
If its nonzero vectors split into $5+10$ or $5+5+5$, an element of
order three in $A_5$ fixes at least three nonzero vectors, by the
subgroup table. Its order-three lift in $H$ has cyclic centralizer of
order three because it fixes no nonzero vector in $W$. This sector
adds at least nine to $c(H)\ge38$. Thus $U$ is transitive off zero and Tannakian by
Lemma~\ref{ext:binary-reductions}. Equivariantization
gives an extension $K$ of $H$ by $U^*$. Its extra nonzero kernel
class gives $c(K)\ge c(H)+3$. A lift of order at least four over
$W\setminus0$ adds another two. Otherwise the preimage of $W$ is
elementary abelian, so conjugation there factors through $A_5$,
and the $240$ elements above $W\setminus0$ form at least four
classes. Thus $c(K)\ge43$, leaving at most one nonunit orbit in the
neutral component for $\operatorname{Rep}(K)$. Its dimension
$E/16^2$ is square and $K$ is perfect. Lemma~\ref{ext:binary-reductions} forces the neutral component to be trivial, so $\mathcal C$
is a twisted double of $K$. Its nonzero kernel is one class of fifteen involutions,
whose centralizer has order $1024$. Since $1024/2=512$ is not a square, Lemma~\ref{ext:sector} adds two
to $c(K)$, giving rank at least $45$.

If $p=3$, the bound $|V'|\le1+6\cdot960$ and $5\mid |V'|-1$
give $|V'|=81$. The elementary abelian group $W$ acts trivially
on this four-dimensional ternary space: it acts diagonally and
its fifteen nontrivial characters are conjugate under $A_5$.
Thus $V'$ is a faithful $A_5$-module. The image $s$ of $g$ is a
noncentral involution and has determinant one. Its action cannot
be scalar; hence it fixes a plane, contradicting $n_g\le8$.

\paragraph{Case B4: the deleted binary module.}\label{ext:neutral-B4}
If $V\setminus0$ has several orbits, their sizes are $5+10$ or
$5+5+5$; the latter exceeds the weight budget. No proper union of
the five- and ten-point orbits has size $2^j-1$ for $1\le j<4$,
so $V$ is irreducible. Its five-point orbit spans $V$ and sums to
zero, identifying it with the deleted binary permutation module.
The radical of its invariant bicharacter is zero or $V$. In the latter case the twist is an invariant character,
hence trivial, giving a Tannakian extension $H$ as above. The dual
module is isomorphic to $V$ by the permutation dot product on the
even-weight space. An element $g$ of order three fixes a plane in $W$;
choose a lift $a$ of order three. Classes above $g$ are the
$C_G(g)$-orbits on $W/(1-g)W$, and $C_G(g)=\langle g\rangle$ acts
trivially via $a$. They have representatives $wa$, $w\in W^g$,
of orders $3,6,6,6$. There are two nonzero kernel classes of involutions, contributing
$2(1+2)=6$ beyond $c(A_5)$. Replacing the order-three class by
classes of orders $3,6,6,6$ contributes a further $3+18=21$.
The classes above $2,5,5$ retain at least their original contributions.
Thus
\[
 c(H)\ge c(A_5)+6+21=20+6+21=47.
\]

If the bicharacter is nondegenerate, the pointed subcategory on $V$
splits off $\mathcal B$ \cite[Theorem~3.13]{DGNO}; in particular
$16\mid\operatorname{rk}\mathcal B$. The five- and ten-point orbits have seven lifts by
Lemma~\ref{ext:deleted-action}. With that value of $u$ in
\eqref{ext:A5rank}, at most six ranks remain before introducing further
orbits. The weight and parity terms therefore permit one additional
size-thirty orbit with at most two free orbits, or at most six
additional free orbits. The former violates
divisibility of the rank by sixteen. In the latter case a nontrivial
complementary factor would have rank sixteen ($16\mid60j$,
$1\le j\le6$), but its
fifteen nonunit labels cannot form free orbits. Hence $\mathcal B$ is pointed on $V$. Lemma~\ref{ext:deleted-action}
gives seven neutral nonunit lifts, so the neutral sector has rank twelve. For $s=(12)(34)$ its sector has
four labels forming a torsor under $V/(1-s)V$. Indeed $V$ acts
transitively, since for two labels $Y,Z$ of grade $s$ a summand of
$Z\otimes Y^*$ carries $Y$ to $Z$. The two crossed braidings show
that $(1-s)V$ stabilizes every label, and it has the required order
four. The element $(13)(24)$ of the centralizer acts on this torsor
by an affine involution. Its linear part exchanges the two coordinates
$x_1+x_2,x_3+x_4$ of the coinvariants in the even-weight realization.
An affine involution with this linear part has two fixed points
(its translation is $(a,a)$). The involution sector therefore
contributes $4+4+2=10$ lifts, six more than the lower estimate in
\eqref{ext:A5rank}. The order-three and both order-five sectors have cyclic centralizers.
Their contributions are $3\cdot4,5,5$, respectively. Thus the rank
in this case is exactly $12+10+12+5+5=44$.


\subsection{The twisted doubles}
If $\mathcal B=\mathrm{Vec}$, $\operatorname{Rep}(G)$ is Lagrangian, so
$\C\simeq\mathcal Z(\mathrm{Vec}_G^\omega)$
\cite[Proposition~4.58 and Theorem~4.66]{DGNO}.
There is one simple per grade before equivariantization. Apply
Lemma~\ref{ext:sector} for cyclic quotients and the $D_8$ calculation
in Lemma~\ref{ext:groups} for the involutions. The ranks are
\[
\begin{array}{c|l}
A_5&5+4+3+5+5=22\\
\mathrm{PSL}_2(7)&6+5+3+4+7+7=32\\
A_6&7+5+9+9+4+5+5=44 .
\end{array}
\]
The doubles are integral and perfect: an invertible has central flux,
hence comes from a linear character of $G$ \cite[Theorem~2.1]{GruenMorrison}.
For the untwisted examples see \cite[discussion before Question~11.8]{BPBurnside},
\cite[\S4.2.5]{CZ}.

For each of the three simple groups $G$ the canonical Lagrangian is unique, since
$G$ has no nontrivial abelian normal subgroup. Braided equivalence
classes thus correspond to $\operatorname{Out}(G)$-orbits in
$H^3(G,\mathbb C^\times)\simeq H^4(G,\mathbb Z)$
\cite[Theorem~1.2 and Corollary~1.5]{NNLagrangian}.
For $A_5$ this group is $C_{30}$ \cite[p.~355]{EckmannHilton}.
Sylow restriction is injective on each primary part by transfer.
The outer involution \cite{AtlasA5} fixes the $2$- and $3$-parts
(the transposition $(45)$ centralizes $(123)$). On the translations
of $\mathbb F_5$ it is represented by $x\mapsto2x$, giving
$2^2=-1\bmod5$ in degree four. Thus there are $(30+6)/2=18$ orbits.

For $\mathrm{PSL}_2(7)=\mathrm{GL}_3(2)$ the cohomology group is
$C_{12}$ \cite[Lemma~4.1]{JFT}. The seven-point permutation
representation has cycle types $(1,3,3)$ and $(1,2,4)$ at orders three
and four. The second elementary symmetric function of its weights
gives $c_2=13\bmod3$ and $23\bmod4$ on these cyclic subgroups.
Thus $c_2$ generates $C_{12}$, and every automorphism fixes it since
the permutation character depends only on element order. All twelve
doubles are distinct, also distinguished by their twist spectra
\cite[\S\S3--4, \href{https://tqft.net/web/research/students/AngusGruen/Modular_Data/168/42/}{entry $(168,42)$}]{GruenMorrison}.
For $A_6$ the cohomology group is $H^4(BA_6,\mathbb Z)\simeq C_{60}$
\cite[calculation for $A_6$, p.~46 of the extended version]{HKY243}.
To compute the outer action, restrict the second Chern class of the
six-point permutation representation. An element of order four has
weights $0,1,2,3,0,2$ modulo four, so $c_2$ restricts to
$23\equiv3\pmod4$. This detects the entire $2$-primary part. Every
automorphism preserves the unique class of order four and hence fixes
that part. At a three-cycle and a product of two disjoint three-cycles
the restrictions of $c_2$ are respectively $2$ and $1$ modulo three.
Interchanging the two three-classes thus negates the $3$-primary part.
Restriction to a Sylow five-group is injective on the $5$-primary
part. Its generator powers $\pm1$ give one five-class and $\pm2$ the
other, so exchanging the classes acts on degree four by
$2^2\equiv-1\pmod5$. The two exchanges occur independently
\cite[Lemma~2.1]{KOZA6}. The orbit counts on $C_4,C_3,C_5$ are therefore
$4,2,3$, giving $24$ braided equivalence classes.

\subsection{The binary icosahedral family}\label{ext:binary}
Suppose $G=A_5$ and $\mathcal B$ is semion, as in
\hyperref[ext:neutral-A1]{Case A1}. Every component of $\mathcal F$ has dimension two
and consists of two invertibles. Hence $\mathcal F=\mathrm{Vec}_E^\omega$
for a central extension $1\to C_2\to E\to A_5\to1$. Applying
\cite[Corollary~3.30]{DMNOWitt} to $\operatorname{Fun}(A_5)$ gives
\[
 \mathcal Z(\mathrm{Vec}_E^\omega)\simeq \C\boxtimes\mathcal B^{\rm rev}.
\]
Since $\C$ is perfect, the pointed part on the left is the displayed semion
factor. A nontrivial linear character of $E$ would give an electric invertible
with twist one. Thus $E$ is perfect, and $E=\Gamma:=\mathrm{SL}_2(5)$
by the Schur cover of $A_5$ \cite{AtlasA5}. The associator on the kernel
is the semion associator, so $[\omega]|_{Z(\Gamma)}\ne0$.

Conversely, for any such class the pointed part $\mathcal S_\omega$ of
$\mathcal Z(\mathrm{Vec}_\Gamma^\omega)$ is semion
\cite[Example~2.1]{BGHNPRW}, \cite[\S4, Example~3]{MasonNgOrbifolds}.
Its centralizer $\mathcal K_\omega$ is therefore perfect, modular and
of dimension $120^2/2=7200$. It is group-theoretical
\cite[Proposition~8.44(i)]{ENOannals}.
The electric $\operatorname{Rep}(A_5)$ centralizes the semion factor,
since its representations are trivial on the central flux. Count in
the full double:
the two central fluxes each contribute nine simples (the Schur
multiplier of $\Gamma$ vanishes), and the seven other centralizers
are cyclic of orders $4,6,6,10,10,10,10$. Tensoring by the semion preserves the image of the flux in $A_5$.
Dividing the class counts by two therefore gives the sector ranks
\[
 \begin{array}{c|ccccc}
 A_5\text{-class}&1&2&3&5\mathrm A&5\mathrm B\\ \hline
 \text{rank}&9&2&6&10&10
 \end{array}
 \qquad 9+2+6+10+10=37.
\]
Kac and Todorov computed the same rank in the level-one icosahedral
orbifold \cite[\S6B, before Table~2]{KacTodorov}. The ordinary centralizer
and character data in their Table~2 give the simple dimensions
$1,2,3,4,5,6,12,20,30$ with multiplicities $1,2,2,2,1,1,20,6,2$,
respectively.

There are exactly thirty-six braided equivalence classes
$\mathcal K_\omega$. Indeed, $H^3(\Gamma,\mathbb C^\times)=C_{120}$,
and its odd residues are precisely the allowed classes
\cite[\S4, Example~3]{MasonNgOrbifolds}.
The charge of $\mathcal K_\omega$ determines the complementary semion:
the product has charge one and the two semion charges are
$\exp(\pm\pi i/4)$ \cite[Lemmas~5.26--5.27]{DMNOWitt}.
Hence equivalence of centralizers is equivalent to equivalence of their
full centers, whose pointed parts are characteristic.
Every abelian normal subgroup of $\Gamma$ lies in $Z(\Gamma)$ by the
simplicity of $A_5$; the restriction condition excludes $Z(\Gamma)$
from the Lagrangian parametrization
\cite[Theorem~1.2]{NNLagrangian}. Thus the canonical Lagrangian is unique,
and de-equivariantization identifies equivalent doubles precisely on
$\operatorname{Out}(\Gamma)$-orbits of the cocycles.
The noninner automorphism acts on $H^4(B\Gamma,\mathbb Z)=C_{120}$ by
multiplication by $49$: it is induced on $S^3/\Gamma$ by a self-map of
degree $49$ \cite[\S4.5]{HKWZ}, and naturality of transgression in
$S^3\to S^3/\Gamma\to B\Gamma$ gives this action on the Euler generator.
The involution $t\mapsto49t$ fixes the twelve odd multiples of five,
so the sixty allowed parameters give $(60+12)/2=36$ orbits.

\subsection{The deleted permutation gaugings}\label{ext:deleted-gauging}

For $G=A_5$ and $\mathcal P$ as in Corollary~\ref{ext:perfect}, a
$G$-gauging means the equivariantization of a faithful
braided $G$-crossed extension inducing the permutation action on $V$.

Such extensions exist. Apply \cite[Theorem~1.1 and Corollary~1.2]{GannonJonesPermutation}
to the $S_5$-permutation action on five copies of the semion category.
Restrict a crossed extension to $G$ and equivariantize. The diagonal
semion has its coordinatewise equivariant structure. Its invertible
lift is unique because $G$ has no nontrivial linear character. It has
order two and twist $i^5=i$, and it centralizes the electric
$\operatorname{Rep}(G)$. Its centralizer is therefore modular, contains
$\operatorname{Rep}(G)$, and has dimension $60^2\cdot16=57\,600$.
After de-equivariantization the neutral component is the centralizer of
the diagonal semion in the five semions, namely $\mathcal P$.
Since $A_5$ is normal in $S_5$, the crossed action of an odd
permutation preserves the restricted construction and the unique diagonal lift. Hence it preserves the resulting gauging, while inducing
the outer automorphism of $G$.

In every gauging, Section~\ref{ext:neutral} gives sector ranks
$12,10,12,5,5$ over
$1,2,3,5\mathrm A,5\mathrm B$, respectively. All labels of a fixed grade have the same degree,
since $V$ acts transitively on them. The degrees in the nonidentity
grades before equivariantization are consequently $2,2,4,4$.
The lift-dimension formula now gives integral degrees, at least $30$
for nonidentity flux. An invertible neutral lift lies over a $G$-fixed
vector, hence over zero, and $A_5$ has only the trivial linear character.
The gauging is therefore perfect.

There are exactly eighteen braided equivalence classes of these gaugings.
First, $\operatorname{Rep}(G)$ is their unique nontrivial Tannakian
subcategory. Any Tannakian subcategory has dimension at most
$\sqrt{57\,600}=240$, so none of its simples has nonidentity flux.
Its image under the braided forgetful functor $\mathcal P^G\to\mathcal P$
is a $G$-stable symmetric subcategory. Irreducibility of $V$ and
nondegeneracy of $q$ make this image trivial. Thus the subcategory lies
in $\operatorname{Rep}(G)$, and simplicity of $G$ proves uniqueness.
De-equivariantization by this subcategory is not pointed, as a
five-cycle grade contains a simple of degree four. The criterion
of \cite[Theorem~7.2]{NNWSubcategories} shows that these gaugings
are not group-theoretical. They are weakly group-theoretical by closure
under extensions and equivariantization \cite[Proposition~4.1]{ENOWGT}.

With $G$, $\mathcal P$ and its categorical action identified, crossed
extensions form a torsor under $H^3(G,\mathbb C^\times)=C_{30}$
\cite[\S2.2]{GannonJonesPermutation}. The categorical action is unique
by Lemma~\ref{ext:deleted-action}. Its changes of coherence over
the identity are measured by $H^1(G,V)$: the comparison on each $g$
is a character of $V$, the composition condition is the cocycle
identity, and modification gives a coboundary. Its vanishing was
proved in Lemma~\ref{ext:deleted-action}.
Every isometry of $(V,q)$ permutes the five nonzero isotropic vectors,
which span $V$. All coordinate permutations are isometries, so
$O(V,q)=S_5$. Since $\operatorname{Rep}(G)$ is intrinsic,
forgetting the identifications therefore divides the torsor by
$N_{S_5}(A_5)=S_5$, acting by conjugation on $G$. The gauging
constructed above provides an origin fixed by the outer involution;
inner automorphisms act trivially. The resulting action on $C_{30}$ is
the usual pullback action, with six fixed classes as computed for the
$A_5$ doubles. Hence there are $(30+6)/2=18$ classes.

This completes the proof of Corollary~\ref{ext:perfect}.

\subsection*{AI use declaration}
GPT-6 Astra and GPT-6.2 Astra assisted with the development and drafting of proofs,
the literature search, and the implementation and checking of the
finite calculations.

\begin{thebibliography}{99}
\bibitem{BDPcore}
S.~Burciu, J.~Dong and S.~Palcoux,
\href{https://hexagonmath.org/2611.00233v3}{\emph{Prime factors of fusion-category dimensions: weak group-theoreticity, rank and decomposition}},
Hexagon preprint 2611.00233v3 (2026).

\bibitem{AK}
A.~G.~Aleksandrov and \'E.~A.~Komissarchik,
\href{https://doi.org/10.1007/BF01089081}{\emph{Simple groups with a small number of conjugacy classes}},
J. Soviet Math. \textbf{37} (1987), no.~2, 917--924.
Russian edition, Zap. Nauchn. Sem. Leningrad. Otdel. Mat. Inst. Steklov.
\textbf{75} (1978), 5--15.

\bibitem{ABPP}
M.~A.~Alekseyev, W.~Bruns, S.~Palcoux and F.~V.~Petrov,
\href{https://doi.org/10.1063/5.0334158}{\emph{Classification of integral
modular data up to rank 13}}, J. Math. Phys. \textbf{67} (2026), no.~8, 081701. Numbering follows the \href{https://arxiv.org/abs/2302.01613v10}{open author version}.

\bibitem{Apostol}
T.~M. Apostol,
\href{https://doi.org/10.1007/978-1-4757-5579-4}{\emph{Introduction to Analytic Number Theory}},
Undergraduate Texts in Mathematics, Springer, New York, 1976.

\bibitem{Atlas}
J.~Berean-Dutcher, C.~C\'ordova, J.~Parra-Martinez, and R.~Tuddenham,
\href{https://arxiv.org/abs/2610.25439}{\emph{The Anyon Atlas}},
arXiv:2610.25439 (2026).

\bibitem{BGNPRW}
P.~Bruillard, C.~Galindo, S.-H.~Ng, J.~Y.~Plavnik, E.~C.~Rowell and
Z.~Wang, \href{https://doi.org/10.1016/j.jpaa.2015.11.010}{\emph{On the classification of weakly integral modular
categories}}, J. Pure Appl. Algebra \textbf{220} (2016), 2364--2388.

\bibitem{BNRWFinite}
P.~Bruillard, S.-H.~Ng, E.~C.~Rowell, and Z.~Wang,
\href{https://doi.org/10.1090/jams/842}{\emph{Rank-finiteness for modular categories}},
J. Amer. Math. Soc. \textbf{29} (2016), no.~3, 857--881.

\bibitem{BN}
S.~Burciu and S.~Natale,
\href{https://doi.org/10.1063/1.4774293}{\emph{Fusion rules of equivariantizations of fusion categories}},
J. Math. Phys. \textbf{54} (2013), no.~1, 013511.

\bibitem{BrownLectures}
K.~S.~Brown, \href{https://pi.math.cornell.edu/~kbrown/papers/cohomology_hangzhou.pdf}{\emph{Lectures on the cohomology of groups}},
in \emph{Cohomology of groups and algebraic K-theory},
Adv. Lect. Math., vol.~12, International Press, 2010, pp.~131--166.

\bibitem{BPBurnside}
S.~Burciu and S.~Palcoux,
\href{https://doi.org/10.1142/S0219199726500070}{\emph{Burnside-type
results for fusion categories}}, Commun. Contemp. Math. (2025), article 2650007.

\bibitem{AtlasGroups}
J.~H. Conway, R.~T. Curtis, S.~P. Norton, R.~A. Parker and
R.~A. Wilson, \emph{Atlas of Finite Groups}, Clarendon Press, Oxford, 1985.

\bibitem{CZ}
R.~Coquereaux and J.-B.~Zuber,
\href{https://doi.org/10.3842/SIGMA.2013.039}{\emph{Drinfeld doubles for finite subgroups of $SU(2)$ and $SU(3)$ Lie groups}},
SIGMA \textbf{9} (2013), 039.

\bibitem{CMSCyclotomic}
F.~Calegari, S.~Morrison and N.~Snyder,
\href{https://doi.org/10.1007/s00220-010-1136-2}{\emph{Cyclotomic integers, fusion categories, and subfactors}},
Comm. Math. Phys. \textbf{303} (2011), 845--896.

\bibitem{CasselsRobinson}
J.~W.~S.~Cassels, \emph{On a conjecture of R.~M.~Robinson about sums
of roots of unity}, J. Reine Angew. Math. \textbf{238} (1969), 112--131.

\bibitem{ConwaySloaneLattice}
J.~H. Conway and N.~J.~A. Sloane,
\href{https://neilsloane.com/doc/Me81.pdf}{\emph{On the enumeration of lattices of determinant one}},
J. Number Theory \textbf{15} (1982), 83--94.

\bibitem{CameronKiyota}
P.~J.~Cameron and M.~Kiyota,
\href{https://doi.org/10.1016/0021-8693(88)90285-2}{\emph{Sharp characters of finite groups}},
J. Algebra \textbf{115} (1988), no.~1, 125--143.

\bibitem{CPodd}
A.~Czenky and J.~Plavnik, On odd-dimensional modular tensor categories,
\emph{Algebra Number Theory} \textbf{16} (2022), 1919--1939,
\href{https://arxiv.org/abs/2007.01477}{corrected version}.

\bibitem{CGP}
A.~Czenky, W.~Gvozdjak and J.~Plavnik,
\href{https://doi.org/10.1016/j.jalgebra.2023.10.021}{\emph{Classification of low-rank odd-dimensional modular categories}},
J. Algebra \textbf{655} (2024), 223--293.

\bibitem{DLNCongruence}
C.~Dong, X.~Lin, and S.-H.~Ng,
\href{https://doi.org/10.2140/ant.2015.9.2121}{\emph{Congruence property in conformal field theory}},
Algebra \& Number Theory \textbf{9} (2015), no.~9, 2121--2166.

\bibitem{DP}
J.~Dong and S.~Palcoux,
\href{https://arxiv.org/abs/2510.10311}{\emph{A new criterion for integral
modular categorification}}, arXiv:2510.10311 (2025).

\bibitem{DGNO}
V.~Drinfeld, S.~Gelaki, D.~Nikshych and V.~Ostrik,
\href{https://doi.org/10.1007/s00029-010-0017-z}{\emph{On braided fusion categories. I}},
Selecta Math. (N.S.) \textbf{16} (2010), no.~1, 1--119.

\bibitem{EckmannHilton}
B.~Eckmann and P.~J.~Hilton,
\href{https://doi.org/10.1007/BF02566849}{\emph{On central group
extensions and homology}}, Comment. Math. Helv. \textbf{46}
(1971), 345--355.

\bibitem{EGNO}
P.~Etingof, S.~Gelaki, D.~Nikshych and V.~Ostrik,
\href{https://doi.org/10.1090/surv/205}{\emph{Tensor Categories}},
Math. Surveys Monogr., vol.~205, Amer. Math. Soc., Providence, RI, 2015.

\bibitem{ENOannals}
P.~Etingof, D.~Nikshych, and V.~Ostrik,
\href{https://doi.org/10.4007/annals.2005.162.581}{\emph{On fusion categories}},
Ann. of Math. \textbf{162} (2005), no.~2, 581--642.

\bibitem{ENOWGT}
P.~Etingof, D.~Nikshych and V.~Ostrik,
\href{https://doi.org/10.1016/j.aim.2010.06.009}{\emph{Weakly
group-theoretical and solvable fusion categories}}, Adv. Math.
\textbf{226} (2011), 176--205.

\bibitem{GNN}
S.~Gelaki, D.~Naidu and D.~Nikshych,
\href{https://doi.org/10.2140/ant.2009.3.959}{\emph{Centers of graded
fusion categories}}, Algebra Number Theory \textbf{3} (2009), 959--990.

\bibitem{GelakiNikshych}
S.~Gelaki and D.~Nikshych,
\href{https://arxiv.org/abs/math/0610726}{\emph{Nilpotent fusion categories}},
Adv. Math. \textbf{217} (2008), no.~3, 1053--1071.

\bibitem{GannonData}
T.~Gannon, \href{https://arxiv.org/abs/math/0103044}{\emph{Modular data:
the algebraic combinatorics of conformal field theory}},
J. Algebraic Combin. \textbf{22} (2005), 211--250.

\bibitem{HornJohnson}
R.~A.~Horn and C.~R.~Johnson, \emph{Matrix Analysis},
Cambridge University Press, Cambridge, 1985.

\bibitem{PrasadSL}
A.~Prasad, \href{https://arxiv.org/abs/0712.4051}{\emph{Representations
of $\mathrm{GL}_2(\mathbb F_q)$ and $\mathrm{SL}_2(\mathbb F_q)$, and
some remarks about $\mathrm{GL}_n(\mathbb F_q)$}}, 2007.

\bibitem{DigneMichel}
F.~Digne and J.~Michel,
\href{https://www.cambridge.org/core/books/representations-of-finite-groups-of-lie-type/3E2A81B39F7BBC7B3B0ECBD0AC0B8E60}{\emph{Representations of Finite Groups of Lie Type}},
London Math. Soc. Student Texts, vol.~21, Cambridge Univ. Press, 1991.

\bibitem{GruenMorrison}
A.~Gruen and S.~Morrison,
\href{https://doi.org/10.1512/IUMJ.2021.70.8309}{\emph{Computing modular data for pointed fusion categories}},
Indiana Univ. Math. J. \textbf{70} (2021), 561--593.

\bibitem{HardyWright}
G.~H.~Hardy and E.~M.~Wright, \emph{An Introduction to the Theory of Numbers},
4th ed., Clarendon Press, Oxford, 1960.

\bibitem{HM}
R.~Heffernan and D.~MacHale,
\href{https://doi.org/10.1353/mpr.2024.a931175}{\emph{Finite groups
with at most nine non-central conjugacy classes}},
Math. Proc. R. Ir. Acad. \textbf{124A} (2024), no.~1, 33--46.
\href{https://hdl.handle.net/10468/16879}{Author manuscript in the University College Cork repository}, \S4, pp.~6--10.

\bibitem{HR}
S.-M.~Hong and E.~C.~Rowell,
\href{https://doi.org/10.1016/j.jalgebra.2009.11.044}{\emph{On the classification of the Grothendieck rings of non-self-dual modular categories}},
J. Algebra \textbf{324} (2010), no.~5, 1000--1015.

\bibitem{JFT}
T.~Johnson-Freyd and D.~Treumann,
\href{https://doi.org/10.3842/SIGMA.2019.059}{\emph{Third homology of
some sporadic finite groups}}, SIGMA \textbf{15} (2019), 059, 38~pp.

\bibitem{JMNR}
C.~Jones, S.~Morrison, D.~Nikshych and E.~C.~Rowell,
\href{https://doi.org/10.1007/s00031-020-09576-2}{\emph{Rank-finiteness for $G$-crossed braided fusion categories}},
Transform. Groups \textbf{26} (2021), no.~3, 915--927.

\bibitem{Kirillov}
A.~Kirillov Jr.,
\href{https://arxiv.org/abs/math/0401119}{\emph{On $G$-equivariant modular categories}},
arXiv:math/0401119 (2004).

\bibitem{KraftSchwarz}
H.~Kraft and G.~W.~Schwarz,
\href{https://doi.org/10.1016/j.jalgebra.2006.12.025}{\emph{Compression of finite group actions and covariant dimension}},
J. Algebra \textbf{313} (2007), 268--291;
\href{https://arxiv.org/abs/math/0609253}{author preprint}.

\bibitem{Landauer}
C.~A.~Landauer,
\href{https://doi.org/10.7907/V2KC-3W98}{\emph{Simple groups with 9, 10, and 11 conjugate classes}},
Ph.D. thesis, California Institute of Technology, 1973.

\bibitem{LPR}
Z.~Liu, S.~Palcoux, and Y.~Ren,
\href{https://doi.org/10.1142/S0129167X23500301}{\emph{Interpolated family of non-group-like simple integral fusion rings of Lie type}},
Internat. J. Math. \textbf{34} (2023), no.~6, 2350030.

\bibitem{NNLagrangian}
D.~Naidu and D.~Nikshych,
\href{https://doi.org/10.1007/s00220-008-0441-5}{\emph{Lagrangian
subcategories and braided tensor equivalences of twisted quantum doubles
of finite groups}}, Comm. Math. Phys. \textbf{279} (2008), no.~3, 845--872.

\bibitem{NSWGauss}
S.-H.~Ng, A.~Schopieray and Y.~Wang,
\href{https://doi.org/10.1007/s00029-019-0499-2}{\emph{Higher Gauss sums of modular categories}},
Selecta Math. (N.S.) \textbf{25} (2019), Paper No.~53.

\bibitem{Nevins}
M.~Nevins, \href{https://doi.org/10.4153/CJM-2005-026-1}{\emph{Branching
rules for principal series representations of $\mathrm{SL}(2)$ over a
$p$-adic field}}, Canad. J. Math. \textbf{57} (2005), 648--672.

\bibitem{NSfs}
S.-H. Ng and P. Schauenburg, \href{https://doi.org/10.1016/j.aim.2006.07.017}{\emph{Frobenius--Schur indicators and exponents of spherical categories}}, Adv. Math. \textbf{211} (2007), 34--71.

\bibitem{NSind}
S.-H. Ng and P. Schauenburg, \href{https://doi.org/10.1007/s00220-010-1096-6}{\emph{Congruence subgroups and generalized Frobenius--Schur indicators}}, Comm. Math. Phys. \textbf{300} (2010), 1--46.

\bibitem{SchopFormal}
A.~Schopieray, \href{https://doi.org/10.1016/j.jalgebra.2020.10.014}
{\emph{Norm, trace, and formal codegrees of fusion categories}},
J. Algebra \textbf{568} (2021), 362--385.

\bibitem{SchopierayFPF}
A.~Schopieray,
\href{https://doi.org/10.1090/conm/813/16287}{\emph{Fixed-point-free fusion automorphisms}},
in \emph{Quantum Symmetries: Tensor Categories, TQFTs, and Vertex Algebras},
Contemp. Math. \textbf{813}, Amer. Math. Soc., 2025, 151--174.

\bibitem{SerreLocal}
J.-P. Serre, \href{https://doi.org/10.1007/978-1-4757-5673-9}{\emph{Local Fields}}, translated by M.~J. Greenberg,
Graduate Texts in Mathematics 67, Springer, 1979.

\bibitem{Stacks}
The Stacks Project Authors, \emph{The Stacks Project},
\href{https://stacks.math.columbia.edu/tag/04GH}{Tag 04GH}.

\bibitem{VeraLopez85}
A.~Vera-L\'opez and J.~Vera-L\'opez,
\href{https://doi.org/10.1007/BF02764723}{\emph{Classification of finite groups according to the number of conjugacy classes}},
Israel J. Math. \textbf{51} (1985), no.~4, 305--338.

\bibitem{VeraLopez86}
A.~Vera-L\'opez and J.~Vera-L\'opez,
\href{https://doi.org/10.1007/BF02766124}{\emph{Classification of finite groups according to the number of conjugacy classes. II}},
Israel J. Math. \textbf{56} (1986), no.~2, 188--221.

\bibitem{VeraSang}
A.~Vera-L\'opez and J.~Sangroniz,
\href{https://doi.org/10.1002/mana.200410508}{\emph{The finite groups with thirteen and fourteen conjugacy classes}},
Math. Nachr. \textbf{280} (2007), 676--694.

\bibitem{Washington}
L.~C.~Washington,
\href{https://doi.org/10.1007/978-1-4612-1934-7}{\emph{Introduction to
Cyclotomic Fields}}, second ed., Grad. Texts in Math., vol.~83,
Springer, New York, 1997.

\bibitem{AtlasA5}
R.~A.~Wilson et al.,
\href{https://brauer.maths.qmul.ac.uk/Atlas/v3/alt/A5/}{\emph{ATLAS of
Finite Group Representations: Alternating group $A_5$}}.

\bibitem{XuDong}
F.~Xu and J.~Dong,
\href{https://doi.org/10.4153/S0008439525101197}{\emph{On perfect weakly group-theoretical modular categories}},
Canad. Math. Bull. \textbf{69} (2026), no.~1, 292--298.

\bibitem{BGHNPRW}
P.~Bruillard, C.~Galindo, T.~Hagge, S.-H.~Ng, J.~Y.~Plavnik, E.~C.~Rowell
and Z.~Wang, \href{https://doi.org/10.1063/1.4982048}{\emph{Fermionic modular
categories and the 16-fold way}}, J. Math. Phys. \textbf{58} (2017), 041704.

\bibitem{DMNOWitt}
A.~Davydov, M.~M\"uger, D.~Nikshych and V.~Ostrik,
\href{https://arxiv.org/abs/1009.2117}{\emph{The Witt group of non-degenerate
braided fusion categories}}, J. Reine Angew. Math. \textbf{677} (2013), 135--177.

\bibitem{HKWZ}
C.~Hayat-Legrand, E.~Kudryavtseva, S.~Wang and H.~Zieschang,
\href{https://rendiconti.dmi.units.it/volumi/32s1/08.pdf}{\emph{Degrees of
self-mappings of Seifert manifolds with finite fundamental groups}},
Rend. Istit. Mat. Univ. Trieste \textbf{32}, suppl.~1 (2001), 131--147.

\bibitem{KacTodorov}
V.~G.~Kac and I.~T.~Todorov,
\href{https://arxiv.org/abs/hep-th/9612078}{\emph{Affine orbifolds and rational
conformal field theory extensions of $W_{1+\infty}$}},
Comm. Math. Phys. \textbf{190} (1997), 57--111.

\bibitem{MasonNgOrbifolds}
G.~Mason and S.-H.~Ng,
\href{https://doi.org/10.21915/BIMAS.2019101}{\emph{Generalized twisted
quantum doubles of a finite group and rational orbifolds}},
Bull. Inst. Math. Acad. Sin. (N.S.) \textbf{14} (2019), 1--13.

\bibitem{GannonJonesPermutation}
T.~Gannon and C.~Jones,
Vanishing of categorical obstructions for permutation orbifolds,
\emph{Comm. Math. Phys.} \textbf{369} (2019), 245--259.
\href{https://doi.org/10.1007/s00220-019-03288-9}{doi:10.1007/s00220-019-03288-9}.
\bibitem{HKY243}
A.~Hoshi, M.-c.~Kang and A.~Yamasaki,
\href{https://arxiv.org/abs/1710.01958}{\emph{Degree three unramified cohomology groups and Noether's problem for groups of order $243$}},
J. Algebra \textbf{544} (2020), 262--301. The link gives the extended version.

\bibitem{KOZA6}
J.~Keum, K.~Oguiso and D.-Q.~Zhang,
\href{https://arxiv.org/abs/math/0408105}{\emph{Extensions of the alternating group of degree $6$ in the geometry of K3 surfaces}},
Eur. J. Combin. \textbf{28} (2007), 549--558.

\bibitem{NNWSubcategories}
D.~Naidu, D.~Nikshych and S.~Witherspoon,
\href{https://doi.org/10.1093/imrn/rnp084}{\emph{Fusion subcategories of representation categories of twisted quantum doubles of finite groups}},
Int. Math. Res. Not. IMRN \textbf{2009} (2009), no.~22, 4183--4219.
\end{thebibliography}
\end{document}

